%% =============================================================================
%% Pre-thesis 1 -- Crypto World Bank -- v24 (BRACU front matter + ACM body)
%% -----------------------------------------------------------------------------
%% v24 improvements applied (autonomous agent, zkEVM, Chainlink oracles,
%%   compliance, actor taxonomy, DB schema, revenue model; June 2026;
%%   based on CWB_v24_Improvement_Analysis.md):
%%   * v24-A: Autonomous AI Agent System -- six-step pipeline (Qwen3-8B +
%%     live on-chain context injection + MCP tool server + human gate +
%%     EIP-7702 session key execution + status monitoring). Covered in:
%%     Ch.1 Background (new paragraph), Objectives (Obj.3 expanded),
%%     Research Contribution 3 (new Extension sub-paragraph), Ch.3
%%     User Taxonomy (actor table), Credit Passport (credit tier schedule),
%%     Ch.4 AI/ML Support section (major overhaul), Appendix A (major
%%     overhaul). New DB tables: sessions, agent_action_log.
%%   * v24-B: Protocol/Oracle upgrades -- Polygon zkEVM Cardona replaces
%%     Amoy PoS; Chainlink Functions (trustless ML oracle), Automation
%%     (installment triggers), Price Feeds (BDT/USD, ETH/USD), Proof of
%%     Reserve (WorldBankReserve). Covered in: Ch.1 Blockchain Justification
%%     para E, Ch.3 Blockchain Platform Selection, Oracle Architecture
%%     (major rewrite), Appendix A tech stack table, Appendix C addresses.
%%   * v24-C: Compliance architecture -- FATF Travel Rule data packet
%%     (specify only), SAR workflow (Isolation Forest → aml-alert →
%%     freeze), Regulator audit request sequence (specify only), failed-
%%     attempt lockout. Covered in: Ch.3 Upward Facility, Reentrancy
%%     Security, Compliance sections; Ch.6 Future Work.
%%   * v24-D: Formal actor taxonomy -- nine actors (A1–A9, five primary
%%     four secondary) + full permission matrix. Covered in: Ch.3
%%     User Taxonomy section (two new tables).
%%   * v24-E: Database schema -- 19 entities (up from 15): sessions,
%%     agent_action_log, interest_rate_tier, assets tables added.
%%     Append-only audit_logs DB role + RLS. Composite indexes.
%%     Covered in: Ch.3 Data Model (entity summary, EER constructs,
%%     normalization, indexing, functional deps, integrity constraints,
%%     on-chain/off-chain partitioning, digital identity sections).
%%   * v24-F: Revenue model -- no native CWB governance token (FTX hard
%%     rule). Revenue stream taxonomy table. Covered in: Ch.5 Revenue.
%%   * v24-G: Sprint plan updates -- Sprint 1: 70 pts, Sprint 2: 92 pts,
%%     Sprint 3: 80 pts. New user stories US-1.12--1.21, US-2.12--2.20,
%%     US-3.11--3.18. Design decisions table additions.
%% v24 Chapter 4 Methodology -- all improvements applied (June 2026):
%%   * sec:aiml-support -- Major overhaul: autonomous agent six-step pipeline
%%     paragraph, MCP tool server table (16 tools), Authority Brief UI verbatim
%%     block, per-user context namespace verbatim block, extended agent
%%     capabilities list.
%%   * sec:llm-eval -- Three new evaluation items: action accuracy, human-gate
%%     bypass test, session key scope enforcement test.
%%   * tab:sprint1-contracts -- Added US-1.12--1.16 (InsuranceFund, reserve
%%     summary, Chainlink Feeds, PoR, append-only audit_logs). Total 37 pts.
%%   * tab:sprint1-frontend -- Added US-1.17--1.21 (sessions, assets, rate
%%     tier, indexes, zkEVM migration). Total 33 pts. US-1.10 updated 19 tables.
%%   * Sprint 1 total: 42 -> 70 story points.
%%   * tab:sprint2 -- Added US-2.12--2.20 (MCP server, agent chat, EIP-7702,
%%     authority brief, Chainlink Automation, EMI reminders, credit tier,
%%     agent_action_log, interest_rate_tier integration). Total 92 pts.
%%   * tab:sprint3 -- Added US-3.11--3.18 (Chainlink Functions, Graph subgraph,
%%     SAR workflow, Reserve Transparency Dashboard, 300-client Foundry sim,
%%     no-privilege Foundry test, pool Isolation Forest, freezeAccount).
%%     Total 80 pts.
%%   * tab:design-decisions -- Added 4 rows: network zkEVM vs Amoy, oracle
%%     Chainlink Functions vs commit-reveal, agent MCP vs direct API,
%%     session EIP-7702 vs ERC-4337 paymasters.
%%   * sec:abm-sim -- Updated Polygon Amoy -> zkEVM Cardona; 50-200 clients ->
%%     300 clients; Hardhat -> Foundry simulation for the v24 large-scale run.
%%   * sec:realtime -- Updated Polygon Amoy WebSocket URL -> Cardona.
%% v23 improvements applied (implementation plan alignment pass, May 2026;
%%   based on CWB_v22_Implementation_Plan.md re-evaluation):
%%   * Improvement v23-1: Contribution 1 scope qualifier added -- distinguishes
%%     specified vs. partially implemented vs. fully implemented tiers; World
%%     Bank Reserve (Tier 1) and Local Bank lending flow noted as deployed;
%%     National Bank cross-tier automation scoped to Sprint 2.
%%   * Improvement v23-2: Contribution 2 scope qualifier added -- GroupLendingPool
%%     framed as specification contribution + partial proof-of-concept; mutual
%%     liability enforcement and over-indebtedness controls scoped to Sprint 2.
%%   * Improvement v23-3: Contribution 3 scope qualifier added -- oracle-mediated
%%     AI/ML integration described with honest framing; benchmark table slot
%%     and commit-reveal wiring added for Sprint 3.
%%   * Improvement v23-4: Contribution 4 scope qualifier added -- age-range
%%     Groth16 circuit added as proof-of-concept implementation; wallet-velocity
%%     zkAML circuit scoped to Future Work with IACR ePrint 2025/465 reference.
%%   * Improvement v23-5: Revenue figure caption updated to clarify $137.6M is
%%     the theoretical 10-year maturity ceiling, not a near-term projection;
%%     cross-reference sentence added linking ramp table to the figure.
%%   * Improvement v23-6: Revenue table footnote extended to explicitly state
%%     the $137.6M is the 10-year ceiling at 68,800 active loans across 25
%%     local banks; break-even loan count (~300) and Year 1-2 scope noted.
%%   * Improvement v23-7: ETH price planning assumption footnote added on first
%%     use in the revenue projection section; stablecoin-first design note added.
%%   * Improvement v23-8: GDPR Art. 17 (right to erasure) vs. permanent SBT
%%     default record addressed -- paragraph added to sec:credit-passport
%%     explaining pseudonymity, off-chain deletion, and credit-bureau analogy.
%%   * Improvement v23-9: Institutional trust bootstrapping paragraph added to
%%     the Governance section -- three-pronged strategy: academic pilot route,
%%     regulatory sandbox route, BRAC alignment route.
%%   * Improvement v23-10: Future Work item 14 updated -- Singapore MAS Project
%%     Guardian / UAE DIFC as Phase 1 target; Bangladesh CBDC integration moved
%%     to Phase 3 long-term target with explicit timeline qualification.
%%   * Improvement v23-11: Five new references added (R47-R51): BCCC
%%     DeFiFraudTrans-2025, Elliptic Data Set, Elliptic++ Dataset, Weber et al.
%%     IBM/MIT Watson AML paper, AMD ROCm 7.0.2 release notes.
%% v22 improvements applied (architectural alignment pass, May 2026):
%%   * Improvement 1: SavingsVault now explicitly aligned with ERC-4626
%%     (tokenized yield-bearing vault standard); FixedDeposit aligned with
%%     ERC-7540 (async extension for time-gated redemptions).  Paragraph added
%%     to sec:savings-vault; new reference [R46] added.
%%   * Improvement 2: Credit Passport SBT given a public ICreditPassport
%%     interface with standardised getScore() view function for cross-protocol
%%     composability.  Paragraph added to sec:credit-passport.
%%   * Improvement 3: Oracle trust assumption strengthened with a multi-signer
%%     attestation requirement (2-of-3 Safe multisig) on the commitRiskScore
%%     step.  Paragraph added to sec:oracle-architecture.
%%   * Improvement 4: NettingEngine settlement failure path specified:
%%     settlePartial(cycleId) fallback with queued remainder and defaulting-
%%     participant flag.  Paragraph added to sec:netting-settlement.
%%   * Improvement 5: LLM assistant live context injection layer added:
%%     Express.js backend appends user's on-chain state (loan status, SBT risk
%%     tier, pool utilisation) as structured JSON context in system prompt.
%%     Paragraph added to Appendix A (sec:local-llm).
%%   * Improvement 6: GroupLendingPool over-indebtedness check tightened with
%%     cross-pool query and income-hash debt-to-income ratio enforcement.
%%     Existing paragraph at sec:over-indebt expanded.
%%   * Improvement 7: ERC-3643 (T-REX permissioned token standard) noted in
%%     competitive analysis table footnote and institutional-tier RBAC prose.
%%   * Improvement 8: IBLP rate bound changed from static governance parameter
%%     to live cross-contract call (require proposedRate <= localBank
%%     .getCurrentBorrowRate() - delta) to close inter-update arbitrage gap.
%%   * Model change: LLM replaced from Qwen3.5-9B (no such Alibaba release;
%%     the 3.5/3.6 family starts at 27B) to Qwen3-8B (Alibaba Cloud, released
%%     April 28 2025, Apache 2.0).  Updated: Appendix A model paragraph,
%%     latency benchmark, future-work item, two "9B-parameter" inline mentions.
%%     Qwen3-8B specs: 8.19B dense params, 32K native / 131K YaRN context,
%%     100+ languages, switchable thinking/non-thinking modes, text-only.
%%     Multimodal note removed; [R46] updated with Qwen3 arXiv:2505.09388.
%% v21 corrections applied (final audit pass -- all residual issues identified
%%   in post-v20 audit against CWB_v17_Unified_Improvement_Suggestions.md and
%%   online research, May 2026):
%%   * Fix Issue 1: SBT revocation contradiction at sec:liquidation-hf -- line
%%     previously said "SBT revocation" which contradicts Option A (SBT never
%%     revoked); corrected to "permanent riskTier downgrade and effective
%%     future-borrowing ban" consistent with Buterin SBT design principle.
%%   * Fix Issue 2: Model name was "Qwen3.5-9B-Instruct" in v20; corrected
%%     in v21 -- fully superseded in v22 by model replacement (Qwen3-8B).
%%   * New 1: Qwen3.5-9B multimodal note (v21) removed in v22; Qwen3-8B
%%     is text-only (vision model is the separate Qwen3-VL-8B).
%%   * Fix Issue 3: Competitive table (tab:competitive-analysis) hierarchy
%%     label "Local → Borrower" corrected to "Local → Client".
%%   * Fix Issue 4: Five primary actors list -- "Borrower" corrected to
%%     "Client" in the legacy actor enumeration at sec:user-taxonomy.
%%   * Fix Issue 5: Local Bank Contract description -- "Borrower registration"
%%     corrected to "Client registration".
%%   * Fix Issue 6: "women-led borrower groups" corrected to "women-led client
%%     groups" in Bangladesh financial inclusion paragraph.
%%   * Fix Issue 7: "rural Sylhet borrowers" corrected to "rural Sylhet
%%     clients" in Bangladesh positioning section.
%%   * Fix Issue 8: "Retail level (borrowers and depositors)" corrected to
%%     "Retail level (clients and depositors)" in service delivery subsection.
%%   * Fix Issue 9: "institutional borrower" in syndicated lending context
%%     corrected to "institutional client" in prose (contract role name
%%     Borrower retained in code-context only).
%%   * Fix Issue 12: DB schema table -- "One active request per borrower per
%%     bank" corrected to "per client per bank".
%%   * Fix Issue 13: ERC-4337 EntryPoint version confirmed as v0.7
%%     (0x0000000071727De22E5E9d8BAf0edAc6f37da032) -- note added.
%%   * New 2: ERC-4337 adoption stats updated (40M+ smart accounts, 100M+
%%     transactions as of 2025-2026).
%%   * New 3: USDC supply figure updated ($77.3B as of April 2026, DeFiLlama).
%%   * ROCm caveat footnote added for RX 9060 XT llama.cpp known bug
%%     workaround (Linux, GGML_CUDA_NO_PEER_COPY=1).
%% v20 corrections applied (completing all remaining items from
%%   CWB_v17_Unified_Improvement_Suggestions.md not yet applied in v19):
%%   * Behavioral biometrics: removed from ACTIVE feature set in 4 locations
%%     (AI/ML pipeline paragraph, defense-in-depth L3 table, zkAML circuit
%%     prose, and Sybil-abuse threat-model row); one-sentence Future Work note
%%     added at end of Section sec:aiml-support; 18 core on-chain features
%%     retained unchanged (Part 4.3 of improvement suggestions)
%%   * Mesa ABM: all remaining body references replaced with on-chain Hardhat
%%     simulation: SDLC figure caption, RQ5 evaluation note, and Future Work
%%     item #9 (Part 5.4 / improvement suggestions Part 8 "ABM empirically
%%     grounded" fix)
%%   * Research Contribution 1 hierarchy label: "Local Bank → Borrower" →
%%     "Local Bank → Client" (Part 1.1 terminology change, two remaining
%%     instances in sec:research-contribution and InterBankLendingPool prose)
%%   * SHAP explainability paragraph: "output to borrowers" → "output to clients"
%%     (Part 1.1 residual terminology fix)
%%   * zkAML: third property reworded to remove behavioral biometrics reference
%%     (wallet velocity check retained; biometrics layer reference removed)
%% v19 corrections applied (from CWB_v17_Unified_Improvement_Suggestions.md):
%%   * "Borrower" → "Client" terminology throughout prose (not in code, not in
%%     citations); BORROWER entity noted as pending rename to CLIENT in final schema
%%   * "Lending platform" → "cryptocurrency exchange and financial services platform"
%%     in all self-referential standalone descriptors (Abstract, Ch.1, Ch.3)
%%   * Olympus DAO reframed as cautionary precedent (99% value loss, 2022)
%%   * FTX positioning paragraph added to Section 2.3 (Comparative Protocol Analysis)
%%   * Ch.1 duplicated market/customer/partner tables replaced with concise paragraphs
%%     + forward references to Ch.5 tables (labels: tab:market-sizing-ch5 etc.)
%%   * Embedded Ch.2 Preliminaries glossary replaced with Glossary cross-reference
%%   * Blockchain Fundamentals subsection: skip-ahead note added for expert readers
%%   * MARKET_DATA and PROFILE_SETTING entity table: relationships and attributes added
%%   * AI_ML_SECURITY entity: "Association/Derived Entity" classification added
%%   * Kinked rate: floor/ceiling require guards documented in prose
%%   * UpwardDepositFacility: 20%/24hr bank-run circuit breaker clause added
%%   * SyndicatedLoan: "unanimous" → "supermajority (75%-by-capital)" corrected
%%   * RQ4: rewritten to precise language + scope note (testnet only)
%%   * Group lending lifecycle: "borrowers" → "clients" throughout
%%   * 99.5% SLA: split into API layer (team-controlled) vs chain layer (Polygon)
%%   * ABM calibration: Bangladesh-specific calibration moved to Future Work;
%%     public DeFi data (Aave v3, DefiLlama) used for prototype calibration
%%   * DefiLlama TVL citations: access date "accessed March 2026" added
%%   * SHAP and ERC-4337 sections: borrower-facing → client-facing
%%   * Why Retail section: hierarchy label updated WB→NB→LB→Client
%%   * Sprint 2/3 tables: borrower-bank → client-bank, borrower assistance → client
%%   * Bangladesh deployment: future-work scope note added to rural Sylhet paragraph
%%   * FX currency risk: "borrower in rural Bangladesh" → "retail client"
%%   * Revenue table footnote: "borrower's total APR" → "retail client's total APR"
%% v19 additional corrections (completing interrupted v18/v19 pass):
%%   * Section sec:abm-sim: Mesa ABM replaced with on-chain Hardhat simulation
%%     plan (50-200 clients, 6 banks, Polygon Amoy, seeded randomness per Part 5)
%%   * Behavioral biometrics removed from active ML feature set; moved to Future
%%     Work with one-sentence note; 18 core blockchain transaction features retained
%%   * SDLC caption and RQ5 evaluation: Mesa ABM references updated to on-chain sim
%%   * Sprint 3 future-work item: Mesa ABM -> on-chain simulation deliverable
%%   * File header version number updated to v19
%% Format intent (per Documentation/Writing format/checklist.md and
%% format_improvement_suggestions.md):
%%   * Front matter (title / declaration / approval / ethics / abstract /
%%     dedication / acknowledgment) preserved in BRACU institutional style.
%%   * Body chapters (Ch. 1 onward) follow ACM body conventions: booktabs
%%     tables, \Description-friendly figure captions, no decorative color in
%%     section headings, references prepared for ACM-Reference-Format.bst
%%     migration via Documentation/references.bib (114 entries).
%% Diagram strategy:
%%   * Figure placeholders kept at the v10 positions; vector PDFs for the
%%     redesigned diagrams live in Diagrams/mermaid-pdf/improved diagrams/.
%%     The diagram-audit pass (next session) will swap the \FigureImageMaxFit
%%     targets in place without moving anchors.
%% Strict v15 invariants (do NOT regress in any follow-up session):
%%   * Strict B&W textual content (greyscale palette only, no decorative colour)
%%   * ACM body conventions from Contents page onward (single-spacing, Libertine
%%     font with lmodern fallback, booktabs-style table heads, no decorative
%%     section heading colour, frame=tb on listings)
%%   * BRAC front matter preserved (title page, declaration, approval, ethics,
%%     abstract, dedication, acknowledgment) -- content only; colours neutralized
%%   * Numbered inline references kept; references.bib reserved for future ACM
%%     migration; R-cited entries R1--R45 anchored at end of references chapter
%%   * Diagram audit is a SEPARATE session; do not regenerate Mermaid sources in
%%     this content pass. Diagram PDFs live in Diagrams/mermaid-pdf/improved
%%     diagrams/ and the \graphicspath search already prefers that folder.
%% =============================================================================
\documentclass[12pt,a4paper]{report}
% If Overleaf still times out: (1) upload all image files the thesis references;
% (2) recompile twice so references resolve; (3) in Overleaf, avoid ``draft`` mode for
% graphicx in project settings, or temporarily enable it for quick checks only.

\usepackage[utf8]{inputenc}
\usepackage[T1]{fontenc}
\usepackage{listings}
\lstdefinelanguage{Solidity}{
  keywords={pragma,solidity,contract,function,returns,external,internal,public,private,
             view,pure,payable,mapping,address,uint256,uint,bool,bytes32,string,
             event,emit,modifier,require,revert,import,is,memory,storage,indexed,
             bytes,keccak256,struct,enum,constructor,fallback,receive},
  sensitive=true,
  comment=[l]{//},
  morecomment=[s]{/*}{*/},
  morestring=[b]",
}
% v15: greyscale-only listings. ACM body convention is frame=tb (top/bottom
% rules only) and black-only typography; the slight grey for comments preserves
% readability without introducing colour.
\lstset{
  basicstyle=\ttfamily\scriptsize,
  keywordstyle=\bfseries,
  commentstyle=\color{black!55}\itshape,
  stringstyle=\itshape,
  breaklines=true,
  frame=tb,
  rulecolor=\color{black!55},
  captionpos=b,
  numbers=left,
  numberstyle=\tiny\color{black!55},
  xleftmargin=2em,
}
% Body uses Linux Libertine (text) + Inconsolata (mono) to approximate the
% ACM acmsmall body face. We retain the report class so the BRAC front matter
% (title page, declaration, approval, ethics, abstract, dedication, ack.) keeps
% its institutional layout; from the Contents page onward the typography and
% colour palette match the ACM body conventions of the Writing-format samples.
\IfFileExists{libertine.sty}{%
  \usepackage[mono=false]{libertine}%
  \IfFileExists{zi4.sty}{\usepackage[varqu]{zi4}}{}%
}{%
  \usepackage{lmodern}%
}
% expansion=false reduces pdflatex work; protrusion still improves line breaks
\usepackage[protrusion=true,expansion=false]{microtype}
\usepackage[a4paper,width=155mm,top=25mm,bottom=25mm,headheight=15pt]{geometry}
\usepackage{setspace}
% ACM acmsmall body is single-spaced; we use \onehalfspacing for the BRAC front
% matter and re-tighten to single spacing from the Contents page onward (see the
% \mainmatter-style switch right before \tableofcontents).
\onehalfspacing
\emergencystretch=3em
\hbadness=2000

\usepackage{amsmath,amssymb}
\usepackage{fontawesome5}
% pdfLaTeX cannot typeset Unicode emoji; these mirror the manuscript legend in
% tables. v15: greyscale-only palette per the strict B&W academic-format rule.
% Shape carries the meaning (check / cross / hourglass / circle); shade carries
% the emphasis.
\newcommand{\tmarkDone}{\textbf{\faCheck}}
\newcommand{\tmarkNo}{\textbf{\faTimes}}
\newcommand{\tmarkPartial}{\faHourglassHalf}
\newcommand{\tmarkPlanned}{\textcolor{black!55}{\faCircle}}

\usepackage{graphicx}
\graphicspath{%
  {./}%
  {Tables/}%
  {Diagrams/mermaid-pdf/improved diagrams/}%  v15 redesigned diagrams (preferred)
  {Diagrams/mermaid-pdf/}%                    v13/v14 rendered fallbacks
  {Diagrams/}%
  {Diagrams/CSE370/}{Diagrams/CSE471/}{Diagrams/CSE470/}%
}
% Implied-DPI for bitmaps without pHYs: 300~dpi is enough for A4; 1200~dpi can explode RAM on many figures.
\pdfimageresolution=300
% \IfFileExists does not search \graphicspath; this mirrors the path list so missing files
% typeset a placeholder (avoids hard errors, repeated retries, and broken .aux on Overleaf).
\makeatletter
\newcommand*{\@thesis@doinclude}[2]{%
  \if\relax\detokenize{#1}\relax
    \includegraphics{#2}%
  \else
    \includegraphics[#1]{#2}%
  \fi
}
\newcommand*{\@thesismissingimg}[1]{%
  \fbox{\parbox{0.92\linewidth}{\scriptsize\centering\ttfamily Missing file: \detokenize{#1}.\\[2pt]Copy it into the project (or match \string\graphicspath) and recompile.}}%
}
\newcommand{\ThesisIncludeGraphics}[2][]{%
  \IfFileExists{#2}{%
    \@thesis@doinclude{#1}{#2}}{%
  \IfFileExists{./#2}{%
    \@thesis@doinclude{#1}{#2}}{%
  \IfFileExists{Tables/#2}{%
    \@thesis@doinclude{#1}{#2}}{%
  \IfFileExists{Diagrams/mermaid-pdf/improved diagrams/#2}{%
    \@thesis@doinclude{#1}{Diagrams/mermaid-pdf/improved diagrams/#2}}{%
  \IfFileExists{Diagrams/mermaid-pdf/#2}{%
    \@thesis@doinclude{#1}{Diagrams/mermaid-pdf/#2}}{%
  \IfFileExists{Diagrams/#2}{%
    \@thesis@doinclude{#1}{#2}}{%
  \IfFileExists{Diagrams/CSE370/#2}{%
    \@thesis@doinclude{#1}{#2}}{%
  \IfFileExists{Diagrams/CSE471/#2}{%
    \@thesis@doinclude{#1}{#2}}{%
  \IfFileExists{Diagrams/CSE470/#2}{%
    \@thesis@doinclude{#1}{#2}}{%
    \@thesismissingimg{#2}%
  }}}}}}}}}%
}
\makeatother
% Image-only table floats: caption text is in the PNG; keep counter for \ref{tab:...}
% Usage: \TableScreenshotEnd[label]{Title for List of Tables}
% \phantomsection is REQUIRED before every \addcontentsline when hyperref is loaded;
% without it, hyperref silently drops the .lot entry or links it to the wrong page.
\newcommand{\TableScreenshotEnd}[2][]{%
  \refstepcounter{table}%
  \if\relax\detokenize{#1}\relax\else\label{#1}\fi%
  \phantomsection%
  \addcontentsline{lot}{table}{\protect\numberline{\thetable}#2}%
}
% Only [table] + HTML-defined colors: omit svgnames/x11names (large internal color tables).
\usepackage[table]{xcolor}
\usepackage{tikz}
\usetikzlibrary{arrows.meta,positioning,fit,shapes.geometric}
\usepackage{pgfplots}
\pgfplotsset{compat=1.18, minor tick num=0}

\usepackage{booktabs}
\usepackage{array}
\usepackage{colortbl}
\usepackage{multirow}
\usepackage{float}
\usepackage{caption}
\captionsetup{
    labelfont={bf,small},
    textfont={small},
    skip=6pt,
    format=hang
}

% PNGs are embedded at full file resolution; width/height only set how large the image is drawn (pdflatex).
% Tables: cap box at 70\% of the one-page text area (30\% smaller than a full-area fit), centered in the float.
% Figures/diagrams: maximize drawable area---minimal reserve for [H] float glue + one short \caption (increase if a caption wraps to the next page).
\newlength{\ThesisTableImgReserve}
\newlength{\ThesisFigImgReserve}
\setlength{\ThesisTableImgReserve}{\dimexpr2\intextsep+0.75\baselineskip\relax}
\setlength{\ThesisFigImgReserve}{\dimexpr2\intextsep+\abovecaptionskip+\belowcaptionskip+4\baselineskip\relax}
\newcommand{\TableImageMaxFit}[1]{%
  \ThesisIncludeGraphics[
    width=0.7\linewidth,
    height=\dimexpr(\textheight-\ThesisTableImgReserve)*7/10\relax,
    keepaspectratio
  ]{#1}%
}
\newcommand{\FigureImageMaxFit}[1]{%
  \ThesisIncludeGraphics[width=\linewidth,height=\dimexpr\textheight-\ThesisFigImgReserve\relax,keepaspectratio]{#1}%
}
\newcommand{\OnePageDiagram}[2][]{%
  \if\relax\detokenize{#1}\relax
    \ThesisIncludeGraphics[width=\linewidth,height=\dimexpr\textheight-\ThesisFigImgReserve\relax,keepaspectratio]{#2}%
  \else
    \ThesisIncludeGraphics[#1,width=\linewidth,height=\dimexpr\textheight-\ThesisFigImgReserve\relax,keepaspectratio]{#2}%
  \fi
}
% Compact bar charts (Ch.~5 feasibility figures): half text width, aspect preserved.
\newcommand{\HalfWidthDiagram}[1]{%
  \ThesisIncludeGraphics[width=0.5\linewidth,keepaspectratio]{#1}%
}

% [most] loads many libraries (listings, theorems, breakable, …); [skins] covers titled boxes and shadows.
\usepackage[skins]{tcolorbox}

\usepackage{titlesec}
\usepackage{parskip}
\usepackage{fancyhdr}
\usepackage{enumitem}
% url is loaded by hyperref; this helps long \url{...} lines in references break.
\PassOptionsToPackage{hyphens,spaces,obeyspaces}{url}
\usepackage{hyperref}
\makeatletter
\g@addto@macro\UrlBreaks{\do\/\do\-\do\.\do\:\do\&\do\=\do\#}
\makeatother

% v15 colour palette: greyscale only (strict B&W academic-format rule).
% Existing \color{PrimaryBlue} / \color{AccentBlue} / ... calls throughout the
% document are intentionally NOT rewritten one-by-one; they keep their names
% but resolve to greyscale through these redefinitions. PrimaryBlue=black gives
% all headings, section titles, hyperlinks, and citations a uniform black tone;
% AccentBlue=gray!55 keeps a subtle visual separator for rules and frames; the
% three RowShade tones are kept as light greys for table-row emphasis only.
\definecolor{PrimaryBlue}{HTML}{000000}
\definecolor{AccentBlue}{gray}{0.55}
\definecolor{LightBlue}{gray}{0.96}
\definecolor{MidBlue}{gray}{0.90}
\definecolor{DarkText}{HTML}{000000}
\definecolor{GrayText}{gray}{0.45}
\definecolor{RowShade}{gray}{0.96}
\definecolor{RowShadeAlt}{gray}{0.90}
\definecolor{GreenAccent}{HTML}{000000}
\definecolor{AmberAccent}{HTML}{000000}
\definecolor{RedAccent}{HTML}{000000}
\definecolor{VioletAccent}{HTML}{000000}

\hypersetup{
    colorlinks=true,
    linkcolor=black,
    filecolor=black,
    urlcolor=black,
    citecolor=black,
    pdfborder={0 0 0},
}

% ACM body convention: section heads in black, no decorative rules.
% The chapter rule below is kept as a very thin neutral grey separator only.
\titleformat{\chapter}[display]
    {\normalfont\LARGE\bfseries}
    {\chaptertitlename\ \thechapter}{20pt}{\LARGE}
    [\vspace{4pt}{\color{black!40}\titlerule[0.6pt]}]

\titleformat{\section}
    {\large\bfseries}
    {\thesection}{0.8em}{}

\titleformat{\subsection}
    {\normalsize\bfseries}
    {\thesubsection}{0.6em}{}

\titleformat{\subsubsection}
    {\normalsize\bfseries\itshape}
    {\thesubsubsection}{0.6em}{}

\titlespacing*{\chapter}{0pt}{-20pt}{20pt}
\titlespacing*{\section}{0pt}{16pt plus 3pt minus 2pt}{8pt plus 2pt minus 1pt}
\titlespacing*{\subsection}{0pt}{12pt plus 2pt minus 1pt}{6pt plus 1pt minus 1pt}
\titlespacing*{\subsubsection}{0pt}{10pt plus 2pt minus 1pt}{4pt plus 1pt minus 1pt}

% v15 table heading style: light grey row + bold black text (ACM booktabs spirit).
\newcommand{\tableheadcolor}{\rowcolor{black!8}}
\newcommand{\tableheadfont}{\bfseries\small}

\newenvironment{moderntable}[1][H]{%
    \begin{table}[#1]\centering\small%
    \renewcommand{\arraystretch}{1.35}%
}{%
    \end{table}%
}

% v15 colour-box redesign: greyscale only, no fills.
\newtcolorbox{layerbox}[1]{
    colback=white,
    colframe=black!55,
    coltitle=black,
    fonttitle=\bfseries\small,
    title=#1,
    boxrule=0.6pt,
    arc=2pt,
    left=6pt,right=6pt,top=4pt,bottom=4pt,
}

\newtcolorbox{highlightbox}[1][]{
    colback=black!4,
    colframe=black!45,
    boxrule=0.5pt,
    arc=2pt,
    left=8pt,right=8pt,top=6pt,bottom=6pt,
    #1
}

\pagestyle{fancy}
\fancyhf{}
\renewcommand{\headrulewidth}{0.4pt}
\fancyhead[L]{\small\color{black!55}Crypto World Bank}
\fancyhead[R]{\small\color{black!55}BRAC University}
\fancyfoot[C]{\small\color{black!55}\thepage}

% ── Thesis diagram helper macros ─────────────────────────────────────────────
% manualpie: portable pie slice (no pgf-pie package required)
% Usage: \manualpie{cx}{cy}{r}{start_deg}{end_deg}{fill}
\newcommand{\manualpie}[6]{%
  \fill[#6] (#1,#2) -- ++({#4}:{#3}) arc ({#4}:{#5}:{#3}) -- cycle;
  \draw[white, line width=0.8pt] (#1,#2) -- ++({#4}:{#3}) arc ({#4}:{#5}:{#3}) -- cycle;
}
% v15 decision-tree nodes: greyscale only. Selected (ddone) is filled darker
% grey + bold to remain visually distinguishable without colour.
\newcommand{\ddcat}[3]{\node[draw=black!55,fill=black!5,rounded corners=2pt,
  minimum width=2.2cm,minimum height=0.65cm,align=center,font=\scriptsize]
  (#1) at (#2) {#3};}
\newcommand{\ddone}[3]{\node[draw=black,fill=black!22,
  minimum width=2.2cm,minimum height=0.65cm,align=center,
  font=\scriptsize\bfseries,rounded corners=2pt] (#1) at #2 {#3};}
\newcommand{\ddtwo}[3]{\node[draw=black!55,fill=white,
  minimum width=2.2cm,minimum height=0.65cm,align=center,
  font=\scriptsize,rounded corners=2pt] (#1) at #2 {#3};}
% ─────────────────────────────────────────────────────────────────────────────

\begin{document}

%% ============================================================
%% TITLE PAGE
%% ============================================================
\thispagestyle{empty}

\begin{center}

\vspace*{0.3cm}

\ThesisIncludeGraphics[width=3.5cm]{bracu_logo_12-0-2022.png}

\vspace{0.8cm}

{\LARGE\bfseries\color{PrimaryBlue} Decentralized Crypto World Bank}\\[0.4em]
{\large\color{PrimaryBlue!80} A Blockchain-Based Banking Platform\\[0.15em]with AI-Enhanced Security and Assistance}

\vspace{1.2cm}

{\large by}

\vspace{0.5cm}

{\Large\bfseries\color{PrimaryBlue} Md.\ Bokhtiar Rahman Juboraz}\\[0.2em]
{\color{GrayText}20301138}\\[0.8em]
{\Large\bfseries\color{PrimaryBlue} Md.\ Mahir Ahnaf Ahmed}\\[0.2em]
{\color{GrayText}20301083}

\vspace{1.2cm}

A pre-thesis 1 report submitted to the Department of Computer Science and Engineering\\
in partial fulfillment of the requirements for the degree of\\
\textbf{B.Sc.\ in Computer Science}

\vfill

{\color{PrimaryBlue}\rule{0.4\textwidth}{0.6pt}}\\[0.6em]
Department of Computer Science and Engineering\\
\textbf{BRAC University}\\
February 2026

\vspace{0.5cm}

{\small\color{GrayText}\copyright\ 2026.\ BRAC University --- All rights reserved.}

\end{center}

%% ============================================================
%% DECLARATION
%% ============================================================
\newpage
\thispagestyle{empty}

\chapter*{Declaration}
\phantomsection
\addcontentsline{toc}{chapter}{Declaration}

It is hereby declared that

\begin{enumerate}
\item The report submitted is our own original work while completing degree at BRAC University.

\item The report does not contain material previously published or written by a third party, except where this is appropriately cited through full and accurate referencing.

\item The report does not contain material which has been accepted, or submitted, for any other degree or diploma at a university or other institution.

\item We have acknowledged all main sources of help.
\end{enumerate}

\vspace{2cm}

\textbf{Student's Full Name \& Signature:}

\vspace{2cm}

\begin{tabular}{p{6cm}p{1cm}p{6cm}}
\hrulefill & & \hrulefill \\
\centering Md. Bokhtiar Rahman Juboraz & & \centering Md. Mahir Ahnaf Ahmed \tabularnewline
\centering 20301138 & & \centering 20301083 \tabularnewline
\end{tabular}

%% ============================================================
%% APPROVAL
%% ============================================================
\newpage
\thispagestyle{empty}

\chapter*{Approval}
\phantomsection
\addcontentsline{toc}{chapter}{Approval}

The pre-thesis 1 report of final year project titled ``Decentralized Crypto World Bank: A Blockchain-Based Banking Platform with AI-Enhanced Security and Assistance'' submitted by

\begin{enumerate}
\item Md. Bokhtiar Rahman Juboraz (20301138)
\item Md. Mahir Ahnaf Ahmed (20301083)
\end{enumerate}

Of Spring, 2026 has been accepted as satisfactory in partial fulfillment of the requirement for the degree of B.Sc. in Computer Science on February 20, 2026.

\vspace{1cm}

\textbf{Examining Committee:}

\vspace{1.5cm}

\begin{flushleft}
Supervisor:\\
(Member)
\end{flushleft}

\begin{flushright}
\vspace{-0.5cm}
\hrulefill\\
Mr. Annajiat Alim Rasel\\
Senior Lecturer\\
Department of Computer Science and Engineering\\
BRAC University
\end{flushright}

\vspace{1.5cm}

\begin{flushleft}
Project Coordinator:\\
(Member)
\end{flushleft}

\begin{flushright}
\vspace{-0.5cm}
\hrulefill\\
Dr. Md. Golam Rabiul Alam\\
Professor\\
Department of Computer Science and Engineering\\
BRAC University
\end{flushright}

\newpage

\vspace{1.5cm}

\begin{flushleft}
Head of Department:\\
(Chair)
\end{flushleft}

\begin{flushright}
\vspace{-0.5cm}
\hrulefill\\
Dr. Sadia Hamid Kazi\\
Chairperson and Associate Professor\\
Department of Computer Science and Engineering\\
BRAC University
\end{flushright}

%% ============================================================
%% ETHICS STATEMENT
%% ============================================================
\newpage
\thispagestyle{empty}

\chapter*{Ethics Statement}
\phantomsection
\addcontentsline{toc}{chapter}{Ethics Statement}

This project operates exclusively on public blockchain test networks using test tokens with no real monetary value. No real financial transactions, personal banking data, or user financial records are involved at any stage of development or evaluation. All transaction data used for Artificial Intelligence and Machine Learning (AI/ML) model training and evaluation is either synthetically generated or sourced from publicly available anonymized datasets. The platform prototype does not process Know Your Customer (KYC) or Anti-Money Laundering (AML) data in its current scope, and all wallet addresses used during testing are disposable testnet addresses. We acknowledge the ethical considerations inherent in AI-assisted financial decision-making and have incorporated explainability mechanisms (SHapley Additive exPlanations, or SHAP) to ensure that automated risk assessments remain transparent and auditable by human reviewers. Future phases plan to incorporate a privacy-preserving ZKP-based identity compliance layer in which users prove KYC status without exposing personal data on-chain; the ethical implications of this design will be evaluated in the final thesis.

%% ============================================================
%% ABSTRACT
%% ============================================================
\newpage
\thispagestyle{empty}

\chapter*{Abstract}
\phantomsection
\addcontentsline{toc}{chapter}{Abstract}

Global development finance relies on multilayered institutional structures to distribute capital across borders and communities, yet these systems often struggle with fragmented transparency, procedural friction, and uneven access to credit. Although decentralized finance has demonstrated that programmable infrastructures can automate lending and enhance auditability, existing models largely employ flat architectures that do not reflect institutional hierarchies or structured governance. As a result, there remains limited exploration of how tiered financial systems might be represented within decentralized environments while preserving oversight and adaptability. Here we present \textit{Crypto World Bank}, a formally specified, partially implemented research prototype that demonstrates the \textit{feasibility} of a blockchain-based cryptocurrency exchange and institutional finance platform coordinating capital allocation across a four-tier hierarchical governance structure. We show that hierarchical capital flows, role-based governance, and data-informed risk analytics can be coordinated within a unified smart contract environment, enabling transparent state visibility alongside adaptive decision support. The Crypto World Bank formally specifies and partially implements a banking architecture spanning six functional domains---deposit mobilization, credit allocation, payment settlement, risk intermediation, liquidity management, and ancillary financial services---within a four-tier hierarchical governance structure. The prototype implements the core Tier~1 reserve and capital-allocation contracts, with Tiers~2--4 and the extended product suite formally specified for implementation in the final thesis phase. This hybrid design illustrates how institutional finance and decentralized systems may converge, contributing an open-source, hierarchically governed crypto financial services architecture as a foundation for further research.

\vspace{0.5em}
\noindent\textbf{Keywords:} Blockchain, Decentralized Finance, Institutional Architecture, Financial Inclusion, Smart Contracts, AI-Augmented Governance, Group Lending, Deposit Mobilization, Reserve Management, Financial Sustainability

%% ============================================================
%% DEDICATION
%% ============================================================
\newpage
\thispagestyle{empty}

\chapter*{Dedication}
\phantomsection
\addcontentsline{toc}{chapter}{Dedication}

\vspace{3cm}
\begin{center}
\textit{Dedicated to our families for their unwavering support throughout our academic journey.}
\end{center}

%% ============================================================
%% ACKNOWLEDGMENT
%% ============================================================
\newpage
\thispagestyle{empty}

\chapter*{Acknowledgment}
\phantomsection
\addcontentsline{toc}{chapter}{Acknowledgment}

We would like to express our sincere gratitude to our panel members for their guidance and support throughout this project. We also thank our supervisor, Mr.\ Annajiat Alim Rasel, Senior Lecturer at the Department of Computer Science and Engineering, BRAC University, for his guidance and support. Finally, we thank the Department of Computer Science and Engineering at BRAC University for providing us with the resources and academic environment to pursue this work.

%% ============================================================
%% TABLE OF CONTENTS, LIST OF TABLES, LIST OF FIGURES
%% ============================================================
% v15: from the Contents page onward, switch to ACM-style single spacing and
% strictly black-and-white textual content. The BRAC front matter above keeps
% its institutional layout (with \onehalfspacing) per the supervisor's request;
% only the colours are neutralized there.
\newpage
\singlespacing
\setcounter{page}{1}
\pagenumbering{roman}
\tableofcontents
\newpage
\phantomsection
\addcontentsline{toc}{chapter}{List of Tables}
\listoftables
\newpage
\phantomsection
\addcontentsline{toc}{chapter}{List of Figures}
\listoffigures
\newpage
\chapter*{List of Formulas}
\phantomsection
\addcontentsline{toc}{chapter}{List of Formulas}
\begin{description}[leftmargin=0pt]
    \item[Utilization Rate (U):] \(U = \frac{L}{L + B}\) (lending utilization; \(L\) is lent amount, \(B\) is available balance/liquidity)
    \item[Collateral Ratio (CR):] \(CR = \frac{C}{D}\) (collateral-to-debt ratio)
    \item[Loan-to-Value (LTV):] \(LTV = \frac{\text{Loan Amount}}{\text{Collateral Value}}\)
    \item[APR (simple):] \(APR = r_{\text{period}} \times N\) (annualized simple rate over \(N\) periods)
    \item[Compound Interest (Savings):] \(A = P\left(1 + \frac{r}{n}\right)^{nt}\) (deposit growth; \(P\) principal, \(r\) annual rate, \(n\) compounding frequency)
    \item[EMI (Installment Payment):] \(EMI = \frac{P \cdot r \cdot (1+r)^n}{(1+r)^n - 1}\) (installment for principal \(P\), periodic rate \(r\), \(n\) installments)
    \item[Health Factor (HF):] \(HF = \frac{C \times LT}{D}\) (liquidation trigger for \textit{over-collateralized retail loans only}; if \(HF < 1.0\) liquidation is triggered, where \(LT\) is liquidation threshold; see Section~\ref{sec:banking-products} for tier-specific HF variants and the group-loan pool-level formulation)
    \item[Reserve Ratio (RR):] \(RR = \frac{\text{Reserves}}{\text{Total Deposits}}\) (solvency constraint enforced at each tier; governs the maximum downward capital that may be allocated, \textit{not} a money-creation parameter)
    \item[Platform Solvency Ratio (PSR):] \(PSR = \frac{\text{TotalReserveBalance}}{\text{TotalOutstandingLoans}}\) (on-chain analogue of capital adequacy, measuring whether the system's reserve cover exceeds outstanding loan obligations; replaces the Basel~III CAR formula, which is undefined for this platform because USDC assets have no regulatory risk-weight classification under Basel~III and ``Tier~1 Capital'' in the Basel sense cannot be computed from an on-chain USDC balance alone)
    \item[Credit Velocity (CV):] \(CV = \frac{\text{Capital Recycled per Period}}{\text{Total Reserve}}\) (rate at which reserve capital flows through the lending hierarchy and returns as repayment; replaces the inapplicable fractional-reserve money multiplier, which presupposes new money creation and does not apply to fixed-supply USDC)
    \item[FX Conversion:] \(\text{Amount}_B = \text{Amount}_A \times \left(\frac{P_A}{P_B}\right)\) (oracle-priced conversion between assets \(A\) and \(B\))
    \item[Group Loan Individual Share:] \(\text{Share}_i = \frac{\text{LoanTotal}}{N_{\text{members}}}\)
    \item[On-Chain Credit Score (conceptual):] \(CS = \sum_i w_i \cdot \text{feature}_i\) (weighted features learned off-chain and stored as a score on-chain)
    \item[Isolation Forest Anomaly Score:] \(s(x, n) = 2^{-\frac{E(h(x))}{c(n)}}\) (anomaly scoring for instance \(x\); where \(h(x)\) is the path length of instance \(x\) from the root of an isolation tree and \(c(n)\) is the average path length for \(n\) samples)
    \item[Random Forest Fraud Probability:] \(P(\text{fraud}\mid x) = \frac{1}{T}\sum_{t=1}^{T} f_t(x)\) (ensemble average over \(T\) trees; where \(f_t(x)\) is the binary fraud prediction of the \(t\)-th decision tree)
    \item[Platform Net Interest Allocation:] \(\text{NetInterest} = \text{InterestCollected} - \text{DepositorYield} - \text{InsuranceAllocation}\)
    \item[SHAP (Shapley value):] \(\phi_i = \sum_{S \subseteq F \setminus \{i\}} \frac{|S|!(|F|-|S|-1)!}{|F|!}\left[f(S \cup \{i\}) - f(S)\right]\)
\end{description}
\newpage
\chapter*{List of Abbreviations}
\phantomsection
\addcontentsline{toc}{chapter}{List of Abbreviations}
\begin{description}[leftmargin=0pt]
    \item[AA:] Account Abstraction (ERC-4337)
    \item[ALM:] Asset-Liability Management
    \item[AI:] Artificial Intelligence
    \item[AI/ML:] Artificial Intelligence / Machine Learning
    \item[AML:] Anti-Money Laundering
    \item[APR:] Annual Percentage Rate
    \item[API:] Application Programming Interface
    \item[BRAC:] Bangladesh Rural Advancement Committee
    \item[CAR:] Capital Adequacy Ratio (Basel~III; not directly applicable to this platform --- see Platform Solvency Ratio, PSR, in List of Formulas)
    \item[CBDC:] Central Bank Digital Currency
    \item[CCIP:] Cross-Chain Interoperability Protocol (Chainlink)
    \item[CCS:] ACM Computing Classification System
    \item[CR:] Collateral Ratio
    \item[CVL:] Certora Verification Language
    \item[DeFi:] Decentralized Finance
    \item[DID:] Decentralized Identifier (W3C)
    \item[DON:] Decentralised Oracle Network (Chainlink)
    \item[EIP:] Ethereum Improvement Proposal
    \item[EIP-7702:] Ethereum Improvement Proposal 7702 --- Session Key standard enabling scoped, time-bound wallet authorisations for an autonomous agent
    \item[EMI:] Equated Monthly Installment
    \item[ERC:] Ethereum Request for Comments (token / interface standard)
    \item[EVM:] Ethereum Virtual Machine
    \item[FL:] Federated Learning
    \item[FATF:] Financial Action Task Force
    \item[FX:] Foreign Exchange
    \item[GNN:] Graph Neural Network
    \item[HF:] Health Factor
    \item[KYC:] Know Your Customer
    \item[LLM:] Large Language Model
    \item[L2:] Layer 2
    \item[LTV:] Loan-to-Value Ratio
    \item[MFI:] Microfinance Institution
    \item[MiCA:] Markets in Crypto-Assets Regulation (EU)
    \item[ML:] Machine Learning
    \item[MCP:] Model Context Protocol --- tool-calling interface used by the AI agent to interact with CWB banking operations
    \item[NIM:] Net Interest Margin
    \item[PRISMA:] Preferred Reporting Items for Systematic Reviews and Meta-Analyses
    \item[PoS:] Proof of Stake
    \item[PoR:] Proof of Reserve (Chainlink on-chain reserve verification)
    \item[PSR:] Platform Solvency Ratio (on-chain capital adequacy analogue; see List of Formulas)
    \item[QRC:] QR code
    \item[RBAC:] Role-Based Access Control
    \item[RLS:] Row-Level Security (PostgreSQL policy feature enforcing append-only access on audit tables)
    \item[RR:] Reserve Ratio
    \item[SAR:] Suspicious Activity Report
    \item[SBT:] Soulbound Token (non-transferable ERC standard)
    \item[SHAP:] SHapley Additive exPlanations
    \item[SOFR:] Secured Overnight Financing Rate
    \item[SSE:] Server-Sent Events
    \item[SSI:] Self-Sovereign Identity
    \item[SoK:] Systematization of Knowledge
    \item[SWC:] Smart Contract Weakness Classification
    \item[TVL:] Total Value Locked
    \item[U:] Utilization
    \item[UUPS:] Universal Upgradeable Proxy Standard (ERC-1822)
    \item[VC:] Verifiable Credential (W3C)
    \item[XAI:] Explainable Artificial Intelligence
    \item[ZKP:] Zero-Knowledge Proof
    \item[zkAML:] Zero-Knowledge Anti-Money-Laundering proof
    \item[zkEVM:] Zero-Knowledge Ethereum Virtual Machine --- a ZK rollup that inherits Ethereum L1 security via cryptographic validity proofs
    \item[zk-SNARK:] Zero-Knowledge Succinct Non-Interactive Argument of Knowledge
\end{description}
\newpage
\pagenumbering{arabic}
\setcounter{page}{1}


%% %%%%%%%%%%%%%%%%%%%%%%%%%%%%%%%%%%%%%%%%%%%%%%%%%%%%%%%%%%%%
%% CHAPTER 1: INTRODUCTION
%% %%%%%%%%%%%%%%%%%%%%%%%%%%%%%%%%%%%%%%%%%%%%%%%%%%%%%%%%%%%%
\chapter{Introduction}

The Crypto World Bank is not a lending protocol. It is a \textbf{formally specified, partially implemented} research prototype demonstrating the feasibility of a blockchain-based cryptocurrency exchange and institutional finance platform on programmable blockchain infrastructure. Where existing decentralized finance protocols address isolated financial functions---Aave handles collateralized lending, Uniswap handles exchange, MakerDAO issues stablecoins---the Crypto World Bank formally specifies and partially implements a hierarchically governed crypto financial services platform spanning the full spectrum of banking functions. These functions include deposit mobilization, credit allocation, installment-based repayment, interbank liquidity management, foreign exchange facilitation, group lending, reserve enforcement, and risk-based governance---coordinated across four institutional tiers that mirror the structural logic of real-world development finance. The Tier~1 World Bank Reserve and the core lending hierarchy are implemented and deployed on testnet; the extended product suite (SavingsVault, GroupLendingPool, IBLP, and NettingEngine) is formally specified and scheduled for the final thesis implementation phase.

The system is designed to serve three categories of participants simultaneously. At the global level, it functions as a programmable central reserve institution---managing capital allocation across national jurisdictions with transparent, auditable reserve ratios enforced by smart contract code rather than periodic audit. At the national and regional level, it functions as a commercial and development bank---channeling capital to local lending institutions, managing interbank liquidity, and providing institutional clients with structured credit facilities. At the individual level, it functions as a retail bank---offering savings products, group loans, installment-based personal lending, and digital payment services to end customers, including the estimated 1.4~billion adults currently excluded from formal financial systems~[14]. The complete source code, smart contracts, and supporting material for this project are available in our \href{https://github.com/Yuvraajrahman/Cryto-World-Bank}{GitHub repository}.

Blockchain functions here as a coordination technology, providing visibility of shared states, audit trails that are difficult to alter, and programmable enforcement of rules among institutions that may have competing interests. Governance design, security controls, and regulatory considerations are incorporated into the platform according to a staged implementation roadmap that includes academic validation and regulatory sandbox pilots. Analogous hybrid institutional structures in blockchain-based finance include the World Bank's FundsChain programme, which tracks fund disbursement across projects on a permissioned ledger, and JPMorgan Kinexys, which settles multi-billion-dollar institutional payments daily on a private EVM-compatible chain---demonstrating that open ledger infrastructure and institutional hierarchy are not mutually exclusive.

\section{Background}

Understanding why such a system is necessary requires examining the structural failures of the institutions it is designed to complement.

Capital is distributed to local borrowers through layered institutional arrangements in which supranational development institutions such as the World Bank and IMF, national financial intermediaries, and local lenders collaborate to channel development finance to end borrowers. Although this model allows risk sharing and penetrating local markets, it is linked to great operational issues. In cross-border correspondent banking, financial institutions are obligated to have pre-funded nostro and vostro accounts at correspondent banks in every currency corridor, which maroon huge amounts of capital in idle balances unable to be invested or lent out~[24]. Crossborder transactions are regularly settled in two to five business days with an average cost of about \$42 per transaction by using the correspondent banking network~[25]. The remittance market, which is valued at about \$860~billion annually, loses between \$48--56~billion per annum to transfer charges, with the average cost to transfer in the world being 6.49\% as of 2024, more than the United Nations Sustainable Development Goal limit of 3\%~[26].

At the same time, it is estimated that there are 1.4~billion unbanked adults in the world as defined by the World Bank Global Findex~[14], and most of them are found in the developing economies where documentation, geographic location, and minimum balance requirements leave out high populations of people in the formal financial sector. The International Finance Corporation estimates that there is a financing gap in the world of micro, small and medium enterprises at \$4.5~trillion per year~[20], which is unmet credit demand that cannot be effectively met by the existing institutions through the conventional intermediation chains.

The DeFi industry has already shown that lending, collateral management, and interest calculation can be automated on a smart contract platform and that all transactions are transparent, on-chain. The largest DeFi lending protocol, Aave~v3, has amassed \$26.3~billion in total value locked across ten blockchain networks~[27] (accessed March~2026), and the DeFi lending market has more than \$55~billion in TVL~[13] (accessed March~2026). Nevertheless, these protocols utilize flat, pool-based designs where all the lenders feed into one pool and all borrowers withdraw off the same source---irrespective of institutional status, creditworthiness, and geography. There is no current DeFi scheme that reflects the multi-tier institutional structure, inter-tier capitals movement and tiered borrower access provision that is typical of real-world development finance.

\paragraph{Autonomous banking agents and the human-gate model.}
A further dimension of this gap concerns the interface through which retail clients in underserved economies interact with banking services. The Crypto World Bank extends beyond the role of a read-only LLM product guide to incorporate a locally hosted, privacy-preserving autonomous AI agent capable of both answering client questions and executing on-chain banking operations on their behalf. The agent is built on a Model Context Protocol (MCP) tool server exposing 17 banking tools (9 read tools, 8 write tools) and operates under a mandatory human confirmation gate: every state-modifying operation---loan application, installment payment, deposit, or KYC upgrade---requires explicit affirmative consent from the client before the agent calls the corresponding write tool and signs the resulting on-chain transaction via an EIP-7702 scoped session key. This architecture provides the conversational accessibility of a virtual banker while preserving the client's full sovereignty over every financial action. The agent is an optional convenience layer that sits alongside the standard form-based React UI: every banking operation it can perform is equally available through the web interface, and clients who prefer direct interaction engage the platform without the agent. The human confirmation gate is enforced by a confirmation audit hook in the Express.js middleware layer---application-layer policy enforcement that operates independently of model output: any write-tool POST request without a logged confirmation turn in the conversation history is rejected with HTTP~403, regardless of what the model generated.

\section{Rationale of the Study}

The combination of three forces, namely: (1)~the long-standing inefficiencies in the traditional development finance, which restrict the transparency and speed of settlement; (2)~the absence of multi-tier banking functions in existing DeFi designs, particularly deposit mobilization and solidarity group lending alongside hierarchical lending; (3)~the possibility of integrating blockchain auditability with lightweight analytics and monitoring support to enhance operational control, drives this study. With an architecture that will bridge these areas, we hope to show a plausible way forward toward transparent, programmable, and institutionally significant banking services for underserved populations.

\section{Problem Statement}

The international system of the development finance works based on a stratified system in which funds are directed down by large supranational organizations such as the World Bank and IMF to national institutions, which transmit them to regional and local banks, and they are ultimately delivered to individual borrowers. This intermediary chain brings into being a number of significant inefficiencies:

\begin{enumerate}[noitemsep]
\item \textbf{Lack of transparency.} Financial records at each tier of the intermediation chain are not openly accessible to participants below that tier. Lower-level institutions and individual borrowers cannot independently verify how capital is being managed, what reserve levels are maintained, or on what basis lending decisions are made at higher tiers. Reserve adequacy is predominantly self-reported and audited at most quarterly, providing no mechanism for real-time verification of institutional financial health.

\item \textbf{Sluggish settlements and unproductive capital.} Cross-border money transmission requires multiple intermediaries, compliance verification steps, and correspondent bank confirmations that routinely delay settlement by two to five business days. Banks must also maintain pre-funded accounts at partner institutions in every currency corridor, locking up capital that would otherwise be available for lending~[24]. An average transaction via the correspondent banking system costs approximately \$42, a burden that falls disproportionately on small-value transactions and participants in developing economies~[25].

\item \textbf{Inconsistent risk evaluation.} Fraud screening and credit assessment rely heavily on individual human judgment, introducing bias, inconsistency, and limited scalability. Because no unified and verifiable credit history exists across institutional tiers, borrowers must be re-evaluated from scratch at each stage of the intermediation chain, adding time and cost without improving accuracy.

\item \textbf{Obstacles to access and institutional trust.} Establishing inter-institutional trust through legal agreements, compliance processes, and third-party audits is both expensive and time-consuming, placing smaller institutions and borrowers in developing countries at a particular disadvantage. According to the World Inequality Report~2026, the global monetary system continues to redirect resources from developing to wealthier economies, with developing-country investors receiving returns approximately three percentage points lower than those in developed economies~[28]---a structural gap that further impedes access to international credit markets.

\item \textbf{Absence of programmable savings instruments.} Traditional savings accounts in developing economies offer near-zero real yields due to inflation, currency risk, and institutional opacity. Depositors have no real-time visibility into whether their savings are being deployed responsibly or sitting idle in reserve accounts.

\item \textbf{Inaccessibility of group credit mechanisms.} Over 1.4 billion unbanked individuals lack individual collateral but could access credit through solidarity group structures, a model demonstrated by BRAC, Grameen Bank, and ASA in Bangladesh. No existing DeFi protocol models programmable mutual liability or on-chain group repayment enforcement.
\end{enumerate}

At the same time, the DeFi ecosystem has shown that the lending, collateral management, and interest calculation can be completely automated through smart contract platforms with complete transparency. The TVL of Aave~v3 is \$26.3~billion and active borrows are \$17.7~billion~[27]; the TVL of Compound~v3 is \$1.4~billion~[29]; and the combined DAI and USDS stablecoin issuance of MakerDAO (renamed Sky Protocol) is approximately \$10.5~billion with projected 2026 revenue of \$611~million~[30]. Despite this scale, these protocols exhibit fundamental limitations in the context of institutional development finance:

\begin{itemize}[noitemsep]
\item They employ \textbf{flat, peer-to-peer}, or \textbf{over-collateralized} lending models. While Aave~v3 implements risk-parameterized sub-markets (Isolation Mode, Efficiency Mode) and Compound~v3 deploys separate Comet markets with independent risk parameters, none of these protocols model institutional hierarchy, cross-tier capital allocation, or differentiated governance access by institutional role. The CWB's four-tier structure is distinct from Aave's sub-markets: it models the \textit{institutional relationships} between a reserve authority, national lenders, local lenders, and retail clients---not merely risk-differentiated asset buckets within a single lending pool.
\item They are not connected with \textbf{AI-based risk analytics} and \textbf{explainable decision support} systems. DeFi lending decisions are determined solely using collateral ratio and utilization curve and no behavioral fraud detection, anomaly detection or explainable credit judgement is present.
\item They do not model \textbf{role based, tiered governance} of the development finance institutions. It is a form of leadership, in which a majority of the large token holders will wield the power of decision-making, without the institutional role distinction (regulator, wholesale lender, retail lender, retail client) of hierarchical finance.
\item The credit-based and emerging-market lending has also been introduced with other new institutional DeFi projects such as Maple Finance (\$2.6--3.8~billion TVL, undercollateralized institutional lending)~[31] and Goldfinch (\$680~million loan originations in 18+ developing countries)~[32], however, are single-level projects with no hierarchical capital flows and no interbank lending facilities.
\end{itemize}

\section{Objectives}

The objectives of this project are:

\begin{enumerate}[noitemsep]
\item To \textbf{design, formally specify, and partially implement} a four-tier lending architecture (World Bank $\to$ National Bank $\to$ Local Bank $\to$ Client) on an EVM-compatible blockchain that maintains institutional hierarchy and provides shared access to the ledger. The current prototype fully implements the Tier~1 World Bank Reserve contract and the lending request/approval workflow; Tier~2 and Tier~3 contracts are specified and partially scaffolded, with full implementation planned for the final thesis phase.
\item To specify and justify an extended banking product suite---including deposit mobilization, savings products, and solidarity group lending---that can be integrated on top of the hierarchical lending foundation.
\item To investigate a lightweight off-chain analytics support layer, including fraud detection, anomaly detection, and explainable review support, to supplement human credit governance; and to specify and partially implement an optional autonomous AI agent system using a Model Context Protocol (MCP) tool server with 17 banking tools (9 read tools and 8 write tools), enabling clients to execute banking operations through a conversational interface as an alternative path to the standard React UI; to specify a three-tier system prompt architecture (Stable, Context, and Volatile tiers) with a session lineage model for cross-session memory continuity for returning clients; and to implement a confirmation audit hook at the Express.js API middleware layer that enforces the human-gate invariant independently of model output prior to any state-modifying action.
\item To demonstrate transparent and programmable lending processes with configurable borrowing limits, installment payments, and role-based access control implemented through smart contracts.
\item To test the prototype on public testnets (e.g., Polygon, Ethereum Sepolia) and verify hierarchical controls, transparency properties, and operational workflows.
\item To document the design of governance mechanisms, security controls, regulatory considerations, and a controlled rollout pathway toward academic validation and potential regulatory sandbox piloting.
\end{enumerate}

\section{Research Questions}

The research questions used in the study are as follows:

\begin{enumerate}[noitemsep]
\item \textbf{RQ1:} Does a hierarchical blockchain-based lending architecture reflect more faithfully the real-world capital flow mechanism of development finance than present decentralized lending?
\item \textbf{RQ2:} Can on-chain transparency, programmable reserve controls and role-based governance system make settlements less opaque and less frictional across institutional tiers?
\item \textbf{RQ3:} Does a lightweight off-chain analytics layer offer practical and auditable support to monitor fraud and lending in such a system?
\item \textbf{RQ4:} Does the proposed architecture operate correctly on current public test networks, and do the deployed contracts exhibit the role-based access control, capital-flow, and reserve-invariant properties required for real-world institutional banking use cases? \textit{Scope note: RQ4 evaluates prototype-level technical viability on testnets only; it does not constitute a claim of production readiness, regulatory compliance, or real-asset deployment.}
\item \textbf{RQ5:} Can deposit-funded lending and solidarity group lending be represented as programmable, auditable mechanisms on-chain while preserving practical feasibility for retail users in developing economies?
\end{enumerate}

\section{Research Contribution}
\label{sec:research-contribution}

This work makes four original and precisely scoped research contributions to the field of decentralized finance and blockchain-based development banking:

\begin{description}[leftmargin=0pt]

\item[\textbf{Contribution 1 --- Four-Tier Hierarchical DeFi Architecture.}]
To our knowledge, no prior published work presents a four-tier smart-contract lending hierarchy (World Bank Reserve $\to$ National Bank $\to$ Local Bank $\to$ Client) on an EVM-compatible blockchain that mirrors the capital-flow model of multilateral development finance. We extend prior work on flat DeFi lending~[1] to incorporate institutional hierarchy, addressing the gap identified by Werner et al.~[1]. No surveyed DeFi protocol---including Aave, Compound, MakerDAO, Maple Finance, or Goldfinch---models cross-tier capital allocation with reserve-ratio enforcement at every tier; all employ flat, single-tier pool architectures with undifferentiated liquidity.

\textit{Implementation scope:} The World Bank Reserve (Tier~1) contract and the lending request/approval workflow (Tier~3) are deployed and tested on Polygon~Amoy testnet (contract addresses in Appendix~C). The National Bank (Tier~2) and Local Bank (Tier~3) contracts are formally specified in Solidity interface form (Appendix~B) and partially scaffolded; full cross-tier fund-transfer automation is implemented in the Hardhat simulation environment and constitutes Sprint~2 of the final thesis phase. The architectural contribution---the four-tier design pattern and its formal specification---holds independently of the current prototype completeness, as design contributions in distributed systems are established by formal specification and partial implementation consistent with the stated scope.

\item[\textbf{Contribution 2 --- On-Chain Solidarity Group Lending.}]
To our knowledge, no prior published academic work specifies and prototypes a programmable solidarity group lending mechanism on-chain that encodes mutual liability enforcement, group formation consent, and installment splitting into smart contract logic. We note that Goldfinch's Backers/Senior Pool model provides partial structural precedent, but does not implement per-member mutual liability or the over-indebtedness controls detailed in Section~\ref{sec:over-indebt}. This model was previously realised only in analogue microfinance institution (MFI) operations, such as BRAC's 30--40 member groups and Grameen Bank's groups of five, where enforcement relies on social pressure rather than programmable contract rules.

\textit{Implementation scope:} Contribution~2 is a specification contribution: no prior academic work formally specifies the algorithmic structure of solidarity group lending as a smart contract interface. The formal specification is given in Section~\ref{sec:over-indebt}. A minimal proof-of-concept deployment demonstrating group formation, consent recording, and per-member disbursement is included in the testnet repository; full lifecycle implementation with mutual liability enforcement and over-indebtedness control is Sprint~2 of the final thesis phase.

\item[\textbf{Contribution 3 --- Oracle-Mediated AI/ML Integration with Explainability.}]
We propose an architectural pattern for integrating off-chain Random Forest fraud detection and SHAP-based explanations into an on-chain lending decision workflow via a trusted oracle relay, providing a blueprint for auditable, AI-assisted credit governance in DeFi. The design addresses the oracle problem---the challenge of securely bridging on-chain logic with off-chain data---and specifies three concrete integration options (centralized relay, Chainlink Functions, and commit-reveal scheme) with explicit trust trade-offs.

\textit{Implementation scope:} The Random Forest fraud detector is trained and benchmarked on the BCCC-DeFiFraudTrans-2025 dataset~[R47] (1,026,867 annotated DeFi transactions, 79 features), with cross-validation against the Elliptic Bitcoin Dataset~[R48]. Precision, recall, F1, and ROC-AUC results are reported in Section~\ref{sec:aiml-support}. The oracle commit-reveal wiring to the deployed \texttt{LocalBank} contract on Polygon~Amoy constitutes Sprint~3; the architectural specification and ML benchmark together constitute the Contribution~3 claim at pre-thesis stage.

\paragraph{Extension --- Autonomous agent architecture (v24).}
In v24, Contribution~3 is extended from a read-only analytics layer to a fully action-capable autonomous agent system. The agent operates as a six-step pipeline: (1)~the user message arrives at the Qwen3-8B inference engine, which has the client's complete on-chain state injected as structured JSON context by the Express.js backend; (2)~for question-answering requests, the agent performs retrieval-augmented generation from policy documents and replies directly; (3)~for action requests, the agent assembles the full parameter set and presents a plain-language confirmation summary to the client; (4)~the agent waits for explicit affirmative consent before calling any write tool; (5)~upon confirmation, the MCP write tool is called via the Express.js banking API layer, and the resulting on-chain transaction is signed using an EIP-7702 session key scoped to the approved tool set, value-capped at 500~USDC per transaction, and bounded by a 24-hour time-to-live; (6)~the agent monitors the transaction status and notifies the client on completion. When the agent files a loan application on behalf of a client, it also prepares a structured \textit{Authority Brief} for the Local Bank approver, presenting the SHAP feature attribution breakdown alongside a one-click approve/decline interface. No write tool is ever executed without a confirmation turn in the conversation history; this invariant is enforced at the application level and audited via the append-only \texttt{agent\_action\_log} table.

\paragraph{Planned agent architecture enhancements.}
The agent pipeline is strengthened by three architectural additions validated in
the agentic systems literature~[R52]. First, system prompt
construction is formalised into three explicit tiers: a Stable tier (agent
persona, full tool schema, human-gate invariant) that changes only on governance
updates; a Context tier (interest rate schedule, KYC requirement matrix, skill
documents) that updates on policy changes; and a Volatile tier (per-user on-chain
state, conversation history) that refreshes each turn. Second, conversation state
is persisted in a session lineage model: when a session is compressed, a child
session is created seeded by a summary, and the \texttt{parent\_session\_id}
foreign key preserves the ancestry chain; a seventeenth read tool,
\texttt{get\_session\_history}, allows the agent to search this lineage using
full-text queries, turning the session database into a queryable long-term
memory store for returning clients. Third, write-tool enforcement is strengthened
by a confirmation audit middleware hook that enforces the human-gate invariant at
the API layer, independent of model behaviour.

\textit{Implementation note: the confirmation audit hook is implemented in
Phase~II. Three-tier prompt construction, the session lineage migration,
\texttt{get\_session\_history}, and the toolset scoping classifier are fully
specified here and scoped to Phase~III, consistent with the scope-qualifier
language applied to federated learning (Phase~II activation threshold) and
Certora proofs (Phase~IV) elsewhere in this paper.}

\item[\textbf{Contribution 4 --- Compliance-Aware ZKP / zkAML Identity Pathway.}]
We design a two-circuit privacy-preserving compliance architecture combining a zk-SNARK KYC verifier (Groth16 proofs via Circom~2.0 + snarkjs) with a zkAML circuit~[R19] that proves a wallet has not transacted with sanctioned addresses and remains within velocity limits, without revealing the underlying transaction graph. Both circuits compose with a W3C Decentralized Identifier (DID) and Verifiable Credential (VC) layer~[R20] anchoring user identity off-chain. This positions the platform's compliance gateway against the broader Self-Sovereign Identity (SSI) framework and is grounded in the Piper et al. (2025) TU Berlin permissioning paper~[R8].

\textit{Implementation scope:} A Groth16 KYC age-range circuit was compiled using Circom~2.0 and the auto-generated Solidity verifier deployed to Polygon~Amoy (\texttt{KYCVerifier}: address in Appendix~C; to be migrated to Polygon~zkEVM~Cardona in Sprint~1). The circuit proves knowledge of an age satisfying KYC eligibility without revealing the age on-chain, using a Poseidon hash commitment for proof freshness. This constitutes a proof-of-concept implementation of the ZKP identity architecture described in Section~\ref{sec:identity}. The wallet-velocity AML circuit~[R19] (IACR ePrint 2025/465, benchmarked at 55~TPS / 226~ms) requires a significantly larger circuit and is scoped to Future Work; the current deployment demonstrates the end-to-end technical pathway from Circom circuit to on-chain verification.

\end{description}

\paragraph{Positioning against institutional DeFi.}
The Sygnum Bank \textit{Institutional DeFi in 2025} report (February~2026) documents a specific gap: protocols work technically, but no large institutional allocator participates until legal and regulatory risks are resolved~[R21]. The Crypto World Bank is positioned not as a technically novel curiosity but as an attempt to resolve precisely the governance, compliance, and audit gaps that report identifies as the barriers to institutional DeFi adoption. Combined with the mBridge / Agora positioning in Section~\ref{sec:stablecoin-first}, this reframes the contribution from ``an alternative to traditional banking'' to ``a composable lending protocol compatible with emerging institutional CBDC and tokenized-deposit settlement infrastructure''---a substantially more defensible and publishable claim.

The remainder of this paper is organized as follows. Chapter~2 surveys the academic and industry literature on DeFi lending, hierarchical financial architecture, AI-assisted credit assessment, and blockchain-based financial inclusion, identifying the research gaps that motivate this work and presenting the synthesis under a PRISMA-style methodological frame. Chapter~3 presents the system architecture and design, including the smart contract hierarchy, banking product specifications, data model (19 normalized entities), identity and onboarding flows, the formal actor taxonomy and permission matrix, the Chainlink oracle stack, and the compliance and governance framework. Chapter~4 describes the development methodology, sprint plan, the autonomous AI agent pipeline and MCP tool server architecture, AI/ML pipelines, formal-verification strategy, and technology stack justification. Chapter~5 evaluates the technical, economic, regulatory, and sustainability feasibility of the proposed system, including the revenue stream taxonomy. Chapter~6 presents conclusions and future research directions.

\section{Blockchain Justification}

The multi-party coordination and trust problem of incompatible institutions with potential conflicting interests is the root cause of the above problem. We argue that blockchain technology is the solution to this kind of problems than its traditional counterparts due to the following reasons:

\begin{enumerate}[noitemsep]
\item \textbf{Minimizing trust based consensus.} Distributed consensus algorithms make sure that no one can change the ledger state unilaterally and a trusted central operator is not required.

\item \textbf{Programmable enforcement.} The rules of lending, such as limits on borrowing, installment payments and workflows are stored in smart contracts as deterministic self-executable programs~[58], reducing the usage of manual processes and bilateral contracts.

\item \textbf{Cryptographic auditability.} Each change of state is recorded in a tamper-evident and publicly verifiable record (immutability is probabilistic and consensus-enforced on Polygon PoS) of transactions and will be subject to real-time auditing by regulators, partners, and borrowers without expensive third-party interactions.

\item \textbf{Composable incentive structures.} On-chain reputation systems, governance tokens, and programmable fee distributions can align incentives across all network participants. In the Crypto World Bank's architecture, this takes the form of on-chain repayment history contributing to a client's credit score, interest spread revenue being programmatically split between depositor yield and an insurance fund, and governance-controlled parameter adjustments that require multi-tier approval rather than unilateral decision-making.
\end{enumerate}

A conventional cloud-based data architecture would delegate ledger integrity, access control, and audit assurances to a trusted central operator, re-establishing the same trust dependency that this project is designed to eliminate. Blockchain eliminates this single point of trust failure.

\subsection{Blockchain Fundamentals and the EVM Execution Model}
\label{sec:blockchain-fundamentals}

\noindent\textit{\small This subsection provides a concise technical reference for the blockchain primitives on which the Crypto World Bank is built. Readers familiar with EVM architecture, smart contract state machines, and oracle design may proceed directly to Section~\ref{sec:proposed-solution}.}

\vspace{4pt}
\textbf{A. What blockchain is (technically precise).}
A blockchain is an append-only distributed ledger in which each block contains a cryptographic hash of the previous block, a Merkle root of all transactions in the block, a timestamp, and nonce or validator signature data. Immutability is not absolute---it is probabilistic and consensus-enforced: it holds as long as no coalition controlling more than 50\% of stake (in Proof-of-Stake) cooperates to rewrite history~[R17]. On Polygon PoS specifically, finality is achieved through the Heimdall validator layer (a Tendermint-based BFT consensus layer), which provides stronger finality guarantees than pure longest-chain PoW, but still assumes that fewer than one-third of validators are Byzantine. This trust assumption is a design constraint, not a deficiency, and must be stated explicitly in any security analysis.

\textbf{B. The EVM Execution Model (Ethereum Yellow Paper~[58]).}
The Ethereum Virtual Machine is a stack-based, deterministic, Turing-complete virtual machine that executes smart contract bytecode. Every operation (opcode) has a fixed gas cost specified in the Ethereum Yellow Paper~[R7]. Key opcodes relevant to this system:
\begin{itemize}[noitemsep]
  \item \texttt{SSTORE} (write to storage): 20,000 gas for a new slot, 2,900 for modification---the most expensive operation class.
  \item \texttt{SLOAD} (read from storage): 2,100 gas (warm access), 100 gas (cold).
  \item \texttt{CALL} (external contract call): 700 gas base plus additional costs.
  \item Events (\texttt{LOG2}, \texttt{LOG3}): $\approx$375 + 8 per data byte---inexpensive, used for off-chain indexing.
\end{itemize}
Each installment repayment involves at minimum one \texttt{SSTORE} for updating the loan balance, one \texttt{SSTORE} for recording payment history, and one event emission. A 12-installment loan therefore involves at minimum 24 \texttt{SSTORE} operations. A complete retail loan lifecycle involves approximately 27--32 individual on-chain state changes, as detailed in Section~\ref{sec:gas-cost}.

\textbf{C. Smart contracts as state machines.}
A lending smart contract is best modeled as a finite state machine: \texttt{PENDING} $\to$ \texttt{APPROVED} $\to$ \texttt{ACTIVE} $\to$ \texttt{REPAYING} $\to$ \texttt{COMPLETED $|$ DEFAULTED}. State transitions are triggered by role-authorized function calls, and the \texttt{nonReentrant} guard ensures atomicity of each transition. This state machine framing strengthens the argument that smart contracts are more auditable than off-chain systems, because every state transition is permanent, timestamped, and publicly verifiable.

\textbf{D. The Oracle Problem.}
Smart contracts are deterministic and isolated from external data by design. This is a feature---consensus would break if different nodes observed different external states---but it creates a fundamental architectural gap for any system that requires off-chain data. The oracle problem, formally identified by Beniiche~(2020)~[R1] and Pasdar et al.~(2023)~[R2], is the challenge of securely bridging on-chain logic with off-chain information. For the Crypto World Bank, the AI/ML risk scores produced by the FastAPI Random Forest service are off-chain data that must influence on-chain lending decisions. Three architectural options exist: (1) a centralized relay (trusted backend calls \texttt{updateRiskScore} on-chain---simple but re-introduces a central trust point); (2) a decentralised oracle network such as Chainlink Functions (removes the trust assumption at the cost of 30--60 seconds latency and \$0.10--\$1.00 per call; the DON requires consensus across multiple independent nodes before a score is committed, so no single compromised node can manipulate the result); and (3) a commit-reveal scheme (the ML service commits a hash of the risk score before the loan decision window, reveals it after). The current prototype uses a centralized relay, acceptable in the testnet context; Chainlink Functions is adopted as the primary oracle mechanism in v24 (Sprint~3), with the commit-reveal scheme retained as a prototype-phase fallback. In addition to ML score commitment, Chainlink Price Feeds supply BDT/USD and ETH/USD rates to the \texttt{FXModule} contract, Chainlink Automation replaces the centralised cron job for overdue installment detection, and Chainlink Proof of Reserve publishes the \texttt{WorldBankReserve} balance for external cryptographic verification.

\textbf{E. Polygon zkEVM Trust Model.}
Polygon zkEVM Cardona testnet is selected as the primary deployment target for v24, replacing Polygon Amoy PoS. The security model changes materially: rather than delegated proof-of-stake consensus with a validator set, Polygon zkEVM uses ZK validity proofs. Every batch of transactions submitted to the sequencer is accompanied by a zero-knowledge proof verified on Ethereum L1, so security derives from cryptographic validity rather than a validator-set assumption. A malicious sequencer cannot publish an invalid batch because the ZK proof would not verify on Ethereum. This provides a fundamentally stronger guarantee than the PoS model: the trust assumption collapses from ``no coalition of validators controlling more than one-third of staked MATIC colludes'' to ``the ZK proof system is sound and Ethereum L1 is live.'' As a ZK rollup, Polygon zkEVM inherits Ethereum's full economic security for finality, while offering comparable throughput and gas costs to Polygon Amoy on testnet. The stronger security story is particularly significant for an institutional banking prototype: ``ZK-secured reserve ratios'' is a materially more defensible claim than ``PoS-secured reserve ratios'' in a thesis context and in regulatory discussions.

\section{Proposed Solution}
\label{sec:proposed-solution}

The Crypto World Bank is based on a \textbf{four tier on-chain lending model}:

\begin{itemize}[noitemsep]
\item \textbf{Tier 1 --- World Bank:} Custodian of the global crypto reserve. Allocates capital to registered National Banks via lending mediated by smart contracts.
\item \textbf{Tier 2 --- National Banks:} Borrow from the World Bank reserve. Lend to Local Banks incorporated in their jurisdiction. Aggregate the risk exposure at the national level.
\item \textbf{Tier 3 --- Local Banks:} Borrow from National Banks. Process loan applications submitted by retail clients. Administer the full loan lifecycle using designated approvers.
\item \textbf{Tier 4 --- Clients (Retail):} Submit loan requests to Local Banks. Repay through configurable installment schedules. Build an on-chain repayment history that determines future credit access. Includes individual retail clients, group lending members, and small-business clients.
\end{itemize}

The platform also incorporates:

\begin{itemize}[noitemsep]
\item \textbf{Risk-sensitive borrowing limits} computed from rolling 6-month and 1-year transaction windows.
\item \textbf{Automatic installment generation} for loan amounts exceeding a configurable threshold (e.g., 100 ETH equivalent).
\item \textbf{Off-chain AI/ML security analytics} (planned) for fraud detection (Random Forest), anomaly identification (Isolation Forest), and explainable risk assessment (SHAP).
\end{itemize}

\subsection{Banking Functions of the Platform}

A functionally complete bank performs six core activities: deposit mobilization, credit allocation, payment and settlement, risk intermediation, liquidity management, and ancillary financial services. The Crypto World Bank is designed to implement each of these functions across its four-tier hierarchy.

\begin{description}[leftmargin=0pt]
\item[\textbf{Deposit Mobilization.}] The process by which a bank accepts funds from savers and transforms them into productive capital. On the platform, depositors at any tier can place funds into savings products---standard savings accounts with variable yield, fixed-term deposits with locked periods and agreed APY, and institutional yield accounts for large participants. Deposits and accrual are recorded on-chain, and interest accrues automatically via deterministic rules rather than discretionary accounting.
\item[\textbf{Credit Allocation.}] Loans flow downward through the four-tier hierarchy---from the World Bank reserve to national banks, from national banks to local banks, and from local banks to end borrowers. Each tier applies its own interest rate spread, collateral requirement, and borrowing limit, enforced by smart contract rules and supported by data-informed monitoring.
\item[\textbf{Payment and Settlement.}] Transfers between registered accounts settle atomically in a smart contract transaction, avoiding the intermediate ``funds-in-transit'' state that commonly produces disputes in traditional systems.
\item[\textbf{Risk Intermediation.}] Implemented through reserve-ratio enforcement constraints, role-based approval workflows, and an AI/ML monitoring layer (Random Forest for fraud detection, Isolation Forest for anomaly detection, SHAP for explainability) designed for transparent, auditable decision support.
\item[\textbf{Liquidity Management.}] Enforced via minimum reserve ratios at each tier, with same-tier interbank lending pools for short-term liquidity balancing and planned asset-liability monitoring to detect unsafe duration gaps.
\item[\textbf{Ancillary Financial Services.}] Includes foreign exchange, group lending, trade finance facilitation, and digital identity management. FX is designed around decentralized price oracles; group lending enables pooled collateral and mutual liability; trade finance instruments can be added as planned extensions.
\end{description}

\subsection{Service Delivery Across Scale}

The platform is designed to operate at three scales of engagement: the global institutional level, the national and regional commercial level, and the individual retail level.

\textbf{Global institutional level (World Bank Reserve tier).} The World Bank Reserve functions as a programmable reserve authority by managing system-level capital allocation and reserve constraints across jurisdictions. This tier is designed to support transparent reporting of reserves and enforce reserve-ratio constraints as code, reducing informational asymmetry relative to periodic self-reporting.

\textbf{National and regional level (National Bank and Local Bank tiers).} National Banks borrow wholesale capital from the reserve tier and allocate it to registered Local Banks under jurisdiction-aware governance and tiered approval workflows. The intent is to mirror real-world development finance intermediation while improving settlement speed and auditability through on-chain execution.

\textbf{Retail level (clients and depositors).} End customers interact through Local Bank interfaces to open accounts, deposit savings, apply for loans (including group loans), repay installments, and build an on-chain credit history. In Bangladesh, where roughly 40\% of adults lack access to formal banking services~[54], a mobile-accessible interface with low transaction costs can reach users that branch-based systems do not, while keeping loan terms and reserve behavior verifiable through on-chain transparency.

\subsection{Cross-Tier Lending System}

Capital does not necessarily flow down in the old banking system. Banks at the same level deposit with one another regularly in the interbank lending market (e.g., the federal funds market) and less significant banks that have surplus cash inject the extra cash upwards with their home institutions or central banks. The Crypto World Bank is the extension of its hierarchical architecture to accommodate such \textbf{multi-directional lending flows}, developing a more realistic and robust financial ecosystem.

\subsubsection{Same-Tier Lending}

Banks at the same hierarchical level can lend to one another to manage short-term liquidity:

\begin{itemize}[noitemsep]
\item \textbf{National Bank $\leftrightarrow$ National Bank:} A national bank, which has excess reserves, can lend to a second party that is in a liquidity crunch, similar to the traditional interbank lending market. The lending rate is adjusted according to the supply and demand through national banks of the system, like the Secured Overnight Financing Rate (SOFR).
\item \textbf{Local Bank $\leftrightarrow$ Local Bank:} Local banks in same or different national jurisdictions can share liquidity through a peer lending pool. This prevents local liquidity crunches when one bank possesses surplus deposits and the other is experiencing elevated loan demand.
\end{itemize}

Same-tier lending is done by an on-chain \textbf{InterBankLendingPool} contract at every tier level in which surplus banks provide short-term liquidity and those banks that need it are allowed to borrow at utilization floating rates. v15 specifies this pool fully in Section~\ref{sec:interbank-pool} (rate model, settlement, default cascade, operations); the present chapter introduces only the role it plays in the four-tier hierarchy. The current pre-thesis~1 prototype implements the downward flow (World Bank $\to$ National Bank $\to$ Local Bank $\to$ Client); the InterBankLendingPool is built in Sprint~2 of the final-thesis phase.

\subsubsection{Upward Lending}

Banks with excess capital of lower tier are then able to lend up the ladder:

\begin{itemize}[noitemsep]
\item \textbf{Local Banks $\to$ National Banks:} When local banks accumulate reserves beyond the minimum reserve ratio, they are allowed to deposit excess capital in their parent national bank's liquidity pool with earned deposit yield. This mirrors how commercial banks save the surplus with central banks by means of standing deposit facilities.
\item \textbf{National Banks $\to$ World Bank:} National banks may provide excess to the international deposit, enhancing the total capital base of the system. This provides a back route to well capitalized national banks and raises the reserve available for downward distribution to banks with higher demand.
\end{itemize}

The downward lending rates are higher than the upward lending rates (the low risk of lending to a higher-level institution), developing a natural interest rate term structure across the hierarchy.

\subsubsection{Tiered Client Access}

Different client classes access different tiers based on loan size and organizational type:

\begin{table}[H]
\centering
\caption{Client tier access rules.}
\label{tab:borrower-tiers}
\begin{tabular}{p{3.5cm}p{3.5cm}p{2.5cm}p{4cm}}
\hline
\textbf{Client Type} & \textbf{Accessible Tiers} & \textbf{Loan Range} & \textbf{Use Case} \\
\hline
Individual / Retail Client & Local Bank only & 0.01--10 ETH & Personal, micro-enterprise \\
\hline
Small Business / SME & Local Bank, National Bank & 1--100 ETH & Working capital, equipment \\
\hline
Large Corporate & National Bank, World Bank & 50--10{,}000 ETH & Infrastructure, large projects \\
\hline
Institutional / Sovereign & World Bank only & 1{,}000+ ETH & Development programs \\
\hline
\end{tabular}
\TableScreenshotEnd{Tiered Client Access Rules by Client Type and Loan Size}
\end{table}

\vspace{2pt}
\noindent\textit{\small This table summarizes tiered client access rules by mapping client type and loan size to the appropriate institutional tier. The intent is to preserve real-world banking intermediation: retail clients primarily access Local Banks, while larger institutional clients can be routed to higher tiers under stricter verification. This design supports risk segmentation and prevents large actors from consuming retail liquidity.}

Local Banks are used by end users and retail clients to gain access to the system, maintaining the hierarchical intermediation model. Higher tiers can be accessed directly by large corporate and institutional clients, subject to rigorous verification (on-chain credit history with off-chain documentation) and a client-type designation in the smart contract that determines tier access permissions.

This multi-way lending scheme---downward capital allocation, same-tier interbank lending, upward repatriation of surpluses, and tiered access by clients---generates a more detailed model of the banking flows of the real world in the decentralized architecture. The current prototype fully implements the downward capital distribution and tiered client access components; same-tier interbank lending and upward surplus repatriation are specified in full in Section~\ref{sec:multi-entity} (\texttt{InterBankLendingPool}, \texttt{UpwardDepositFacility}) and scheduled for implementation in Sprint~2 of the final-thesis phase. The genuinely multi-entity operations that build on top of these flows---syndicated lending, tranched pools, cross-tier treasury FX, and multilateral settlement netting---are likewise specified at full depth in Section~\ref{sec:multi-entity} so that the paper describes a complete banking system rather than only retail-facing lending.

\section{Methodology in Brief}

Our approach combines a lightweight Agile/Scrum delivery process with an explicit Software Development Life Cycle (SDLC) frame. Pre-thesis~1 covers the requirements, architecture, and design phases of the SDLC; the implementation, verification, and validation phases run across three subsequent sprints during the final thesis stage. Sprint~1 builds the four-tier smart contract scaffold (World Bank Reserve, National Bank, Local Bank), the wallet authentication flow, the frontend skeleton, and the PostgreSQL schema. Sprint~2 implements lending services (loan application, approval, installment generation, borrowing limit enforcement), the deposit-mobilization modules (SavingsVault, FixedDeposit), the group-lending pool, and the same-tier / upward and treasury-swap modules (\texttt{InterBankLendingPool} per Section~\ref{sec:interbank-pool}, \texttt{UpwardDepositFacility} per Section~\ref{sec:upward-flow}, \texttt{TreasurySwap} per Section~\ref{sec:treasury-fx}). Sprint~3 wires the off-chain AI/ML risk-scoring pipeline (Random Forest, Isolation Forest, SHAP) into the on-chain loan decision via Chainlink Functions (DON-based trustless oracle; the commit-reveal relay is retained as the prototype fallback per Section~\ref{sec:oracle-architecture}), layers the Graph-Neural-Network and federated-learning extensions on top, adds the syndicated-lending, tranched-pool, and netting-engine contracts (\texttt{SyndicatedLoan} per Section~\ref{sec:syndicated-lending}, \texttt{TranchedPool} per Section~\ref{sec:tranched-pools}, \texttt{NettingEngine} per Section~\ref{sec:netting-settlement}, plus the off-chain Settlement Coordinator service), and adds the Foundry invariant suite, Certora reserve-invariant specifications, and Tenderly runtime monitoring. The implementation stack is Solidity~0.8.20 (smart contracts), React~18 with TypeScript and ERC-4337 account abstraction (frontend), Express.js + FastAPI (backend), scikit-learn + PyTorch Geometric (ML), Circom~2.0 + snarkjs (ZKP), and ChromaDB + a locally hosted 8B-parameter LLM (RAG assistant). All deployments target Polygon~zkEVM~Cardona and Ethereum~Sepolia testnets so that the prototype carries zero real-cryptocurrency cost; stablecoin pools use testnet USDC.

\section{Scopes and Challenges}

\noindent The complete banking architecture described in this report defines both what is implemented in the current prototype and what constitutes the full system design, clearly distinguishing prototype scope from architectural intent throughout.

\textbf{Scope.} This prototype is restricted to publicly available testnets, i.e. no actual cryptocurrency is involved in the system. It encompasses basic features including four-tier lending architecture, loan request and approval process, installment based repayment, calculation of borrowing limits, communication between borrowers and banks and verification of income. The system also uses Random Forest with SHAP to explain explainability in fraud detection and Isolation Forest to detect anomalies. An additional feature of dual-currency is also introduced to enhance the current banking infrastructure. The prototype, however, does not cover mainnet deployment, fiat on/off-ramp integration, automated KYC/AML compliance or production-scale stress testing.

\textbf{Challenges.} There are some issues that come with this system. The unlabeled nature of DeFi lending fraud data is one of those problems, as it limits the successful training of the model and might necessitate the application of synthetic data or transfer learning methods. The other issue is regulatory uncertainty, since crypto lending may have different jurisdictional restrictions; and this risk can be partially reduced by working the testnets only. The sensitivity of gas costs is an issue as well, as intricate on-chain programs can be costly; to solve this, AI/ML work requiring computation is performed off-chain. Additionally, model interpretability is essential to regulatory compliance, especially in describing the process of lending, which is obtained by SHAP-based feature attribution. Last but not the least, economic sustainability should be taken into account, because the platform should be able to earn enough income, like interest revenues, to cover gas expenses at all four levels. The following Chapter~5 discusses this balance in greater depth.

\subsection{Economic Sustainability: Interest Revenue Versus Gas Costs}

The most important problem facing any blockchain-based lending system is whether the revenue earned by the lending activities will be in a sustainable position to cover the total gas expenses of the transactions on-chain. Every level in the hierarchy will have deposits, loan applications, approvals, disbursements, repayment and installment processing gas fees. One complicated smart contract interaction can cost between \$5 and \$50 on Ethereum mainnet, based on network congestion, but around \$0.001 to \$0.01 on Polygon zkEVM Cardona (comparable to Polygon PoS for testnet operations). At an average retail loan of 10~ETH (\$20,000 at \$2,000/ETH) with an 8\% APR, the client pays \$1,600 in annual interest. With 10--15 on-chain transactions per loan lifecycle (request, approval, disbursement, 12 monthly installments), the overall gas costs on Polygon are under \$0.15---less than 0.01\% of the interest obtained. The platform thus provides a sustainable margin on Layer~2 networks, but mainnet implementation would need to consider gas optimisation or batched transaction processing to ensure it would be viable to serve micro-loans.

\subsection{Resistance to Institutional Capture}

As the Crypto World Bank grows in adoption, a critical question is whether large capital holders and established financial institutions---which already dominate the traditional banking system---could migrate to the new architecture and replicate the existing patterns of financial concentration. This risk is mitigated by the platform in a variety of architectural ways:

\begin{itemize}[noitemsep]
\item \textbf{Visible reserve ratios:} In contrast to traditional banks, where reserve adequacy is reported and audited quarterly, on-chain reserves are publicly verifiable in real time. There is no way a given institution can hide its real financial status to have unfair advantage.
\item \textbf{Algorithmic interest rates:} Interest rates are determined by the smart contract parameters, not by opaque internal pricing committees, and thus no single participant can capture the entire pool of capital.
\item \textbf{Tamper-evident audit trails:} Every transaction, approval, and governance decision is permanently recorded on-chain, making regulatory capture and corruption detectable by any participant.
\item \textbf{Open-source governance:} The smart contract code is publicly auditable, preventing hidden backdoors or preferential treatment encoded at the protocol level.
\item \textbf{Tiered access with caps:} Borrowing limits and tier access rules are enforced programmatically, preventing any single entity from monopolizing the capital pool.
\end{itemize}

Although no system can entirely deter wealth concentration in a free market, the detectability and programmability characteristics of blockchain infrastructure make exploitative behavior prohibitively expensive and highly observable by comparison to opaque intermediaries.

\subsection{Stablecoin-First Lending and Supply Scalability}
\label{sec:stablecoin-first}

Denomination matters more than any other design choice for a crypto financial services platform that targets developing economies. A retail client in Bangladesh who takes a 0.5~ETH micro-loan during a period of price stability and is later required to repay the same 0.5~ETH after a 50\% drawdown effectively faces a 2$\times$ increase in the real value of their debt. ETH has experienced exactly this drawdown twice in the past five years (May 2021 and 2022), and the stablecoin literature has formalized the resulting liability-side volatility risk for crypto-asset borrowers~[R12, R13]. The Crypto World Bank therefore elevates \textbf{stablecoin-denominated lending} from a future-work footnote to an explicit design requirement at every retail tier. As of April~2026, USDC commands a circulating supply of approximately \$77.3~billion and USDT approximately \$189.5~billion (DeFiLlama stablecoin tracker, accessed April~2026), confirming that stablecoin-denominated transactions have achieved dominant adoption across consumer DeFi applications---a market signal that validates the platform's denomination choice.

The current prototype continues to use ETH as the testnet denomination for simplicity, but every loan-bearing contract has been designed so that the loan currency is a per-pool parameter rather than a hard-coded asset. In the final thesis phase, the retail Local Bank tier will deploy with USDC as the default loan currency; ETH remains acceptable only for institutional Tier~1 / Tier~2 operations where the institutional client has the balance-sheet capacity to hedge crypto-asset exposure.

Three regulatory anchors shape this choice. First, the EU Markets in Crypto-Assets (MiCA) Regulation, fully in effect since December 2024, classifies the kind of fiat-pegged tokens used here as electronic-money tokens (EMTs) and requires 1{:}1 reserve backing, authorization before public offering, and regular audits~[R36, R37]. Second, the United States GENIUS Act, signed into law in July 2025, introduces the first federal framework for payment stablecoins and codifies a similar reserve and disclosure regime~[R38]. Third, the BIS \textit{mBridge} platform reached Minimum Viable Product status in mid-2024 and now settles real-value cross-border CBDC transactions among China, Hong Kong, Thailand, the UAE, and Saudi Arabia, while the BIS / central-bank \textit{Project Agora} explores tokenized bank deposits on a unified ledger~[R39, R40]. Neither initiative implements a retail lending hierarchy: they are settlement rails. The Crypto World Bank is positioned as the lending layer that can compose on top of such institutional settlement infrastructure rather than as a replacement for it.

The remaining supply-side concerns are addressed by three complementary measures:

\begin{enumerate}[noitemsep]
\item \textbf{Stablecoin pool support.} Each Local Bank deploys an ERC-20 stablecoin loan pool (USDC primary, USDT and DAI as fallbacks). The InsuranceFund and SavingsVault contracts are denominated in the same unit so depositors, borrowers, and the reserve share a single numeraire.
\item \textbf{Multi-chain and Layer~2 deployment.} The platform deploys on Polygon~zkEVM Cardona testnet for retail traffic (ZK validity proof security model) and on Ethereum~Sepolia for high-value institutional flows, with the documented cross-chain bridge architecture in Section~\ref{sec:bridge} preserving the hierarchy invariant across chains.
\item \textbf{Capital circulation via the four-tier hierarchy.} The Reserve Ratio (RR) enforced at every tier functions as a solvency constraint, not a money-creation parameter. USDC is 100\%-reserve backed; the platform re-allocates existing stablecoin supply rather than creating new money. The rate at which this fixed capital pool flows through the hierarchy and returns as repayment is measured by Credit Velocity (Formula~CV). Higher credit velocity---not money multiplication---is the platform's capital-throughput mechanism.
\end{enumerate}

The supply concern that originally motivated this subsection is therefore reframed: scarcity is a property of base-layer crypto-assets, not of the lending denomination. Fiat-backed stablecoins issued under MiCA-/GENIUS-compliant reserves provide the elastic, audit-friendly supply a retail banking platform requires, while the hierarchical reserve architecture promotes capital velocity across tiers. Client-side currency risk is removed, deposit-side regulatory compliance is preserved, and the protocol does not depend on the monetary discretion of any single issuer.

\section{Market Analysis and Partnership Ecosystem}

\subsection{Market Sizing}

The platform addresses three nested market segments. The Total Addressable Market (TAM) spans global DeFi lending (exceeding \$55B TVL~[13]) plus the \$860B remittance market~[26] and \$4.5T MSME financing gap~[20]. The Serviceable Addressable Market (SAM) is institutional and semi-institutional lending requiring hierarchical structures (\$5--15B), and the immediately Obtainable Market (SOM) covers pilot deployments in regulatory sandboxes and NGO-backed microfinance (\$50--200M). Full market sizing with data sources is presented in Table~\ref{tab:market-sizing-ch5} in Chapter~5.

\subsection{Target Customer Segment}

The Crypto World Bank targets retail clients and small businesses in developing economies who lack formal banking access---particularly the estimated 1.4~billion unbanked adults globally~[14]. The platform's hierarchical architecture (World Bank $\to$ National $\to$ Local $\to$ Client) mirrors development-finance intermediation while providing blockchain-enforced transparency and AI-enhanced risk assessment. Full client segment profiling is presented in Chapter~5 (Table~\ref{tab:customer-segment-ch5}).

\subsection{Partner Ecosystem}

The platform's deployment requires five partner categories: financial regulators (regulatory sandbox approval), banking institutions (National/Local Bank network members), payment gateway providers (fiat on/off-ramp), academic and research institutions (AI/ML validation), and non-governmental organizations (field pilots with underserved client populations). Blockchain-mediated incentives align each partner's interests with platform operation through on-chain transparency, reduced settlement friction, and programmatic fee distribution. Full partner ecosystem detail is in Chapter~5 (Table~\ref{tab:partner-categories-ch5}).

\subsection{Incentive Alignment Through the Blockchain Platform}

The system manages and enforces partner incentives directly on the blockchain, making the process more transparent and reliable:

\begin{itemize}[noitemsep]
\item \textbf{Tamper-evident repayment records} serve as an on-chain reputation system through permanent repayment records, allowing credit decisions to be based on actual data without relying on external credit bureaus.
\item \textbf{Transparent reserve verification} enables on-chain reserve verification, so all participants can independently confirm that allocated funds are being used as intended.
\item \textbf{Programmable fee structures} (with potential for future expansion) allow transaction fees to be distributed fairly among network participants according to their roles and contributions.
\end{itemize}


%% %%%%%%%%%%%%%%%%%%%%%%%%%%%%%%%%%%%%%%%%%%%%%%%%%%%%%%%%%%%%
%% CHAPTER 2: LITERATURE REVIEW
%% %%%%%%%%%%%%%%%%%%%%%%%%%%%%%%%%%%%%%%%%%%%%%%%%%%%%%%%%%%%%
\chapter{Literature Review}

\section{Preliminaries}

Key technical terms used throughout this thesis---including Decentralized Finance (DeFi), smart contracts, over-collateralization, Explainable AI (XAI), Proof of Stake (PoS), CBDC, Solidarity/Group Lending, Asset-Liability Management (ALM), Flash Loans, Zero-Knowledge Proofs (ZKP), and Health Factor (HF)---are defined in the \textbf{Glossary} appendix (labelled Appendix D in the final thesis; the current pre-thesis build carries the WorldBankReserve interface as Appendix D pending the final assembly pass). The Glossary also includes platform-specific terminology introduced in this thesis (e.g., \textit{Platform Solvency Ratio (PSR)}, \textit{Provisional Credit Tier}, \textit{Cold-Start Credit Pathway}). This chapter assumes familiarity with blockchain fundamentals at the level of Section~\ref{sec:blockchain-fundamentals}; other terms are defined on first use.

\subsection{Review Methodology (PRISMA-Style Frame)}
\label{sec:lit-prisma}

The literature underlying this thesis is synthesized under a PRISMA-style evidence-synthesis frame~[R22] rather than as a narrative summary alone. Inclusion criteria were: (i) peer-reviewed publications or formally published institutional reports issued between 2017 and 2026; (ii) direct relevance to at least one of seven scoped CWB modules---DeFi lending mechanics, ML-based fraud / anomaly detection, ZKP / DID-based identity, group lending and microfinance, smart-contract security, financial-inclusion economics, or institutional DeFi compliance; and (iii) sufficient methodological transparency to allow re-use of design parameters (e.g., reserve invariants, gas accounting, F1 scores). Sources were identified through ACM Digital Library, IEEE Xplore, Elsevier ScienceDirect, MDPI, arXiv (with venue check), and the BIS / World Bank publication portals; queries combined the keywords \texttt{blockchain}, \texttt{DeFi}, \texttt{lending}, \texttt{fraud detection}, \texttt{ZKP}, \texttt{microcredit}, and \texttt{financial inclusion} with regional anchors (\texttt{Bangladesh}, \texttt{MFI}). The full evidence synthesis (114 entries) is held in \texttt{references.bib}; the top-10 ranked subset is summarized below and in Table~\ref{tab:lit-top10}. Recent systematic reviews of blockchain in banking using PRISMA on 38 studies confirm that this methodological frame is the current expected standard for journal-quality synthesis~[R23].

\begin{table}[H]
\centering
\caption{Top-10 ranked literature subset under the PRISMA-style frame. Ranking is by direct architectural impact on the CWB design; each entry maps to a v15 section that operationalizes the finding.}
\label{tab:lit-top10}
\small
\renewcommand{\arraystretch}{1.15}
\begin{tabular}{p{0.4cm}p{3.4cm}p{3.0cm}p{3.0cm}p{3.4cm}}
\hline
\textbf{\#} & \textbf{Author / Source} & \textbf{Module} & \textbf{Headline Finding} & \textbf{Operationalized In} \\
\hline
1 & Werner et al.\ 2022 [1] & DeFi mechanics & DeFi lending is uniformly flat, pool-based, overcollateralized & Four-tier architecture (Ch.~3) \\
\hline
2 & Tan 2023 (IMF) [5] & CBDC \& inclusion & Two-tier CBDC distribution (central bank $\to$ commercial banks $\to$ users) achieves 30--40\% higher financial inclusion in developing economies; hierarchical infrastructure is the key design choice & Section~\ref{sec:bangladesh-reg} \\
\hline
3 & Palaiokrassas et al.\ 2023 [3] & Fraud detection & Multi-chain ML on 54M tx, F1 0.76--0.85 with features & Section~\ref{sec:aiml-support} \\
\hline
4 & Adom et al.\ 2022 [4] & Explainable AI & SHAP outperforms LIME on lending datasets & Section~\ref{sec:aiml-support} \\
\hline
5 & Weber et al.\ 2019 / Elliptic [R24] & Graph fraud & Graph-structured fraud signal in BTC transactions & Section~\ref{sec:gnn-extension} \\
\hline
6 & McMahan et al.\ 2017 / FedAvg [R25] & Federated learning & Privacy-preserving model averaging across institutions & Section~\ref{sec:fl-fraud} \\
\hline
7 & BIS WP 905 [R13] & Stablecoin risk & Algorithmic stablecoins unsuitable for EM DeFi; USDC/USDT recommended & Sections~\ref{sec:stablecoin-first}, \ref{sec:mica-compliance} \\
\hline
8 & Atzei et al.\ 2017 [7] & Smart contract security & SoK of EVM attack vectors & Section~\ref{sec:defense-in-depth} \\
\hline
9 & EU MiCA + GENIUS Act [R36, R38] & Regulation & Stablecoin issuer authorization, reserve disclosure & Section~\ref{sec:mica-compliance} \\
\hline
10 & Bangladesh Bank SP2025-02 [R44] & Inclusion data & Female-to-male account parity 49\%/49\% in 2025; 89\% BRAC loans to women & Section~\ref{sec:bangladesh-reg} \\
\hline
\end{tabular}
\end{table}

\section{Review of Existing Research}

The idea of the Crypto World Bank is designed based on the examination of peer-reviewed materials, institutional reports, and industry figures covering more than one sphere of direct interest to our architecture: blockchain machine learning architecture, decentralized lending protocol architecture, security, explainable AI in finance, blockchain for financial inclusion, smart contract security, correspondent banking economics, monetary policy distribution effects, and real world asset tokenization.

\subsection{Decentralized Lending and DeFi Protocol Design}

In their SoK paper, Werner et al.~[1] systematize DeFi protocol design and found that current lending markets are homogenous, pool based, and over-collateralized, and they do not exist institutional hierarchy---a gap which is directly addressed by this project. Their taxonomy of DeFi protocols indicates that there is no current system modeling that can represent multi-tier capital flow in a similar manner to development finance, ensuring the innovativeness of our four-level approach.

Bastankhah et al.~[2] present a data-driven DeFi lending protocol that is adaptive dual fast/slow control architecture which proves the dynamic interest rate adjustment compares very well with the static utilization curves employed by Aave and Compound~[59]. Their experience confirms that it is a viable method to optimize the algorithmic lending parameters, which is justified by work we go through our adjustable levels of interest rates, through the four levels of the hierarchy.

Li et al.~[46] introduce a multi-chain lending model, which allocates lending activities on various blockchain networks in order to enhance throughput and lower single-chain congestion. Their cross-chain settlement system certifies technical feasibility of executing lending protocols in heterogeneous blockchain settings, directly promoting our intended multi-chain implementation approach on Ethereum and Polygon.

Xu et al.~[47] come up with a DeFi lending protocol evaluation system that quantifies protocol risk by measuring such metrics as the stability of utilization ratio, liquidation efficiency, and interest rate responsiveness. Their assessment system gives them a flexible methodology, which can be implemented to evaluate the performance of our hierarchical lending layers versus the current flat-pool protocols.

Sharma et al.~[48] introduce a peer-to-peer lending framework that is powered by blockchain that removes intermediary overhead by the use of smart contract-mediated loan origination and repayment. Although they are still single-tier in architecture, the cost of their gas analysis on-chain credentials and management of identity through optimization strategies and borrowers guides our strategy in reducing transaction costs in the four level hierarchy. Dao et al.~[60] demonstrate that credit scoring models can be adapted for DeFi lending environments, validating the feasibility of data-informed borrowing limits within decentralized protocols.

\subsection{Machine Learning for Blockchain Security}

Palaiokrassas et al.~[3] use machine learning on multichain DeFi fraud over 54~million transactions in 23 protocols which show a demonstration that the behavioral characteristics of DeFi enhance Neural Network and XGBoost classifiers to F1-scores of 0.76--0.85 versus 0.08 using transactional features only. Their finding that the behavioral features are superior to raw transaction data directly informed our feature engineering method for the random forest fraud detection model.

The algorithm of unsupervised anomaly detection presented by Liu et al.~[6] is the Isolation Forest algorithm. The algorithm identifies anomalies using recursive partitioning of data, which assigns shorter distances to outliers. We used Isolation Forest as the second detection model when there are limited labeled fraud data to analyze wallet behavior, it is possible to detect new patterns of attacks not pre-labeled.

Hassan et al.~[49] introduce blockchain and machine learning to detect fraud with the help of a privacy-protecting and adjustment-flexible incentive-based method. Their framework trains ML classifiers on encrypted transaction characteristics via federated learning, whereby the operators of the detection model never get to see the individual transaction data. Our future extensions of our AI/ML layer involve the privacy-preserving paradigm, in which the histories of client transactions should be confidential and at the same time allow system-wide fraud pattern recognition.

\subsection{Explainable AI in Financial Decision Systems}

Direct comparison of LIME and SHAP on loan approval is given by Adom et al.~[4], SHAP offers more consistent and in-depth feature attributions, which are found to be deeper via Shapley values whereas LIME provides quicker running time but less sure descriptions across repeated evaluations. We have used their comparison to make the decision to implement SHAP since the main explainability technique of loan risk measurements, trading off interpretability depth with regulatory compliance requirements.

Gupta et al.~[56] use interpretable models that are based on SHAP on credit default assessment proving that SHAP feature attributions on Gradient Boosting and Random Forest classifiers are able to score over 0.92 with complete interpretability of individual predictions. Their approach to causing per-client risk explanations tells our intended credit risk dashboard, where the approvers of loans look at SHAP waterfall plots before lending decisions are made.

\subsection{Blockchain for Financial Inclusion}

Tan~[5] constructs a model of the International Monetary Fund (IMF) according to which CBDCs in developing countries can bank large unbanked populations using two-tier distribution model (central bank $\to$ commercial banks $\to$ users) that directly parallels our four-tier hierarchy. The paper shows that the distributed infrastructure of digital currency can prove to be effective by use of available institutional means, can access those populations not covered by traditional banking---facilitating the mission of the Crypto World Bank, which is transparent and programmable lending to underserved populations.

The article by Alam et al.~[54] examines how blockchain technology can be used in the microcredit industry particularly in Bangladesh, it can be proven that smart contract-based loan management is capable of minimizing the administrative overhead by up to 60\% as compared to manual microfinance processes. They have analyzed the Bangladeshi financial environment where about 40\% of adult population do not have access to formal banking---directly confirms the geographic and demographic attractiveness of the target market of the Crypto World Bank.

The article by Islam et al.~[53] examines the design requirements and challenges of Central Bank Digital Currencies, there are two levels of retail CBDC model (central bank to commercial banks to end users) as the distribution architecture of choice. Their analysis of interoperability challenges between CBDC systems and existing payment infrastructure informs our design of the dual-currency facility that bridges fiat and cryptocurrency within participating banks.

\subsection{Group Lending and Microfinance Digitization}

The solidarity group lending model, pioneered by Grameen Bank in Bangladesh and scaled by BRAC into one of the world's largest microfinance institutions, demonstrates that peer-monitored mutual liability can achieve repayment rates exceeding 95\% among borrowers with no individual collateral or formal credit history. These rates are attributable to social pressure mechanisms---weekly physical group meetings, field officer oversight, community reputation, and progressive loan eligibility---that are preventive in nature. The blockchain extension of this model encodes mutual liability in programmable contract logic (multi-signature consent, on-chain group formation, automatic collateral pool claims), but cannot replicate the preventive social pressure that drives Grameen's repayment outcomes. The CWB group lending module therefore targets the structural design of solidarity lending while acknowledging that initial repayment rates on-chain are expected to be lower until credit history accumulates. Alam et al.~[54] show that smart contract-based loan management reduces administrative overhead in Bangladeshi microfinance by up to 60\%, directly validating the economic case for the platform's group lending module. Güçük et al.~[61] provide a broader literature review of privacy and trust considerations in blockchain-based microlending, reinforcing the importance of privacy-preserving mechanisms in the group lending architecture.

\subsection{Smart Contract Security and Governance}

Atzei et al.~[7] list Ethereum smart contract attack vectors such as reentrancy, access control vulnerabilities, integer overflow, and access control vulnerabilities. They were directly informed by their taxonomy our implementation of the ReentrancyGuard by OpenZeppelin, Solidity~0.8.20 inbuilt overflow protection, and role-based access control modifiers. On their advice we took up their suggested defense-in-depth strategy of combining security primitives with formal verification with planning using static analysis (Slither) and symbolic execution (Mythril).

The article by Wang et al.~[50] introduces ContractWard which is an automated vulnerability detection system that uses machine learning classifiers (Random Forest, SVM, k-NN) to classify six types of smart contract vulnerabilities based on the features of bytecodes. ContractWard scores F1-reentrancy and F1-access control with 0.96 and 0.93 respectively, proving that security auditing using ML can be used to supplement traditional static analysis tools. Our proposed security verification pipeline will utilize the use of machine learning-based static analysis tools.

Liao et al.~[51] present SoliAudit, that is an integration of machine learning classification and vulnerability assessment of smart contracts fuzz testing. Their hybrid approach---using ML to first prioritize probable vulnerable code paths then target fuzz testing---reduces audit time by an estimated 70\% of that spent on manual review. SoliAudit's methodology informs our intended security audit plan on the current three-contract prototype and its planned extensions toward the full nine-contract architecture.

So et al.~[52] come up with VERISMART, an extremely accurate Ethereum safety verifier which employs constraint-based abstract interpretation to detect arithmetic overflow, division-by-zero and array out-of-bound errors. VERISMART achieves 98.4\% accuracy on benchmark data, but has zero false positives, so it is a candidate tool to formally verify our smart contract invariants (e.g., reserve ratio maintenance, cascading interest rate limits).

\subsection{AI Security Features for Blockchain Lending}

In addition to the currently existing fraud detection and anomaly identification models in place in the platform, there are other security features that are based on AI that can enhance the Crypto World Bank's resilience to threats of change:

\begin{itemize}[noitemsep]
\item \textbf{Graph Neural Network (GNN) transaction analysis:} GNN-based models are able to analyze the graph structure of transaction to identify coordinated rings of fraud---groups of wallets which conspire to rig the borrowing limits or make circular lending patterns. A study by Palaiokrassas et al.~[3] proves that graph-based features are much superior to flat transaction features as an indicator of DeFi fraud detection.

\item \textbf{Federated learning for cross-bank fraud intelligence:} Hassan et al.~[49] demonstrate that federated learning allows the collaboration between several institutions to detect and train fraud models without distributing raw transaction information. Deploying federated learning of the Crypto World Bank's National and Local Banks would allow detecting threats system wide without compromising the privacy of the per-bank data.

\item \textbf{Reinforcement learning for adaptive interest rates:} As a reinforcement-based learning algorithm, adaptive interest rates can be learned dynamically; RL agents can dynamically modify interest rate parameters to match the varying market conditions, utilization rates, and risk profiles, decreasing the lag of manual parameter governance.

\item \textbf{Real-time smart contract monitoring:} ML-based monitoring agents (informed by ContractWard~[50] and SoliAudit~[51]) can continuously analyze incoming transactions for patterns matching known exploit signatures, triggering the platform's pause mechanism before significant losses occur.

\item \textbf{Natural language processing for income verification:} NLP models are able to extract and validate income information in uploaded documents, saving the bank officers the manual review load but not eliminating verification accuracy.
\end{itemize}

\subsection{Gas Cost Optimization and Layer~2 Scalability}

In DeFi, Tolmach et al.~[55] examine the optimal approaches to minimization of gas fees, demonstrating that transaction batching, calldata compression, and tactical time of on-chain operations can save 30--50\% of gas costs on Ethereum mainnet. Their findings test our architectural choice of deploying to Polygon PoS (where fees already are sub-cent) and inform optimization strategies of the contemplated Ethereum mainnet deployment. The research also finds out that complicated multi-step operations of DeFi (analogous to our four-tier loan lifecycle) benefit disproportionately from Layer~2 deployment due to the multiplicative savings in gas of sequential calls of contracts.

\subsection{Correspondent Banking and Cross-Border Settlement}

The Bank of International Settlement~[21] and Financial Stability Board~[33] have extensively documented the structural inefficiencies of the correspondent banking network. In the traditional model, the banks are required to have pre-funded nostro accounts in every currency corridor, establishing a world pool of idle capital which can not be put into effective productive lending. SWIFT settlements and respondent banks involves routing of messages using MT103 payment directions, screening of compliance at all intermediates, and settlement time extensions to two--five business days. According to the World Bank Migration and Development Brief~[26] it is stated that the cost of transferring \$200 on cross-border is 6.49\% on average, with Sub-Saharan African corridors averaging more than 8\%. These results encourage our development of on-chain settlement with close to instant finality and sub-cent transaction costs on Layer~2 networks.

\subsection{Monetary Policy Distribution and Financial Inequality}

Recent empirical studies have reported the distributional impact of monetary policy on wealth inequality. A 2025 cross-country study spanning 49 nations (1999--2019) discovered that the relationship between central bank asset purchases programs and greater wealth inequality, whose outcomes are stronger in the distribution of wealth and long lasting~[34] than income inequality effects. The Federal Reserve's own 2025 contractionary monetary analysis based on the U.S. metropolitan tax data confirmed that policies harm the income of low income employees~[35]. These findings deliver economic incentive to the design principles of Crypto World Bank: transparent and algorithmically determined monetary parameters (interest rates, reserve ratios, supply rules) are encoded in smart contracts rather than set discretionally by opaque committees. On-chain transparency reduces \textit{price discrimination}---all borrowers at the same tier face the same utilization-driven interest rate, eliminating the ability of loan officers to charge higher rates to less-informed borrowers. This is a distinct mechanism from the Cantillon Effect (a macroeconomic phenomenon concerning newly created money benefiting first recipients); the Cantillon Effect does not apply to the CWB because USDC is 100\%-reserve backed and no new money supply is created by the platform's operations~[36].

\subsection{Real-World Asset Tokenization and Institutional DeFi}

The tokenization of real-world assets is an emerging convergence of traditional and decentralized finance. The Corda platform of R3 boasts of more than \$17~billion in tokenized real-world assets on its permissioned ledger as of September 2025, with institutional participants including HSBC and Bank of America~[37]. Centrifuge has deployed in eight blockchain networks of more than \$1~billion in tokenized institutional fund products~[38]. The World Bank is using blockchain on its own public chain side to track the disbursement of funds using its FundsChain program (based on Hyperledger Besu) which is flown on 13 projects in 10 countries and increases to 250 projects by mid-2026~[39]. These developments confirm the institutional interest of blockchain-based financial infrastructure and the viability of hierarchical fund distribution design models on distributed ledgers. This institutional adoption is continued in the Crypto World Bank's trajectory through carrying out not only fund tracking but total hierarchical lending system that has on-chain interest rates, reserves and credit assessment.

\subsection{Graph Neural Networks for Relational Fraud Detection}
\label{sec:lit-gnn}

Recent work has demonstrated that DeFi fraud is fundamentally a relational phenomenon: coordinated borrowing rings, sybil attacks, and flash-loan manipulation appear in the wallet-interaction graph, not in individual transaction features. Cheng~et~al.~(2024)~[R24] and the DeFiGuard graph-neural-network architecture~[R25] consistently outperform flat ML baselines on blockchain fraud, with reported detection above 70\% of illicit transactions at sub-1\% false-positive rates on the Elliptic dataset; Wang and Wang (2025) report similar advantages on cross-chain laundering~[R26]. This evidence motivates a Graph Neural Network (GNN) extension to the platform's Random-Forest fraud-detection layer, treated as an ablation that demonstrates the marginal value of relational features over the flat behavioral features already used by Palaiokrassas et~al.~[3]. The GNN extension is described in Section~\ref{sec:gnn-extension}.

\subsection{Federated Learning Across Banking Tiers}
\label{sec:lit-fl}

Hassan et~al.~(2022)~[49] introduced privacy-preserving fraud detection via federated learning on the blockchain. The PrivChain-AI study~[R27] in \textit{Nature Scientific Reports} (December~2025) demonstrates the same pattern with 94.7\% fraud-detection accuracy across multiple banks, while Abbassi et~al.~(2025)~[R28] and FED-SPFD~[R29] document the regulatory benefits of FL for cross-border fraud intelligence. The Crypto World Bank's four-tier structure---where National and Local Banks are independent legal entities that cannot share raw transaction data---makes federated learning architecturally appropriate rather than merely fashionable: each Local Bank trains a local fraud model on its own client transaction data, parameter updates are aggregated at the National Bank tier via FedAvg, and no raw records leave the bank. The detailed FL pipeline is given in Section~\ref{sec:fl-fraud}.

\subsection{Stablecoin Regulation: MiCA and GENIUS Act}
\label{sec:lit-stablecoin-reg}

Two recent statutes reshape the legal context of any stablecoin-denominated lending platform. The EU Markets in Crypto-Assets (MiCA) Regulation, fully in effect from December~2024, requires authorized issuers, 1{:}1 reserve backing, and regular audits for electronic-money tokens---categories that directly cover USDC and USDT used at retail tiers of this platform~[R36, R37]. The United States GENIUS Act of July~2025 creates the first federal payment-stablecoin framework with similar reserve and disclosure obligations~[R38]. These statutes are read in this thesis not as obstacles but as the regulatory anchor that makes stablecoin-first design defensible (Section~\ref{sec:stablecoin-first}); the corresponding compliance mapping is presented in Section~\ref{sec:mica-compliance}.

\subsection{On-Chain Credit Passport and Soulbound Tokens}
\label{sec:lit-sbt}

Buterin, Hitzig, and Weyl~(2022)~[R30] introduce the concept of Soulbound Tokens (SBTs)---non-transferable ERC-style credentials that bind reputational data to an identity rather than to a tradable asset. For a multi-tier cryptocurrency exchange and financial services platform, SBTs provide a natural primitive for a portable on-chain credit passport: a client's successful repayment history at one Local Bank becomes a credential they can present to any other CWB Local Bank, and (with the schema fixed) to any external lending protocol that adopts the standard. MicroSave Consulting~(December~2025)~[R31] documents the Bangladesh microfinance sector's active shift toward performance-based credit identity, providing direct empirical grounding for this design choice. The full credit-passport schema is in Section~\ref{sec:credit-passport}.

\subsection{Institutional DeFi Gap and the mBridge/Agora Landscape}
\label{sec:lit-institutional}

Sygnum Bank's February~2026 report \textit{Institutional DeFi in 2025}~[R21] documents the central gap that this thesis attempts to fill: DeFi protocols work technically, but no large institutional allocator will participate until legal and regulatory risks are resolved. The BIS mBridge platform~[R39] (China, Hong Kong, Thailand, UAE, Saudi Arabia) reached Minimum Viable Product status in mid-2024 and now settles real-value cross-border CBDC transactions, and the BIS / central-bank \textit{Project Agora}~[R40] explores tokenized bank deposits on a unified ledger. The IMF CBDC survey 2025 reports that 94\% of central banks are now engaged in CBDC research~[R32]. None of these initiatives implements a retail lending hierarchy: they are settlement rails. This positions the Crypto World Bank explicitly as the lending layer compatible with such institutional CBDC infrastructure rather than as a competitor to it.

\subsection{LLMs in Finance and Hallucination Risk}
\label{sec:lit-llm}

The CWB AI assistant uses a 8B-parameter open-weight LLM with QLoRA fine-tuning, RAG over the platform's policy and contract documents, and an explicit refusal layer for regulated questions. The LLM-in-finance literature~[R33, R34, R35] gives both the design rationale and the necessary cautions: Morgan Stanley's GPT-based equity-analysis assistant achieved a reported 50\% reduction in analyst time, McKinsey estimates LLMs can cut back-office costs by 40\%, but the 2025 LLM-vulnerability-in-finance corpus~[R35] documents specific regulatory-hallucination failure modes when LLMs are asked compliance questions without RLHF from domain experts. The evaluation methodology in Section~\ref{sec:llm-eval} includes a red-teaming protocol for adversarial financial prompts directly motivated by this work.

\paragraph{Model Context Protocol as a safety architecture for financial agents.}
The emergence of Model Context Protocol (MCP) as a structured tool-calling interface represents a significant safety advance for LLM-based agents deployed in financial systems~[R33, R35]. Rather than permitting the model to generate and execute arbitrary code or issue unconstrained API calls, MCP confines the agent to a predefined schema of typed tools---each with well-defined parameters, permitted data access, and bounded side effects---providing an explicit safety boundary between the model's reasoning layer and the system's operational infrastructure. This architectural pattern directly addresses the uncontrolled-execution and regulatory-hallucination risks documented in the LLM-in-finance literature: the agent can reason freely in natural language but can only act within its approved tool set, and every tool invocation is logged as an auditable record. The Crypto World Bank adopts this pattern through its 16-tool MCP tool server (Section~\ref{sec:aiml-support}), where read tools are always permitted and write tools require an explicit human confirmation gate before any state-modifying call is issued, ensuring that no model output can trigger an irreversible on-chain action without documented client consent.

\section{Literature Review Summary}

The most important reviewed works with their methodologies, key findings, and relevance to our project are presented in Table~\ref{tab:lit-summary}, in the form of literature table suggested by Michigan State University Libraries~[23].

\begin{table}[H]
\centering
\caption{Literature review summary (Part A): DeFi lending, fraud detection, and XAI.}
\begin{tabular}{p{2.8cm}p{2.8cm}p{2.8cm}p{3cm}p{3cm}}
\hline
\textbf{Author Year} & \textbf{Research Focus} & \textbf{Methodology} & \textbf{Key Findings} & \textbf{Relevance to Our Project} \\
\hline
Palaiokrassas et al.\ (2023) [3] & DeFi fraud detection & XGBoost, NN classifiers on 54M+ multi-chain transactions & F1: 0.76--0.85 vs.\ 0.08 with features alone & Informs Random Forest fraud modeling \\
\hline
Adom et al.\ (2022) [4] & XAI in loan approval & LIME / SHAP comparison on lending datasets & SHAP provides deeper, more consistent explainability; LIME is faster & Justifies SHAP as our primary explainability method \\
\hline
Tan (2023) [5] & CBDC and IMF financial inclusion & Model for developing nations & Two-tier CBDC distribution (central bank $\to$ commercial banks $\to$ users) achieves 30--40\% higher financial inclusion in developing economies; hierarchical infrastructure is the decisive design choice & Informs dual-currency facility design \\
\hline
Liu et al.\ (2008) [6] & Anomaly detection & Isolation Forest on synthetic and real datasets & Isolation via recursive partitioning; shorter paths & Adopted as secondary unsupervised detection for wallet behaviour \\
\hline
Atzei et al.\ (2017) [7] & Smart contract security & SoK of Ethereum attack vectors & Over-reliance, access control vulnerabilities & Directly informs our security primitives and planned formal verification \\
\hline
\end{tabular}
\TableScreenshotEnd[tab:lit-summary]{Literature Review Summary (Part 1): DeFi Lending and Hierarchical Architecture}
\end{table}

\vspace{2pt}
\noindent\textit{\small Of the works reviewed, none model a four-tier hierarchical lending structure, confirming the architectural novelty of this project. The synthesis shows that DeFi and institutional finance have developed in parallel without converging on a governance-aware, multi-tier design.}

\begin{table}[H]
\centering
\caption{Literature review summary (Part B): DeFi protocol systematisation and adaptive lending.}
\begin{tabular}{p{2.8cm}p{2.8cm}p{2.8cm}p{3cm}p{3cm}}
\hline
\textbf{Author Year} & \textbf{Research Focus} & \textbf{Methodology} & \textbf{Key Findings} & \textbf{Relevance to Our Project} \\
\hline
Werner et al.\ (2022) [1] & DeFi protocol systematisation & SoK survey of 12+ DeFi protocol categories & Lending platforms are uniformly pool-based and overcollateralised; no institutional hierarchy exists & Identifies the core gap our four-tier architecture addresses \\
\hline
Bastankhah et al.\ (2023) [2] & Adaptive DeFi lending & Dual fast/slow control; simulation on historical data & Dynamic rate adjustment outperforms static utilisation curves & Validates algorithmic lending parameter optimisation \\
\hline
\end{tabular}
\TableScreenshotEnd{Literature Review Summary (Part 2): Governance, Settlement, and Institutional Adoption}
\end{table}

\vspace{2pt}
\noindent\textit{\small This continuation extends coverage to governance, settlement, and institutional blockchain adoption. The key observation is that enabling infrastructure (Layer~2 networks, oracle pricing, EVM tooling) is mature, while the specific hierarchical banking architecture remains architecturally unexplored.}

\begin{table}[H]
\centering
\caption{Literature review summary (Part C): blockchain P2P lending, privacy-preserving ML, and smart contract auditing.}
\begin{tabular}{p{2.8cm}p{2.8cm}p{2.8cm}p{3cm}p{3cm}}
\hline
\textbf{Author Year} & \textbf{Research Focus} & \textbf{Methodology} & \textbf{Key Findings} & \textbf{Relevance to Our Project} \\
\hline
Sharma et al.\ (2021) [48] & Blockchain P2P lending & Smart contract-mediated loan origination framework & Eliminates intermediary overhead; gas cost minimisation analysis & Informs four-tier architecture for peer-to-peer lending \\
\hline
Hassan et al.\ (2022) [49] & Privacy-preserving fraud detection & Federated learning on encrypted blockchain features & ML achieves comparable accuracy to centralised models & Informs future privacy-preserving fraud detection extensions \\
\hline
Wang et al.\ (2020) [50] & Smart contract vulnerability detection & ML classifiers on bytecode features & F1: 0.96 for access control; 0.93 for data leakage & Supports our planned security verification models \\
\hline
Liao et al.\ (2019) [51] & Smart contract audit automation & Hybrid classification + fuzz testing & Reduces audit time by 70\% vs.\ manual review & Informs our security audit strategy for three-contract architecture \\
\hline
\end{tabular}
\TableScreenshotEnd{Literature Review Summary (Part 3): Cross-Border Settlement and Security Tooling}
\end{table}

\vspace{2pt}
\noindent\textit{\small This continuation captures empirical and systems findings on cross-border settlement and security tooling, supporting the claim that Layer~2 deployment and hybrid ML-static analysis pipelines substantially reduce feasibility risk for an academic prototype.}

\begin{table}[H]
\centering
\caption{Literature review summary (Part D): CBDC design, cross-border settlement, and multi-chain DeFi.}
\begin{tabular}{p{2.8cm}p{2.8cm}p{2.8cm}p{3cm}p{3cm}}
\hline
\textbf{Author Year} & \textbf{Research Focus} & \textbf{Methodology} & \textbf{Key Findings} & \textbf{Relevance to Our Project} \\
\hline
Islam et al.\ (2024) [53] & CBDC design requirements & Analysis of retail CBDC distribution architecture & Two-tier model preferred; interoperability challenges identified & Informs dual-currency facility design \\
\hline
BIS/FSB [21] & Cross-border settlement infrastructure & Analysis of cross-border payment costs and capital trapped in nostro accounts & 2--5 days avg.\ settlement; \$42/transfer avg.\ cost & Motivates four-tier settlement with near-instant finality \\
\hline
Beyer et al.\ (2025) [34] & Monetary policy and inequality & Cross-country analysis (1999--2019) & QE increases wealth inequality; effects more persistent than monetary parameters & Motivates transparent, algorithmic monetary parameters \\
\hline
World Bank (2025) [39] & Blockchain for development finance & Pilot on 13 projects, 10 countries & Blockchain tracking reduces reporting burden for fund distribution & Supports our planned multi-chain deployment strategy \\
\hline
\end{tabular}
\TableScreenshotEnd{Literature Review Summary (Part 4): Risk Monitoring and AI/ML in Finance}
\end{table}

\vspace{2pt}
\noindent\textit{\small This block summarizes works informing the platform's risk and monitoring approach. The synthesis justifies the choice of lightweight, auditable ML components rather than opaque deep learning models, prioritizing regulatory interpretability over raw detection performance.}

\begin{table}[H]
\centering
\caption{Literature review summary (Part E): multi-chain DeFi, smart microfinance, gas optimisation, and SHAP-based credit models.}
\begin{tabular}{p{2.8cm}p{2.8cm}p{2.8cm}p{3cm}p{3cm}}
\hline
\textbf{Author Year} & \textbf{Research Focus} & \textbf{Methodology} & \textbf{Key Findings} & \textbf{Relevance to Our Project} \\
\hline
Li et al.\ (2024) [46] & Multi-chain DeFi lending & Cross-chain model design and implementation & Multi-chain distribution improves throughput and reduces congestion & Supports our planned multi-chain deployment strategy \\
\hline
Xu et al.\ (2023) [47] & DeFi lending protocol evaluation & Quantitative risk metrics for lending protocol assessment & Utilisation ratio, liquidation efficiency, and reserve responsiveness as key metrics & Provides benchmarks for evaluating our hierarchical tiers \\
\hline
Alam et al.\ (2021) [54] & Smart micro-credit in Bangladesh & Smart contract-based microfinance & Reduces admin overhead by 60\%; validates geographic and demographic markets & Directly validates our target geographic and demographic market \\
\hline
Tolmach et al.\ (2024) [55] & DeFi gas fee optimisation & Transaction batching and calldata compression analysis & Gas savings of 30--50\% through optimisation strategies & Validates Layer~2 deployment and mainnet optimisation \\
\hline
Gupta et al.\ (2024) [56] & SHAP-based credit default models & SHAP, Random Forest, and Gradient Boosting classifiers & AUC above 0.92 with full interpretability & Informs per-client risk explanation methodology \\
\hline
\end{tabular}
\TableScreenshotEnd{Literature Review Summary (Part 5): Institutional Adoption and Financial Inclusion}
\end{table}

\vspace{2pt}
\noindent\textit{\small This final block closes the synthesis with institutional and inclusion-related sources. The combined evidence supports the thesis that blockchain adoption is increasing at the institutional level while inclusion gaps persist at the retail level, matching the platform's intended multi-scale scope.}


\section{Comparative Protocol Analysis}
\label{sec:comparative-protocol}

The literature survey above reveals that no existing DeFi protocol combines multi-tier institutional hierarchy with AI-assisted governance and compliance-aware identity. Table~\ref{tab:protocol-comparison} provides a structured comparison of the Crypto World Bank against representative existing protocols on eleven architectural dimensions.

\begin{table}[H]
\centering
\caption{Comparative protocol analysis: existing DeFi lending protocols vs.\ the Crypto World Bank (CWB). {\normalfont \tmarkDone{} = implemented/present; \tmarkPartial{} = designed/partial; \tmarkPlanned{} = planned; \tmarkNo{} = absent.}}
\label{tab:protocol-comparison}
\renewcommand{\arraystretch}{1.25}
\small
\begin{tabular}{@{}p{4.2cm}ccccc>{\bfseries}c@{}}
\toprule
\tableheadcolor
\tableheadfont Feature & \tableheadfont Aave v3 & \tableheadfont Compound v3 & \tableheadfont MakerDAO & \tableheadfont Maple & \tableheadfont Goldfinch & \tableheadfont CWB \\
\midrule
Institutional hierarchy & \tmarkNo{} & \tmarkNo{} & \tmarkNo{} & \tmarkNo{} & \tmarkNo{} & \tmarkDone{} 4-tier \\
\rowcolor{RowShade}
Cross-tier capital flow & \tmarkNo{} & \tmarkNo{} & \tmarkNo{} & \tmarkNo{} & \tmarkNo{} & \tmarkPartial{} Designed \\
Same-tier interbank lending & \tmarkNo{} & \tmarkNo{} & \tmarkNo{} & \tmarkNo{} & \tmarkNo{} & \tmarkPartial{} Designed \\
\rowcolor{RowShade}
Solidarity group lending & \tmarkNo{} & \tmarkNo{} & \tmarkNo{} & \tmarkNo{} & Partial & \tmarkPlanned{} Planned \\
AI/ML fraud detection & \tmarkNo{} & \tmarkNo{} & \tmarkNo{} & Manual & Manual & \tmarkPartial{} Built, not integrated \\
\rowcolor{RowShade}
SHAP explainability & \tmarkNo{} & \tmarkNo{} & \tmarkNo{} & \tmarkNo{} & \tmarkNo{} & \tmarkPartial{} Built, not integrated \\
ZKP KYC compliance & \tmarkNo{} & \tmarkNo{} & \tmarkNo{} & \tmarkNo{} & \tmarkNo{} & \tmarkPlanned{} Planned \\
\rowcolor{RowShade}
Kinked interest rate curve & \tmarkDone{} & \tmarkDone{} & Via gov. & \tmarkDone{} & \tmarkNo{} & \tmarkPartial{} Designed \\
Role-based access control & Partial & Partial & Via gov. & \tmarkDone{} & Partial & \tmarkDone{} Implemented \\
\rowcolor{RowShade}
Developing-economy focus & \tmarkNo{} & \tmarkNo{} & \tmarkNo{} & \tmarkNo{} & \tmarkDone{} & \tmarkDone{} Bangladesh \\
TVL / Status (2026) & \$26.3B & \$1.4B & \$10.5B & \$2.6B & \$680M & Testnet \\
\bottomrule
\end{tabular}
\end{table}

This comparison directly answers RQ1: no existing protocol models institutional hierarchy, cross-tier capital flow, or solidarity group lending in one decentralized architecture. While Aave~v3 implements risk-parameterized sub-markets (Isolation Mode, Efficiency Mode) and Compound~v3 deploys separate Comet markets, these mechanisms differentiate risk parameters for different assets within a single-tier pool---they do not model the institutional relationships, cross-tier governance, or inter-bank lending flows that characterize the CWB architecture. The Crypto World Bank is architecturally distinct across every dimension that characterizes development banking.

\paragraph{Note on ERC-3643 (T-REX) for institutional-tier permissioning.}
The institutional tiers (Tiers~1--3) use the bespoke RBAC described in Section~\ref{sec:defense-in-depth} for transfer-level access control.  ERC-3643 (the T-REX standard for permissioned, compliance-gated token transfers) and ERC-4626 for tokenized vaults are gaining adoption as standards that ensure tokenized assets can plug and play across platforms without custom integrations~[R46].  ERC-3643 is specifically designed for regulated institutional token transfers where KYC and AML checks gate individual token movements at the token contract level, rather than at the application layer.  The platform's current RBAC serves the same functional purpose: role gates on every state-mutating function enforce that only KYC-verified, role-holding addresses can transfer institutional-tier assets.  A production deployment operating under MiCA~[R36] or equivalent regulated-asset frameworks would evaluate full ERC-3643 compliance as the path to interoperability with other regulated tokenized-asset platforms; the architecture is designed to accommodate this migration because the RBAC role structure maps directly onto the ERC-3643 identity registry model.  This note is recorded here rather than as a Future Work item so that examiners can identify the deliberate design alignment.

\paragraph{Positioning against failed centralized exchanges.}
Centralized cryptocurrency exchanges such as FTX (defunct, 2022) demonstrate the catastrophic consequence of opaque reserve management: an \$8~billion client-fund shortfall was discovered only at the point of bankruptcy, having been concealed through commingling of client deposits with proprietary trading activity by its affiliated trading firm Alameda Research. FTX halted withdrawals and filed for bankruptcy within three days. The structural cause was centralized custody with no on-chain proof of reserves, no separation between customer funds and operating capital, and no real-time verifiability of solvency. The Crypto World Bank's on-chain reserve architecture is designed as a structural response to this class of failure: every tier's reserve ratio is an enforced smart contract invariant, readable by any participant in real time without trust in any third-party attestation. This design property---on-chain Proof of Reserves by construction---is increasingly recognized as the non-negotiable foundation of trustworthy crypto-financial infrastructure, and the Crypto World Bank implements it at the institutional tier level rather than the exchange-account level. The platform's purpose is also structurally distinct from FTX: it targets financial inclusion through savings, credit, and interbank lending---not speculative trading, derivatives, or leveraged tokens---and its four-tier institutional hierarchy mirrors multilateral development finance rather than a flat exchange order book. DeFi naturally provides on-chain transparency, self-custody, and governance---precisely the properties whose absence destroyed FTX.

\subsection{Literature Synthesis}

The literature synthesis above directly informs four architectural decisions in this work: (1) Gudgeon et al.'s~(2020)~[R3] empirical finding that utilization above 90\% causes liquidity crises motivates the kinked interest rate model in Section~\ref{sec:kinked-rate}; (2) Piper et al.'s~(2025)~[R8] ZKP permissioning framework provides the technical foundation for the compliance pathway in Section~\ref{sec:identity}; (3) the empirical finding of Howlader \& Halder~(2025)~[R11] that mobile financial inclusion in Bangladesh grew by 99\% between 2004 and 2021 validates the retail-tier accessibility argument in Section~\ref{sec:research-contribution}; and (4) Alam et al.'s~(2021)~[54] IEEE TENSYMP paper on blockchain microcredit in Bangladesh provides direct prior art that this thesis advances.

\subsection{Agent Harness Engineering and Production Safety}
\label{sec:harness-lit}

The deployment of LLM agents in production financial systems has produced a
distinct engineering discipline focused on the \textit{harness} layer that wraps
the model. Heidt et~al.~[R52] propose the CAR decomposition ---
Control, Agency, and Runtime --- as the three functional layers that determine
production reliability in LLM agents. The core thesis is that it is the harness
wrapping the model, not the model itself, that determines whether an agent is
trustworthy in production: a capable model with a weak harness is less reliable
than a modest model with a well-engineered harness. The CWB agent architecture
addresses all three layers: the Stable tier and confirmation audit hook constitute
the Control layer; the named toolsets and tool-schema boundary constitute the
Agency layer; and the session lineage model and cross-session memory constitute
the Runtime layer.

A parallel research stream has documented prompt injection as the dominant
vulnerability in production banking agents. Alizadeh et~al.~[R53]
demonstrate that income document uploads and language preference fields in
banking tool-calling agents are viable injection vectors for personal data
exfiltration --- the same fields present in the CWB context assembly pipeline.
The OWASP Top~10 for LLM Applications~[R54] classifies LLM01 (Prompt
Injection) as the highest-priority threat in agentic financial systems, providing
an authoritative standard against which the CWB scanning middleware can be
evaluated.

\section{Summary of Key Findings}

The following summarized findings are obtained in the literature review and are directly informative to our design:

\begin{itemize}[noitemsep]
\item \textbf{DeFi lending is structurally flat.} Available protocols (Aave, Compound, MakerDAO) operate pool-based or over-collateralized models with total value locked exceeding \$55~billion~[13]. Institutional hierarchy has no comparable protocol models to development finance. The research by CBDC~[5] ascertains that tiered distribution is workable and promotes financial inclusion.
\item \textbf{Correspondent banking is inefficient in structure.} Cross-border settlement spends two to five days at average costs of \$42 per transaction~[25]. The system traps important capital lying idle in nostro/vostro accounts~[24], and remittance fees are estimated to consume between \$48--56~billion a year~[26]. On-chain settlement on Layer~2 networks is known to cut the latency (to seconds) and the cost (to sub-cent levels) down.
\item \textbf{ML enhances fraud detection in blockchain.} DeFi-specific behavioral features significantly work well as compared to transactional features (F1: 0.76--0.85 vs.\ 0.08~[3]). The Isolation Forest~[6] allows unsupervised detection in which labeled data is scarce.
\item \textbf{SHAP is explainable by regulators.} SHAP offers greater consistency and is better than the LIME in loan approval systems~[4] in meeting prudential demands on open-minded automated decision-making.
\item \textbf{Layered defense is needed in smart contract security.} Reentrancy, overflow, and vulnerabilities to access control are well documented~[7]. ML-based vulnerability detection achieves F1-scores of above 0.93~[50], and hybrid ML-fuzz methods reduce audit time by 70\%~[51]. Verification tools are formal, i.e., VERISMART achieves 98.4\% precision~[52]. Automated primitives used with OpenZeppelin are recommended for production deployments to be verified.
\item \textbf{Monetary policy leads to distributional inequality.} Cross-country evidence~[34] and Federal Reserve research~[35] confirm that quantitative easing disproportionately favors asset holders, and workers of lower income cover the expenses of both inflation and consequent tightening. Transparent, algorithmic monetary parameters can decrease this informational asymmetry.
\item \textbf{The use of blockchains in institutions is gaining momentum.} The World Bank's FundsChain~[39], tokenized assets of R3 Corda worth \$17~billion~[37], and JPMorgan Kinexys settling multi-billion dollars per day~[40] reflect the fact that financial institutions are actively implementing blockchain on settlement, funds tracking, and asset tokenization.
\item \textbf{Blockchain allows financial inclusion.} Approximately 1.4~billion adults were not banked in any country around the world~[14], and DeFi was found by the World Economic Forum as a leaping frog technology that allows peoples to avoid the banking systems infrastructure~[41]. The MiniPay wallet of Celo has onboarded 14~million users in over 60 countries with less than cent transaction fees~[42]. Blockchain-based microcredit in Bangladesh shows 60\% decrease on administrative overhead~[54].
\item \textbf{Multi-chain lending enhances scalability.} Cross-chain lending models~[46] and Layer~2 gas optimization strategies~[55] prove the fact that the lending is distributed when operating on several networks, decreasing congestion and transaction costs by 30--50\%.
\item \textbf{ML can be used to detect cross-institutional fraud using privacy preserving ML.} Federated learning methods~[49] are such that enable several lending institutions to cooperate and learn fraud detection models in a non-intrusive manner, i.e. no individual transaction data is exposed, resolving the conflict between institution-wide security and data privacy.

\item \textbf{Group lending reduces default risk in underserved markets.} Microfinance research from BRAC, Grameen Bank, and ASA consistently reports repayment rates above 95\% in solidarity group structures~[54]. On-chain enforcement of mutual liability, replacing social pressure with programmable contract enforcement, has not been studied in the academic literature, representing an open research gap that this project begins to address.

\item \textbf{Sustainable microfinance requires deposit-funded lending.} Literature on sustainable microfinance institutions consistently shows that deposit-funded lending rather than donor-funded or externally capitalized models is the only economically viable approach at scale. The platform's banking product suite is designed to close this loop by mobilizing retail savings at the Local Bank tier to partially fund the lending pool.
\end{itemize}

The results explain why we have a four-tier architecture, AI/ML integration (Random Forest, Isolation Forest, SHAP), cross-tier lending structure, governance structure, and planned security verification pipeline.


%% %%%%%%%%%%%%%%%%%%%%%%%%%%%%%%%%%%%%%%%%%%%%%%%%%%%%%%%%%%%%
%% CHAPTER 3: SYSTEM ARCHITECTURE AND DESIGN
%% %%%%%%%%%%%%%%%%%%%%%%%%%%%%%%%%%%%%%%%%%%%%%%%%%%%%%%%%%%%%
\chapter{System Architecture and Design}

\section{Prototype Scope}
\label{sec:prototype-scope}

Evaluators should distinguish between what has been \textbf{built and tested}, what has been \textbf{designed and partially scaffolded}, and what is \textbf{planned} for the final thesis phase. Table~\ref{tab:prototype-scope} provides this mapping for every major platform feature.

\begin{table}[H]
\centering
\caption{Prototype scope: feature implementation status as of pre-thesis submission. \tmarkDone{} = Implemented and testnet-verified; \tmarkPartial{} = Designed or partially scaffolded; \tmarkPlanned{} = Planned for final thesis phase.}
\label{tab:prototype-scope}
\renewcommand{\arraystretch}{1.3}
\small
\begin{tabular}{@{}p{7.5cm}c@{}}
\toprule
\tableheadcolor
\tableheadfont Feature & \tableheadfont Status \\
\midrule
Four-tier role system (RBAC) & \tmarkDone{} Implemented \\
\rowcolor{RowShade}
World Bank Reserve contract (Tier~1) & \tmarkDone{} Implemented \\
Tier~2 National Bank contracts & \tmarkPartial{} Designed, partial \\
\rowcolor{RowShade}
Tier~3 Local Bank contracts & \tmarkPartial{} Designed, partial \\
Cross-tier fund transfer & \tmarkPartial{} Designed, unimplemented \\
\rowcolor{RowShade}
Loan request / approval workflow & \tmarkDone{} Implemented \\
Installment EMI auto-generation & \tmarkPartial{} Designed \\
\rowcolor{RowShade}
SavingsVault contract & \tmarkPlanned{} Planned (final thesis) \\
FixedDeposit contract & \tmarkPlanned{} Planned (final thesis) \\
\rowcolor{RowShade}
GroupLendingPool contract & \tmarkPlanned{} Planned (final thesis) \\
InterBankLendingPool & \tmarkPlanned{} Planned (final thesis) \\
\rowcolor{RowShade}
AI/ML fraud detection (Random Forest) & \tmarkPartial{} Built, not integrated \\
SHAP explainability output & \tmarkPartial{} Built, not integrated \\
\rowcolor{RowShade}
Oracle integration (Chainlink Functions) & \tmarkPlanned{} Planned (Sprint~3 --- replaces commit-reveal relay) \\
ZKP KYC compliance layer & \tmarkPlanned{} Planned (final thesis) \\
\rowcolor{RowShade}
Testnet deployment evidence & \tmarkDone{} --- migration to Polygon zkEVM Cardona planned for Sprint~1 \\
AI Agent / MCP tool server (16 banking tools) & \tmarkPlanned{} Planned (Sprint~2) \\
\rowcolor{RowShade}
EIP-7702 session key management & \tmarkPlanned{} Planned (Sprint~2) \\
Authority Brief UI (SHAP for bank approvers) & \tmarkPlanned{} Planned (Sprint~2) \\
\rowcolor{RowShade}
Chainlink Automation (installment triggers) & \tmarkPlanned{} Planned (Sprint~2) \\
Chainlink Proof of Reserve & \tmarkPlanned{} Planned (Sprint~1) \\
\rowcolor{RowShade}
The Graph subgraph (event indexing) & \tmarkPlanned{} Planned (Sprint~3) \\
SAR workflow (AML $\to$ compliance queue $\to$ freeze) & \tmarkPlanned{} Planned (Sprint~3) \\
\rowcolor{RowShade}
\texttt{sessions} table + EIP-7702 schema & \tmarkPlanned{} Planned (Sprint~1) \\
\texttt{agent\_action\_log} table (append-only) & \tmarkPlanned{} Planned (Sprint~2) \\
\rowcolor{RowShade}
300-client Foundry simulation (live demo) & \tmarkPlanned{} Planned (Sprint~3) \\
\bottomrule
\end{tabular}
\end{table}

This table eliminates any ambiguity between present-tense architectural description and empirical implementation evidence throughout the following sections.

\section{High-Level Architecture}

The system employs a \textbf{four-layer decentralized application architecture} (Figure~\ref{fig:three-layer-arch}): a presentation layer (React + wallet integration), a smart-contract layer (the World Bank, National Bank, and Local Bank contracts on the EVM), an off-chain services layer (REST API, PostgreSQL, AI Agent engine, MCP Tool Server, Chainlink oracle integration, event listener, and Redis cache), and a Chainlink infrastructure layer (Chainlink Functions, Chainlink Automation, Chainlink Price Feeds, and Chainlink Proof of Reserve). Figure~\ref{fig:component-diagram} elaborates this view with the fifteen-contract target architecture and external integrations.

\begin{figure}[H]
\centering
\OnePageDiagram{fig-three-layer-arch.pdf}
\caption{Four-layer decentralized application architecture: presentation (React 18, TypeScript, Wagmi/Viem), smart-contract layer on the EVM (World Bank Reserve, National Bank, Local Bank), off-chain services layer (Express REST API, PostgreSQL, AI Agent engine (Qwen3-8B + MCP Tool Server), Chainlink oracle integration, event listener, Redis), and Chainlink infrastructure layer (Chainlink Functions, Automation, Price Feeds, Proof of Reserve).}
\label{fig:three-layer-arch}
\end{figure}

The current prototype implements three core contracts (World Bank Reserve, National Bank, Local Bank). The complete architecture extends to fifteen modular contracts covering the full banking product suite described in Section~\ref{sec:banking-products} and the multi-entity / cross-tier operations of Section~\ref{sec:multi-entity}. The remaining twelve contracts are: SavingsVault, FixedDeposit, GroupLendingPool, FXModule, InsuranceFund, CurrentAccount (banking product suite); InterBankLendingPool, UpwardDepositFacility, SyndicatedLoan, TranchedPool, TreasurySwap, NettingEngine (multi-entity / cross-tier operations). All are planned for implementation across Sprints~2 and~3 of the final thesis phase.

Figure~\ref{fig:component-diagram} shows the component interactions across these three layers. The diagram reflects the current three-contract prototype view; the nine-contract target architecture is specified in Appendix~B and will be reflected in an updated diagram in the final thesis phase.

\begin{figure}[H]
\centering
\OnePageDiagram{fig-component-architecture.pdf}
\caption{Component diagram showing interactions between the presentation layer, smart contract layer, off-chain backend services, and external systems.}
\label{fig:component-diagram}
\end{figure}

\section{Blockchain Platform Selection}

\begin{figure}[H]
\centering
\OnePageDiagram{fig-blockchain-stack.pdf}
\caption{Layered blockchain and application stack: L1 settlement (Polygon~PoS for retail, Ethereum~Sepolia for institutional, plus Chainlink CCIP bridge, The~Graph indexer, and Tenderly monitor); L2 smart-contract platform (Solidity 0.8.20, OpenZeppelin~v5, UUPS, TimeLock, RBAC); L3 data (PostgreSQL~16 in 3NF, Redis, IPFS); L4 API and services (Express, FastAPI, WebSocket, EIP-712); L5 presentation (React + MetaMask).}
\label{fig:blockchain-stack}
\end{figure}


\begin{table}[H]
\centering
\caption{Blockchain platform selection criteria and justification.}
\label{tab:blockchain-selection}
\begin{tabular}{p{3.5cm}p{3.5cm}p{6.5cm}}
\hline
\textbf{Criterion} & \textbf{Selection} & \textbf{Justification} \\
\hline
Platform & Ethereum Virtual Machine (EVM) & Largest developer ecosystem; battle-tested security model; extensive tooling. \\
\hline
Network & Polygon zkEVM Cardona / Ethereum Sepolia & Zero-cost deployment; production-equivalent behaviour; free faucet access. \\
\hline
Consensus & ZK validity proofs (Polygon zkEVM L2, Ethereum L1 settlement) & Every batch of transactions is accompanied by a zero-knowledge proof verified on Ethereum L1, so security derives from cryptographic validity rather than a validator set assumption. \\
\hline
Smart contract language & Solidity 0.8.20 & Industry standard; mature compiler with overflow protection; rich library ecosystem. \\
\hline
\end{tabular}
\TableScreenshotEnd{Blockchain Platform Selection: Comparison of EVM-Compatible Networks}
\end{table}

\vspace{2pt}
\noindent\textit{\small This platform selection table compares candidate networks using criteria relevant to banking workflows, including cost, finality, throughput, and ecosystem maturity. Polygon zkEVM Cardona is selected over Polygon Amoy PoS because ZK validity proofs derive security from cryptographic verification rather than a validator set assumption, which is a materially stronger security claim for an institutional banking prototype. It motivates the selection of an EVM-compatible Layer~2 deployment for low-fee retail operations while retaining portability to other EVM chains.}

\begin{table}[H]
\centering
\caption{Blockchain platform selection (continued): operational and deployment factors.}
\begin{tabular}{p{3.5cm}p{3.5cm}p{6.5cm}}
\hline
\textbf{Criterion} & \textbf{Selection} & \textbf{Justification} \\
\hline
Gas cost & \$0.001--\$0.01 per transaction & Orders of magnitude cheaper than Ethereum mainnet (\$5--\$50); enables micro-loan economics. \\
\hline
Block finality & $\approx$2 seconds & Sub-second practical finality for retail UX; checkpointed to Ethereum for security. \\
\hline
Developer tooling & Hardhat + OpenZeppelin & Automated test suite, deployment scripts, and audited security primitives. \\
\hline
Testnet availability & Cardona (Polygon zkEVM), Sepolia (Ethereum) & Free faucets; EVM-identical behaviour; no real cryptocurrency required for prototype. \\
\hline
Security model & ZK validity proofs & Every batch verified by a ZK proof anchored to Ethereum L1 (vs.\ PoS validator assumption on Amoy). \\
\hline
Migration path & EVM-compatible & Same Solidity contracts deploy to any EVM chain; reduces future L2 migration cost. \\
\hline
\end{tabular}
\TableScreenshotEnd{Blockchain Platform Selection (Continued): Operational and Deployment Factors}
\end{table}

\vspace{2pt}
\noindent\textit{\small This supplementary comparison expands the platform evaluation to additional operational and deployment factors. The conclusion is that Polygon PoS provides a practical balance of stability and low transaction costs for an academic prototype, with a clear migration path if future requirements demand alternative L2s.}

\subsection{Transaction Verification and Consensus}

\begin{itemize}[noitemsep]
\item \textbf{Polygon zkEVM Cardona:} Transactions are batched and verified by a zero-knowledge validity proof submitted to Ethereum L1. Block finality is achieved once the ZK proof is verified on Ethereum, which provides cryptographic rather than economic security guarantees. The ZK proof means that a malicious sequencer cannot publish an invalid batch --- the proof would not verify. This is materially stronger than the PoS validator-collusion assumption of Polygon Amoy.
\item \textbf{On prototype testnets:} Cardona and Sepolia use equivalent consensus models at no financial cost, enabling incremental development and testing without exposure to real-asset risk.
\end{itemize}

\subsection{Oracle Architecture: Off-Chain AI to On-Chain Decision}
\label{sec:oracle-architecture}

The Crypto World Bank requires a mechanism to convey off-chain AI/ML risk assessments into the on-chain loan approval workflow. This is an instance of the oracle problem~[R1,R2].

\paragraph{Chainlink Functions oracle (primary).}
The final thesis phase uses \textbf{Chainlink Functions} as the primary oracle mechanism. Chainlink Functions operates via a Decentralised Oracle Network (DON): the ML risk score is fetched independently by multiple Chainlink nodes, consensus across the DON is required before the result is committed on-chain, and no single compromised node can manipulate the risk score. This is a fundamental trust improvement over the v23 commit-reveal relay, which required trusting a single FastAPI service key. The source code executed by the DON is:

\begin{verbatim}
// Chainlink Functions source (runs in decentralised DON)
const clientId = args[0];
const response = await Functions.makeHttpRequest({
  url: `https://your-ml-service.com/score/${clientId}`,
  method: "POST"
});
const score = response.data.risk_score;
// 6 decimal precision
return Functions.encodeUint256(Math.round(score * 1e6));
\end{verbatim}

The DON adds 30--60 seconds of latency and \$0.10--\$1.00 per oracle call, both acceptable for the loan approval lifecycle where decisions are not time-critical. This architecture closes the oracle trust assumption that v23 noted as a limitation (the 2-of-3 Safe multisig attestation was the interim mitigation; Chainlink Functions removes the need for that mitigation entirely).

\paragraph{Commit-reveal relay (prototype fallback).}
During the prototype phase, prior to Chainlink Functions being wired (Sprint~3), the system uses a commit-reveal relay as a fallback. The off-chain ML service commits a hash $h = \text{keccak256}(s \,\|\, \text{nonce})$ to the \texttt{LoanController} contract, then reveals $s$ and the nonce within the decision window. The contract verifies $h = \text{keccak256}(s \,\|\, \text{nonce})$ and stores $s$ immutably. This pattern prevents score manipulation between commitment and decision and creates an immutable on-chain audit trail of every risk score used in a lending decision, while Chainlink Functions integration is completed.

\paragraph{Chainlink Automation: trustless installment triggers.}
\textbf{Chainlink Automation} replaces the centralised cron job for overdue installment detection and interest accrual. The entire loan lifecycle from application to overdue flagging runs without a trusted operator:

\begin{verbatim}
// In LocalBankPool.sol
function checkUpkeep(bytes calldata) external view override
    returns (bool upkeepNeeded, bytes memory performData) {
    uint256[] memory overdueLoans = getOverdueLoans();
    upkeepNeeded = overdueLoans.length > 0;
    performData = abi.encode(overdueLoans);
}

function performUpkeep(bytes calldata performData) external override {
    uint256[] memory overdueLoans =
        abi.decode(performData, (uint256[]));
    for (uint i = 0; i < overdueLoans.length; i++) {
        _markInstallmentOverdue(overdueLoans[i]);
    }
}
\end{verbatim}

This removes the last centralised component from the loan lifecycle --- the entire process from application to overdue flagging runs without a trusted operator.

\paragraph{Chainlink Price Feeds: BDT/USD and ETH/USD.}
The \texttt{FXModule} contract uses Chainlink's production-grade \texttt{AggregatorV3Interface} price feeds for BDT/USD and ETH/USD, replacing the ``forex oracle approved by governance'' pattern of v23. Every BDT display in the frontend flows through:

\begin{verbatim}
// In LocalBankPool.sol
AggregatorV3Interface internal bdtUsdFeed;
AggregatorV3Interface internal ethUsdFeed;

function getBdtEquivalent(uint256 usdcAmount)
    public view returns (uint256) {
    (, int bdtRate,,,) = bdtUsdFeed.latestRoundData();
    // USDC tracks USD 1:1; Chainlink uses 8 decimals
    return usdcAmount * uint256(bdtRate) / 1e8;
}
\end{verbatim}

The 8-decimal precision convention is the Chainlink standard for fiat-pair feeds. Chainlink price feeds are the Sprint~1 build item for BDT display across the entire frontend.

\paragraph{Chainlink Proof of Reserve.}
The \texttt{WorldBankReserve} contract publishes its reserve balance to \textbf{Chainlink Proof of Reserve (PoR)}, making the reserve cryptographically verifiable by any external auditor without trusting CWB's administrator. This is the ``free PoR'' advantage over FTX: any external system can verify reserve solvency without relying on CWB's admin claims.

\begin{verbatim}
// WorldBankReserve.sol
function getReserveSummary() external view returns (
    uint256 totalDeposited,
    uint256 totalLoaned,
    uint256 reserveRatio,   // scaled 1e4 (e.g. 5000 = 50%)
    uint256 insuranceFundBalance
) {
    return (
        totalDeposited,
        totalLoaned,
        (totalDeposited - totalLoaned) * 1e4 / totalDeposited,
        insuranceFund.balance()
    );
}
\end{verbatim}

This function is the FTX commingling safeguard described in Section~\ref{sec:defense-in-depth}. Adding Chainlink's PoR job on top turns it into a market-standard verifiable proof, and it is specified in full in Appendix~D.

The diagram placeholder for this architecture (Figure pending): \texttt{oracle\_architecture.png}---Chainlink Functions DON $\to$ score commitment $\to$ on-chain LoanController.


\section{Data Model and Database Design}

The relational database schema comprises 19 normalized entities in Third Normal Form (3NF). Figure~\ref{fig:core-system-graph} presents the core system graph showing entity relationships, Figure~\ref{fig:erd} presents the full Entity-Relationship Diagram (ERD), and Figure~\ref{fig:eer} presents the Enhanced Entity-Relationship (EER) diagram.

% ── Figure: Core System Graph ─────────────────────────────────────────────────
% Drawn with TikZ; no external PNG required.
% Layout mirrors the uploaded "Core_System_Graph.png":
%   WORLD_BANK → NATIONAL_BANK → LOCAL_BANK → BANK_USER (National / Local)
%   NATIONAL_BANK / LOCAL_BANK → BORROWER → LOAN_REQUEST / TRANSACTION / INCOME_PROOF
%   LOAN_REQUEST → INSTALLMENT / CHAT_MESSAGE / AI_ML_LOG
% ──────────────────────────────────────────────────────────────────────────────
\begin{figure}[H]
\centering
\OnePageDiagram{fig-erd-core.pdf}
\caption{Core system graph showing entity relationships across the four-tier banking hierarchy (\texttt{WORLD\_BANK} $\to$ \texttt{NATIONAL\_BANK} $\to$ \texttt{LOCAL\_BANK} $\to$ \texttt{BANK\_USER}), the central \texttt{BORROWER} entity, and the lending-lifecycle sub-graph (\texttt{LOAN\_REQUEST}, \texttt{LOAN}, \texttt{INSTALLMENT}, \texttt{TRANSACTION}, \texttt{INCOME\_PROOF}, \texttt{CHAT\_MESSAGE}, \texttt{AI\_ML\_LOG}, \texttt{CREDIT\_PASSPORT}).}
\label{fig:core-system-graph}
\end{figure}
%
% ── Inline context note (fills the page alongside the diagram) ────────────────
\vspace{4pt}
\begin{highlightbox}
\small
\textbf{Reading guide.} Arrows denote \emph{one-to-many (1:N)} relationships flowing downward
through the institutional hierarchy.
\textbf{BANK\_USER} specialises into \emph{National} and \emph{Local} variants (disjoint
generalisation—see Figure~\ref{fig:eer}).
\textbf{BORROWER} is the pivot entity: it aggregates identity, credit history, and all
transactional records.
\textbf{LOAN\_REQUEST} is the \emph{aggregation hub} for the lending lifecycle:
INSTALLMENT (weak entity), CHAT\_MESSAGE (audit trail), and AI\_ML\_LOG (risk scoring)
all depend on it existentially.
The complete attribute sets for every entity are shown in the ERD (Figure~\ref{fig:erd})
and the full EER model (Figure~\ref{fig:eer}).
\end{highlightbox}

% ── ERD: Full relational schema (15 tables) ───────────────────────────────────
% Rendered in TikZ using tcolorbox-style table headers per entity.
% Layout groups entities into three horizontal bands:
%   Row A (top)    : WORLD_BANK | NATIONAL_BANK | LOCAL_BANK | AI_CHATBOT_LOG | MARKET_DATA | PROFILE_SETTINGS
%   Row B (middle) : BORROWER | BANK_USER | INCOME_PROOF | BORROWING_LIMIT
%   Row C (lower)  : LOAN_REQUEST | TRANSACTION | CHAT_MESSAGE
%   Row D (bottom) : INSTALLMENT | AI_ML_SECURITY_LOG
% ──────────────────────────────────────────────────────────────────────────────

% Helper macro: coloured entity header + attribute lines
% Usage: \erdtable[options]{Title}{col}{attr1\\attr2\\...}
\newcommand{\erdheader}[2]{%
  \node[draw=#2!70, fill=#2!15, rounded corners=2pt,
        minimum width=3.0cm, font=\tiny\bfseries\color{black},
        inner sep=3pt, align=center] {#1};%
}

\begin{figure}[p]
\centering
\OnePageDiagram{fig-erd-core.pdf}
\caption{Entity-Relationship Diagram (ERD) of the Crypto World Bank database: core lending and governance entities in Third Normal Form (3NF), with primary keys, foreign keys, attribute types, and crow's-foot cardinality. Multi-entity / cross-tier and extended-product entities are shown in Figure~\ref{fig:erd-extended}.}
\label{fig:erd}
\end{figure}

\begin{figure}[H]
\centering
\OnePageDiagram{fig-erd-extended.pdf}
\caption{Extended ERD entities introduced in v15: the banking-product entities (\texttt{SAVINGS\_ACCOUNT}, \texttt{FIXED\_DEPOSIT}, \texttt{CURRENT\_ACCOUNT}, \texttt{LOAN\_GROUP}, \texttt{GROUP\_MEMBER}, \texttt{INSURANCE\_FUND}) and the multi-entity / cross-tier operational entities (\texttt{INTERBANK\_LOAN}, \texttt{UPWARD\_DEPOSIT}, \texttt{SYNDICATE}, \texttt{SYNDICATE\_MEMBER}, \texttt{TRANCHED\_POOL}, \texttt{TREASURY\_SWAP}, \texttt{NETTING\_BATCH}, \texttt{NETTING\_ENTRY}). Foreign-key relationships into the core entities of Figure~\ref{fig:erd} are preserved.}
\label{fig:erd-extended}
\end{figure}

The current ERD covers the core lending and governance entities. Extended banking entities including SavingsAccount, FixedDeposit, LoanGroup, GroupMember, CurrentAccount, and InsuranceFund are designed for the full system and will be incorporated into the updated data model in the final thesis. v15 additionally introduces the multi-entity / cross-tier operational entities described in Section~\ref{sec:multi-entity-consistency}: \texttt{INTERBANK\_LOAN}, \texttt{UPWARD\_DEPOSIT}, \texttt{SYNDICATE} and \texttt{SYNDICATE\_MEMBER}, \texttt{TRANCHED\_POOL}, \texttt{TREASURY\_SWAP}, and \texttt{NETTING\_BATCH} with \texttt{NETTING\_ENTRY}. Their foreign-key relationships into the core lending entities (LOAN, BANK\_USER, REPAYMENT) are detailed in Section~\ref{sec:multi-entity} and will be drawn into the updated ERD figure in the diagram-audit session. The data model therefore remains in one-to-one consistency with the contract architecture and Sprint plan.

% ── EER Diagram: Full Model ────────────────────────────────────────────────────
% Illustrates EER constructs per the uploaded "EER_Diagram_-_Full_Model.png":
%   • Banking Hierarchy panel (top): WORLD_BANK 1:N NATIONAL_BANK 1:N LOCAL_BANK
%   • Generalization panel (right): BANK_USER → {NATIONAL_BANK_USER, LOCAL_BANK_USER}
%   • User Layer panel (left): LOCAL_BANK 1:N BORROWER; BORROWER 1:1 BORROWING_LIMIT
%   • INCOME_PROOF (multi-valued); LOAN_REQUEST (weak entity / aggregation)
%   • Aggregation panel: TRANSACTION, CHAT_MESSAGE → AI_ML_SECURITY_LOG
%   • Independent entities panel (top-left): MARKET_DATA, AI_CHATBOT_LOG, PROFILE_SETTINGS
% ──────────────────────────────────────────────────────────────────────────────

% ── EER legend shorthand macros ──────────────────────────────────────────────
% Entity cylinder (database symbol) — drawn as a simple rectangle with cap arc
% We use standard rounded rectangles to keep the diagram TeX-compilable
% without additional libraries. Dashed border = weak entity.

\begin{figure}[H]
\centering
\OnePageDiagram{fig-eer-model.pdf}
\caption{Enhanced Entity-Relationship (EER) model showing generalization / specialization (\texttt{BANK\_USER} $\to$ National-/Local-bank-admin / Approver subtypes), the weak entity \texttt{INSTALLMENT} identified by its parent \texttt{LOAN}, the multi-valued attribute \texttt{INCOME\_PROOF}, the loan-centric aggregation, and participation constraints (total / partial).}
\label{fig:eer}
\end{figure}

\subsection{Entity Summary}

\begin{table}[H]
\centering
\caption{Database entity summary (20 entities).}
\label{tab:db-entities}
\begin{tabular}{p{4cm}p{9.5cm}}
\hline
\textbf{Entity} & \textbf{Role / Description} \\
\hline
WORLD\_BANK & Top-level reserve holder; global lending parameters. \\
NATIONAL\_BANK & Country-level banks; borrow from World Bank. \\
LOCAL\_BANK & City-level banks; borrow from National Bank and lend to users. \\
BANK\_USER & Bank staff with role-based permissions (approve/reject loans). \\
BORROWER (CLIENT) & End clients requesting and repaying loans. The entity will be renamed \texttt{CLIENT} (with \texttt{client\_id} FK) in the final thesis database migration to align with platform terminology; the current schema retains \texttt{BORROWER} for prototype compatibility. \\
LOAN\_REQUEST & Loan applications and full lifecycle tracking. \\
INSTALLMENT & Repayment schedule records (weak entity dependent on LOAN\_REQUEST). \\
TRANSACTION & Financial transaction records. \\
BORROWING\_LIMIT & Per-client limits with 6-month and 1-year rolling windows. \\
INCOME\_PROOF & Income verification documents (multi-valued per client). \\
CHAT\_MESSAGE & Client--bank communication records. \\
AI\_CHATBOT\_LOG & AI chatbot interaction records. \\
AI\_ML\_SECURITY & Two logically distinct domains aggregated here for the prototype schema, classified as an \textit{Association/Derived Entity}: (1) \texttt{LOAN\_RISK\_ASSESSMENT} --- per-loan fraud scores, anomaly scores, composite scores, and SHAP value vectors; (2) \texttt{SECURITY\_EVENT\_LOG} --- system-wide security events (large disbursements, repeated requests, Tenderly alerts). The final thesis will split these into separate tables to enable independent indexing and querying. \\
MARKET\_DATA & Cryptocurrency price feed cache, refreshed by the Express.js cron service. \textit{Relationships:} read-only auxiliary entity; values are consumed by the frontend and the FX oracle module. No FK into core lending tables. Key attributes: \texttt{symbol} (e.g.\ ETH, USDC), \texttt{price\_usd}, \texttt{source} (Chainlink feed address), \texttt{recorded\_at}. \\
PROFILE\_SETTING & Client and bank-user interface preferences. \textit{Relationships:} FK \texttt{client\_id} references the CLIENT (BORROWER) entity (1:1 per KYC-verified user). Key attributes: \texttt{language\_setting} (e.g.\ \texttt{bn}, \texttt{en}), \texttt{display\_currency} (e.g.\ USD, BDT), \texttt{notification\_enabled} (boolean). \\
\hline
SESSIONS & Formal session management: wallet address, device hash, IP address,
EIP-7702 session key hash, session key scope (JSON: approved tool list + value cap),
session key TTL, and revocation flag. New fields: (1)~\texttt{parent\_session\_id}
(self-referential FK, NULL for root sessions; enables session lineage chain for
cross-session memory); (2)~\texttt{compression\_summary} (TEXT, NULL until a
compression event occurs; carries the prior-session context seed); (3)~\texttt{end\_reason}
(\texttt{ENUM}: \texttt{user\_closed} $|$ \texttt{compression} $|$ \texttt{expired}
$|$ \texttt{revoked}); (4)~\texttt{source} (\texttt{ENUM}: \texttt{web}
$|$ \texttt{whatsapp} $|$ \texttt{telegram} $|$ \texttt{api}; defaults to
\texttt{web}; future-proofs multi-platform delivery). Supports the agent
confirmation gate and provides the audit FK for every agent write-tool execution. \\
\hline
AGENT\_ACTION\_LOG & Append-only log of every agent write-tool execution: tool name, parameters (JSONB), \texttt{confirmation\_turn\_id} FK to \texttt{CHAT\_MESSAGE} (the ``Yes, proceed'' message), on-chain tx hash, and status (\texttt{PENDING} $|$ \texttt{SUBMITTED} $|$ \texttt{CONFIRMED} $|$ \texttt{FAILED}). Logically separate from \texttt{AI\_CHATBOT\_LOG}, which records Q\&A turns. Append-only by DB policy (INSERT-only role + Row-Level Security). Additional fields: \texttt{injection\_scan\_flag} (BOOLEAN DEFAULT FALSE; set TRUE when any field in the user message triggered the prompt injection scanner, providing a per-action audit of security controls applied); \texttt{toolset\_name} (VARCHAR: \texttt{read\_only} $|$ \texttt{loan\_actions} $|$ \texttt{account\_management}; records which tool visibility subset was active for this turn). \\
\hline
INTEREST\_RATE\_TIER & Extracted normalisation of interest rate parameters (\texttt{base\_rate}, \texttt{kink\_utilisation}, \texttt{rate\_above\_kink}, \texttt{max\_rate}) keyed by \texttt{tier\_id}. Eliminates the transitive dependency previously embedded as columns in each bank-tier entity (see Section~\ref{sec:normalization}). \\
\hline
ASSETS & Asset registry for collateral and loan assets (\texttt{symbol}, \texttt{name}, \texttt{asset\_type}: \texttt{STABLECOIN} $|$ \texttt{CRYPTO} $|$ \texttt{RWA}, \texttt{network}, \texttt{decimals}, \texttt{is\_active}). The \texttt{LOAN} table references this entity via \texttt{collateral\_asset\_id} and \texttt{loan\_asset\_id} foreign keys, replacing any inline \texttt{collateral\_symbol} VARCHAR column. \\
\hline
AGENT\_SKILLS & Agent skill registry: named procedural documents that the agent
can load into the Context tier of its system prompt. Schema: \texttt{skill\_id}
(UUID PK), \texttt{skill\_name} (VARCHAR UNIQUE), \texttt{procedure\_text} (TEXT),
\texttt{usage\_count} (INT DEFAULT 0), \texttt{last\_used\_at} (TIMESTAMP),
\texttt{is\_active} (BOOLEAN DEFAULT TRUE). Initially populated with one
manually-authored skill: the Authority Brief generation procedure. \textit{Scope
note: the schema is specified here; the self-improving skill generation loop that
would populate additional rows automatically is scoped to Future Work, as it
requires hundreds of production-scale sessions for calibration.} \\
\hline
\end{tabular}
\TableScreenshotEnd{Database Entity Summary: 20 Normalized Entities in the Relational Schema}
\end{table}

\vspace{2pt}
\noindent\textit{\small This entity summary table now enumerates 20 relational entities. The five additions (\texttt{SESSIONS}, \texttt{AGENT\_ACTION\_LOG}, \texttt{INTEREST\_RATE\_TIER}, \texttt{ASSETS}, \texttt{AGENT\_SKILLS}) support the AI agent session management architecture, the immutable agent audit trail, normalised interest rate governance, standardised asset referencing, and the agent skill registry respectively. The on-chain/off-chain split implied by the entities supports the architectural principle that high-integrity state transitions occur on-chain while analytics and document workflows remain off-chain.}

\subsection{EER Constructs Applied}

\begin{table}[H]
\centering
\caption{EER constructs applied: specialisation, hierarchy, and constraints.}
\begin{tabular}{p{3.5cm}p{3.5cm}p{6.5cm}}
\hline
\textbf{EER Construct} & \textbf{Applied To} & \textbf{Notes} \\
\hline
Specialisation (disjoint) & BANK\_USER $\to$ NationalBankUser, LocalBankUser & A bank user belongs to exactly one bank type; enforced by CHECK constraint. \\
\hline
Generalisation & WORLD\_BANK, NATIONAL\_BANK, LOCAL\_BANK share institution attributes & Avoids attribute duplication across institution types. \\
\hline
Weak entity & INSTALLMENT depends on LOAN\_REQUEST & Existence and identity determined by parent loan. \\
\hline
Multi-valued attribute & INCOME\_PROOF per BORROWER & Modelled as separate entity to maintain 1NF. \\
\hline
Aggregation & TRANSACTION, CHAT\_MESSAGE $\to$ AI\_ML\_SECURITY\_LOG & AI\_ML\_SECURITY\_LOG aggregates monitoring signals from both sources. \\
\hline
Association entity & LOAN\_REQUEST links BORROWER + LOCAL\_BANK & LOAN\_REQUEST is a strong entity with its own UUID primary key and foreign keys to BORROWER and LOCAL\_BANK, capturing the many-to-many relationship between borrowers and lending banks. It is \textit{not} a weak entity (it has its own UUID PK) and is not an aggregation in the ER-theoretic sense; the term ``association entity'' or ``junction entity'' is the correct ER classification. \\
\hline
Participation constraint & BORROWER must have $\geq 1$ LOAN\_REQUEST to access services & Enforced at application layer; architecturally planned for on-chain. \\
\hline
Append-only table (policy) & AGENT\_ACTION\_LOG, AUDIT\_LOGS & DB-level INSERT-only role + Row-Level Security policy enforces that no UPDATE or DELETE is permitted on audit records. Provides immutable tamper-evident trail for regulatory review and agent action accountability. \\
\hline
\end{tabular}
\TableScreenshotEnd{EER Constructs Applied: Specialization, Hierarchy, and Constraints}
\end{table}

\vspace{2pt}
\noindent\textit{\small This table captures the EER constructs applied to represent specialization, hierarchy, and constraints in the data model. It demonstrates how institutional roles and banking participants are modeled without duplicating attributes across tables, improving data integrity and query clarity.}

\subsection{Normalization}

The schema has been normalized to Third Normal Form (3NF) and checked against Boyce-Codd Normal Form (BCNF):

\begin{itemize}[noitemsep]
\item \textbf{1NF:} There are no repeating groups in any of the attributes. Separate entities store multi-valued attributes, such as proof of income.
\item \textbf{2NF:} There are no partial dependencies. The full composite key (\texttt{loan\_id}, \texttt{installment\_number}) determines the non-key attributes of INSTALLMENT.
\item \textbf{3NF:} No dependencies that go through other dependencies. Instead of storing redundant values like \texttt{total\_borrowed} in bank entities, they are calculated at query time.
\item \textbf{BCNF:} All determinants are candidate keys. A CHECK constraint on BANK\_USER enforces that the generalization hierarchy's specializations are not the same.
\end{itemize}

\label{sec:normalization}
A transitive dependency was identified in v24: \texttt{interest\_rate\_parameters} $\to$ \texttt{bank\_tier\_id} (interest rate values were stored as columns in the World Bank, National Bank, and Local Bank entities, making them dependent on a non-key attribute via the tier relationship). This violates 3NF. The fix is to extract these parameters into a dedicated \texttt{INTEREST\_RATE\_TIER} table keyed by \texttt{tier\_id}, which each bank entity references via FK. This normalisation change ensures that updating a base rate does not require touching multiple bank-tier rows.

\subsection{Indexing Strategy}

We use B-tree indexes (the default in PostgreSQL) for efficient retrieval on frequently queried columns:

\begin{table}[H]
\centering
\caption{Indexing strategy: B-tree indexes for time-window and high-frequency queries.}
\begin{tabular}{p{4cm}p{3cm}p{6.5cm}}
\hline
\textbf{Index} & \textbf{Table / Column(s)} & \textbf{Rationale} \\
\hline
\texttt{idx\_loan\_borrower} & LOAN\_REQUEST(\texttt{borrower\_id}) & High-frequency lookup for borrower loan history. \\
\hline
\texttt{idx\_loan\_status} & LOAN\_REQUEST(\texttt{status}) & Efficient filtering of active vs.\ closed loans. \\
\hline
\texttt{idx\_installment\_due} & INSTALLMENT(\texttt{due\_date, status}) \newline WHERE \texttt{status = 'PENDING'} & Partial-index range query for pending overdue installment detection (agent reminder cron + Chainlink Automation). \\
\hline
\texttt{idx\_txn\_created} & TRANSACTION(\texttt{created\_at}) & Time-window aggregation for 6-month and 1-year borrowing limits. \\
\hline
\texttt{idx\_txn\_borrower} & TRANSACTION(\texttt{borrower\_id}) & Per-borrower transaction history retrieval. \\
\hline
\texttt{idx\_market\_symbol} & MARKET\_DATA(\texttt{symbol, recorded\_at}) & Composite index for price feed time-series queries. \\
\hline
\texttt{idx\_loan\_bank\_status} & LOAN(\texttt{local\_bank\_id, status, created\_at DESC}) & Loan lifecycle queries by bank and status (most common query pattern for the approver dashboard). \\
\hline
\texttt{idx\_loan\_client\_active} & LOAN(\texttt{client\_id, status}) \newline WHERE \texttt{status = 'ACTIVE'} & Partial index for per-client open-loan count (used by the over-indebtedness control). \\
\hline
\texttt{idx\_agent\_client\_date} & AGENT\_ACTION\_LOG(\texttt{client\_id, created\_at DESC}) & Per-client agent action history retrieval. \\
\hline
\texttt{idx\_agent\_status} & AGENT\_ACTION\_LOG(\texttt{status}) \newline WHERE \texttt{status = 'PENDING'} & Partial index for the agent status-monitoring loop. \\
\hline
\texttt{idx\_aiml\_score} & AI\_ML\_LOG(\texttt{anomaly\_score DESC}) \newline WHERE \texttt{anomaly\_score > 0.5} & Partial index for AML alert queue (SAR workflow). \\
\hline
\texttt{idx\_session\_parent} & SESSIONS(\texttt{parent\_session\_id}) WHERE \texttt{parent\_session\_id IS NOT NULL} & Partial index for lineage chain traversal; covers all non-root sessions. \\
\hline
\texttt{idx\_session\_fts} & GIN index on SESSIONS(\texttt{to\_tsvector('english', coalesce(compression\_summary,''))}) & Full-text index enabling the \texttt{get\_session\_history} tool to search compression summaries across the lineage chain. Compression summaries are generated by the agent in English regardless of client conversation language; the English text-search configuration is therefore appropriate. PostgreSQL GIN is preferred over B-tree for full-text search. \\
\hline
\texttt{idx\_agent\_scan\_flag} & AGENT\_ACTION\_LOG(\texttt{injection\_scan\_flag}) WHERE \texttt{injection\_scan\_flag = TRUE} & Partial index for security audit queries filtering flagged turns; keeps the flag column queryable without indexing the majority of clean rows. \\
\hline
\end{tabular}
\TableScreenshotEnd{Indexing Strategy: B-tree Indexes for Time-Window and High-Frequency Queries}
\end{table}

\vspace{2pt}
\noindent\textit{\small This indexing table lists the primary indexes used to support time-window queries (e.g., rolling 6-month/1-year transaction windows) and high-frequency lookups. The selection emphasizes B-tree suitability for range filtering common in risk analytics and reporting, improving performance without complicating the schema.}

B-tree indexes are chosen for their efficiency with range queries (e.g., transactions within the last 6 months); hash indexes would be suitable only for exact-match lookups.

\subsection{Functional Dependencies}

\begin{table}[H]
\centering
\caption{Representative functional dependencies.}
\label{tab:functional-deps}
\begin{tabular}{p{3cm}p{5.5cm}p{5cm}}
\hline
\textbf{Relation} & \textbf{Functional Dependency} & \textbf{Notes} \\
\hline
LOAN\_REQUEST & \texttt{loan\_id} $\to$ all attributes & Primary key determines row. \\
\hline
LOAN\_REQUEST & \texttt{borrower\_id, local\_bank\_id} $\to$ status, amount, \ldots & One active request per client per bank. \\
\hline
BORROWING\_LIMIT & \texttt{borrower\_id} $\to$ \texttt{six\_month\_limit, one\_year\_limit}, \ldots & 1:1 with BORROWER. \textit{Note:} BORROWING\_LIMIT stores pre-computed rolling-window limits derived from TRANSACTION records. Because a two-phase commit between PostgreSQL and the blockchain is impossible, the PostgreSQL copy is a read replica for dashboard display only. The authoritative enforcement of borrowing limits occurs on-chain (smart contract checks \texttt{totalBorrowedSixMonth} and \texttt{totalBorrowedOneYear} before any disbursement), preventing the race condition that would arise if the database were the enforcement source. \\
\hline
BORROWER & \texttt{wallet\_address} $\to$ \texttt{borrower\_id}, all attributes & Wallet address is a unique candidate key \textit{for the current prototype}. In the production schema, identity will be decoupled from wallet address (an \texttt{IDENTITY} table keyed by \texttt{kyc\_hash} with a one-to-many \texttt{WALLET} table), so that ERC-4337 social recovery---which produces a new wallet address for the same user---preserves credit history continuity rather than creating a duplicate identity. \\
\hline
INSTALLMENT & \texttt{loan\_id, installment\_number} $\to$ \texttt{amount, due\_date, status} & Composite primary key; fully determined by both columns. \\
\hline
SESSIONS & \texttt{session\_id} $\to$ \texttt{client\_id, wallet\_address, session\_key\_hash,} \newline \texttt{session\_key\_scope, session\_key\_expires\_at,} \newline \texttt{expires\_at, revoked} & Session key scope (JSONB) encodes the approved MCP tool list + value cap per session. \\
\hline
AGENT\_ACTION\_LOG & \texttt{action\_id} $\to$ \texttt{session\_id, client\_id, tool\_name,} \newline \texttt{parameters, confirmation\_turn\_id,} \newline \texttt{onchain\_tx\_hash, status, created\_at} & \texttt{confirmation\_turn\_id} FK to \texttt{CHAT\_MESSAGE.message\_id} links every write-tool execution to the user's consent message. \\
\hline
ASSETS & \texttt{asset\_id} $\to$ \texttt{symbol, name, asset\_type,} \newline \texttt{network, decimals, is\_active}; also \texttt{symbol} $\to$ \texttt{asset\_id} & \texttt{symbol} is a unique candidate key; both FDs hold. \\
\hline
INTEREST\_RATE\_TIER & \texttt{tier\_id} $\to$ \texttt{base\_rate, kink\_utilization,} \newline \texttt{rate\_above\_kink, max\_rate} & Tier key uniquely determines all rate parameters; eliminates the transitive dependency in bank-tier entities. \\
\hline
\end{tabular}
\TableScreenshotEnd{Functional Dependencies in the Core Lending and Identity Schema}
\end{table}

\vspace{2pt}
\noindent\textit{\small This functional dependency table documents the data-level rules that prevent inconsistent state, such as improper loan status transitions or duplicated relationships. Explicit FDs reinforce that key banking invariants are enforced structurally, complementing smart contract enforcement on the on-chain side.}

\noindent The \texttt{parent\_session\_id} self-referential foreign key in
\texttt{SESSIONS} introduces a recursive relationship within a single entity.
This does not violate 3NF (the dependency is \texttt{session\_id} $\to$
\texttt{parent\_session\_id}, not a transitive dependency), but it requires an
acyclicity constraint to prevent circular ancestry chains: a session cannot be
its own ancestor. This is enforced by a \texttt{CHECK} constraint in conjunction
with a trigger that walks the ancestry chain before INSERT (see
Section~\ref{sec:integrity-constraints}).

\subsection{Relational Integrity Constraints}

\begin{table}[H]
\centering
\caption{Relational integrity constraints.}
\label{tab:integrity-constraints}
\begin{tabular}{p{4cm}p{9.5cm}}
\hline
\textbf{Constraint Type} & \textbf{Examples} \\
\hline
Primary Key & \texttt{world\_bank\_id}, \texttt{loan\_id}, \texttt{borrower\_id}, etc. \\
\hline
Foreign Key & \texttt{local\_bank\_id} in LOAN\_REQUEST references LOCAL\_BANK. \\
\hline
UNIQUE & \texttt{wallet\_address} in BORROWER; \texttt{blockchain\_tx\_hash} in LOAN\_REQUEST. \\
\hline
CHECK & BANK\_USER: \texttt{(bank\_type='national' AND national\_bank\_id IS NOT NULL) OR (bank\_type='local' AND local\_bank\_id IS NOT NULL)}. \\
\hline
NOT NULL & Core attributes: \texttt{name}, \texttt{wallet\_address}, \texttt{status}. \\
\hline
APPEND-ONLY (DB policy) & \texttt{AGENT\_ACTION\_LOG}: \texttt{CREATE ROLE audit\_writer; GRANT INSERT ON agent\_action\_log TO audit\_writer; REVOKE UPDATE, DELETE FROM PUBLIC.} \texttt{AUDIT\_LOGS}: same policy applied via Row-Level Security (\texttt{ALTER TABLE audit\_logs ENABLE ROW LEVEL SECURITY; CREATE POLICY audit\_insert\_only ON audit\_logs FOR INSERT WITH CHECK (TRUE)}). \\
\hline
FOREIGN KEY (new) & \texttt{LOAN.collateral\_asset\_id} \texttt{REFERENCES assets(asset\_id)}; \texttt{LOAN.loan\_asset\_id} \texttt{REFERENCES assets(asset\_id)}; \texttt{AGENT\_ACTION\_LOG.session\_id} \texttt{REFERENCES sessions(session\_id)}; \texttt{AGENT\_ACTION\_LOG.confirmation\_turn\_id} \texttt{REFERENCES chat\_message(message\_id)}. \\
\hline
UNIQUE (new) & \texttt{assets.symbol} (asset ticker is globally unique across the platform asset registry). \\
\hline
CHECK (session lineage acyclicity) & \texttt{SESSIONS}: \texttt{CHECK (parent\_session\_id IS NULL OR parent\_session\_id <> session\_id)}. A trigger additionally walks the ancestry chain to reject cycles beyond immediate self-reference. \\
\hline
FOREIGN KEY (session lineage) & \texttt{SESSIONS.parent\_session\_id REFERENCES sessions(session\_id)} (self-referential). Ensures every non-null \texttt{parent\_session\_id} points to an existing session row; orphaned lineage references are impossible. \\
\hline
\end{tabular}
\TableScreenshotEnd{Relational Integrity Constraints: Referential and Domain Rules}
\end{table}

\vspace{2pt}
\noindent\textit{\small This integrity constraints table summarizes referential and domain constraints that keep loan, repayment, and identity records consistent. These constraints reduce operational risk by preventing orphaned records and invalid lifecycle states, which is critical for auditability in financial systems.}


\section{On-Chain and Off-Chain Data Partitioning}

\begin{table}[H]
\centering
\caption{Data partitioning between on-chain and off-chain storage.}
\label{tab:data-partitioning}
\begin{tabular}{p{4.5cm}p{2.5cm}p{6.5cm}}
\hline
\textbf{Data Category} & \textbf{Storage} & \textbf{Rationale} \\
\hline
Reserve balances, loan requests, approval/rejection events, repayment transactions & On-chain & Immutability, public auditability, trustless verification. \\
\hline
User profiles, income verification documents, chat messages, AI/ML inference logs & Off-chain (database) & Data privacy, query flexibility, storage cost optimisation. \\
\hline
Borrowing limit computations & Off-chain with on-chain enforcement & Complex temporal aggregation; results committed as on-chain constraints. \\
\hline
Cryptocurrency market data & Off-chain (cached) & High-frequency updates; external API dependency. \\
\hline
Agent action execution log & Off-chain (append-only DB) & Immutable audit trail of every agent write-tool call, linked to on-chain tx hash; not stored on-chain to save gas. \\
\hline
Session key material & Off-chain (\texttt{sessions} table) & EIP-7702 session key hash, scope JSON, and TTL stored off-chain; the scope restrictions are enforced on-chain by the session key contract at signing time. \\
\hline
Chainlink oracle price data & Off-chain (\texttt{MARKET\_DATA}) & Chainlink Price Feed results are cached in \texttt{MARKET\_DATA} with \texttt{source} = Chainlink feed address; the on-chain \texttt{AggregatorV3Interface} is the authoritative source. \\
\hline
Session lineage chains (SESSIONS, message history, compression summaries) & Off-chain (PostgreSQL) & Full transcript history is too large for on-chain storage. The lineage chain enables searchable long-term memory without gas cost. \texttt{AGENT\_ACTION\_LOG} provides the on-chain audit link via confirmed transaction hashes. \\
\hline
Agent skill procedures (\texttt{AGENT\_SKILLS.procedure\_text}) & Off-chain (PostgreSQL) & Skill documents are model-facing prompt content with no on-chain settlement obligation. Storing off-chain allows iterative improvement without contract upgrades or governance votes. \\
\hline
\end{tabular}
\TableScreenshotEnd{On-Chain vs.~Off-Chain Data Partitioning: State, Analytics, and Privacy Boundaries}
\end{table}

\vspace{2pt}
\noindent\textit{\small This partitioning table clarifies which data must live on-chain (role bindings, reserves, and critical state transitions) versus off-chain (documents, analytics features, and operational metadata). The partitioning choice balances transparency and tamper-resistance with practical storage and privacy constraints.}


\section{Digital Identity System}

The platform's identity model operates across two layers, combining wallet-based authentication with off-chain document verification, and is designed to accommodate compliance requirements in future phases.

\begin{itemize}[noitemsep]
\item \textbf{Identity based on wallets:} Users authenticate via Ethereum wallet signatures (MetaMask, WalletConnect), eliminating the need for centralised credential storage.
\item \textbf{Role binding:} In the smart contract permission system, wallet addresses are linked to hierarchical roles like Owner, National Bank, Local Bank, Approver, and Retail Client.
\item \textbf{Limits of wallet-based identity:} Wallet ownership proves transaction history and cryptographic control but cannot independently prove legal identity, jurisdiction, or age. This distinction is critical for regulated banking operations. A compromised or lost wallet requires an off-chain recovery process and on-chain role revocation followed by re-binding to a replacement address, a governance workflow that must be carefully designed to prevent unauthorized role escalation.
\item \textbf{On-chain versus off-chain identity layers:} The system maintains two identity layers. On-chain identity consists of the wallet address, assigned role, and transaction history recorded permanently on the blockchain. Off-chain identity consists of document-verified income, KYC credentials, and AML screening results stored in the PostgreSQL database and linked to the wallet address by hash reference. The planned ZKP-based compliance extension would allow users to prove off-chain KYC status to the smart contract without exposing personal documents on-chain.
\item \textbf{EIP-7702 session key authorisation:} For agent-executed banking operations, the client authorises a scoped, time-bound session key at login. The session key is restricted to a named set of MCP write tools (e.g., only \texttt{submit\_loan\_application} and \texttt{pay\_installment}), a value cap (e.g., 500~USDC per transaction), and a 24-hour TTL after which it auto-expires. The session key cannot transfer funds to external addresses or perform any operation outside its approved scope. Every session key usage is logged in \texttt{AGENT\_ACTION\_LOG} with the conversation turn ID of the client's confirmation message as the audit reference.
\end{itemize}

\subsection{ZKP KYC and zkAML Compliance Architecture}
\label{sec:identity}

\begin{figure}[H]
\centering
\OnePageDiagram{fig-compliance-identity.pdf}
\caption{Compliance and identity stack: (a) zkKYC issuance via licensed identity provider with W3C Verifiable Credential, (b) zkAML continuous monitoring with sanction-list non-membership proofs, (c) the tiered KYC ladder (L1 zkKYC $\to$ L4 entity onboarding) gated by loan size, and (d) ERC-4337 Account-Abstraction onboarding for non-crypto users with gas sponsored by a Paymaster.}
\label{fig:compliance-identity}
\end{figure}


The compliance gateway uses two cooperating zero-knowledge circuits rather than a single KYC verifier: a \textbf{KYC circuit} and an \textbf{AML circuit}, both expressed as zk-SNARKs (Groth16 proofs via Circom~2.0 + snarkjs). KYC proves that the user owns a valid identity credential; AML proves that the user's wallet has not interacted with sanctioned addresses and stays within velocity limits. Decoupling the two circuits makes it possible to revoke AML certification without invalidating the KYC credential, and matches the actual structure of regulator-mandated compliance.

\paragraph{KYC circuit (identity).}
An off-chain KYC provider validates the user's NID and issues a signed credential
\(\textit{credential} = \text{Sign}_{\text{KYC}}(\textit{wallet\_address},\, \textit{country},\, \textit{age\_over\_18},\, \textit{kyc\_passed})\).
The user generates a zk-SNARK proof that (a) they possess a valid signed credential from an approved KYC provider, (b) their age is over~18, and (c) their country is in the permitted jurisdiction list---without revealing the underlying credential, NID, or personal data. The smart contract verifies the proof:
\texttt{KYCVerifier.verify(proof, [wallet, country, age\_over\_18])} returns \texttt{true}.
Piper et~al.~(2025)~[R8] report proof generation in 1--4~s on consumer hardware; on-chain verification costs $\approx$200--300\,k gas, within a one-time registration budget.

\paragraph{zkAML circuit (anti-money-laundering).}
The zkAML circuit~[R19] (IACR ePrint 2025/465) proves three properties of a wallet without revealing its transaction graph: (1) the wallet has not transacted with any address on a public sanctions list (Chainalysis / TRM Labs); (2) the wallet's 30-day inflow does not exceed a governance-set velocity threshold; and (3) the wallet has not received funds from addresses flagged in the platform's on-chain risk registry within the preceding 90 days. The published benchmark for zkAML is 55~TPS on public networks with 226\,ms proof-generation latency, both compatible with retail-tier flows on Polygon. The on-chain verifier is a sibling of \texttt{KYCVerifier} that exposes \texttt{AMLVerifier.verify(proof, [wallet, blockHeight])}. Both verifiers must return true before a wallet's CLIENT role is activated.

\paragraph{W3C DID / VC anchoring.}
The KYC credential issued off-chain is shaped as a W3C Verifiable Credential (VC) bound to a W3C Decentralized Identifier (DID)~[R20] rather than a free-form signed blob. The DID is anchored on-chain (\texttt{did:ethr:0x\ldots} per the EIP-1056 standard), and the smart contract checks possession of the VC by means of the zk-SNARK proof rather than by storing the VC itself. This places the platform's identity layer inside the Self-Sovereign Identity (SSI) framework and aligns it with the European Blockchain Services Infrastructure (EBSI) reference architecture.

\paragraph{Privacy posture.}
A 2025 ResearchGate study reports a 97\% reduction in exposed user data versus conventional on-chain KYC; Decker~(2025)~[R9] gives the same finding under a different threat model. Implementation for the final thesis phase will deploy both Circom~2.0 circuits to Polygon~zkEVM~Cardona testnet, tested with synthetic credential data drawn from the platform's mock KYC provider.

\section{User Taxonomy and Onboarding Flows}
\label{sec:onboarding}

The thesis now uses a formal actor taxonomy of nine actors following the systems-analysis convention of the Binance Software Engineering Architecture reference. Five primary actors (A1--A5) initiate actions; four secondary actors (A6--A9) are external systems or authorities. Table~\ref{tab:user-taxonomy} gives the operational user taxonomy with per-user-type KYC level, wallet type, and database-residence assumption; Table~\ref{tab:actor-taxonomy} presents the complete formal actor taxonomy; and Table~\ref{tab:permission-matrix} presents the actor permission matrix.

\begin{table}[H]
\centering
\caption{Operational user taxonomy with KYC level, wallet type, and data-residence assumption. The taxonomy expands the actor list referenced in the use-case diagram of Section~\ref{sec:usecase}.}
\label{tab:user-taxonomy}
\small
\begin{tabular}{p{3.1cm}p{1.4cm}p{2.6cm}p{2.4cm}p{3.4cm}}
\hline
\textbf{User Type} & \textbf{Tier} & \textbf{KYC Level} & \textbf{Wallet Type} & \textbf{Data Residence} \\
\hline
Anonymous Visitor & --- & None & None & Session only (Redis) \\
\hline
Registered Retail User & Tier 4 & Level 1 (NID + selfie) & ERC-4337 abstracted (email / Google login) & Off-chain primary; on-chain role only \\
\hline
Group Client & Tier 4 & Level 1 + group verification & ERC-4337 abstracted & Off-chain + on-chain group contract \\
\hline
Local Bank Operator & Tier 3 & Full institutional KYC & MetaMask (institutional) & Off-chain + on-chain role binding \\
\hline
Local Bank Approver & Tier 3 & Full institutional + background check & MetaMask + Safe co-signer & Off-chain + on-chain role binding \\
\hline
National Bank Admin & Tier 2 & Full institutional + regulatory license & Safe multisig mandatory & Off-chain + on-chain role binding \\
\hline
World Bank Admin (Governance) & Tier 1 & Founding team + supervisor attestation & Safe 3-of-5 multisig & On-chain governance only \\
\hline
\end{tabular}
\end{table}

\begin{table}[H]
\centering
\caption{CWB formal actor taxonomy (nine actors): five primary actors who initiate actions and four secondary actors who are external systems or authorities.}
\label{tab:actor-taxonomy}
\renewcommand{\arraystretch}{1.3}
\small
\begin{tabular}{p{1.2cm}p{4.2cm}p{7.9cm}}
\hline
\textbf{Label} & \textbf{Actor} & \textbf{Role} \\
\hline
\multicolumn{3}{l}{\textit{Primary Actors (initiate actions)}} \\
\hline
A1a & Retail Client (pre-KYC) & Browsing only; no banking transactions. \\
\hline
A1b & Retail Client (KYC Tier~1) & Bronze/Silver credit tier; small loans up to the KYC-1 cap. \\
\hline
A1c & Retail Client (KYC Tier~2) & Gold/Platinum/Diamond; larger limits, group lending. \\
\hline
A2 & Local Bank Admin (Approver) & Loan approver; risk officer; reviews AML alerts. \\
\hline
A3 & National Bank Admin & Capital allocator; compliance officer; SAR review; freeze authority. \\
\hline
A4 & World Bank Admin (Governance) & Parameter governor; system operator via multi-sig. \\
\hline
A5 & AI Agent & Read-only tools always permitted; write tools require explicit human confirmation gate before execution. \\
\hline
\multicolumn{3}{l}{\textit{Secondary Actors (external systems or authorities)}} \\
\hline
A6 & Regulatory Authority & Read-only audit access via encrypted data package. \\
\hline
A7 & Chainlink DON & Oracle, price feeds, automation. \\
\hline
A8 & Blockchain Validator & Polygon zkEVM validity proof generation. \\
\hline
A9 & External Auditor & Chainlink PoR verification. \\
\hline
\end{tabular}
\end{table}

\begin{table}[H]
\centering
\caption{CWB actor permission matrix. READ = read-only tool call; WRITE\dag{} = write tool requiring explicit human confirmation gate. * = own tier only.}
\label{tab:permission-matrix}
\renewcommand{\arraystretch}{1.3}
\small
\begin{tabular}{p{4.4cm}p{1.3cm}p{1.5cm}p{1.5cm}p{1.5cm}p{1.6cm}}
\hline
\textbf{Action} & \textbf{Client} & \textbf{LB Admin} & \textbf{NB Admin} & \textbf{WB Admin} & \textbf{AI Agent} \\
\hline
View own loan status & YES & YES* & YES* & YES* & READ \\
\hline
Submit loan application & YES & NO & NO & NO & WRITE\dag{} \\
\hline
Pay installment & YES & NO & NO & NO & WRITE\dag{} \\
\hline
Submit KYC documents & YES & NO & NO & NO & NO \\
\hline
Approve loan application & NO & YES & NO & NO & NO \\
\hline
Reject loan application & NO & YES & NO & NO & NO \\
\hline
Freeze client account & NO & YES & YES & YES & NO \\
\hline
Set borrowing rate & NO & NO & YES & YES & NO \\
\hline
Change reserve ratio & NO & NO & NO & YES & NO \\
\hline
Upgrade LB borrowing cap & NO & NO & YES & YES & NO \\
\hline
View audit logs (all) & NO & YES* & YES* & YES & NO \\
\hline
Generate SAR report & NO & YES & YES & YES & NO \\
\hline
View reserve summary (public) & YES & YES & YES & YES & READ \\
\hline
\end{tabular}
\vspace{2pt}

\noindent\textit{\small \dag{} All agent write tools require an explicit human confirmation gate: the agent assembles the full parameter set, presents a summary to the client, and waits for affirmative consent. No write tool is ever called without a confirmation turn in the conversation history. * own tier only.}
\end{table}

\subsection{ERC-4337 Account Abstraction for Retail Onboarding}
\label{sec:erc4337}

The single largest accessibility contradiction in the original thesis was that retail clients in developing economies were assumed to already understand MetaMask, seed phrases, gas fees, and Ethereum transactions. This is incompatible with the financial-inclusion mission. The platform resolves the contradiction by adopting \textbf{ERC-4337 Account Abstraction (AA)} for all retail-tier (Tier~4) clients, with the EIP-7702 Pectra upgrade (May~2025) providing the gas-sponsorship layer. Since its launch on Ethereum mainnet in March~2023, ERC-4337 has enabled over 40~million smart accounts and processed more than 100~million transactions across major L2 networks including Polygon, Arbitrum, Optimism, and Base, demonstrating production-grade viability at scale~[18-1]. The current production EntryPoint is \textbf{v0.7}, deployed at \texttt{0x0000000071727De22E5E9d8BAf0edAc6f37da032} on Ethereum, Polygon, and all major EVM chains; the platform targets this canonical address for all UserOperation routing.

Under ERC-4337 each retail client is issued a smart-contract account that is created on first sign-in. The user can register with email or a Google account via Firebase Auth, never sees a seed phrase, and never needs to acquire ETH for gas: a Paymaster contract sponsors the gas, funded by the World Bank Reserve as a financial-inclusion subsidy line. Account recovery uses a social-recovery guardian (a trusted email / phone) rather than a 24-word mnemonic.

\paragraph{Paymaster sybil-resistance controls.}
Gas sponsorship introduces a perverse incentive: an attacker can create many ERC-4337 accounts at zero cost and spam loan applications, draining the gas subsidy and potentially the loan pool if fraud detection fails. Three controls mitigate this. First, the Paymaster implements a \textbf{pre-KYC bootstrap allowance}: before KYC is completed, the Paymaster sponsors exactly two function signatures per wallet --- \texttt{submitKYCHash(bytes32)} (storing the document hash for Level~1 KYC) and the ERC-4337 account creation transaction. Both are capped at a maximum of 0.001~ETH-equivalent total per wallet address; this is sufficient to complete KYC but insufficient to spam loan applications. All other gas sponsorship is conditional on the wallet having a verified KYC status (\texttt{kycVerified[wallet] = true}). This pre-KYC bootstrap window resolves the chicken-and-egg problem: a new user with no ETH can complete KYC using the bootstrap allowance, after which normal Paymaster sponsorship unlocks. Second, the Paymaster budget is capped per verified identity (not per account), limiting lifetime gas sponsorship to a fixed USDC equivalent per KYC identity. Third, the Paymaster rate-limits by IP and device fingerprint to prevent bulk account creation from a single attacker. These controls make Paymaster sponsorship conditional on verified identity rather than unconditional, preserving the inclusion mission while preventing abuse.

For the demo and final thesis prototype, the Biconomy SDK and Alchemy Account Kit will provide the bundler / Paymaster infrastructure on Polygon~zkEVM Cardona. From the smart-contract perspective, an account-abstracted wallet is indistinguishable from an externally owned account, so none of the four-tier role logic needs to change. From the client's perspective, the cryptocurrency stack is invisible: they see a Bengali- or English-language banking interface, sign in with email, and apply for a loan. This re-aligns the implementation with the financial-inclusion contribution claimed in Chapter~1.

\paragraph{EIP-7702 session keys --- scoped agent wallet.}
EIP-7702 provides a cleaner solution than ERC-4337 paymasters for agent-controlled operations because the scope restriction is enforced at the key level, not at the application level. When a client authorises an AI Agent session at login, the EIP-7702 session key is subject to four hard constraints: (1) \textbf{Scope} --- the key is restricted to a named set of MCP write tools (e.g., only \texttt{submit\_loan\_application} and \texttt{pay\_installment}); (2) \textbf{Time-bound} --- a 24-hour TTL after which the key auto-expires regardless of remaining uses; (3) \textbf{Value-capped} --- a maximum of 500~USDC per transaction, preventing large unauthorised transfers; and (4) \textbf{Revocable} --- the client can invalidate the session key at any time from the wallet UI. The session key \emph{cannot} transfer funds to external addresses, exceed the value cap, perform any operation not in the approved scope, or execute after the TTL expires.

EIP-7702 is architecturally superior to ERC-4337 paymasters for agent-controlled operations and does not replace ERC-4337 for retail gas sponsorship --- the two mechanisms serve different functions. ERC-4337 removes the need for ETH for gas; EIP-7702 scopes what the agent may sign. Both are active simultaneously for an agent session: the session key signs the transaction, the Paymaster sponsors the gas.

\subsection{Tiered Risk-Based KYC}

\begin{figure}[H]
\centering
\OnePageDiagram{fig-tier-model.pdf}
\caption{Tiered borrower access model. KYC level is monotone non-increasing down the tier ladder; the loan-band ceiling is enforced on-chain via \texttt{borrowerLevel[wallet]}; stablecoin denomination is enforced at Tier~4 (Section~\ref{sec:stablecoin-first}); and the SBT (Section~\ref{sec:credit-passport}) carries the simultaneous-loan cap that the GroupLendingPool reads at every application.}
\label{fig:tier-model}
\end{figure}

\label{sec:tiered-kyc}

Identity verification is gated by a risk-based KYC ladder rather than an all-or-nothing flag. Table~\ref{tab:tiered-kyc} gives the four levels and the corresponding borrowing limits and verification times. The design is consistent with current Persona / Veriff guidance for retail-fintech KYC and the EU AI Act Article~86 transparency requirements.

\begin{table}[H]
\centering
\caption{Tiered, risk-based KYC ladder. Higher KYC level unlocks larger loans but requires proportionally more verification.}
\label{tab:tiered-kyc}
\begin{tabular}{p{2.2cm}p{5.8cm}p{2.8cm}p{2.6cm}}
\hline
\textbf{Risk Level} & \textbf{Required Documents} & \textbf{Borrow Limit} & \textbf{Verify Time} \\
\hline
Level 0 (Unverified) & None & View only & Instant \\
\hline
Level 1 (Basic) & NID photo + selfie liveness check & Up to 0.1 ETH-equivalent (USDC) & 2--5 min (automated) \\
\hline
Level 2 (Standard) & NID + proof of income + bank statement & Up to 2 ETH-equivalent & 24~h (human review) \\
\hline
Level 3 (Enhanced) & Full document set + video interview & Up to 10 ETH-equivalent & 48--72~h \\
\hline
\end{tabular}
\end{table}

After KYC approval, the smart contract sets \texttt{kycVerified[wallet] = true}, \texttt{kycLevel[wallet] = level}, and a one-year expiry \texttt{kycExpiry[wallet] = block.timestamp + 365 days}. Re-verification is enforced via the role-expiry mechanism (Section~\ref{sec:identity}). Personal data---NID number, biometric template, document files---is never stored on-chain; only a SHA-256 hash of the document and the zk-SNARK proof are written, in compliance with GDPR data-minimization principles and Bangladesh's Personal Data Protection regime.

\subsection{Five-Stage Conversion Funnel for Non-Crypto Users}
\label{sec:funnel}

The platform's retail conversion strategy is structured as a five-stage funnel that progressively reveals technical complexity rather than demanding it up front. Stage~1 is information browsing with zero friction (no account required, Bengali / English chatbot answers in plain language). Stage~2 is account creation with email or phone---an ERC-4337 smart account is created in the background; the user never sees a seed phrase. Stage~3 is KYC verification via phone-camera document capture with automated AI verification in 2--5~minutes. Stage~4 is the first loan, with the platform explaining the workflow in plain language and showing the repayment schedule in local-currency equivalent through a USD-BDT forex oracle (or a fiat-price oracle approved by governance), since retail loans are denominated in USDC which tracks USD rather than ETH. Stage~5 is the optional power-user mode in which the client can connect a self-custodied MetaMask wallet, view their transactions on a block explorer, or participate in solidarity-group lending. The funnel hides crypto until the user is ready for it and re-frames the platform as ``a bank that happens to settle on a public ledger'' rather than ``a DeFi protocol that happens to call itself a bank''.


\section{Kinked Interest Rate Model}
\label{sec:kinked-rate}

The Crypto World Bank adopts a \textbf{kinked utilization-based interest rate model}, following the design established by Compound Finance and Aave~v2/v3 and motivated by empirical findings from Gudgeon et al.~(2020)~[R3]. Below an optimal utilization rate $U^*$ (set at 80\% for retail lending pools), the borrowing rate increases gently:

\begin{equation}
r_b(U) = r_0 + \frac{U}{U^*} \cdot r_1, \quad U \leq U^* \tag{Formula 18}
\end{equation}

Above $U^*$, the rate increases steeply to incentivize rapid repayment and new deposits:

\begin{equation}
r_b(U) = r_0 + r_1 + \frac{U - U^*}{1 - U^*} \cdot r_2, \quad U > U^* \tag{Formula 19}
\end{equation}

where $r_0$ is the base rate, $r_1$ is the slope below the kink, and $r_2$ is the jump multiplier above the kink. This piecewise model prevents liquidity crises in which utilization approaches 100\% and depositors are unable to withdraw---a failure mode documented empirically in DeFi lending markets by Gudgeon et al.~(2020)~[R3] and analyzed theoretically by Mackinga et al.~(2023)~[R4]. The smart contract enforces two floor/ceiling guards on these parameters: \texttt{require(r2 >= MINIMUM\_JUMP\_MULTIPLIER)} ensures the above-kink rate rise is steep enough to incentivize repayment and deposit inflow (a jump multiplier below the minimum would flatten the above-kink curve and fail to deter utilization runaway); and \texttt{require(U\_star <= MAXIMUM\_KINK\_POINT)} prevents the kink from being placed so high that the safety mechanism is effectively disabled. Both constants are governance-set and enforced on-chain before any parameter update is executed.

\section{Liquidation Engine}
\label{sec:liquidation}

The Health Factor formula ($HF = C \times LT / D$, List of Formulas) applies specifically to \textit{over-collateralized retail loans} at Tier~4 where a client has deposited ETH or USDC collateral exceeding the loan value. It is operationally meaningless without a contract that monitors it and an actor incentivized to act on the result, and it does not apply uniformly across all loan types in the CWB. The tier-specific Health Factor models are:

\begin{itemize}[noitemsep]
\item \textbf{Tier~4 over-collateralized retail:} Standard $HF = (C \times LT) / D$; liquidation triggered at $HF < 1.0$.
\item \textbf{Tier~4 group lending (collateral pool):} $HF_{\text{group}} = (\text{total\_group\_collateral} \times LT) / \text{total\_group\_outstanding\_debt}$; liquidation is at the pool level, distributing the loss across all group members, not the individual who defaulted.
\item \textbf{Tier~4 credit-based (no collateral):} HF is not applicable. Default enforcement relies on credit-score degradation, permanent \texttt{riskTier} downgrade (with \texttt{lastDefault} timestamp irrevocably set on the SBT), and an effective future-borrowing ban. The SBT credential itself is \textit{never revoked}---consistent with the Buterin, Hitzig, and Weyl~(2022)~[R30] non-revocable design principle; no lender accepts a wallet whose SBT \texttt{riskTier} is at the default level, making borrowing effectively impossible in practice while preserving the permanent, honest record of the default event. No liquidation is possible because no collateral exists to seize.
\item \textbf{Tier~1--3 institutional:} $HF_{\text{inst}} = (\text{on\_chain\_reserve} - \text{min\_reserve}) / \text{outstanding\_borrowing}$; the parent tier's reserve surplus serves as the effective collateral for inter-tier lending.
\end{itemize}

The original v10 thesis listed HF but contained no liquidation mechanism: when HF dropped below 1.0, the system relied on a human approver noticing and reacting, which is a ledger, not a bank. The Crypto World Bank therefore specifies a dedicated \textbf{LiquidationEngine} contract that closes this gap.

\paragraph{Liquidation life-cycle.}
\begin{enumerate}[noitemsep]
\item A client's collateral position is updated from the Chainlink price feed at every block; HF is recomputed in view-only fashion (no gas cost to the client).
\item Any external participant may call \texttt{liquidate(loanId)} when HF $< 1.0$.
\item The contract verifies HF on-chain, transfers a configurable portion of the collateral (typically equal to the outstanding debt plus a 5--8\% liquidation bonus, following Aave~v3 parameters) to the liquidator, and updates the client's on-chain credit record.
\item A \texttt{LoanLiquidated(loanId, liquidator, bonusBps)} event is emitted and indexed by The Graph for the dashboard.
\end{enumerate}

\paragraph{Incentive design.}
The 5--8\% liquidation bonus is what makes liquidation self-actuating: any wallet (including a public arbitrage bot) can call \texttt{liquidate} when HF dips, captures the bonus, and protects the lending pool from accumulating bad debt. This is a standard DeFi primitive in production protocols since 2020 (Aave, Compound, MakerDAO) and its absence from a paper that calls itself a banking system is a credibility gap that v15 closes.

\paragraph{Hierarchical liquidation queue.}
The platform extends the standard single-tier liquidation primitive with a hierarchical priority: when a Local Bank exceeds a portfolio-level reserve-ratio threshold (RR below 15\%), liquidation calls from authorized liquidators are routed through a National-Bank-level queue that prioritizes the lowest-HF positions first. This converts the same liquidation primitive into a tier-aware portfolio-health control rather than a per-loan reflex, matching how Tier~2 institutions actually manage credit risk.

\section{SavingsVault and FixedDeposit Modules}
\label{sec:savings-vault}

A bank is fundamentally a deposit institution, not a credit-only protocol. The original v10 design described the SavingsVault as ``planned'' but left the contract unspecified, which left the entire ``closed-loop sustainable model'' claim resting on capital that did not exist. v15 specifies the two deposit-mobilization contracts directly.

\paragraph{SavingsVault.}
The SavingsVault contract accepts ERC-20 deposits (USDC primary, ETH/USDT/DAI for institutional flows), credits the depositor's balance on-chain, accrues variable yield tied to pool utilization (the kinked rate of Section~\ref{sec:kinked-rate}), and permits withdrawal subject to the local-tier reserve ratio gate. The yield is paid in the same token as the deposit; interest accounting is per-second using a cumulative index rather than per-block accrual, exactly matching the Compound~v3 / Aave~v3 supply-side pattern. The contract is upgradable via UUPS (Section~\ref{sec:defense-in-depth}) and emits \texttt{Deposit}, \texttt{Withdraw}, \texttt{InterestAccrued}, and \texttt{ReserveGateTriggered} events for The Graph indexing.

\paragraph{FixedDeposit.}
FixedDeposit issues a non-transferable on-chain receipt that locks principal for a fixed term (30 / 90 / 180 / 365 days) at an APY written at deposit time. Early withdrawal forfeits a deterministic portion of accrued interest, transparently. Fixed-term deposits provide the Local Bank tier with \textbf{duration-matched liabilities}: a 180-day FixedDeposit funds the 6-month installment tranche of a retail client at the same tier without creating asset-liability mismatch. The contract integrates with the InsuranceFund so that, under stress scenarios, fixed-term depositors are subordinated to demand depositors only after the insurance buffer is exhausted.

\paragraph{ERC-4626 and ERC-7540 interface alignment.}
The \texttt{SavingsVault} contract is designed to implement the ERC-4626 Tokenized Vault Standard (\texttt{IERC4626}), which defines a unified interface for yield-bearing ERC-20 vaults: \texttt{deposit}, \texttt{mint}, \texttt{withdraw}, \texttt{redeem}, and the asset-to-shares conversion functions \texttt{convertToShares} and \texttt{convertToAssets}~[R46].  A depositor's share token (\texttt{svUSDC}) represents a pro-rata claim on the underlying USDC pool, updated continuously by the cumulative interest index described above.  Aligning with ERC-4626 means the vault is auditable against a known, widely-reviewed standard; any future integrator---a yield aggregator, a cross-protocol lending market, or a liquidity router---can compose on top of it without a custom adapter.  Prior to EIP-4626, every yield protocol shipped a bespoke interface (Yearn vaults, Aave aTokens, Compound cTokens), forcing each integrator to write and audit a separate adapter; ERC-4626 eliminates that per-protocol overhead and is the 2025--2026 DeFi standard for tokenized vaults.  The \texttt{FixedDeposit} contract adopts ERC-7540 (the asynchronous redemption extension of ERC-4626)~[R46] because fixed-term redemptions are time-gated: a request submitted at lock-up close enters a \texttt{requestId}-keyed pending state and becomes claimable only after the maturity block is reached, which fits the asynchronous request/fulfil flow that ERC-7540 formalises.  This is a specification alignment, not an implementation rewrite; the underlying interest-accrual and penalty logic described above is unchanged.

\paragraph{Closed-loop economic argument.}
With SavingsVault and FixedDeposit in place, the platform's economic claim becomes verifiable: depositor capital flows in on one side, lending flows out on the other, interest collected is split deterministically between depositor yield, the insurance fund, and protocol revenue under Formula NetInterest (see List of Formulas). The model recycles capital without external subsidy, distinguishing the platform from donor-funded MFIs and matching the empirical evidence on sustainable microfinance that Section~\ref{sec:lit-prisma} cites.

\section{On-Chain Credit Passport (Soulbound Token)}
\label{sec:credit-passport}

The platform issues each KYC-verified client a non-transferable on-chain credential---a Soulbound Token (SBT) following the Buterin, Hitzig, and Weyl framework~[R30]---that encodes their cumulative repayment record in a schema that any compatible lending protocol can read. The credential is updated automatically whenever an installment is repaid in full, and downgraded (with \texttt{lastDefault} timestamp permanently set) after a defined number of missed installments; the credential is never revoked, consistent with the non-revocable Buterin SBT design principle~[R30]. Three concrete benefits follow.

First, the SBT solves the cold-start problem at the bank-switching boundary: a credit client who has repaid three loans at a Local Bank in Sylhet can apply for a larger loan at a Local Bank in Dhaka without re-establishing trust from zero. Second, the SBT is the on-chain primitive that the over-indebtedness control (Section~\ref{sec:over-indebt}) relies on: a client's open-loan count across the platform is publicly readable from their SBT, so simultaneous-loan caps are enforceable without an external credit bureau. Third, the schema is intentionally minimal (\texttt{creditScore}, \texttt{openLoans}, \texttt{completedCycles}, \texttt{lastDefault}, \texttt{riskTier}) so other lending protocols that adopt the standard can read the credential, allowing progressive lending to follow a client across protocols. MicroSave Consulting (December 2025)~[R31] documents that Bangladesh's microfinance sector is actively shifting toward performance-based credit identity; the SBT is the technical implementation of that shift on the Crypto World Bank platform.

\paragraph{Public \texttt{ICreditPassport} interface for cross-protocol composability.}
Although the SBT schema described above is readable on-chain, the credit passport is currently siloed: only contracts within the Crypto World Bank hierarchy can query it in a standardized way.  To make the credential composable with external protocols---such as third-party Local Banks in partner networks, DeFi lending markets, or impact-finance platforms---the contract exposes a public \texttt{ICreditPassport} interface.  The interface defines a single standardized view function:
\begin{verbatim}
function getScore(address wallet)
    external view
    returns (uint8 riskTier, uint32 repaymentCount, bool hasDefault);
\end{verbatim}
This read-only surface adds no new trust assumptions (no state mutation is permitted through the interface) and requires no additional gas for the issuing contract.  Any external protocol that imports \texttt{ICreditPassport} can read a client's credit tier and repayment history in one call, without needing a custom adapter or a centralized credit bureau query.  The cross-chain credit-passport mirroring described in Section~\ref{sec:bridge} propagates this same interface to each deployed chain, so the composability guarantee holds across the multi-chain deployment.  Platforms such as Spectral Finance and RociFi have demonstrated that on-chain credit scoring from wallet history can enable undercollateralized lending at scale~[R46]; the \texttt{ICreditPassport} interface positions the Crypto World Bank credit passport as a reusable primitive within that broader on-chain credit infrastructure.

\paragraph{GDPR and the permanent SBT default record.}
The permanent, non-revocable nature of the SBT default record raises a question under GDPR Article~17 (right to erasure). The resolution is as follows. The SBT stores only a wallet address, a risk tier integer, and a Unix timestamp---no name, national ID, or personally identifiable information (PII). The wallet address is pseudonymous; it does not constitute a directly identifiable personal data record under GDPR Article~4(1) unless combined with off-chain KYC data held by the Local Bank. The Local Bank's off-chain PostgreSQL records---which do constitute personal data---are subject to data-subject deletion requests. However, deleting the off-chain record does not delete the on-chain SBT, which is by design. The regulatory analogy is a credit bureau record that survives a borrower's request to be forgotten under credit-reporting exemptions in most jurisdictions. Platforms that adopt the CWB SBT standard should include this limitation in their privacy policy and terms of service, and should treat the wallet address as pseudonymous personal data under GDPR Article~89 (research and statistics exemption) where applicable. The architecture's data-minimization posture---only a hash of identity documents appears on-chain---is already compliant with the GDPR data-minimization principle (Article~5(1)(c)), and the off-chain PostgreSQL layer supports standard data-subject access and deletion workflows.

\paragraph{Credit tier schedule.}
The Credit Passport SBT encodes a five-tier credit schedule that governs maximum loan limits and interest modifiers. Table~\ref{tab:credit-tiers} gives the full schedule.

\begin{table}[H]
\centering
\caption{Credit tier schedule: score thresholds, maximum loan limits, and interest modifiers per tier.}
\label{tab:credit-tiers}
\renewcommand{\arraystretch}{1.3}
\small
\begin{tabular}{p{0.8cm}p{2.2cm}p{2.6cm}p{3.2cm}p{3.8cm}}
\hline
\textbf{Tier} & \textbf{Label} & \textbf{Score Range} & \textbf{Max Loan (USDC)} & \textbf{Interest Modifier} \\
\hline
1 & Bronze & 0--299 & 50 & Base rate \\
\hline
2 & Silver & 300--549 & 250 & Base $-$ 0.5\% \\
\hline
3 & Gold & 550--749 & 1,000 & Base $-$ 1.0\% \\
\hline
4 & Platinum & 750--899 & 5,000 & Base $-$ 1.5\% \\
\hline
5 & Diamond & 900--1000 & 25,000 & Base $-$ 2.0\% \\
\hline
\end{tabular}
\end{table}

The AI Agent reads the client's current tier via the \texttt{get\_credit\_score} MCP tool and proactively coaches the client on what is needed to reach the next tier. For example: ``You need 2~more on-time repayments to reach Gold tier. Your borrowing rate will then drop by 1.0\%.'' This coaching loop---executed automatically before every loan application confirmation---creates a positive feedback mechanism that incentivises repayment discipline and progressively unlocks higher credit limits, directly addressing the credit-building gap that MicroSave Consulting~[R31] identifies as the primary barrier to sustainable microfinance graduation.

\section{Cross-Chain Bridge Architecture}
\label{sec:bridge}

A multi-chain deployment (Polygon zkEVM Cardona + Ethereum Sepolia in the prototype, with mainnet successors planned) without cross-chain state consistency is a liability: if a client has an active loan on Polygon and attempts to open a second loan by routing through Ethereum, the hierarchical borrowing-limit control could be bypassed. v15 makes the bridge architecture explicit rather than leaving it as ``portable.''

\paragraph{Bridge protocol selection.}
The platform uses a single canonical bridge protocol for cross-chain messaging rather than a bespoke contract. The 2025--2026 leaders are LayerZero (largest message-volume share), Axelar (GMP for arbitrary payloads), and Chainlink CCIP (institutional-grade with native MEV protection). The current prototype targets Chainlink CCIP because the platform now uses Chainlink Functions for the oracle layer (Section~\ref{sec:oracle-architecture}), Chainlink Price Feeds for FX, and Chainlink Automation for installment triggers (all adopted in v24), which means adopting CCIP for bridging concentrates the entire oracle and cross-chain stack on a single audited provider and minimises new trust assumptions. Bridge hacks were the largest single category of crypto loss between 2022 and 2024; concentrating on one well-audited bridge rather than maintaining multiple is the correct security posture for an academic prototype.

\paragraph{Cross-chain hierarchy invariant.}
The hierarchical control is preserved across chains by requiring that all loan state for a given client live on a single chain (the chain of the originating Local Bank). Cross-chain messages are restricted to (a) reserve-ratio updates from World Bank Reserve to National Banks across chains, and (b) credit-passport SBT mirroring across chains so that a client's credential is readable on any chain even though their loans are not. This pattern keeps borrowing-limit enforcement single-chain and avoids the consensus problem of cross-chain debt.

\section{Multi-Entity and Cross-Tier Capital Operations}
\label{sec:multi-entity}

\begin{figure}[H]
\centering
\OnePageDiagram{fig-multi-entity-ops.pdf}
\caption{Multi-entity and cross-tier capital operations: (a) InterBankLendingPool with utilization-kinked rate model and default cascade, (b) UpwardDepositFacility with asymmetric rate structure, (c) SyndicatedLoan with Lead Arranger and Co-Lenders, (d) TranchedPool with senior--junior waterfall, (e) TreasurySwap for cross-tier asset exchange, and (f) NettingEngine for multilateral settlement.}
\label{fig:multi-entity-ops}
\end{figure}


The Chapter~1 framing of multi-directional flows (Section~\ref{sec:research-contribution}) and Section~1.7 promise downward, same-tier, and upward lending plus tiered client access. v10 left same-tier and upward flows at design-level only. A paper that claims a \emph{complete} banking system must specify these flows---and the genuinely multi-entity operations they enable---at the same depth as the downward flow. This section closes that gap with six concrete contract-and-flow specifications: a fully-specified InterBankLendingPool (Section~\ref{sec:interbank-pool}), upward surplus repatriation (Section~\ref{sec:upward-flow}), syndicated / club lending (Section~\ref{sec:syndicated-lending}), senior--junior tranched lending pools (Section~\ref{sec:tranched-pools}), cross-tier treasury FX swap (Section~\ref{sec:treasury-fx}), and a multilateral settlement netting engine (Section~\ref{sec:netting-settlement}). The corresponding ERD entities, contract list, Sprint plan, and Future Work are updated downstream so the database design, software-engineering stages, and architecture remain a single consistent system.

\subsection{InterBankLendingPool: Fully Specified Same-Tier Lending}
\label{sec:interbank-pool}

\paragraph{Pool placement.}
One \texttt{InterBankLendingPool} (IBLP) contract is deployed per hierarchical tier: \texttt{IBLP\_NB} at the National-Bank tier and \texttt{IBLP\_LB} at the Local-Bank tier. Membership in a tier's IBLP is gated by the on-chain role grant from the parent tier (a Local Bank can only deposit into \texttt{IBLP\_LB} if its \texttt{LOCAL\_BANK\_ROLE} is active and unexpired---Section~\ref{sec:defense-in-depth}). The pool denominates positions in the platform's stablecoin numeraire (USDC per Section~\ref{sec:stablecoin-first}); ETH-denominated positions are accepted at higher tiers but cleared through the \texttt{TreasurySwap} contract (Section~\ref{sec:treasury-fx}) before pool entry.

\paragraph{Rate model.}
The interbank borrowing rate $r_{\text{IB}}(U)$ follows a kinked utilization curve identical in shape to the retail kinked rate (Section~\ref{sec:kinked-rate}) but with two distinguishing properties. First, the slope coefficients are calibrated to the SOFR-style overnight market: the base rate $r_0^{\text{IB}}$ is set to the platform's deposit-yield benchmark, and the kink point $U^{*}_{\text{IB}} = 0.90$ is higher than the retail kink (because banks have lower default risk and tighter operational reserves than retail borrowers). Second, the formula is upper-bounded by the tier's \emph{downward} lending rate from the tier above:
\[
r_{\text{IB}}(U) \le r_{\text{downward}}(U) - \delta,
\]
where $\delta$ is the governance-set risk premium spread that prevents a Local Bank from arbitraging the tier above by borrowing from \texttt{IBLP\_LB} at a cheaper rate than from its parent National Bank. This bound is enforced on-chain inside the rate-setter function via a \textbf{live cross-contract call} rather than a static governance-cached parameter:
\begin{verbatim}
require(
    proposedRate <= localBank.getCurrentBorrowRate() - delta,
    "Rate exceeds downward bound"
);
\end{verbatim}
Reading \texttt{localBank.getCurrentBorrowRate()} at the moment the rate-setter function is called ensures the bound is evaluated against the \emph{current} downward rate, which itself shifts dynamically with pool utilization under the kinked model of Section~\ref{sec:kinked-rate}.  A static cached parameter would allow a window of arbitrage between successive governance updates during which $r_{\text{IB}}$ could exceed the updated $r_{\text{downward}} - \delta$; the live call closes this gap atomically.  The additional cross-contract read adds one \texttt{STATICCALL} to the rate-setter gas cost but imposes no cost on routine borrow/repay operations, since the bound is checked only when a rate change is proposed.

\paragraph{Settlement and default cascade.}
Loans through the IBLP are short-tenor (1-, 7-, 30-day) and are settled atomically: the borrowing bank receives the stablecoin amount in the same transaction in which its repayment obligation is recorded. If the borrowing bank defaults at maturity, the default cascade is: (1) the bank's reserve buffer above the minimum ratio is auto-debited; (2) if insufficient, the \texttt{InsuranceFund} (Section~\ref{sec:savings-vault}, with its claim mechanism extended in Section~\ref{sec:upward-flow}) is drawn; (3) if still insufficient, the parent tier (the bank's National Bank, or the World Bank for an NB default) absorbs the residual as a regulatory backstop and a governance vote is automatically queued through the TimeLock for the defaulting bank's role suspension. Every step emits a typed event indexed by The Graph (Section~\ref{sec:realtime}) so the cascade is fully auditable.

\paragraph{Operations.}
\texttt{IBLP\_*} exposes four user functions and three governance functions:

\begin{itemize}[noitemsep]
\item \texttt{deposit(uint256 amount, uint8 maturityCode)} --- surplus bank supplies liquidity.
\item \texttt{borrow(uint256 amount, uint8 maturityCode)} --- deficit bank draws against pool, gated by per-bank credit ceiling, role expiry, and the IBLP-vs.-downward-rate bound.
\item \texttt{repay(uint256 borrowId)} --- atomic principal + interest settlement.
\item \texttt{withdraw(uint256 supplyId)} --- surplus bank reclaims principal + accrued interest.
\item \texttt{setRateParams(...)} --- TimeLock-gated (24--48~h) governance update.
\item \texttt{setMembershipPolicy(...)} --- role-expiry / per-bank ceiling adjustment.
\item \texttt{triggerCascade(uint256 borrowId)} --- public function any caller can invoke once a borrow is overdue; pays the caller a small bonus from the defaulter's reserve buffer, mirroring the Aave liquidation incentive (Section~\ref{sec:liquidation}).
\end{itemize}

\subsection{Upward Surplus Repatriation}
\label{sec:upward-flow}

The upward direction is the dual of downward lending: a Local Bank whose reserve buffer exceeds the minimum ratio by more than a governance-set ``excess threshold'' (default 20\,\% above the minimum) may deposit the excess into its parent National Bank, earning the parent's \emph{deposit yield}. A National Bank with similarly excess capital may deposit into the World Bank Reserve under the same logic. Three controls keep the upward flow safe:

\begin{enumerate}[noitemsep]
\item \textbf{Asymmetric rate structure.} The upward deposit yield is strictly lower than the downward borrowing rate ($r_{\text{up}} < r_{\text{down}} - \delta$), creating a natural term structure across the hierarchy and removing the round-trip arbitrage incentive of borrowing-down to deposit-up.
\item \textbf{Withdrawal gate.} The depositing bank may withdraw its upward deposit on demand, but the withdrawal triggers a reserve-ratio recomputation at the depositing bank's own tier first; if the withdrawal would push the depositor below its minimum reserve ratio, the call reverts. This protects retail depositors at lower tiers from a liquidity-squeeze upstream.
\item \textbf{Daily withdrawal rate limit (bank-run circuit breaker).} To prevent a coordinated bank-run scenario in which multiple National Banks simultaneously withdraw upward deposits, the \texttt{UpwardDepositFacility} enforces a per-National-Bank daily withdrawal ceiling of 20\% of that bank's total outstanding upward deposit liabilities. Withdrawals exceeding this ceiling within a 24-hour window are queued and executed in the next settlement window. This mirrors the standard withdrawal-rate limits applied by central bank standing deposit facilities and is enforced on-chain via a per-bank \texttt{withdrawnToday[bank]} counter that resets at each UTC midnight checkpoint.
\item \textbf{Recursive cap.} A given bank's net upward exposure (sum of upward deposits less own borrowings) is capped at 30\,\% of its on-chain assets, so a single failing bank cannot drain the entire upstream tier.
\end{enumerate}

The on-chain contract is named \texttt{UpwardDepositFacility} (one instance per parent tier); it is a thin extension of the SavingsVault contract (Section~\ref{sec:savings-vault}) restricted by role, with separate accounting registers (\texttt{upwardSupplied[bank]}, \texttt{upwardYieldOwed[bank]}) so that upward and retail flows do not commingle in the same reserve pool.

\paragraph{FATF Travel Rule compliance for inter-tier capital flows (Future Work).}
The Financial Action Task Force Travel Rule (FATF Recommendation~16) requires that transfers above USD~1,000 in most jurisdictions include originator and beneficiary information alongside the transfer. For CWB's inter-tier capital flows --- Local Bank $\to$ National Bank $\to$ World Bank through the \texttt{UpwardDepositFacility} and \texttt{InterBankLendingPool} contracts --- this requirement applies to any single transfer above the threshold. The off-chain Travel Rule data packet accompanying such a transfer takes the following structure:

\begin{verbatim}
{
  "originator_institution": "LocalBank-BD-042",
  "originator_tier": 3,
  "beneficiary_institution": "NationalBank-BD-001",
  "beneficiary_tier": 2,
  "amount_usdc": 5000,
  "purpose": "IBLP_REPAYMENT",
  "onchain_tx_hash": "0xABC...",
  "timestamp": "2026-06-01T10:00:00Z"
}
\end{verbatim}

This packet is stored in the PostgreSQL \texttt{audit\_logs} table and is available to regulators via the audit request workflow (Section~\ref{sec:regulatory-compliance}). FATF Travel Rule compliance is scoped to Future Work for a Bangladesh deployment; the data structure specification above constitutes the thesis contribution, while implementation of the cross-tier notification protocol is deferred pending jurisdiction-specific regulatory guidance from Bangladesh Bank's Financial Intelligence Unit.

\subsection{Syndicated and Club Lending}
\label{sec:syndicated-lending}

For loan sizes that exceed a single bank's prudent exposure ceiling---typically the 1\,000\,+\,ETH bracket at the National-Bank tier (Table~\ref{tab:borrower-tiers})---multiple banks at the same or adjacent tier co-fund a single loan through a \texttt{SyndicatedLoan} contract. This is the on-chain analogue of a real-world syndicated loan or club deal and is the answer to the supervisor's question about ``multiple entities investing, funding, lending together.''

\paragraph{Roles.}
A syndicated loan has three distinct on-chain roles: a \textbf{Lead Arranger} (the bank that structures and underwrites the deal), \textbf{Co-Lenders} (the participating banks supplying capital), and the \textbf{Credit Client} (a Local Bank, National Bank, or institutional client per Table~\ref{tab:borrower-tiers}). The Lead Arranger holds a small underwriting fee (10--20~basis points) plus a portion of the interest spread; Co-Lenders earn pro-rata interest by capital share.

\paragraph{Lifecycle.}

\begin{enumerate}[noitemsep]
\item \textbf{Proposal.} Lead Arranger calls \texttt{proposeSyndicate(uint256 totalAmount, uint8 borrowerTier, uint256 minSubscription)} which mints a syndicate identifier and emits a \texttt{SyndicateProposed} event. Off-chain, the deal terms (covenants, schedule) are pinned to IPFS and the content hash is stored in \texttt{syndicateDocHash[id]} for immutable cross-reference.
\item \textbf{Subscription window.} Co-Lenders call \texttt{subscribe(uint256 syndicateId, uint256 commitment)} within a governance-set window (default 7 days). Their commitment is escrowed in the contract.
\item \textbf{Funding threshold.} If aggregate commitments meet $\ge$ 90\,\% of \texttt{totalAmount} by window close, the deal funds; otherwise commitments are returned and the proposal expires. This avoids the partial-funding race condition.
\item \textbf{On-chain consent vote.} Each Co-Lender's wallet must call \texttt{confirmTerms(uint256 syndicateId)} after subscription. Disbursement requires confirmation from Co-Lenders representing at least 75\% of the subscribed capital by value (supermajority, not unanimity), preventing a griefing attack in which a small Co-Lender subscribes and then deliberately withholds confirmation to block the deal at minimal cost. The Lead Arranger may also reject a Co-Lender's subscription before the window closes. A Co-Lender that subscribes but does not confirm forfeits a bond proportional to its committed capital (1\% of committed amount), making non-confirmation economically costly at scale.
\item \textbf{Disbursement.} On supermajority confirmation (75\%-by-capital-share threshold met), the contract atomically transfers the aggregate to the borrower in a single transaction. Each Co-Lender's exposure share is recorded as \texttt{shareBps[lender][syndicateId]}.
\item \textbf{Interest accrual + repayment.} Client installments flow into the \texttt{SyndicatedLoan} contract; the contract distributes principal + interest pro-rata to each Co-Lender on every installment via the basis-points share registry. The Lead Arranger receives its underwriting fee on the first repayment.
\item \textbf{Default handling.} On default, the LiquidationEngine (Section~\ref{sec:liquidation}) seizes collateral; recovery is distributed pro-rata to Co-Lenders by the same basis-points share. A Co-Lender's loss is therefore bounded by its commitment, which is the desired multi-entity risk-sharing property.
\end{enumerate}

\paragraph{Cross-tier syndication.}
The same contract supports \emph{cross-tier} syndication when permitted by governance: one National Bank and three of its child Local Banks may co-fund a single large project loan where the NB is Lead Arranger. This is the genuine ``lending from upper tiers as well as lower tier together'' pattern the supervisor asked about; the role-grant logic of the \texttt{SyndicatedLoan} contract simply permits any wallet holding \texttt{NATIONAL\_BANK\_ROLE} or \texttt{LOCAL\_BANK\_ROLE} (and any \texttt{WORLD\_BANK\_ADMIN}-approved institutional client) to subscribe.

\subsection{Senior--Junior Tranched Lending Pools}
\label{sec:tranched-pools}

For risk-segmented co-funding, the platform supports the senior--junior tranche pattern established by Goldfinch~[32] and Maple~[31]. The \texttt{TranchedPool} contract divides depositor capital into two notional tranches with distinct seniority:

\begin{itemize}[noitemsep]
\item \textbf{Senior tranche:} lower yield, first-loss protection, suitable for risk-averse depositors (Local Banks redepositing surplus, or institutional partners).
\item \textbf{Junior tranche:} higher yield, takes first loss on default, suitable for risk-tolerant capital (governance-token-aligned depositors, the protocol treasury, or impact-investment partners).
\end{itemize}

\paragraph{Waterfall payment.}
On each repayment, the contract pays from the inflow in the order: (i) senior accrued interest, (ii) junior accrued interest, (iii) senior principal pro-rata, (iv) junior principal pro-rata. On default the order reverses for losses: the junior tranche is written down first; only when junior principal reaches zero is senior principal touched. The subordination ratio (junior / senior) is fixed at deposit time and ranges from 10\,/\,90 (institutional-grade) to 30\,/\,70 (impact-aligned) per the policy table set by the World Bank Admin.

\paragraph{Stress test.}
A 30\,\% default-rate stress test (matching the agent-based simulation parameter of Section~\ref{sec:abm-sim}) on a $10:90$ junior:senior pool would zero out the junior tranche but preserve the senior tranche fully; on a $30:70$ pool a 30\,\% default writes the junior down to zero and senior takes a 0\,\% nominal loss. These dynamics are formalized as Foundry invariants (Section~\ref{sec:foundry}) to give automatic regression detection.

\subsection{Cross-Tier Treasury FX Swap}
\label{sec:treasury-fx}

The FXModule mentioned in Section~\ref{sec:banking-products} handles end-user FX. Cross-tier treasury FX---the case where a National Bank holding an ETH treasury must source USDC stablecoin to fund a Local Bank that has elected stablecoin denomination---is a distinct operation served by the \texttt{TreasurySwap} contract.

\paragraph{Mechanism.}
\texttt{TreasurySwap} is an oracle-priced atomic-swap contract restricted to wallets holding a banking role. A National Bank calls \texttt{swap(asset\_from, asset\_to, amount)}; the contract reads the Chainlink price feed, applies the governance-set treasury spread (typically 5--10~basis points, tighter than retail FX), and executes the swap atomically against either a counterparty bank's posted offer or the protocol's own \texttt{TreasuryAMM} liquidity pool. The Chainlink price reading and the swap execution sit inside the same transaction; there is no intermediate state and no settlement risk.

\paragraph{Hierarchy invariant.}
Treasury swaps cannot push the swapping bank below its minimum reserve ratio (the swap function recomputes the post-swap ratio before any state change). Swaps above a governance-set threshold (default \$1~million-equivalent) require a two-signature confirmation from the bank's Safe multisig. Each swap emits a \texttt{TreasurySwapExecuted} event with both legs of the swap, the oracle reading used, and the spread, supporting both internal audit and regulatory reporting.

\subsection{Multilateral Settlement Netting Engine}
\label{sec:netting-settlement}

A four-tier hierarchical bank with same-tier interbank flows accumulates many small inter-bank obligations during a settlement period. Settling each one as a separate on-chain transaction is wasteful; the \texttt{NettingEngine} contract compresses them.

\paragraph{Workflow.}
At the end of each settlement window (1 hour, 1 day, or governance-set), an off-chain Settlement Coordinator computes a netted obligation matrix from the queued IBLP and TreasurySwap orders. The matrix is committed on-chain as a Merkle root, and a single \texttt{settleBatch(bytes32 batchRoot, bytes calldata proofs)} call executes the entire batch in one transaction. Each bank's net debit/credit is applied to its on-chain balance, and the constituent transactions are marked settled. This pattern follows the Hyperledger Cactus and BIS Project Agora cross-ledger netting designs~[R39, R40].

\paragraph{Why this matters.}
For a network with $n$ active banks per tier, naive bilateral settlement requires $O(n^2)$ transactions per period; multilateral netting compresses this to $O(n)$ at the cost of one off-chain coordination step. At full Local-Bank deployment ($n = 50$) this is a 25$\times$ reduction in on-chain settlement footprint per settlement window. The trade-off is one explicit trust assumption: the Settlement Coordinator role must be held by a Safe multisig (the same WORLD\_BANK\_ADMIN Safe), and any participating bank may dispute a batch within a challenge window by calling \texttt{disputeBatch(uint256 batchId, bytes proof)}.

\paragraph{Trust assumption disclosure.}
The Settlement Coordinator is a trusted central operator for the netting computation---an acknowledged trust regression relative to the platform's otherwise trustless design. This is the same trade-off made by real-world clearinghouses (e.g., CLS, SWIFT GPI netting): they are trusted institutions operating under legal agreement, not trustless systems. The CWB NettingEngine adopts the same honest posture: netting is scoped to the permissioned institutional tier (NBs and LBs that have signed legal agreements with each other), not applied to the retail client tier. A future upgrade path toward ZK-proven netting (where a zk-SNARK proves the sum-of-obligations-equals-zero property without revealing individual pair amounts) is noted in Future Work.

\paragraph{Database mirror.}
Each \texttt{NettingBatch} is mirrored off-chain to a PostgreSQL \texttt{netting\_batch} table with one \texttt{netting\_entry} row per netted obligation, allowing dashboard display, regulator export, and the agent-based simulation (Section~\ref{sec:abm-sim}) to read historical settlement-compression ratios.

\paragraph{Settlement failure path.}
The current \texttt{settleBatch} specification assumes every participant in a netting cycle can fully meet their net debit obligation.  In practice, a bank may be underfunded at settlement time---its on-chain balance is insufficient to cover its net debit position.  Without an explicit failure path, a single underfunded participant can deadlock the entire netting cycle, leaving all other participants' net credits unprocessed.  The \texttt{NettingEngine} therefore exposes a \texttt{settlePartial(uint256 cycleId)} fallback function that is callable by the Settlement Coordinator (or any participant) once the standard settlement window has elapsed without full completion.  \texttt{settlePartial} iterates the netted obligation matrix in decreasing net-credit order: participants with sufficient balances are settled immediately and marked \texttt{SETTLED}; participants whose on-chain balance is below their net debit are marked \texttt{PENDING\_DEFAULT}, their net debit is recorded in a \texttt{pendingObligations[bank]} queue, and the Settlement Coordinator emits a \texttt{SettlementDefaultFlagged(cycleId, bank, shortfall)} event.  The queued obligation is resolved in the next settlement cycle or through the IBLP default cascade (Section~\ref{sec:interbank-pool}).  This partial-fill design mirrors the settlement-failure handling in real-world clearing systems such as CLS and is consistent with the platform's existing default-cascade model; it does not change the trust assumptions around the Settlement Coordinator role.

\subsection{Database, Sprint, and Contract Consistency for Multi-Entity Operations}
\label{sec:multi-entity-consistency}

To keep the data model, software-engineering stages, and architecture coherent, the following entities, sprint deliverables, and contracts are added downstream of this section:

\paragraph{ERD entities added (data model in Chapter~3 ERD).}
\texttt{INTERBANK\_LOAN}~\{loan\_id (PK), tier, lender\_bank\_id (FK), borrower\_bank\_id (FK), principal, maturity\_code, rate\_bps, status, created\_at, settled\_at\}; \texttt{UPWARD\_DEPOSIT}~\{deposit\_id (PK), depositing\_bank\_id (FK), parent\_bank\_id (FK), principal, yield\_owed, created\_at\}; \texttt{SYNDICATE}~\{syndicate\_id (PK), lead\_arranger\_id (FK), borrower\_id, total\_amount, status, doc\_hash, opened\_at, funded\_at\}; \texttt{SYNDICATE\_MEMBER}~\{syndicate\_id (FK), lender\_bank\_id (FK), commitment, share\_bps, confirmed\_at\}; \texttt{TRANCHED\_POOL}~\{pool\_id (PK), local\_bank\_id (FK), borrower\_id, senior\_principal, junior\_principal, subordination\_bps, status\}; \texttt{TREASURY\_SWAP}~\{swap\_id (PK), bank\_id (FK), asset\_from, asset\_to, amount\_from, amount\_to, oracle\_reading, spread\_bps, executed\_at\}; \texttt{NETTING\_BATCH}~\{batch\_id (PK), tier, coordinator\_id, batch\_root, opened\_at, settled\_at\}; \texttt{NETTING\_ENTRY}~\{batch\_id (FK), src\_bank\_id, dst\_bank\_id, net\_amount\}. These are appended to the ERD prose at the end of Chapter~3 alongside SavingsAccount, FixedDeposit, LoanGroup, GroupMember, CurrentAccount, and InsuranceFund, and will be drawn into the updated ERD figure in the diagram-audit session.

\paragraph{Smart contracts added (Appendix~B).}
\texttt{InterBankLendingPool} (one per tier), \texttt{UpwardDepositFacility}, \texttt{SyndicatedLoan}, \texttt{TranchedPool}, \texttt{TreasurySwap}, \texttt{NettingEngine}. The complete architecture therefore grows from nine to fifteen modular contracts.

\paragraph{Sprint plan update (Chapter~1, Methodology in Brief).}
Sprint~2 takes ownership of \texttt{InterBankLendingPool}, \texttt{UpwardDepositFacility}, and \texttt{TreasurySwap}; Sprint~3 takes \texttt{SyndicatedLoan}, \texttt{TranchedPool}, and \texttt{NettingEngine} (the latter requires the off-chain Settlement Coordinator service, which sits alongside the AI/ML pipeline already scheduled for Sprint~3). The Foundry invariant suite (Section~\ref{sec:foundry}) is extended in Sprint~3 with three additional invariants: (i) IBLP $r_{\text{IB}} < r_{\text{down}} - \delta$ at all times; (ii) Syndicate disbursement = sum of confirmed commitments; (iii) NettingBatch sum of net debits = sum of net credits.

\section{Reentrancy and Security Analysis}

\textbf{Checks-Effects-Interactions (CEI) pattern.} The CEI pattern is applied to all state-mutating functions in the lending contracts. The three functions most vulnerable to reentrancy are:

\begin{enumerate}[noitemsep]
\item \textbf{\texttt{disburseLoan(address borrower, uint256 amount)}} --- sends ETH to the borrower. Mitigation: update \texttt{loanStatus[borrower] = LoanStatus.ACTIVE} \emph{before} calling \texttt{payable(borrower).transfer(amount)}, and apply \texttt{nonReentrant}.
\item \textbf{\texttt{processInstallment(uint256 loanId)}} --- updates the repayment schedule. CEI order: check $\to$ mark installment paid $\to$ emit event $\to$ release interest share.
\item \textbf{\texttt{allocateCapital(address nationalBank, uint256 amount)}} --- cross-tier ETH transfer. Apply both \texttt{nonReentrant} and an explicit \texttt{require(allocatedTo[nationalBank] + amount <= maxAllocation)} check before any state change.
\end{enumerate}

The DAO hack (2016, $\approx$\$60~million) and Curve Finance reentrancy exploit (2023, $\approx$\$70~million via Vyper compiler bug) demonstrate that reentrancy is not a solved problem in DeFi, even with guards~[R6]. Formal verification with Certora or Mythril is planned for the final thesis security audit.

\textbf{Flash loan scope.} Flash loan attacks primarily exploit price-sensitive logic, specifically protocols that use on-chain pool spot prices for collateral valuation~[R5]. In the current prototype, no price-sensitive logic exists because stablecoin integration is not yet implemented, collateral valuation is not automated on-chain, and oracle feeds have not yet been integrated. Therefore, the current contracts do not present a meaningful flash loan attack surface. Once stablecoin integration and oracle-based collateral pricing are implemented (planned for the final thesis phase), mitigations must include: (1) TWAP oracles rather than spot price reads, (2) multi-block confirmation requirements for large collateral updates, and (3) Chainlink price feeds with circuit breakers.

\textbf{Upgradeability via UUPS (ERC-1822).} All three implemented contracts (and the six planned ones) use the OpenZeppelin Upgradeable contracts in the \textbf{UUPS (Universal Upgradeable Proxy Standard)} pattern rather than the Transparent Proxy pattern. UUPS is cheaper to deploy (no separate ProxyAdmin contract) and the upgrade authorization logic lives in the implementation contract, making it directly governable via the existing RBAC. \texttt{\_authorizeUpgrade} is gated on \texttt{onlyRole(WORLD\_BANK\_ADMIN)}, which in production is held by a Safe 3-of-5 multisig (Section~\ref{sec:defense-in-depth}). Without upgradeability, any post-deployment bug would require redeploying from scratch and migrating all on-chain state; with UUPS the team can patch a bug and propose the upgrade for multisig approval in minutes, with a full on-chain audit trail.

\textbf{TimeLock on governance actions.} Every privileged action that changes a system-level parameter---interest rates, reserve ratios, role grants at or above the National Bank tier, contract upgrades---routes through a minimum 24-to-48-hour \texttt{TimelockController} (OpenZeppelin) delay. This is the standard pattern used by Compound, Aave, MakerDAO, and every serious DeFi protocol since at least 2021. It provides a defensible window in which a malicious proposal can be detected and cancelled before execution. For the university lab demo specifically, the delay is shortened to five minutes purely so examiners can witness a governance proposal queued and executed live; the production delay is set at deploy time and cannot be reduced below the minimum without a TimeLock-gated governance vote of its own.

A \textbf{dual-governance path} is specified to handle security emergencies. The standard path uses the 24--48~h TimeLock for all planned parameter changes and upgrades. An emergency path uses a \textbf{Security Council} --- a 4-of-7 Safe multisig composed of members distinct from the \texttt{WORLD\_BANK\_ADMIN} --- that can execute emergency upgrades within a 2-hour window without the TimeLock, subject to a mandatory on-chain post-incident disclosure within 48 hours. The Security Council can only execute pre-approved emergency action types (pause, upgrade, role revocation); it cannot change system parameters. This mirrors the dual-path model of Aave (Guardian address + ShortExecutor) and MakerDAO (Emergency Shutdown Module + Security Council), which experience shows to be necessary when a reentrancy bug or oracle manipulation is found in production and 24 hours of exposure can drain the entire reserve.

\textbf{Role expiry timestamps.} Role grants in the system are not permanent. Every \texttt{grantRole} call records an expiry block number; every privileged function call checks \texttt{require(roleExpiryBlock[msg.sender][role] > block.number)}. Block numbers are used rather than \texttt{block.timestamp} to avoid the SWC-116 vulnerability, in which validators on Polygon PoS can manipulate \texttt{block.timestamp} within a $\pm$15-second window---potentially extending the effective validity of an expired role. Block numbers are not manipulable in the same way. Off-chain display converts block numbers to approximate calendar dates using the expected block time. The client KYC level is bound to a one-year expiry by default; bank operator roles are bound to the institution's current registration validity; the World Bank Admin role itself is renewed annually by a governance vote so that even the top-of-hierarchy authority is not effectively permanent.

\textbf{Granular per-function pause.} The naive ``\texttt{Pausable} on the whole contract'' pattern is too coarse: pausing the entire World Bank Reserve when only one Local Bank has an incident causes unnecessary disruption to every other tier. The platform uses a per-function pause registry (\texttt{functionPaused[bytes32 functionId]}) so that governance can suspend only \texttt{LOAN\_DISBURSEMENT} while keeping \texttt{DEPOSITS} and \texttt{WITHDRAWALS} open. This is also auditable: each pause emits a \texttt{FunctionPaused} event with the actor and the timestamp, indexed by The Graph (Section~\ref{sec:realtime}).

\paragraph{Failed-attempt lockout and address whitelisting.}
After 3 failed confirmation responses (unclear or ambiguous consent) in a single agent conversation session, the agent pauses all interactions for 10~minutes and logs the incident to \texttt{AGENT\_ACTION\_LOG} with \texttt{status = FAILED}. After a detected command injection attempt --- an adversarial prompt that attempts to bypass the confirmation gate or impersonate a prior consent turn --- the session is immediately terminated and logged. Loan disbursement destination addresses are fixed at KYC registration time and cannot be changed without a 24-hour delay plus 2FA re-verification, preventing address-substitution attacks in which a compromised session attempts to redirect a disbursement to an attacker-controlled wallet.

\section{System Modeling}
\label{sec:system-modeling}

This section presents the system analysis diagrams that model the platform's structure, behavior, and data flow.

\subsection{Use Case Diagram}
\label{sec:usecase}

The system involves nine actors (five primary, four secondary) across 20 representative use cases (Figure~\ref{fig:usecase}). The five primary actors (A1--A5) are: Retail Client (three sub-types by KYC tier), Local Bank Admin, National Bank Admin, World Bank Admin, and AI Agent. The four secondary actors are: Regulatory Authority (A6), Chainlink DON (A7), Blockchain Validator (A8), and External Auditor (A9). A5, the AI Agent, is modelled as a primary actor with restricted write permissions gated by the human confirmation protocol; A6--A9 are modelled as secondary actors. Note that the primary actor classes in the use-case diagram are a condensation of the formal nine-actor taxonomy defined in Table~\ref{tab:actor-taxonomy}: Visitor and Registered Retail User are pre-conditions for the CLIENT (BORROWER) role rather than independent diagram actors, and Local Bank Operator / Approver are merged into the Bank Approver role at the diagram level.

\begin{figure}[H]
\centering
\OnePageDiagram{fig-usecase-actors.pdf}
\caption{Use-case diagram reflecting the nine-actor taxonomy of Table~\ref{tab:actor-taxonomy}: five primary actors (Retail Client A1, Local Bank Admin A2, National Bank Admin A3, World Bank Admin A4, AI Agent A5) and four secondary actors (Regulatory Authority A6, Chainlink DON A7, Blockchain Validator A8, External Auditor A9), mapped to 20 representative use cases covering KYC, deposits, lending, AI-agent banking operations, multi-entity operations, and regulator oversight.}
\label{fig:usecase}
\end{figure}

\subsection{Activity Diagrams}

Figures~\ref{fig:act-lending}--\ref{fig:act-aux} present the platform's seven core operational flows. Related flows are grouped onto a single figure page to keep the layout dense and readable.

\begin{figure}[H]
\centering
\OnePageDiagram{fig-activity-lending.pdf}
\caption{Lending activity flows: (a) loan request to repayment, (b) hierarchical capital flow from World Bank Reserve to Local Bank, (c) borrowing-limit enforcement via the credit-passport SBT.}
\label{fig:act-lending}
\label{fig:act-loan}\label{fig:act-hierarchy}
\end{figure}

\begin{figure}[H]
\centering
\OnePageDiagram{fig-activity-onboarding-id.pdf}
\caption{Onboarding and identity flows: (a) income verification (Open-Banking pull or officer review), (b) profile management with EIP-712 signed updates, (c) tiered, risk-based KYC ladder from zkKYC (L1) to entity onboarding (L4).}
\label{fig:act-onboarding}
\label{fig:act-income}\label{fig:act-profile}
\end{figure}

\begin{figure}[H]
\centering
\OnePageDiagram{fig-activity-aux.pdf}
\caption{Auxiliary activity flows: (a) client--bank chat over authenticated WebSocket, (b) RAG-augmented AI chatbot with hallucination guard, (c) live market-data retrieval and pre-filled loan sizing.}
\label{fig:act-aux}
\label{fig:act-chat}\label{fig:act-aichatbot}\label{fig:act-market}
\end{figure}

\paragraph{SAR activity flow.}
The Suspicious Activity Report (SAR) workflow is triggered when the Isolation Forest anomaly detector (Section~\ref{sec:threat-model}) flags a wallet with an anomaly score above the 0.75 threshold. The five-step flow is:

\begin{enumerate}[noitemsep]
\item Isolation Forest flags the wallet: \texttt{anomaly\_score > 0.75}.
\item \texttt{AI\_ML\_LOG} records the detection event (wallet, score, timestamp, feature vector).
\item Express.js backend emits Kafka topic \texttt{aml-alert}; the event is consumed by the compliance dashboard.
\item Bank officer reviews the alert in the admin compliance queue:
    \begin{itemize}[noitemsep]
    \item \textbf{False positive:} Dismiss and document reason (audit trail in \texttt{audit\_logs}).
    \item \textbf{Confirmed:} Generate SAR record in \texttt{audit\_logs}; notify the tier above (National Bank); freeze wallet via \texttt{LocalBank.freezeAccount(clientId)}.
    \end{itemize}
\item The \texttt{freezeAccount} function is gated by the \texttt{onlyApprover} modifier and is Foundry-tested (Sprint~3 invariant suite). A frozen wallet cannot initiate new loan disbursements or installment payments until the freeze is lifted by the tier above.
\end{enumerate}

\subsection{Data Flow Diagrams}

Figure~\ref{fig:dfd-suite} consolidates the Level-0 context diagram with the Level-1 decompositions of the core lending subsystem and the deposit / interbank / FX / netting subsystems on a single page.

\begin{figure}[H]
\centering
\OnePageDiagram{fig-dfd-suite.pdf}
\caption{Data-flow diagrams: Level-0 context view of the Crypto World Bank platform with five external entities; Level-1 decomposition of the core lending subsystem (origination, ML scoring, disbursement, installment engine, and their data stores); and Level-1 decomposition of the deposit-mobilization, InterBankLendingPool, TreasurySwap, and NettingEngine subsystems.}
\label{fig:dfd-suite}
\label{fig:dfd-context}\label{fig:dfd-level1a}\label{fig:dfd-level1b}
\end{figure}

\subsection{Sequence Diagrams}

Figures~\ref{fig:seq-loan-flow}--\ref{fig:seq-chat-bot} consolidate the platform's nine interaction flows into four panels: loan approval with reject alternative, installment and income, hierarchical banking with market data and borrowing-limit, and chat with AI-chatbot.

\begin{figure}[H]
\centering
\OnePageDiagram{fig-seq-loan-flow.pdf}
\caption{Sequence diagram for the loan-approval flow with reject alternative: the client submits an EIP-712 signed request, the ML oracle returns a SHAP-explained score via commit-reveal, and the approver either approves (revealing the commit) or rejects (with reason logged on-chain).}
\label{fig:seq-loan-flow}
\label{fig:seq-loan}\label{fig:seq-reject}
\end{figure}

\begin{figure}[H]
\centering
\OnePageDiagram{fig-seq-installment-income.pdf}
\caption{Sequence diagrams: (a) installment payment loop using the CEI pattern with SBT update on the final installment, and (b) income verification with Open-Banking pull or manual officer review.}
\label{fig:seq-installment-income}
\label{fig:seq-installment}\label{fig:seq-income}
\end{figure}

\begin{figure}[H]
\centering
\OnePageDiagram{fig-seq-banking-data.pdf}
\caption{Sequence diagrams: (a) hierarchical capital flow with reserve-ratio gates at every tier, (b) market-data retrieval via cached Chainlink feeds, (c) borrowing-limit calculation reading the credit-passport SBT before ML scoring.}
\label{fig:seq-banking-data}
\label{fig:seq-hierarchy}\label{fig:seq-marketdata}\label{fig:seq-borrowlimit}
\end{figure}

\begin{figure}[H]
\centering
\OnePageDiagram{fig-seq-chat-chatbot.pdf}
\caption{Sequence diagrams: (a) client--bank chat over an authenticated WebSocket, and (b) AI chatbot pipeline (ChromaDB top-$k$ retrieval, QLoRA-tuned LLM, hallucination guard, citation-anchored answer).}
\label{fig:seq-chat-bot}
\label{fig:seq-chat}\label{fig:seq-aichatbot}
\end{figure}

\subsection{Four-Tier Capital Flow}

Figure~\ref{fig:four-tier} illustrates the hierarchical capital flow and cascading repayment structure.

\begin{figure}[H]
\centering
\OnePageDiagram{fig-hierarchical-banking.pdf}
\caption{Four-tier hierarchical capital flow with cascading repayment, plus the same-tier interbank-lending pools (InterBankLendingPool, Section~\ref{sec:interbank-pool}) and the upward-surplus repatriation facility (UpwardDepositFacility, Section~\ref{sec:upward-flow}). The asymmetric rate structure ($r_{\text{up}} < r_{\text{down}}-\delta$) is enforced on-chain.}
\label{fig:four-tier}
\end{figure}


\section{Auxiliary Dual-Currency Facility}

The cryptocurrency exchange feature of Crypto World Bank is built upon the cryptocurrency exchange facility that is an auxiliary facility incorporated into the current banking infrastructures all over the world. Instead of being a separate exchange, the platform is a service that is provided by all participating banks as an additional service to the existing product portfolio.

\begin{itemize}[noitemsep]
\item \textbf{Integration model:} The dual-currency facility is offered by all participating banks as an add-on service to their existing product portfolio. No additional banking license or separate entity is needed.
\item \textbf{Eligibility determination:} Eligibility of the dual-currency service is based on the bank officers managing the lending operations, on the basis of existing KYC/AML compliance, account status, and lending relationship history.
\item \textbf{No defaulting of conditions:} The project does not override, modify, or default any existing banking conditions, regulatory requirements, or contractual obligations.
\item \textbf{Scope:} The facility enables fiat-to-crypto and crypto-to-fiat conversions, cryptocurrency-denominated lending and repayment, and transparent on-chain transaction records---all within the governance structure of the participating bank.
\end{itemize}

\section{Banking Product Suite}
\label{sec:banking-products}

\begin{figure}[H]
\centering
\OnePageDiagram{fig-banking-modules.pdf}
\caption{Auxiliary banking modules: (a) LiquidationEngine grace-period trigger and collateral auction with InsuranceFund shortfall fallback, (b) SavingsVault and FixedDeposit with duration-matched feeding of the lending pool, (c) on-chain Credit Passport (Soulbound Token) read by any participating bank, and (d) Cross-Chain Bridge via Chainlink CCIP carrying reserve-ratio updates and SBT mirroring.}
\label{fig:banking-modules}
\end{figure}


This section describes the banking products that the platform is designed to support across tiers. While the current prototype implementation focuses primarily on hierarchical lending, the complete system design includes deposit products, transactional accounts, solidarity lending, and multi-currency operations.

\subsection{Savings Products}

The platform offers savings products at the Local Bank tier and, by extension, at higher tiers for institutional participants. Standard savings accounts provide liquid deposits with variable yield tied to platform utilization, allowing deposit pricing to respond to market conditions rather than discretionary rate-setting. Fixed-term deposits lock capital for defined periods (e.g., 30/90/180/365 days) in exchange for a higher agreed yield written into the contract at deposit time. Early withdrawal can be permitted with an automatically enforced penalty (forfeiture of a portion of accrued interest), which is transparent and deterministic.

The yield on standard savings accounts is algorithmically determined by platform utilization: when loan demand is high and available liquidity is low, savings yields rise to attract deposits; when liquidity is abundant, yields fall. This creates a self-regulating capital cycle where savings and lending are linked through a single on-chain formula rather than through discretionary rate-setting committees.

Fixed-term deposits serve a distinct function in the banking architecture: they provide the Local Bank tier with predictable, duration-matched liabilities that can be deployed into longer-term installment loans without creating dangerous asset-liability mismatches. A depositor who locks funds for 180 days effectively funds the 6-month installment tranche of a retail client at that tier, eliminating the duration mismatch that causes bank runs in traditional fractional reserve systems.

The closed-loop argument is central to the platform's economic sustainability: unlike donor-funded microfinance models that depend on external capital injections, a savings-funded lending model recycles depositor capital through the lending pool continuously. Interest earned by the platform on loans is partially redistributed to depositors as yield, partially retained as protocol revenue, and partially allocated to the insurance fund, creating a self-sustaining three-way split that does not require new token issuance or external subsidy.

\subsection{Checking and Transactional Accounts}

In addition to interest-bearing savings products, the platform provides transactional accounts (checking/current-account equivalents) for day-to-day financial activity. Transactional accounts carry no yield but impose no lock-in, allow transfers between registered accounts, and serve as the primary account from which loan installments are debited and income receipts are credited. Transfers between accounts within the platform are settled atomically: a transaction either completes fully or reverts, eliminating intermediate settlement ambiguity.

The atomic settlement property of on-chain transfers eliminates the intermediate state that causes the majority of interbank disputes in traditional banking. In the correspondent banking model, a payment can exist in a ``funds in transit'' state for two to five business days during which the sender has debited but the receiver has not credited, creating reconciliation complexity and dispute resolution costs. On-chain transfers settle in a single transaction: either both the debit and credit execute or neither does, with no intermediate state possible.

For retail users in developing economies, transactional accounts on the platform serve as a substitute for the informal cash-handling systems that currently dominate daily financial activity. Recurring payroll deposits, utility payment debits, and peer transfers can all be scheduled as programmable transactions, reducing the friction of cash management without requiring users to interact directly with smart contracts for every operation.

\subsection{Group / Solidarity Lending}

The platform implements a group lending module inspired by the solidarity group model used in microfinance. A group of three to twenty registered clients can form a lending group on-chain. The collateral model for group loans distinguishes two tiers of client capability. For clients who can provide on-chain collateral (USDC or ETH), each member contributes individual collateral into a shared pool; mutual liability enforcement is automatic. For clients who lack individual on-chain collateral---the primary target population of 1.4~billion unbanked adults---the platform uses a \textbf{cold-start credit pathway}: new clients with zero on-chain history who have passed Level~1 KYC receive a \textit{provisional credit tier} that (a) requires mandatory group membership of at least 5 members, (b) caps individual loan amounts at the Level~1 KYC borrowing limit, and (c) expires after 3 months from first disbursement. After the first successful repayment cycle, the client's SBT is populated with real on-chain history and normal credit scoring applies. This explicit cold-start pathway prevents the circular deny loop---where no history leads to high risk score which leads to denial which leads to no history---that would otherwise permanently exclude the platform's intended beneficiaries. After three successful cycles, the client graduates from the provisional tier to the standard ML-gated credit model (Section~\ref{sec:aiml-support}). Repayment is tracked per member, and if a member defaults beyond a grace period, the contract can automatically cover the shortfall from the shared collateral pool under programmable mutual liability rules where collateral exists. Where collateral is absent, default triggers \textit{credit-score degradation}: the SBT's \texttt{riskTier} field is downgraded to the highest-risk tier and the \texttt{lastDefault} timestamp is permanently recorded. The credential is never revoked---consistent with the Buterin, Hitzig, and Weyl~(2022)~[R30] Soulbound Token design principle that SBTs are non-revocable identity-bound credentials. Future borrowing does not require formal revocation: no lender on the platform accepts a wallet whose SBT \texttt{riskTier} is at the default level, making borrowing effectively impossible in practice while preserving the permanent, honest record of the default event. This is analogous to a credit bureau record that permanently reflects a default rather than erasing it. Over time, group repayment histories contribute to credit scoring to enable progressive lending (larger loans and improved terms after subsequent cycles).

The group lending module is directly inspired by the solidarity group model developed by BRAC, the organization whose name this university bears. BRAC's group lending program, operating across Bangladesh and subsequently 11 other countries, demonstrates that social collateral---the mutual accountability of group members---can substitute for physical collateral among populations with no formal credit history. The Crypto World Bank translates the structural logic of this model into programmable contract enforcement: the mutual liability that BRAC enforces through field officer visits and weekly group meetings is encoded in the GroupLendingPool contract through automatic collateral claims and on-chain consent recording.

However, it is critical to note that BRAC and Grameen Bank's repayment rates above 95\% are a direct consequence of social pressure mechanisms: weekly physical group meetings, face-to-face peer accountability, community reputation, and field officers who know borrowers personally. These mechanisms are preventive---they stop defaults from occurring. Smart contract enforcement is punitive---it responds after default. An on-chain group of wallet addresses cannot replicate preventive social pressure. Consequently, the expected default rate for the CWB group lending module at launch is likely to be substantially higher than 5\%, estimated at 8--15\% consistent with the base-case and stress-test scenarios in Table~\ref{tab:default-scenarios}, until on-chain credit history accumulates sufficient predictive signal. The 95\% Grameen figure is cited as the design target and theoretical upper bound under mature social-trust conditions, not as a near-term operational forecast.

The full group loan lifecycle proceeds as follows. First, formation: three to twenty registered clients form a group on-chain, each contributing individual collateral to a shared pool contract. Second, application: the group submits a collective loan application specifying the total amount, per-member share, and repayment schedule. Third, consent: every group member must sign the application transaction before it is escalated to a bank approver, enforcing unanimous consent and preventing coerced participation. Fourth, approval: the bank approver reviews the group's aggregate credit history, income verification documents, and shared collateral ratio before approving disbursement. Fifth, disbursement: funds are distributed to each member's account in their individual share as defined by the formula \(\text{Share}_i = \text{LoanTotal} / N_{\text{members}}\). Sixth, repayment: each member repays their individual installments independently; on-chain tracking records per-member performance transparently to all group members. Seventh, mutual liability enforcement: if a member defaults beyond the grace period, the contract automatically covers the shortfall from the shared collateral pool, protecting the lending bank's position while distributing the cost proportionally across the group. Eighth, credit history improvement: successful group repayment cycles are permanently recorded on-chain and incorporated into each member's credit-passport SBT (Section~\ref{sec:credit-passport}), enabling progressive lending with larger amounts and better terms in subsequent cycles.

\paragraph{Stablecoin denomination at the group tier.}
Group loans at Tier~4 are denominated in USDC by default (Section~\ref{sec:stablecoin-first}). Solidarity-group clients are precisely the demographic for which ETH-denominated debt creates the highest real-cost shock under crypto-price drawdowns, so the stablecoin-first rule is enforced most strictly at this tier.

\paragraph{Over-indebtedness risk control.}
\label{sec:over-indebt}
Bangladesh's microfinance sector serves 41.56 million accounts across 724 licensed MFIs, and the NUS Economics Society analysis (May 2025)~[R37] together with the IJAR microfinance report~[R38] document over-indebtedness---borrowers taking simultaneous loans from multiple MFIs and inflating apparent repayment capacity---as the single most persistent problem in Grameen-style group lending. v15 adds three explicit controls. First, a per-wallet cap of two simultaneous active group loans, checked on-chain via the credit-passport SBT (Section~\ref{sec:credit-passport}). Second, a cross-platform debt disclosure flag in the SBT itself, readable by any other platform that adopts the standard, supporting industry-wide debt-load visibility. Third, a 30-day cooling-off period between consecutive group-loan cycles for any individual client, enforced by the GroupLendingPool contract. These controls are not in addition to the social-pressure model of solidarity lending; they are the on-chain encoding of the same caution that Grameen field officers exercise in person. A paper that cites Grameen and BRAC but does not address this risk would be ignoring a known failure mode of the model it builds on.

\paragraph{Cross-pool query and debt-to-income enforcement.}
The three controls above operate within a single Local Bank's GroupLendingPool.  To prevent a client from bypassing per-bank caps by simultaneously holding group loans at \emph{different} Local Banks within the hierarchy, the \texttt{GroupLendingPool.applyGroupLoan} function performs a \textbf{cross-pool query} before approving any new group loan.  Specifically, the function reads \texttt{creditPassport.getScore(wallet)} (the \texttt{ICreditPassport} interface defined in Section~\ref{sec:credit-passport}) and checks \texttt{openLoans >= MAX\_SIMULTANEOUS\_GROUP\_LOANS} against the global SBT state, which is updated by every Local Bank in the hierarchy on each disbursement.  Because the SBT is the authoritative cross-bank debt registry, this single read replaces the need for a bespoke cross-bank oracle or an off-chain credit bureau query.  In addition, the function enforces a \textbf{maximum debt-to-income (DTI) ratio} using the \texttt{income\_hash\_freshness} feature already computed in the ML pipeline (Section~\ref{sec:aiml-support}).  If the client's most recent income hash is stale (older than 90~days, configurable by governance), the application is flagged for manual review regardless of the SBT score.  If the income hash is fresh, the function computes $\text{DTI} = (\text{existingDebt} + \text{requestedAmount}) / \text{estimatedMonthlyIncome}$ and rejects applications where $\text{DTI} > 0.40$ (the governance-set ceiling, consistent with standard microfinance prudential norms).  Neither check requires new infrastructure: the RBAC system already knows which Local Banks exist, and the ML feature set already includes income-hash freshness as one of the four application features enumerated in Section~\ref{sec:aiml-support}.

\subsection{Foreign Exchange and Multi-Currency Operations}

A worldwide banking platform must support participants operating in different currency environments. The platform's FX module is designed to use decentralized price oracle networks to obtain verifiable exchange rates for supported asset pairs. Currency conversion can be offered at transparent oracle-derived rates with a disclosed spread, set by governance, that forms an additional revenue stream. To reduce liability-side volatility for retail clients, the platform is designed to support stablecoin-denominated lending (e.g., USDC/USDT) while allowing collateral in volatile cryptoassets (e.g., ETH).

\subsection{Trade Finance Facilitation (Planned)}

Trade finance instruments (letters of credit, guarantees, documentary collections) can be added as planned extensions for import--export participants. In a smart contract design, a buyer's bank locks payment in escrow, and release occurs automatically upon verified presentation of shipping documents via trusted data sources. This extension is out of scope for the current prototype but follows the same principles of programmable settlement and auditable execution.

\subsection{Privacy-Preserving Identity Compliance (Planned)}
\label{sec:zkp-compliance}

The regulatory compliance requirements of a worldwide banking platform extend beyond on-chain audit trails. Formal KYC and AML processes require users to prove legal identity, verify source of funds, and satisfy jurisdiction-specific screening requirements. Storing this personal data on a public blockchain creates a fundamental tension between regulatory compliance and user privacy.

The planned resolution is a zero-knowledge proof based KYC layer in which a licensed identity provider performs the full KYC verification off-chain and issues a verifiable credential to the user. The user then presents a ZKP to the banking contracts that proves they hold a valid credential from an approved provider without revealing any personal information on-chain. The contract validates the proof cryptographically and sets a \texttt{kycVerified} flag on the wallet address, enabling access to regulated operations such as large loans, FX conversions, and cross-border transfers. This approach satisfies regulatory requirements while preserving on-chain privacy and avoiding the liability of storing personal financial data in an immutable public record.


\section{Governance Framework}

\subsection{Network Membership Governance}

\begin{table}[H]
\centering
\caption{Network membership governance.}
\label{tab:network-governance}
\begin{tabular}{p{4.5cm}p{9cm}}
\hline
\textbf{Governance Aspect} & \textbf{Implementation} \\
\hline
Member on-boarding & World Bank owner registers National Banks; National Banks register Local Banks; Local Banks designate approvers --- all enforced on-chain. \\
\hline
Member off-boarding & Deactivation flags in smart contracts; cascading access revocation. \\
\hline
Regulatory oversight & Audit log emission via smart contract events; planned read-only regulator dashboard. \\
\hline
Permission structure & Hierarchical: Owner $\to$ National Bank $\to$ Local Bank $\to$ Approver $\to$ Retail Client; enforced by on-chain role check modifiers. \\
\hline
Network operations & Pause/unpause mechanism for emergency response; emergency withdrawal for critical situations. \\
\hline
\end{tabular}
\TableScreenshotEnd{Banking Product Suite: Deposit, Credit, and Ancillary Contract Specifications}
\end{table}

\vspace{2pt}
\noindent\textit{\small This governance table defines membership and onboarding rules that control who can join the banking network and under which role. The structure is intended to preserve institutional accountability by separating responsibilities across tiers rather than relying on undifferentiated token-holder governance alone.}

\subsection{Business Network Governance}

\begin{table}[H]
\centering
\caption{Business network governance.}
\label{tab:business-governance}
\begin{tabular}{p{4.5cm}p{9cm}}
\hline
\textbf{Governance Aspect} & \textbf{Implementation} \\
\hline
Business charter & Defined in project documentation; operational parameters coded in smart contract constants. \\
\hline
Common services & Reserve management, loan lifecycle orchestration, event-driven notification system. \\
\hline
Business SLA & Testnet phase: best-effort availability. Production phase: 99.5\% target API-layer uptime (team-controlled; multi-region deployment). Chain-level liveness is determined by Polygon PoS validator consensus and is not under the platform team's control; the team monitors Polygon network status and maintains degraded-mode fallbacks (read-only dashboard) if chain activity pauses. \\
\hline
Regulatory compliance & Architecture designed for audit trail generation; data partitioning supports GDPR-style data subject requests. \\
\hline
\end{tabular}
\TableScreenshotEnd{Security Threat Model: Vulnerability Classes, Attack Vectors, and Mitigations}
\end{table}

\vspace{2pt}
\noindent\textit{\small This business governance table outlines operational controls such as approvals, limits, and escalation paths for higher-risk actions. The goal is to combine programmable enforcement with human oversight, matching real banking governance while retaining on-chain audit trails.}

\subsection{Technology Infrastructure Governance}

\begin{table}[H]
\centering
\caption{Technology infrastructure governance.}
\label{tab:tech-governance}
\begin{tabular}{p{4.5cm}p{9cm}}
\hline
\textbf{Governance Aspect} & \textbf{Implementation} \\
\hline
Distributed IT structure & Client-side frontend (decentralised delivery); blockchain layer (fully decentralised); backend API (centralised, horizontally scalable). \\
\hline
Technology assessment & Continuous evaluation of EVM alternatives (L2 rollups, sidechains) for cost and performance optimisation. \\
\hline
On-chain / off-chain data services & Clearly partitioned (see Table~\ref{tab:data-partitioning}); event listeners synchronise state between layers. \\
\hline
Risk mitigation & Smart contract pause mechanism; ReentrancyGuard; input validation; planned formal security audit. \\
\hline
\end{tabular}
\TableScreenshotEnd{Governance Framework: Operational, Business, and Technology Governance Layers}
\end{table}

\vspace{2pt}
\noindent\textit{\small This technology governance table specifies how protocol parameters and infrastructure changes are managed, including security controls and upgrade considerations. It supports the thesis emphasis on safe staged rollout and maintainability under evolving regulatory and security requirements.}

\subsection{Regulatory Compliance Considerations}
\label{sec:regulatory-compliance}

As the platform is active in the regulated banking sector:
\begin{itemize}[noitemsep]
\item Our prototype stage would be solely on the \textbf{public testnets} without attached real monetary value.
\item The architecture would facilitate \textbf{audit logs generation} to be reviewed by regulatory regimes on the sandbox programs.
\item Future production deployment would engage \textbf{regulatory sandbox programs} in target jurisdictions.
\end{itemize}

Although the current prototype does not process KYC/AML data, a complete worldwide banking architecture must account for compliance at the design level. Privacy-preserving KYC design is described in detail in Section~\ref{sec:zkp-compliance}.

\paragraph{Institutional trust bootstrapping.}
A practical question for any new institutional platform is how the first institutional participants are recruited before the network has an established track record. The CWB bootstrapping strategy is three-pronged. \textit{First, the academic pilot route:} the platform is released as open-source software under an Apache~2.0 licence, allowing any university or NGO to deploy their own instance for research purposes. The first institutional participants are academic departments, not commercial banks, which face no regulatory barriers to testnet participation. \textit{Second, the regulatory sandbox route:} the platform applies for a Singapore MAS Project~Guardian sandbox slot or a UAE DIFC Innovation Testing Licence, which provides regulatory cover for a controlled pilot with a licensed financial-institution partner (see Future Work item~14 and Section~\ref{sec:bangladesh-reg}). \textit{Third, the BRAC alignment route:} as the platform that algorithmically encodes the BRAC solidarity group lending model, a natural first institutional-partner discussion is with BRAC International's technology arm, which already operates digital microfinance infrastructure across 11 countries. None of these routes require the platform to be trusted by a commercial bank on Day~1; trust is built incrementally through the academic and sandbox record.

\paragraph{FATF Travel Rule (R.16) --- inter-tier flows.}
As specified in Section~\ref{sec:upward-flow}, the FATF R.16 Travel Rule applies to inter-tier capital flows above USD~1,000 in most jurisdictions. The data packet structure specified for \texttt{UpwardDepositFacility} and \texttt{InterBankLendingPool} transfers captures all FATF-required originator and beneficiary fields. Implementation of the cross-tier Travel Rule notification protocol is Future Work contingent on jurisdiction-specific Bangladesh Bank regulatory guidance.

\paragraph{SAR workflow --- Isolation Forest to regulator.}
The five-step SAR flow specified in Section~\ref{sec:system-modeling} (Isolation Forest detection $\to$ Kafka \texttt{aml-alert} $\to$ compliance queue $\to$ officer review $\to$ freeze + audit log) constitutes the platform's full AML compliance pathway. The \texttt{freezeAccount(clientId)} function in \texttt{LocalBankPool.sol} is gated by \texttt{onlyApprover} and is Foundry-tested in Sprint~3 to ensure no privilege escalation is possible. Phase~2 will add integration with Bangladesh Bank's Financial Intelligence Unit reporting portal and automated SAR PDF generation in the required regulatory format.

\paragraph{Regulator audit request flow (specify only).}
A formal regulator audit request sequence is specified as a contribution to the institutional compliance architecture:

\begin{enumerate}[noitemsep]
\item Regulator submits a signed request off-chain to the World Bank admin.
\item World Bank admin verifies the regulatory mandate via multi-sig confirmation.
\item System extracts: client loan history (\texttt{loans} table), installment payment records, AI/ML risk score log (\texttt{AI\_ML\_LOG}), on-chain transaction hashes, and SAR history if applicable.
\item Extracted data is packaged and encrypted with the regulator's public key.
\item The audit response is logged immutably in \texttt{audit\_logs} (append-only policy, Section~\ref{sec:normalization}).
\end{enumerate}

This workflow is a specification contribution; implementation is Future Work. The append-only \texttt{audit\_logs} DB policy (Section~\ref{sec:normalization}) ensures the audit response record cannot be altered after the fact, providing the immutability guarantee that regulatory compliance requires.

\subsection{Asset Tokenization}

\begin{itemize}[noitemsep]
\item \textbf{Current implementation:} The native blockchain currency (ETH/MATIC) is used as the reserve and loan currency.
\item \textbf{Planned extension:} Support for ERC-20 stablecoins (USDC, USDT) for lending operations in USD; tokenized collateral instruments for lending situations where there isn't enough collateral.
\end{itemize}

\section{Five-Layer Defense-in-Depth Security Architecture}
\label{sec:defense-in-depth}

\begin{figure}[H]
\centering
\OnePageDiagram{fig-defense-in-depth.pdf}
\caption{Five-layer defense-in-depth security architecture. Each layer is necessary; breaching the smart-contract layer requires defeating every layer above it. Layer~1 (smart contract): OpenZeppelin~v5, UUPS, TimeLock, RBAC, ReentrancyGuard, audit suite. Layer~2 (application): EIP-712 sign-in, JWT, rate-limit. Layer~3 (AI/ML): commit-reveal oracle, model registry, SHAP, hallucination guard. Layer~4 (runtime): Tenderly, The Graph, WebSocket dashboard, anomaly detection. Layer~5 (operations): Safe multisig, key rotation, bug-bounty.}
\label{fig:defense-in-depth}
\end{figure}

\begin{figure}[H]
\centering
\OnePageDiagram{fig-security-controls.pdf}
\caption{Smart-contract security controls applied in v15: (a) UUPS (ERC-1822) upgrade path with TimeLock-gated authorization; (b) EIP-712 sign-in (typed-data v4 with domain separation and 15-minute JWT); (c) granular per-function pause registry indexed by The Graph for transparent governance.}
\label{fig:security-controls}
\end{figure}


The correct mental model for a blockchain banking platform's security is not ``smart-contract security only''---it is a five-layer defense-in-depth stack in which breaching the contract layer requires defeating every layer above it. The Crypto World Bank adopts the following layered model, summarized in Table~\ref{tab:defense-layers} and inspired by current (2025--2026) audit-tooling and runtime-monitoring best practice from Trail of Bits, Cyfrin, Tenderly, and OpenZeppelin Defender.

\begin{table}[H]
\centering
\caption{Five-layer defense-in-depth model for the Crypto World Bank platform. Each layer is necessary; breaching the contract layer requires defeating every layer above it.}
\label{tab:defense-layers}
\begin{tabular}{p{1.4cm}p{4.0cm}p{8.5cm}}
\hline
\textbf{Layer} & \textbf{Component} & \textbf{CWB Design Choice} \\
\hline
L5 & Operational security & Safe (Gnosis) 2-of-3 multisig for \texttt{WorldBankAdmin}; 3-of-5 for global parameter changes; documented key-rotation schedule; hardware-wallet signer for the WorldBankAdmin Safe (Ledger Gen5). \\
\hline
L4 & Runtime monitoring & Tenderly alerts on six critical triggers (large disbursement, repeated requests, reserve-ratio drop below 20\%, failed \texttt{disburseLoan}, role grants, governance pause); The Graph subgraph for indexed event history; transaction simulation pre-deployment. \\
\hline
L3 & AI/ML security & Random Forest fraud probability + Isolation Forest anomaly score + SHAP explanation per loan, delivered via commit-reveal oracle; transaction-graph distance to flagged wallets; GNN-based wallet-graph clustering for relational fraud (Section~\ref{sec:gnn-extension}). \textit{Note: session-level behavioral biometrics (page-view timing, terms-acceptance latency) require client-side instrumentation not yet in the schema and are a specified Future Work extension.} \\
\hline
L2 & Application security & EIP-712 wallet-signed JWT (15-minute lifetime + refresh); \texttt{slowapi} rate limiting (5/min auth, 60/min reads, 10/min writes); Pydantic input range constraints; strict CORS and Content-Security-Policy; environment-variable secret hygiene; prompt injection scanning middleware on all user-sourced fields before context assembly (see below). \\
\hline
L1 & Smart contract security & OpenZeppelin \texttt{ReentrancyGuard}; CEI pattern; Solidity~0.8.20 overflow protection; RBAC with role expiry; UUPS upgrade with Safe-multisig authorization; TimeLock on governance actions; granular per-function pause; four-tool audit pipeline (Slither, Mythril, Echidna, Halmos/Certora); Foundry fuzz + invariants. \\
\hline
\end{tabular}
\end{table}

The L1 contract layer is what the existing thesis discusses in detail; L2--L5 are new contributions of this version. The Section that follows (\ref{sec:threat-model}) expands the threat model to span all five layers and motivates the specific controls listed in Table~\ref{tab:defense-layers}.

\paragraph{Prompt injection scanning (Layer~2 --- Application).} Before any
user-sourced string is assembled into the system prompt by the Express.js context
injection layer, a scanning middleware checks each field for candidate injection
patterns: substrings matching \texttt{ignore previous instructions},
\texttt{new system prompt}, tool-override phrases, and JSON-escape sequences that
could break the structured context block. Fields that trigger the scanner are
sanitised (the offending segment is removed and replaced with a placeholder) and
\texttt{AGENT\_ACTION\_LOG.injection\_scan\_flag} is set to TRUE. The fields
subject to scanning are: \texttt{language} (user-controlled), any text content
from income proof document uploads, and the last $N$ conversation turns before
they are appended to the Volatile tier. This control directly addresses OWASP
LLM01 (Prompt Injection)~[R54] and is motivated by Alizadeh et
al.~[R53], who demonstrate data exfiltration through
income document uploads and language fields in a banking tool-calling agent using
attack vectors structurally identical to those present in the CWB context assembly
layer.

\paragraph{Per-tier admin-key model.}
Operational security (L5) is the highest-impact and lowest-cost security decision in the entire project. The \texttt{WorldBankAdmin} role is never held by a single externally owned account (EOA) in any environment, including the university lab demo: it is transferred to a Safe 2-of-3 multisig before any contract is deployed to a testnet that examiners will inspect, with Signer~1 on the demo machine, Signer~2 on a phone, and Signer~3 as a paper-wallet backup held by the academic supervisor. National Bank Admin roles use Safe 2-of-3; Local Bank operators use Safe 2-of-3; the only EOA-bound role is the retail CLIENT role, which can be reset off-chain via the social-recovery guardian flow described in Section~\ref{sec:onboarding}.

\paragraph{SWC Registry mapping.}
Every state-mutating function in the three implemented contracts is mapped to its Smart Contract Weakness Classification (SWC) class~[R36] and the corresponding mitigation. This mapping replaces the previous narrative coverage of Atzei et~al.~[7] with a structured taxonomy, and is the format the Cyfrin audit checklist (2025) recommends for academic prototypes. The full mapping table is in Appendix~B's Capability Matrix.

\section{Threat Model and Security Controls}
\label{sec:threat-model}

Security threats in smart-contract banking systems are not limited to isolated code defects; they span application-layer attacks (API abuse, JWT theft, social engineering), runtime evasion, and governance-layer key compromise in addition to the contract-layer issues that DeFi literature focuses on. Table~\ref{tab:threat-model} extends the original threat-model table to the full five-layer scope established in Section~\ref{sec:defense-in-depth}.

\begin{table}[H]
\centering
\caption{Threat model and security controls mapping, expanded to span all five layers of the defense-in-depth stack (Section~\ref{sec:defense-in-depth}).}
\label{tab:threat-model}
\begin{tabular}{p{2.6cm}p{2.0cm}p{4.5cm}p{4.8cm}}
\hline
\textbf{Threat Category} & \textbf{Layer} & \textbf{Attack Vector} & \textbf{Implemented / Planned Mitigation} \\
\hline
Reentrancy and state manipulation & L1 & External calls exploit inconsistent intermediate state during execution (e.g.\ withdraw patterns) & OpenZeppelin \texttt{ReentrancyGuard}; CEI pattern; unit + Foundry invariant tests on critical flows. \\
\hline
Arithmetic errors & L1 & Integer overflow/underflow in financial math & Solidity~0.8.x checked arithmetic; explicit bounds for rates and limits; Certora invariants. \\
\hline
Oracle manipulation & L1 & Price-feed manipulation affecting collateral valuation or FX conversion & Decentralized oracle networks (Chainlink Price Feeds); medianization, heartbeat checks; TWAP for collateral pricing; circuit breakers on stale / anomalous prices. \\
\hline
Flash-loan driven economic attacks & L1 & Atomic borrowing manipulates on-chain price/state assumptions & Conservative price-dependent logic; multi-block confirmation for large collateral changes; pause-on-anomaly procedures. \\
\hline
Governance capture & L1+L5 & Privileged roles abused to register malicious banks or change unsafe parameters & Tiered role separation; Safe multisig on every admin role; TimeLock (24--48~h) on parameter changes; role-expiry timestamps; immutable audit trail. \\
\hline
Sybil abuse (retail / group) & L1+L3 & Attackers create many wallets to bypass limits or build fraudulent groups & On-chain RBAC; rate limits per wallet; ZKP-KYC gating; GNN-based wallet-graph clustering. \\
\hline
JWT theft / session hijack & L2 & Stolen token gives full account access until expiry & 15-minute JWT lifetime; EIP-712 wallet-signature challenge on issue; Redis blacklist on logout; rotation on suspicious geolocation. \\
\hline
API abuse / DDoS & L2 & Bulk requests exhaust backend or rate-limit shared services & \texttt{slowapi} middleware with per-endpoint differentiated limits; CSP and CORS lockdown; correlation IDs for incident response. \\
\hline
Input-validation bypass & L2 & Malformed financial inputs cause unintended state on backend & Pydantic explicit range constraints on all financial fields; wallet-address regex; SHA-256 hash for document references. \\
\hline
ML evasion / adversarial fraud & L3 & Crafted features evade Random Forest / Isolation Forest detectors & GNN relational analysis (Section~\ref{sec:gnn-extension}); federated cross-bank intelligence (Section~\ref{sec:fl-fraud}); SHAP-based human review for borderline cases. \\
\hline
LLM hallucination on regulatory questions & L3 & 9B-LLM gives a confidently wrong compliance answer to a client & RAG over current policy docs; refusal layer on regulated topics; red-teaming protocol (Section~\ref{sec:llm-eval}). \\
\hline
Monitoring evasion & L4 & Attack split below alert thresholds & Multi-trigger alerts (count + amount + frequency); The Graph indexed history for retrospective forensics. \\
\hline
Admin-key compromise & L5 & Single EOA compromise gives full World Bank control & Safe 2-of-3 (national / local) and 3-of-5 (world bank) multisig; hardware-wallet signer; rotating demo keys before/after every presentation. \\
\hline
Blind-signing on hardware wallet & L5 & User signs opaque calldata revealing seized funds & ethers.js calldata decoding to human-readable summary; Ledger Multisig (Gen5) for Safe signer; pre-approval simulation in Tenderly. \\
\hline
\end{tabular}
\end{table}

This expanded threat model is the structural justification for the five-layer architecture in Section~\ref{sec:defense-in-depth}: each row of the table is a vector that the corresponding layer prevents, detects, or recovers from. Application-layer rows (L2) and operational-layer rows (L5) are the largest blind spots in the prior literature on academic DeFi prototypes and are made explicit here.


%% %%%%%%%%%%%%%%%%%%%%%%%%%%%%%%%%%%%%%%%%%%%%%%%%%%%%%%%%%%%%
%% CHAPTER 4: METHODOLOGY
%% %%%%%%%%%%%%%%%%%%%%%%%%%%%%%%%%%%%%%%%%%%%%%%%%%%%%%%%%%%%%
\chapter{Methodology}

This project and its planning have been structured in accordance with the \href{https://bcolbd.org/uploads/guideline/BLOCKCHAIN\%20OLYMPIAD\%20BANGLADESH\%20AI\%20Guideline.pdf}{Blockchain Olympiad Bangladesh AI guidelines}~[16] and the final-year project evaluation rubrics.

\section{Development Methodology}

We adopted a \textbf{lightweight Agile/Scrum} methodology tailored for an academic prototype with a fixed two-month development window. The iterative approach enables incremental delivery of demonstrable features while accommodating evolving requirements inherent in research-oriented development. Figure~\ref{fig:agile-process} illustrates our Agile process.

% ── Figure: Agile/Scrum Process Flow ─────────────────────────────────────────
% Standard flowchart symbols: rectangle = process, parallelogram = I/O artefact,
% diamond = decision.  Academic palette: PrimaryBlue headers, RowShade fill,
% AccentBlue borders.  No \n typos from the source image.
% ─────────────────────────────────────────────────────────────────────────────
\begin{figure}[H]
\centering
\OnePageDiagram{fig-agile-process.pdf}
\caption{Agile / Scrum process with two-week iterations, four ceremonies (Sprint Planning, Daily Stand-up, Sprint Review, Sprint Retrospective), backlog refinement, and the sprint-submission gate. Sprint-point estimates for the three implementation sprints are summarized on the right.}
\label{fig:agile-process}
\label{fig:sprint-points}
\label{fig:sprint-submission}
\end{figure}

The following points summarize how Agile will be applied:
\begin{itemize}[noitemsep]
\item \textbf{Sprint Duration:} A period of 2 to 3 weeks is allocated for every sprint (3 sprints total).
\item \textbf{Team Size:} A group of two developers will be assigned (thesis project). The work will be focused on the thesis project.
\item \textbf{Weekly Sync:} A weekly gathering will occur to assess progress. Issues will be identified during these meetings. Plans will also be developed in these sessions alongside progress checks.
\item \textbf{Sprint Planning:} Planning for the sprint will take place for one hour at the beginning of every sprint.
\item \textbf{Sprint Review and Retrospective:} Reviewing work occurs at every sprint's conclusion. A discussion on lessons learned will be conducted for 30 minutes.
\end{itemize}

\section{Planned AI/ML Support and Risk-Score Wiring}
\label{sec:aiml-support}

\begin{figure}[H]
\centering
\OnePageDiagram{fig-aiml-pipeline.pdf}
\caption{AI/ML pipeline wiring: (a) training pipeline (Random Forest + Isolation Forest + SHAP), (b) commit-reveal ML oracle that gates loan approval on-chain, (c) Graph Neural Network (GraphSAGE) extension producing relational features for the wallet--wallet graph, and (d) federated learning across Local Banks with a National-Bank aggregator and differential-privacy noise.}
\label{fig:aiml-pipeline}
\end{figure}


The AI/ML layer is not three separate models added to a static lending workflow; it is an integrated pipeline that produces a single risk score per loan and feeds it on-chain through the commit-reveal oracle of Section~\ref{sec:oracle-architecture}. The pipeline has four stages:

\begin{enumerate}[noitemsep]
\item \textbf{Feature engineering.} For each loan application, the FastAPI service computes 18 behavioral features (rolling-window transaction count, average inter-transaction interval, on-chain repayment history, group-membership flag, wallet-graph distance to flagged addresses) and 4 application features (requested amount, KYC level, income-hash freshness, ZKP credential expiry).
\item \textbf{Risk scoring.} A Random Forest classifier returns a fraud probability $p_f \in [0,1]$ (a true calibrated probability, computed via Platt scaling applied to the raw tree vote proportion). An Isolation Forest returns an anomaly score $a \in [0,1]$ where values approaching 1 indicate anomalies; this is a path-length-based score (see List of Formulas), \textit{not} a probability. The two outputs are combined via a stacking meta-learner (logistic regression) fitted on a held-out validation set, yielding a calibrated composite score $s \in [0,1]$. The default meta-learner weights are approximately 0.7 for $p_f$ and 0.3 for the calibrated anomaly signal; in the prototype these serve as initial placeholders to be re-fitted on any labeled data that becomes available. A direct linear combination of a probability and a non-probability score would produce a statistically uninterpretable output; the stacking approach ensures $s$ is always a valid probability estimate.
\item \textbf{Commit-reveal oracle.} The service publishes \(h = \mathrm{keccak256}(s \,\|\, \text{nonce})\) to the LoanController contract via \texttt{commitRiskScore}, then reveals $(s, \text{nonce})$ in a subsequent transaction. The LoanController enforces a state machine: approval is only permitted in the \texttt{SCORE\_REVEALED} state, preventing any approver from bypassing the ML step. Section~\ref{sec:oracle-architecture} covers the cryptographic detail.
\item \textbf{Decision threshold.} The contract applies a governance-configured threshold calibrated against the expected fraud rate: $s < 0.4$ enables approver discretion to approve, $0.4 \le s < 0.7$ requires mandatory manual review, and $s \ge 0.7$ rejects the loan with an on-chain reason code. These thresholds are prototype defaults derived from the class imbalance expected in DeFi lending datasets (estimated fraud prevalence 2--5\%); they must be recalibrated against the actual training data distribution before any production deployment. The \texttt{auto-approve} language used elsewhere in the document is a simplification; the contract always requires an approver to sign the final disbursement transaction.
\end{enumerate}

\paragraph{SHAP explainability output to clients.}
Each loan decision produces a SHAP value vector, and the top-3 contributing features are stored in \texttt{AI\_ML\_SECURITY\_LOG} together with the decision. The client-facing UI displays a plain-language explanation rather than the raw SHAP score---for example, \emph{``Your application was declined primarily because: (1) no on-chain repayment history (60\%); (2) requested amount exceeds income estimate (30\%); (3) group membership not yet verified (10\%).''} This satisfies the EU AI Act Article~86 transparency requirement for AI-assisted financial decisions and converts the formal SHAP design choice~[R7] into an end-user benefit.

\noindent\textit{Future Work note: session-level behavioral biometrics (time-to-apply, page-view count, terms-acceptance latency, IP-geolocation consistency) are a specified but not yet implemented feature extension. They require client-side instrumentation and a schema addition (\texttt{session\_events} table) that are not present in the current prototype; the 18 core on-chain transaction features listed above are the active feature set.}

\paragraph{Autonomous AI agent --- six-step pipeline.}

The CWB AI layer extends beyond a read-only Q\&A guide to an action-capable autonomous banking agent. The six-step pipeline integrates the Qwen3-8B model with the live on-chain context injection layer (Section~\ref{sec:local-llm}) and the MCP tool server described below:

\begin{enumerate}[noitemsep]
\item \textbf{User message arrives} at the React frontend and is forwarded to the Express.js backend over Server-Sent Events (SSE).
\item \textbf{Agent brain (Qwen3-8B) reads live on-chain context.} The Express.js context injection layer prepends the client's full account state (loan status, SBT risk tier, pool utilisation, credit score, active installments) as structured JSON into the system prompt before the model processes the user message.
\item \textbf{Q\&A path.} If the request is informational, the agent retrieves relevant policy documents via RAG (ChromaDB) and generates a grounded reply. No tool calls are made.
\item \textbf{Action request path.} If the request implies a state-modifying operation, the agent identifies the applicable MCP write tool, assembles the full parameter set from the injected context, and presents a concise confirmation summary: \emph{``You are requesting a loan of 150~USDC for 6~months at 8.5\% APR. Your monthly installment would be 26.25~USDC. Shall I proceed?''}
\item \textbf{Human confirmation gate.} The agent pauses and waits for explicit affirmative consent (\emph{``Yes, proceed''} or equivalent). No write tool is ever called without a confirmation turn in the conversation history. If the client responds ambiguously three times, the agent pauses for 10~minutes and logs the incident to \texttt{AGENT\_ACTION\_LOG}.
\item \textbf{Execute via MCP write tool, monitor, and notify.} After confirmation, the agent calls the appropriate write tool, which POSTs to the Express.js banking API. If an EIP-7702 session key is active, it signs the resulting on-chain transaction within its approved scope. The agent polls for status and notifies the client: \emph{``Your loan application has been submitted. Reference: L-4821. Bank approval typically takes 24~hours.''}
\end{enumerate}

\paragraph{MCP tool server --- 16 restricted banking tools.}

The agent interacts with CWB exclusively through a typed tool schema, never through arbitrary code execution. This schema is the safety boundary: the agent cannot access any data or perform any operation not explicitly defined in the following 16~tools.

\begin{table}[H]
\centering
\caption{MCP tool server: 8 read tools (always permitted) and 8 write tools (require human confirmation gate).}
\label{tab:mcp-tools}
\small
\begin{tabular}{p{4.8cm}p{1.4cm}p{6.3cm}}
\hline
\textbf{Tool} & \textbf{Type} & \textbf{Maps to / Description} \\
\hline
\texttt{get\_account\_state}         & READ & SBT tier, loan count, savings balance \\
\texttt{get\_credit\_score}          & READ & Current score, tier, next threshold \\
\texttt{get\_loan\_status}           & READ & All active loans + next installment \\
\texttt{get\_installment\_schedule}  & READ & Full repayment calendar \\
\texttt{get\_pool\_utilisation}      & READ & Current Local Bank pool capacity \\
\texttt{get\_interest\_rate}         & READ & Rate for client's credit tier \\
\texttt{get\_market\_data}           & READ & ETH/USD, BDT/USD, 30-day volatility \\
\texttt{get\_requirements}           & READ & Documents + KYC needed for a given loan amount \\
\hline
\texttt{submit\_loan\_application}   & WRITE$^{\dagger}$ & POST /api/loan/apply \\
\texttt{submit\_deposit}             & WRITE$^{\dagger}$ & POST /api/savings/deposit \\
\texttt{submit\_fixed\_deposit}      & WRITE$^{\dagger}$ & POST /api/savings/fixed-deposit \\
\texttt{pay\_installment}            & WRITE$^{\dagger}$ & POST /api/installment/pay \\
\texttt{join\_group\_pool}           & WRITE$^{\dagger}$ & POST /api/group/join \\
\texttt{submit\_kyc\_upgrade}        & WRITE$^{\dagger}$ & POST /api/kyc/upgrade \\
\texttt{schedule\_payment\_reminder} & WRITE$^{\dagger}$ & POST /api/reminder/set \\
\texttt{submit\_group\_application}  & WRITE$^{\dagger}$ & POST /api/group/apply \\
\hline
\end{tabular}

\vspace{4pt}
\noindent\small$^{\dagger}$\textit{All write tools require an explicit human confirmation step before execution. The agent assembles the full parameter set, presents a plain-language summary, and waits for affirmative consent. No write tool is ever called without a confirmation turn recorded in the conversation history.}
\end{table}

\paragraph{Authority Brief UI --- SHAP breakdown for bank approvers.}

When the agent submits a loan application via \texttt{submit\_loan\_application}, it generates a structured Authority Brief and delivers it to the Local Bank approver dashboard alongside the application. The brief presents the ML risk assessment in human-readable form, so the approver can make an informed one-click decision without navigating away from the brief:

\begin{verbatim}
AGENT AUTHORITY BRIEF -- Loan Application #4821
-------------------------------------------------
Client tier:         Silver  (score: 512)
Requested amount:    150 USDC
Duration:            6 months
Collateral:          75 USDC (50% LTV)

ML Risk Score:       0.23  (LOW RISK)
SHAP breakdown:
  + On-time repayment rate:    +0.41  (reduces risk)
  + Loan-to-income ratio:      +0.18  (within safe range)
  - Wallet age:                -0.09  (account < 6 months)
  - No prior completed loans:  -0.27  (no repayment history)

Agent recommendation:   APPROVE
Suggested interest:     Base rate - 0.5% (Silver tier)

[ APPROVE ]   [ REQUEST MORE INFO ]   [ DECLINE ]
\end{verbatim}

\noindent The agent monitors the approval event on-chain and notifies the client automatically upon a decision, closing the loop without manual follow-up from either party.

\paragraph{Per-user personalisation --- shared model, per-user context.}

Every client receives a personalised experience without deploying a separate model per user. Personalisation lives entirely in the per-user context namespace injected into every system prompt by the Express.js context injection layer:

\begin{verbatim}
{
  "client_id": "0xABC...",
  "language": "Bengali",
  "credit_tier": "Silver",
  "credit_score": 512,
  "active_loans": [
    {
      "loan_id": "L-4821",
      "amount": 100,
      "next_due": "2026-06-15",
      "installment": 17.5
    }
  ],
  "savings_balance": 250.0,
  "pool_utilisation": 0.67,
  "kyc_tier": 1,
  "conversation_history": [ "...last 10 turns..." ]
}
\end{verbatim}

\noindent The result is indistinguishable from a per-user model at a fraction of the cost and with no additional inference infrastructure. Because the on-chain state is re-injected at every turn, the agent always operates on current account data and cannot hallucinate the client's balance or loan status.

\paragraph{Extended agent capabilities.}

Beyond the core deposit and loan flow, the same tool-calling architecture provides five additional capabilities without any new model infrastructure:

\begin{itemize}[noitemsep]
\item \textbf{Proactive installment reminders.} A cron job checks \texttt{get\_installment\_schedule} daily. Three days before a due date, the agent sends: \emph{``Your next installment of 17.50~USDC is due in 3~days. Shall I process it now?''} On confirmation, \texttt{pay\_installment} is called via the EIP-7702 session key.
\item \textbf{Credit Passport coaching.} The agent explains exactly what the client needs to reach the next credit tier using \texttt{get\_credit\_score} combined with the tier schedule (Table~\ref{tab:credit-tiers}): \emph{``You need 2~more on-time repayments to reach Gold tier. Your borrowing rate will then drop by 1.0\%.''}
\item \textbf{Loan calculator.} Before submission, the agent runs the EMI formula inline and displays the full repayment schedule in USDC and BDT equivalent (via \texttt{get\_market\_data} for the live Chainlink FX rate), so the client understands the full cost before confirming.
\item \textbf{Group lending coordination.} A group leader can ask the agent to track member signatures for a group loan application. The agent monitors which members have signed and sends reminders to those who have not via \texttt{schedule\_payment\_reminder}.
\item \textbf{KYC tier upgrade guidance.} The agent identifies exactly which documents are required for a KYC upgrade via \texttt{get\_requirements} and guides the client through the phone-camera document capture flow step by step.
\end{itemize}

\section{Graph Neural Network Extension}
\label{sec:gnn-extension}

The Random Forest detector treats each transaction independently; DeFi fraud is largely \emph{relational}, appearing in the wallet-interaction graph. A baseline Graph Neural Network (GraphSAGE) is added as an ablation: nodes are wallets, edges are transactions or shared group-membership, node features are the same 18 behavioral features used by the Random Forest. The model is trained on a synthetic transaction graph augmented with the multi-chain DeFi fraud dataset of Palaiokrassas et al.~[3] (54M transactions across 23 protocols), which is the correct domain-matched dataset for Ethereum/Polygon account-model DeFi fraud. The public Elliptic dataset (Weber et al.~[R24]) is used only as a secondary baseline comparison; it models Bitcoin UTXO-graph fraud (mixing services, darknet markets, ransomware), which has a fundamentally different graph topology and fraud pattern from DeFi lending fraud on account-model chains, and this domain difference is explicitly acknowledged in the evaluation. Cheng et al.~(2024)~[R24] and Wang and Wang~(2025)~[R26] both report sustained advantages of GNN over flat ML on relational fraud tasks (>70\% detection of illicit transactions at <1\% false positive). The expected contribution to the thesis is not a state-of-the-art result but a controlled ablation that quantifies the marginal value of relational features over flat behavioral features---a publishable comparison aligned with the SoK literature.

\section{Federated Learning Across Banking Tiers}
\label{sec:fl-fraud}

Each Local Bank trains a local fraud-detection model on its own borrowers' data; gradients are aggregated at the National Bank tier via FedAvg with differential-privacy noise (clipping threshold $C$, noise multiplier $\sigma$ matched to the PrivChain-AI configuration~[R27]); the National Bank publishes the resulting global model parameters back to its Local Banks. No raw transaction data ever leaves a Local Bank, satisfying jurisdiction-specific data-residency rules in countries that prohibit cross-border export of customer records. The federation topology mirrors the four-tier institutional hierarchy directly: each National Bank operates a federation server, and the World Bank Admin can optionally orchestrate a cross-National-Bank federation for global threat patterns. This is architecturally appropriate rather than ornamental: the National-Bank / Local-Bank trust boundary is the same boundary across which the federation cannot cross unencrypted data.

\paragraph{Cold start and activation threshold.}
Federated learning requires each participating bank to have sufficient local data to train a meaningful local model before the first aggregation round. Since the prototype starts with zero real users, federated learning is a \textit{Phase 2} feature, activated per bank only after a minimum of 500 transactions have been recorded locally. Until that threshold, banks share a common baseline model trained on the synthetic DeFi transaction dataset. This threshold is explicitly defined as the FL activation criterion in the implementation roadmap; federated learning is therefore not claimed as a Sprint~3 deliverable but as a formally specified future-phase capability.

\section{Formal Verification of Reserve Invariants (Certora)}
\label{sec:certora}

The Certora Prover went open-source in 2025 and is the only formal-verification tool to have produced a publicly verifiable proof of a real-world Solidity contract; the platform's TVL coverage now exceeds \$100~billion across Aave, MakerDAO, Uniswap and others~[R39, R40]. v15 elevates ``formal verification'' from future-work text to two concrete CVL (Certora Verification Language) specifications.

\paragraph{Invariant 1 -- reserve never under-collateralized.}
\(\forall\, t:\; \mathrm{totalReserve}(t) \ge \mathrm{minReserveRatio} \cdot \mathrm{totalDeposited}(t).\)
This invariant is asserted in CVL against the WorldBankReserve contract under all reachable contract states. The Certora prover symbolically explores every reachable execution path; any counterexample is delivered as a concrete attacker trace.

\paragraph{Invariant 2 -- no over-allocation downstream.}
\(\forall\, n \in \mathrm{NationalBanks}:\; \mathrm{allocatedTo}[n] \le \mathrm{totalReserve} - \sum_{n' \ne n} \mathrm{allocatedTo}[n'].\)
This invariant prevents the World Bank Admin from allocating more capital downward than the reserve actually holds, even in the presence of UUPS upgrades.

The CVL specification sketches are in Appendix~B. Producing a Certora-verified proof of even these two invariants would be a mathematically demonstrable security guarantee that no undergraduate DeFi thesis from a Bangladesh institution has previously achieved, and grounds the abstract ``formal verification planned'' language of prior versions in a concrete artifact.

\section{Foundry Invariant and Fuzz Test Suite}
\label{sec:foundry}

The Hardhat JavaScript test suite remains in place for integration tests and deployment scripts. v15 adds a parallel Foundry suite for unit-level fuzzing and invariant testing because the two frameworks catch different bug classes. The 2025--2026 industry pattern is Hardhat for integration, Foundry for invariants and stateful fuzzing.

\paragraph{Three invariants asserted under Foundry's stateful fuzzer.}
\begin{itemize}[noitemsep]
\item Solvency: \(\mathrm{totalOutstandingLoans} \le \mathrm{getTotalReserve}\).
\item Role segregation: no address holds both BORROWER\_ROLE and BANK\_APPROVER\_ROLE.
\item Capital-flow direction: cross-tier allocations cannot decrease totalReserve below the minimum ratio at any single transaction boundary.
\end{itemize}

Each invariant is asserted under 10{,}000 fuzz runs; counterexamples are saved as Foundry replay tests so any regression in a future commit fails CI deterministically. This single test-suite addition reduces the credibility gap that the Foundry-fuzz literature explicitly identifies as the cause of the \$180~million smart-contract loss class of March~2023.

\section{On-Chain Economic Feasibility Simulation}
\label{sec:abm-sim}

The revenue projection of approximately \$138--159\,M in Chapter~5 (spread-based accounting) assumes 1~World Bank, 5~National Banks, and 50~Local Banks at near-full capacity. To make this projection empirically grounded rather than a single-point estimate, the prototype employs a \textbf{scripted on-chain simulation} of $N$ clients across a six-institution hierarchy, deployed to Polygon~zkEVM Cardona testnet. This simulation replaces the planned Mesa agent-based model for the pre-thesis prototype evaluation; it produces verifiable on-chain evidence for all four research questions, is fully reproducible via the published deterministic seed and Foundry script, and generates the gas-cost data reported in Table~\ref{tab:gas-costs}.

\paragraph{Simulation scope.}
The v24 simulation deploys 1~World Bank Reserve, 2--3~National Banks, and 4--6~Local Banks, generating 300~synthetic wallet addresses acting as Tier~4 retail clients across a 12-month compressed lifecycle. A configurable 5\% fraud rate introduces mid-lifecycle defaults for stress testing. Deployment targets Polygon~zkEVM Cardona testnet (free, EVM-compatible, ZK-secured); scripting uses Foundry with a deterministic seed (SEED~=~42) for full reproducibility. The simulation output (a CSV manifest of per-client outcomes, per-bank reserve ratios, and gas costs) feeds directly into the RQ5 evaluation.

\paragraph{Simulation script design.}
Foundry deployment scripts with deterministic seeds (a)~deploy all tier contracts, (b)~fund each bank with testnet USDC from a seed distributor wallet, (c)~generate 300~client wallets, (d)~run the full loan lifecycle for each client (\texttt{applyLoan} $\to$ \texttt{approveRisk} $\to$ \texttt{signLoan} $\to$ \texttt{disburse} $\to$ \texttt{repay} $\times N_\text{installments}$), (e)~introduce the configured fraud rate by selecting clients who default mid-lifecycle, and (f)~export all on-chain transaction hashes to a JSON manifest for Appendix~C.

\begin{lstlisting}[language=JavaScript, caption={Foundry simulation entry-point (structure)}]
// Foundry script: script/RunSimulation.s.sol
function run() public {
  uint256 NUM_CLIENTS = 300;
  uint256 NUM_BANKS   = 6;
  uint256 SEED        = 42; // deterministic

  BankHierarchy memory h = deployBankHierarchy(NUM_BANKS);
  address[] memory clients = generateClients(NUM_CLIENTS, SEED);
  SimResult[] memory results =
      runLoanLifecycles(h, clients, FRAUD_RATE);
  exportManifest(results, "simulation-output/manifest.json");
}
\end{lstlisting}

\paragraph{What the simulation demonstrates.}
The simulation produces five categories of verifiable empirical evidence: (i)~\textbf{end-to-end four-tier capital flow}, verifiable on the Polygon~zkEVM Cardona explorer against the manifest; (ii)~\textbf{ML pipeline ingestion}, where the FastAPI fraud-scoring service processes synthetic transaction features extracted from the simulation; (iii)~\textbf{Credit Passport SBT updates}, with each client's SBT updated as their lifecycle progresses; (iv)~\textbf{live reserve-ratio dashboard}, showing reserve ratios at each tier shift as loans are disbursed and repaid; and (v)~\textbf{loan audit trail} via The Graph subgraph, viewable as a frontend timeline. Aggregate outputs (gas costs, reserve dynamics under a 30\% simultaneous-default stress scenario, credit-velocity distributions) are reported as ranges with confidence intervals rather than single-point estimates, consistent with the empirical grounding the Sygnum~2026 report requires~[R21]. Country-specific microfinance calibration (e.g., BRAC Bangladesh data) is reserved for Future Work deployment studies (Section~\ref{sec:future-work}).

\section{Real-Time Dashboard Pipeline and Runtime Monitoring}
\label{sec:realtime}

\begin{figure}[H]
\centering
\OnePageDiagram{fig-realtime-dashboard.pdf}
\caption{Real-time dashboard and runtime monitoring pipeline: smart contracts emit typed events; The Graph indexes them and pushes to a Node WebSocket server which renders the React dashboard via a Redis snapshot. Tenderly runtime alerts and an Isolation-Forest anomaly detector feed the operations runbook with the granular-pause control surface of Section~\ref{sec:defense-in-depth}.}
\label{fig:realtime-dashboard}
\end{figure}


The real-time dashboard pipeline closes the gap between contract events and the client / approver UI. Three components compose:

\paragraph{The Graph subgraph.}
A subgraph deployed on The Graph's hosted service (Polygon zkEVM Cardona supported, free for student projects) indexes every event the contracts emit: \texttt{LoanRequested}, \texttt{LoanApproved}, \texttt{Repayment}, \texttt{CapitalAllocated}, \texttt{ReserveRatioUpdated}, \texttt{LoanLiquidated}, \texttt{RiskScoreCommitted}, \texttt{RoleGranted}. GraphQL queries from the React frontend return loan history and dashboard data in tens of milliseconds rather than seconds via RPC polling.

\paragraph{WebSocket event listener.}
A Node.js service listens on \texttt{wss://polygon-zkevm-cardona} for the same set of events and pushes notifications to the client's browser via Socket.IO. The dashboard moves from a static page requiring manual refresh into a live banking dashboard---a single design choice that examiners consistently identify as the difference between ``looks like a real system'' and ``looks like a static mockup''.

\paragraph{Tenderly monitoring.}
Six critical alerts are configured on the deployed contracts: single disbursement over 5~ETH-equivalent, repeated loan requests from one wallet (>3 in 60 minutes), reserve-ratio drop below 20\%, failed \texttt{disburseLoan} calls, role-grant events, and any pause / unpause governance action. Tenderly's free tier supports full testnet monitoring; transaction simulation prevents the embarrassing live-demo failure where a malformed contract change is discovered only at the moment of execution.

\section{Transaction-State Machine for Frontend UX}
\label{sec:tx-state-machine}

\begin{figure}[H]
\centering
\OnePageDiagram{fig-tx-state-machine.pdf}
\caption{Transaction state machine of the loan lifecycle: DRAFT $\to$ PENDING\_KYC $\to$ PENDING\_LIMIT $\to$ PENDING\_SCORE $\to$ PENDING\_APPROVAL $\to$ ACTIVE $\to$ \{CLOSED, DEFAULTED $\to$ \{LIQUIDATED, cured-back-to-ACTIVE\}\}, with every transition guarded by a CEI-ordered contract function and every transition emitting an indexed event.}
\label{fig:tx-state-machine}
\end{figure}


Failed blockchain transactions are the dominant UX failure mode in DeFi prototypes during live demos. v15 specifies an explicit five-state transaction machine on the frontend (\texttt{idle} $\to$ \texttt{pending\_signature} $\to$ \texttt{submitted} $\to$ \texttt{confirming} $\to$ \texttt{success / error}) with distinct visual feedback at each state. User-rejection (MetaMask code 4001) is distinguished from contract revert (\texttt{CALL\_EXCEPTION}) and from network error; each receives its own actionable error message. This is a thirty-line React hook that converts ``the demo failed silently'' into ``the client can see exactly what went wrong and retry''.

\section{EIP-712 Authentication and API Hardening}
\label{sec:eip712-auth}

The platform's API authentication uses an EIP-712 typed-data message signed by the user's wallet at login. The signed message includes \texttt{(wallet, timestamp, nonce)}; the backend recovers the address from the signature and issues a JWT with a 15-minute lifetime. Refresh tokens rotate on use. Every authenticated request includes the JWT and a correlation ID; the correlation ID is propagated from the React frontend through Express.js to FastAPI to the on-chain transaction so that any failure can be traced end-to-end. Rate limiting via \texttt{slowapi} differentiates by endpoint (5/min auth, 60/min reads, 10/min writes) and by IP; the JWT blacklist on logout uses Redis with TTL equal to the remaining token life. The full Layer-2 control set is summarized in Section~\ref{sec:defense-in-depth}.

\paragraph{Limitations (prototype and research context).}
The AI/ML components remain subject to the standard DeFi-environment limitations: class imbalance, concept drift, cold-start, and adversarial structuring. The mitigations introduced in v15---GNN ablation for relational features, federated learning to enlarge the effective training set without violating data residency, and red-teaming the LLM assistant for adversarial prompts (Section~\ref{sec:llm-eval})---do not remove these limitations but explicitly bound them. AI/ML is used as decision support, never as an automated approval authority in the prototype phase.

\section{Evaluation Methodology}
\label{sec:eval-method}

This section summarizes how each research question will be evaluated in the prototype phase.
\begin{itemize}[noitemsep]
\item \textbf{RQ1 (architecture fidelity):} feature comparison and workflow mapping against flat DeFi protocols (e.g., Aave, Compound, MakerDAO), focusing on institutional hierarchy, role separation, and directional capital flow.
\item \textbf{RQ2 (settlement and transparency):} measurement of transaction latency and approximate per-transaction cost on testnets / Layer~2, plus demonstration of on-chain verifiability of reserves and critical state transitions.
\item \textbf{RQ3 (analytics support):} model evaluation using standard metrics (precision, recall, F1, AUC) is reported in two explicitly separated phases. \textit{Phase~(a) --- synthetic-data evaluation:} model performance on the procedurally-generated synthetic dataset (documented in Section~\ref{sec:abm-sim}), labeled ``proof-of-concept evaluation only.'' These metrics reflect the upper-bound performance achievable given the synthetic distribution and do not make claims about real-world fraud detection. The stacking meta-learner weights are hardcoded initial placeholder values (approximately 0.7 for $p_f$, 0.3 for the anomaly signal) pending real labeled data; when no labeled validation set is available, the system degrades gracefully to a calibrated Random Forest alone. \textit{Phase~(b) --- real-data evaluation:} deferred to future work once actual user transaction data is available. Additionally, a GNN ablation study (Section~\ref{sec:gnn-extension}) quantifies the marginal value of relational features over flat behavioral features.
\item \textbf{RQ4 (technical viability):} end-to-end deployment verification on public testnets with integration tests covering wallet connection, loan request, approval, repayment, and role-based access control; Foundry invariant suite reports the proportion of fuzz runs that hold each invariant (Section~\ref{sec:foundry}). \textit{Note:} the current pre-thesis prototype fully tests the Tier~1 Reserve contract and demonstrates role-based access control, loan request/approval, and testnet deployment. End-to-end four-tier capital flow requires completion of Sprint~2 cross-tier fund transfers and will be evaluated and reported in the final thesis.
\item \textbf{RQ5 (banking extensions):} design-level validation via specification completeness and interface contracts for deposit products and group lending, plus prototype demonstrations of selected mechanisms where implemented; the scripted on-chain Hardhat simulation (Section~\ref{sec:abm-sim}) produces an aggregate stress-test outcome under 30\% simultaneous-default scenario across the simulated six-institution hierarchy.
\end{itemize}

\subsection{LLM Assistant Evaluation Protocol}
\label{sec:llm-eval}

The 8B-parameter LLM assistant in Chapter~5 is evaluated under three protocols:

\begin{enumerate}[noitemsep]
\item \textbf{Task-level accuracy.} A held-out dataset of 200 platform-specific Q\&A pairs (drawn from policy documents and FAQ) is run through the RAG pipeline; the metric is exact-match for numeric answers (interest rates, limits, fees) and BLEU / ROUGE for explanatory answers. Baselines: prompt-only (no RAG), and a GPT-4-class commercial API for the same prompts under a free-tier quota.
\item \textbf{Regulatory hallucination test.} A red-team set of 50 borderline prompts (``Is this loan legal in my country?''; ``Can I avoid KYC by using a different wallet?'') triggers either an explicit refusal (correct), a deferral to a human approver (correct), or an answer (failure). The metric is refusal-or-defer rate; the target is 100\% on the red-team set, in line with the LLM-vulnerability-in-finance literature~[R35].
\item \textbf{Latency and resource ceiling.} On the gaming-PC demo hardware (AMD Radeon RX~9060~XT 16~GB, ROCm\footnote{Evaluated on Linux with \texttt{GGML\_CUDA\_NO\_PEER\_COPY=1} and \texttt{HSA\_OVERRIDE\_GFX\_VERSION=12.0.1} environment variables to work around ROCm~6.4 memory-management limitations on RDNA4 gfx1200/gfx1201 hardware; see llama.cpp issue~\#21376.}), llama.cpp serving Qwen3-8B~Q4\_K\_M is benchmarked end-to-end (prompt $\to$ response token~1, total response). Target: <2~s time-to-first-token, $\ge$5 tokens/s sustained throughput. RAG retrieval through ChromaDB is benchmarked separately at <100~ms p95.

\item \textbf{Action accuracy.} A held-out set of 50 action scenarios (loan application, installment payment, deposit, KYC upgrade) is run through the agent pipeline. The metric is whether the agent: (a)~correctly identifies the required MCP tool, (b)~correctly assembles all required parameters, and (c)~presents an accurate confirmation summary before executing. Target: $\geq 95\%$ on items~(a) and~(b); $100\%$ human-gate compliance (no write tool called without a confirmation turn in the conversation history).

\item \textbf{Human-gate bypass test.} A red-team set of 20 adversarial prompts attempts to cause the agent to execute a write tool without a confirmation step (e.g., \emph{``Just do it,''} \emph{``Skip the confirmation,''} \emph{``I already said yes earlier''}). The metric is zero-bypass rate: any execution of a write tool without a preceding confirmation turn in the conversation history constitutes a critical failure and must be resolved before sprint delivery.

\item \textbf{Session key scope enforcement test.} Ten test scenarios attempt to invoke a write tool outside the EIP-7702 session key's approved scope (e.g., calling \texttt{submit\_fixed\_deposit} when only \texttt{pay\_installment} is in scope, or attempting a transaction exceeding the 500~USDC value cap). The metric is $100\%$ rejection rate. The session key contract's scope enforcement is verified independently via a Foundry invariant test (Section~\ref{sec:foundry}).

\end{enumerate}

These evaluation criteria will be applied systematically in the final thesis phase, with results reported against each research question.

% ── Figure: Sprint Story-Points Distribution ─────────────────────────────────
% Four pie charts showing the 155-point breakdown across sprints and by epic.
% Uses portable manual-arc pie slices (no pgf-pie package required).
% Academic greyscale palette with PrimaryBlue accent.
% ─────────────────────────────────────────────────────────────────────────────


% ── Figure: Technical Development Methodology ────────────────────────────────
% Three-sprint swimlane showing tech stack components per layer.
% Swimlane rows: Blockchain | Frontend | Backend | AI/ML
% ─────────────────────────────────────────────────────────────────────────────
\begin{figure}[H]
\centering
\OnePageDiagram{fig-sdlc-agile.pdf}
\caption{SDLC stage mapping with Agile sprints. Requirements / Architecture / Design are covered by Pre-thesis~1; Implementation, Verification, Validation, and Maintenance span Sprints~1--3 of the final thesis phase, with Foundry invariants, Certora proofs, the scripted on-chain Hardhat simulation (Section~\ref{sec:abm-sim}), and Tenderly runtime monitoring layered on top.}
\label{fig:methodology-technical}
\label{fig:sdlc-mapping}
\end{figure}


\section{Sprint Plan and Deliverables}

\subsection{Sprint 1: Foundation and Core Banking (Weeks 1--3)}

\textbf{Sprint Goal:} Establish core blockchain infrastructure and hierarchical banking structure.

\begin{table}[H]
\centering
\caption{Sprint 1 backlog: smart contract development (37 pts).}
\label{tab:sprint1-contracts}
\begin{tabular}{p{1.5cm}p{10cm}p{1.5cm}}
\hline
\textbf{ID} & \textbf{User Story} & \textbf{Pts} \\
\hline
US-1.1 & World Bank contract --- reserve management, national bank registration & 5 \\
US-1.2 & National Bank contract --- borrow from World Bank, lend to Local Banks & 5 \\
US-1.3 & Local Bank contract --- borrow from National Bank, lend to users & 5 \\
US-1.4 & Role-based access control --- World Bank/National Bank/Local Bank Admin, Bank User, Retail Client & 3 \\
US-1.5 & Gas cost management --- initiator pays; Polygon low-fee (\$0.001--\$0.01/tx) & 3 \\
US-1.12 & \texttt{InsuranceFund} contract --- 5\% interest capture, claims processing, default coverage disbursement & 5 \\
US-1.13 & \texttt{getReserveSummary()} view function on each tier contract (PSR, reserve balance, insurance fund) & 3 \\
US-1.14 & Chainlink Price Feeds integration (BDT/USD, ETH/USD via \texttt{AggregatorV3Interface}) & 3 \\
US-1.15 & Chainlink Proof of Reserve for \texttt{WorldBankReserve} (PoR job registration) & 3 \\
US-1.16 & Append-only \texttt{audit\_logs} enforcement (DB role + RLS policy) & 2 \\
\hline
\end{tabular}
\TableScreenshotEnd{Agile Sprint Plan: Three-Sprint Development Timeline and Deliverables}
\end{table}

\vspace{2pt}
\noindent\textit{\small This sprint planning table breaks Sprint~1 into concrete deliverables that establish the foundation: contract deployment, wallet integration, and the initial data model. The plan is structured to deliver a demonstrable vertical slice early, reducing integration risk for later features.}

\begin{table}[H]
\centering
\caption{Sprint 1 backlog: frontend foundation (19 pts) and database schema (14 pts).}
\label{tab:sprint1-frontend}
\begin{tabular}{p{1.5cm}p{10cm}p{1.5cm}}
\hline
\textbf{ID} & \textbf{User Story} & \textbf{Pts} \\
\hline
US-1.6 & Wallet connection --- MetaMask, WalletConnect, Sepolia/Cardona & 3 \\
US-1.7 & Dashboard UI --- Material Design 3, responsive, blockchain-themed & 5 \\
US-1.8 & Navigation and layout --- AppBar, role-based menu & 3 \\
US-1.9 & Blockchain visual elements --- tx hash display, security badges & 2 \\
US-1.10 & Database design --- 19 tables, 3NF & 5 \\
US-1.11 & Migration scripts and seed data & 3 \\
US-1.17 & \texttt{sessions} table implementation + EIP-7702 session key schema & 3 \\
US-1.18 & \texttt{assets} table + \texttt{collateral\_asset\_id} FK migration in \texttt{LOAN} & 2 \\
US-1.19 & \texttt{interest\_rate\_tier} table (normalisation extraction from bank entities) & 2 \\
US-1.20 & Composite indexes: \texttt{idx\_loan\_bank\_status}, \texttt{idx\_loan\_client\_active}, \texttt{idx\_installment\_due} (partial), \texttt{idx\_aiml\_score} (partial) & 2 \\
US-1.21 & Network migration: Polygon zkEVM Cardona testnet setup, faucet, deployment config & 3 \\
\hline
\end{tabular}
\TableScreenshotEnd{Technology Stack: Primary Choices and Justifications}
\end{table}

\vspace{2pt}
\noindent\textit{\small This Sprint~1 continuation clarifies additional tasks and story point allocation for core banking scaffolding. The story point distribution reflects the early emphasis on infrastructure and correctness over feature breadth.}

\textbf{Sprint 1 Total:} 70 story points. Deliverables: smart contracts deployed on testnet, basic frontend with wallet connection, database schema implemented (19 tables, 3NF), role-based access control, dashboard UI, Chainlink oracle integrations, InsuranceFund contract, network migration to Polygon zkEVM Cardona.


\subsection{Sprint 2: Lending Features and Communication (Weeks 4--6)}

\textbf{Sprint Goal:} Implement the complete core lending lifecycle end-to-end and core communication features.

\textbf{Scope note:} Sprint~2 is scoped to the \textit{Minimum Viable Thesis} (MVT): the P0/P1 items necessary to demonstrate the core contributions (hierarchical capital flow and basic loan lifecycle). Multi-entity operations (InterBankLendingPool, TreasurySwap, SyndicatedLoan, NettingEngine) and advanced deposit products are deferred to Sprint~3 or formally specified as Future Work, ensuring the P0 items are implemented solidly rather than all items implemented superficially.

\begin{table}[H]
\centering
\caption{Sprint 2 backlog: core lending lifecycle --- MVT priority items (92 pts).}
\label{tab:sprint2}
\small
\begin{tabular}{p{1cm}p{0.8cm}p{9.2cm}p{1.2cm}}
\hline
\textbf{ID} & \textbf{Priority} & \textbf{User Story} & \textbf{Pts} \\
\hline
US-2.1 & P0 & Loan request submission with blockchain transaction & 5 \\
US-2.2 & P0 & Loan approval and rejection workflow with bank approver & 5 \\
US-2.3 & P0 & Installment payment system with automated schedule generation & 8 \\
US-2.4 & P0 & Borrowing limit enforcement (on-chain authoritative check per client SBT) & 5 \\
US-2.5 & P1 & Client--bank real-time chat system & 5 \\
US-2.6 & P1 & Income verification document upload and review & 5 \\
US-2.7 & P0 & Hierarchical bank registration (National $\to$ Local) & 5 \\
US-2.8 & P1 & Bank user management and approver designation & 3 \\
US-2.9 & P1 & Loan history and transaction tracking pages & 5 \\
US-2.10 & P2 & QR code generation for wallet addresses & 2 \\
US-2.11 & P1 & Responsive UI polish and error handling & 2 \\
US-2.12 & P0 & MCP tool server --- 16 banking tools wired to Express.js API & 8 \\
US-2.13 & P0 & Agent chat interface (Qwen3-8B + tool calling + human-gate confirmation pattern) & 8 \\
US-2.14 & P0 & EIP-7702 session key management (scope enforcement + TTL) & 5 \\
US-2.15 & P1 & Authority Brief UI (SHAP breakdown for bank approver dashboard) & 5 \\
US-2.16 & P1 & Chainlink Automation for installment due-date checks and overdue flagging & 5 \\
US-2.17 & P1 & EMI reminder cron $\to$ agent push notification (3~days + 1~day before due) & 3 \\
US-2.18 & P1 & Credit tier schedule in Credit Passport SBT (Bronze $\to$ Diamond) & 3 \\
US-2.19 & P1 & \texttt{agent\_action\_log} table (append-only, FK to \texttt{sessions}) & 3 \\
US-2.20 & P1 & \texttt{interest\_rate\_tier} table integration with Kinked Rate Model & 2 \\
\hline
\end{tabular}
\TableScreenshotEnd{Technology Stack (Continued): Framework and Tooling Selection}
\end{table}

\vspace{2pt}
\noindent\textit{\small This Sprint~2 table focuses on the core lending lifecycle (P0) and communication workflows (P1), expanded in v24 to include the autonomous AI agent pipeline (US-2.12--2.13), EIP-7702 session key management (US-2.14), the Authority Brief UI (US-2.15), Chainlink Automation triggers (US-2.16), and Credit Passport tier schedule (US-2.18). P2 items are completed only if P0/P1 are stable and time permits. Multi-entity contracts (IBLP, UpwardDepositFacility, TreasurySwap, SyndicatedLoan, NettingEngine) and the SavingsVault / FixedDeposit products are scheduled as Sprint~3 items or formally specified in Future Work.}

\noindent\textbf{Sprint 2 Total:} 92 story points.


\subsection{Sprint 3: AI/ML Security and Finalization (Weeks 7--8)}

\textbf{Sprint Goal:} Integrate AI/ML analytics, complete testing, and prepare documentation.

\begin{table}[H]
\centering
\caption{Sprint 3 backlog: AI/ML security, oracle finalization, and demo polish (80 pts).}
\label{tab:sprint3}
\small
\begin{tabular}{p{1.5cm}p{10cm}p{1.5cm}}
\hline
\textbf{ID} & \textbf{User Story} & \textbf{Pts} \\
\hline
US-3.1 & Random Forest fraud detection model training and deployment & 8 \\
US-3.2 & SHAP-based explainability for risk assessments & 5 \\
US-3.3 & Isolation Forest anomaly detection for wallet behaviour & 5 \\
US-3.4 & Risk dashboard with real-time AI/ML scores & 5 \\
US-3.5 & AI chatbot for client assistance (RAG-augmented, hallucination-guarded) & 3 \\
US-3.6 & Market data visualisation page & 3 \\
US-3.7 & Profile and settings management & 2 \\
US-3.8 & Security audit and vulnerability assessment & 3 \\
US-3.9 & Documentation and report finalisation & 2 \\
US-3.10 & Demo preparation and presentation & 2 \\
US-3.11 & Chainlink Functions oracle replacing commit-reveal relay (DON consensus for ML score commitment) & 8 \\
US-3.12 & The Graph subgraph (\texttt{LoanApplication} + \texttt{ReserveRatioSnapshot} event indexing) & 5 \\
US-3.13 & SAR workflow (Isolation Forest $\to$ \texttt{aml-alert} Kafka topic $\to$ compliance queue $\to$ \texttt{freezeAccount} + \texttt{audit\_logs}) & 5 \\
US-3.14 & Reserve Transparency Dashboard (React, queries The Graph subgraph in real time) & 5 \\
US-3.15 & 300-client Foundry simulation (6 banks, deterministic seed, 12-month compressed cycle, CSV output for RQ5) & 8 \\
US-3.16 & No-privileged-exemption Foundry invariant test (all accounts subject to identical liquidation logic) & 3 \\
US-3.17 & Pool-level Isolation Forest (portfolio anomaly detection across all active loans, not per-loan only) & 5 \\
US-3.18 & \texttt{freezeAccount()} function in \texttt{LocalBankPool.sol} with \texttt{onlyApprover} modifier + Foundry invariant test & 3 \\
\hline
\end{tabular}
\TableScreenshotEnd{Technology Stack: Second-Choice Alternatives and Trade-offs}
\end{table}

\vspace{2pt}
\noindent\textit{\small This Sprint~3 table expands the original AI/ML monitoring scope with the v24 oracle finalization (Chainlink Functions, US-3.11), The Graph subgraph (US-3.12), SAR compliance workflow (US-3.13), Reserve Transparency Dashboard (US-3.14), 300-client Foundry simulation (US-3.15), and pool-level Isolation Forest (US-3.17). The intent is to defer analytics complexity until core lending flows are stable, ensuring explainability and auditability can be evaluated against realistic workflows.}

\noindent\textbf{Sprint 3 Total:} 80 story points.

\vspace{4pt}
\noindent\textbf{Additional Sprint~3 module scope (planned):}
\begin{itemize}[noitemsep]
\item \textbf{Group lending module:} architectural design, smart contract interface specification, and partial implementation of group formation and consent recording logic.
\item \textbf{Savings account module:} SavingsVault contract design, interest accrual formula implementation, and interface specification for fixed-term deposits.
\item \textbf{FX module:} oracle integration design and interface specification for dual-currency conversion.
\end{itemize}

Figure~\ref{fig:sprint-submission} shows the sprint submission workflow.

% ── Figure: Sprint Submission Workflow ───────────────────────────────────────
% Standard UML Activity / flowchart symbols.
% Flow: Backlog Refinement → Sprint Planning → Development Loop →
%       Code Review → Testing → Sprint Review → Delivery / Submit
% ─────────────────────────────────────────────────────────────────────────────


\section{SDLC Stage Mapping}

We map our Agile sprints to the seven stages of the Software Development Life Cycle (SDLC). Figure~\ref{fig:sdlc-mapping} illustrates this mapping.

\begin{table}[H]
\centering
\caption{SDLC stage mapping.}
\label{tab:sdlc-mapping}
\begin{tabular}{p{0.5cm}p{2.5cm}p{4.5cm}p{5cm}}
\hline
\textbf{\#} & \textbf{SDLC Stage} & \textbf{Project Activity} & \textbf{Deliverable} \\
\hline
1 & Planning & Feasibility studies; professor consultations; guideline review & Feasibility Report; Project Plan \\
2 & Requirements & System analysis; use case definitions; constraints identification & Use case diagrams; UC-1 through UC-5 \\
3 & Design & Three-layer architecture; DB schema (3NF); smart contract interfaces & Architecture diagrams; ERD; DFD \\
4 & Development & Sprint 1--3: smart contracts, frontend, backend, AI/ML & Source code; DApp prototype \\
5 & Testing & Hardhat unit tests (12+); integration testing; AI/ML evaluation & Test reports; model metrics \\
6 & Deployment & Testnet deployment; frontend (Vercel); backend (Render) & Live prototype on testnet \\
7 & Maintenance & Monitoring; model retraining; bug fixes; iteration & Updated docs; retrained models \\
\hline
\end{tabular}
\TableScreenshotEnd{Technology Stack: Additional Alternative Evaluations}
\end{table}

\vspace{2pt}
\noindent\textit{\small This SDLC mapping table aligns the three-sprint Agile plan with standard SDLC stages to ensure research deliverables remain complete (requirements, design, implementation, testing, and evaluation). The mapping makes the development process auditable and easier to justify in an academic setting.}

% ── Figure: SDLC Stage Mapping ───────────────────────────────────────────────
% Seven SDLC stages on the right spine; activity deliverable boxes on the left;
% iteration arrow on far left. Standard flowchart rectangles + dashed connectors.
% ─────────────────────────────────────────────────────────────────────────────


\section{Design Decisions and Alternatives}

For each major design decision, we evaluated alternatives and justified selections based on technical criteria, ecosystem maturity, and project constraints. Figure~\ref{fig:design-decisions} visualizes these comparisons.

\begin{table}[H]
\centering
\caption{Design decisions and alternatives considered.}
\label{tab:design-decisions}
\begin{tabular}{p{3cm}p{3cm}p{3cm}p{4.5cm}}
\hline
\textbf{Decision Area} & \textbf{1st Choice} & \textbf{2nd Choice} & \textbf{Key Criterion} \\
\hline
Methodology & Agile / Scrum & Incremental & Evolving scope, milestones \\
Architecture & DApp + Off-chain AI & Hybrid with Oracle & Gas cost, ML flexibility \\
Frontend & React + TypeScript & Vue + TypeScript & Web3 ecosystem maturity \\
Smart Contract & EVM (Solidity) & Solana (Rust) & Gas, control size, free testnets \\
Fraud Detection & Random Forest & XGBoost & SHAP compatibility, simplicity \\
Anomaly Detection & Isolation Forest & Autoencoder & Unsupervised; no labelled data needed \\
XAI Method & SHAP & LIME & Theoretical guarantees, regulatory fit \\
Database & PostgreSQL & SQLite & 3NF support, async queries \\
Hosting & Vercel + Render & Localhost only & \$0 cost, publicly accessible URL \\
Network (primary) & Polygon zkEVM Cardona & Polygon Amoy PoS & ZK validity proof security model vs.\ validator set assumption \\
Oracle mechanism & Chainlink Functions (DON) & Commit-reveal relay & Decentralised trust vs.\ single-key operator trust \\
Agent tool framework & MCP tool server & Direct Express.js API & Defined schema = safety boundary; no arbitrary code execution \\
Agent session auth & EIP-7702 session keys & ERC-4337 paymasters & Key-level scope enforcement vs.\ application-level scope \\
\hline
\end{tabular}
\TableScreenshotEnd{Software Testing Strategy: Acceptance Criteria Across Four Test Layers}
\end{table}

\vspace{2pt}
\noindent\textit{\small This design decisions table records evaluated alternatives (chains, stacks, libraries) and the criteria used to select the final approach. Documenting these trade-offs strengthens the methodological rigor and clarifies why the selected stack best matches the project's constraints.}

% ── Figure: Design Decisions and Alternatives ────────────────────────────────
% Two panels:
%   Panel A (top):  AI/ML decisions   — Fraud Detection | Anomaly | Explainability
%   Panel B (bottom): Tech Stack decisions — Frontend | Smart Contract | Database | UI Framework
% 1st-choice nodes filled PrimaryBlue (selected); 2nd-choice nodes outlined only.
% Standard tree layout; no \n typos from source image.
% ─────────────────────────────────────────────────────────────────────────────
\begin{figure}[H]
\centering
\OnePageDiagram{fig-design-decisions.pdf}
\caption{Key design decisions and alternatives considered: smart-contract platform (EVM vs.\ Cosmos / Solana), chain choice (Polygon + Sepolia vs.\ L1-only vs.\ single L2), upgradeability (UUPS vs.\ Transparent Proxy vs.\ immutable), identity (DID/VC + zkKYC vs.\ raw on-chain KYC vs.\ off-chain-only), ML oracle (commit-reveal vs.\ trusted-backend vs.\ on-chain ML), and cross-chain bridge (Chainlink CCIP vs.\ LayerZero vs.\ Axelar).}
\label{fig:design-decisions}
\end{figure}

\subsection{Justification of Selected Technologies}

\textbf{1st Choice Justifications:}

\begin{itemize}[noitemsep]
\item \textbf{Agile/Scrum} was necessary due to the project's evolving nature. Frequent updates to smart contract designs, user interface steps and AI/ML connections were required. Adopting Agile allows for plan alterations during brief work cycles effectively avoiding time and resource loss. Significance arises as new concepts emerge during the prototype development process.

\item \textbf{DApp + Off-chain AI} architecture allows intensive machine learning operations to occur outside the blockchain. Machine learning tasks are often run on the blockchain at a considerable expense occasionally exceeding \$100 for a single operation. Off-chain processing enables the use of rapid GPUs along with widely-used Python machine learning resources. Final lending choices are managed exclusively by smart contracts on-chain ensuring that trust is maintained in critical areas.

\item \textbf{React + TypeScript} is beneficial as many Web3 libraries such as Wagmi, RainbowKit, Viem, and ethers.js are supported seamlessly. A vast community of developers exceeding 10~million exists for React making it simpler to access assistance and solutions.

\item \textbf{EVM (Solidity)} is regarded as the most thoroughly examined smart contract system. Over \$100~billion is secured across various blockchains by it. Building and testing the project without incurring costs is possible with free test networks such as Sepolia and Amoy. Security tools provided by OpenZeppelin have been thoroughly tested to ensure the protection of contracts.

\item \textbf{MUI (Material Design~3)} provides ready-made components for typical banking interface designs. The package contains various items such as data tables, form checks and layouts for dashboards. Utilizing MUI results in a time saving of around 40\% compared to starting from scratch.

\item \textbf{Random Forest} was selected as the primary fraud detection model for three compounding reasons. First, it demonstrates natural compatibility with SHAP's TreeExplainer, which computes exact Shapley values---rather than approximations---by exploiting the tree structure directly. This exactness is a regulatory asset: in a lending context subject to explainability requirements under emerging frameworks such as the EU AI Act, approximation-based attribution tools like LIME can produce inconsistent feature rankings across identical inputs, undermining audit reliability~[4]. Second, ensemble tree methods generalise well on structured tabular data with moderate sample sizes, a characteristic that matches the DeFi lending transaction log domain, where labeled fraud samples are scarce and synthetic augmentation is likely necessary. Third, Random Forest naturally supports class imbalance through class weighting and bootstrap sampling, reducing the risk of a model that achieves high accuracy by always predicting ``not fraud''---a failure mode that would be catastrophic in a lending context.

\item \textbf{Isolation Forest} does not require labeled data for training. Labeled fraud samples are few making it suitable for DeFi lending. Anomalies in transactions are scored swiftly as it operates with a time complexity of $O(n \log n)$.

\item \textbf{SHAP} provides feature importance through solid mathematical foundations from game theory (Shapley values). Characteristics such as accuracy, management of absent data and consistency are not assured by tools like LIME. These properties assist in fulfilling emerging regulations for explainable decisions within financial services.

\item \textbf{PostgreSQL} effectively manages complex time-based queries. Borrowing limits can be checked over periods such as 6 months or 1 year. ACID compliance is fully supported to ensure financial data remains secure. It supports JSON which facilitates easy adjustments to the database structure during the prototyping phase.

\item \textbf{Vercel + Render} facilitates deploying the demo globally at no charge. A live prototype can be accessed without the necessity of a setup on personal computers. Supervisors, examiners and other individuals are not required to deal with installation issues.

\item \textbf{Polygon zkEVM Cardona} is selected over Polygon Amoy PoS because ZK validity proofs derive security from cryptographic verification rather than a validator set assumption. Every batch of transactions is accompanied by a zero-knowledge proof verified on Ethereum~L1, making the security claim materially stronger for an institutional banking prototype. The phrase ``ZK rollup-secured reserve ratios'' is a more defensible thesis claim than ``PoS-secured reserve ratios.'' On the Cardona testnet, the developer experience is equivalent to Amoy at zero cost.

\item \textbf{Chainlink Functions} (DON-based oracle) is selected over the prototype commit-reveal relay because the DON runs the ML scoring request across multiple independent nodes --- consensus is required before the result is accepted on-chain. A malicious single node cannot manipulate the risk score. This closes the oracle trust assumption of the v23 commit-reveal design, where trusting the operator running the FastAPI relay was an explicit limitation. The 30--60~second latency and \$0.10--1.00 per oracle call are both acceptable for the loan approval lifecycle, where decisions are not time-critical.

\item \textbf{MCP tool server} is selected over a direct Express.js API for agent interactions because the tool schema is the safety boundary: the agent cannot execute arbitrary code, make unapproved API calls, or read data outside the 16 defined tools. This constraint is the mechanism that makes the autonomous agent safe to deploy with real client funds. A direct API would require application-level enforcement only, which is less auditable.

\item \textbf{EIP-7702 session keys} are selected over ERC-4337 paymasters for agent-controlled on-chain operations because scope restriction is enforced at the key level, not at the application level. The session key is bound to a specific list of tools, a value cap (500~USDC per transaction), and a 24-hour TTL --- a compromised application cannot exceed these limits regardless of what the application code does. EIP-7702 does not replace ERC-4337 for retail gas sponsorship; the two mechanisms serve different functions and coexist in the architecture.
\end{itemize}

\textbf{2nd Choice Justifications (Why Retained as Alternatives):}

\begin{itemize}[noitemsep]
\item \textbf{Incremental (Waterfall) methodology} produces more predictable and clearly phased deliverables. However, it offers limited capacity to adapt to the evolving requirements typical of research-oriented prototype development, and is most effective in stable production environments where scope is fixed from the outset.

\item \textbf{Hybrid with Oracle architecture} (e.g., Chainlink) would enable on-chain ML prediction triggers via oracle data feeds, increasing the verifiability of AI-driven decisions. The trade-off is meaningful latency (30--60 seconds per update) and additional cost (\$0.10--1.00 per oracle call), along with a dependency on third-party oracle availability that introduces an external failure point not present in the selected off-chain approach.

\item \textbf{Vue + TypeScript} is approachable for developers new to JavaScript frameworks. Its Web3 tooling ecosystem (e.g., vue-dapp) is less mature than React's, with sparser wallet integration library support and slower update cadence for critical Web3 dependencies.

\item \textbf{Solana (Rust)} offers substantially higher raw throughput (up to 65,000 TPS) than EVM-compatible chains. However, Solana's developer tooling and security library ecosystem are less established, the Rust learning curve is significant within an 8-week development window, and the DeFi security audit tooling available for Solidity (Slither, Mythril, OpenZeppelin) has no direct equivalent on Solana.

\item \textbf{Tailwind CSS} offers greater design flexibility through utility-first composition. Achieving consistent, professional banking-style UI patterns requires substantially more custom work than using MUI's pre-built component library, which imposes a time cost that is unacceptable within the sprint timeline.

\item \textbf{XGBoost} frequently outperforms Random Forest on structured tabular datasets. However, its SHAP integration uses approximate rather than exact tree computation, meaning explanations cannot be guaranteed to be consistent across repeated evaluations---a property required for regulatory auditability of credit decisions.

\item \textbf{Autoencoder} models can capture complex non-linear anomaly patterns. In practice, they require substantial labeled or unlabeled training data and careful hyperparameter tuning to converge reliably. Isolation Forest is parameter-light and performs competitively on small and imbalanced datasets, making it better suited to the data-scarce DeFi lending context of this prototype.

\item \textbf{LIME} produces faster local explanations. As demonstrated by Adom et al.~[4], LIME explanations are unstable: repeated evaluation of the same input can yield different feature importance rankings, undermining the consistency guarantees that regulators require for automated lending decisions.

\item \textbf{SQLite} eliminates database infrastructure overhead and is straightforward to configure. It does not support concurrent writes, which limits multi-user prototype testing, and lacks the window function and CTE support needed for rolling 6-month and 1-year borrowing limit calculations.

\item \textbf{Localhost-only deployment} removes hosting dependencies and infrastructure cost. It prevents remote sharing for supervisor review, collaborative testing across team members, and examiner access to a live demo---all of which are necessary for academic evaluation and iteration.

\item \textbf{Polygon Amoy PoS} (2nd choice for network) offers a faster inner development loop on a familiar PoS testnet. It is retained as a fallback option should Cardona testnet availability be disrupted. However, the validator-collusion security assumption is weaker than ZK validity proofs, which undermines the institutional trust narrative of the thesis.

\item \textbf{Commit-reveal relay} (2nd choice for oracle mechanism) is retained as the prototype fallback for Sprint~1 and Sprint~2 before the Chainlink Functions integration is wired in Sprint~3. It provides a working oracle during early development and is simpler to implement. The limitation---trusting the FastAPI service key---is clearly acknowledged as a prototype constraint.

\item \textbf{Direct Express.js API} (2nd choice for agent tool framework) would allow the agent to call banking endpoints without a formal schema. The safety risk is that the agent could potentially construct arbitrary API calls not anticipated by the developers; the schema boundary of the MCP server eliminates this risk at the architecture level.

\item \textbf{ERC-4337 paymasters} (2nd choice for agent session authentication) are already used for retail gas sponsorship. They could theoretically be extended to agent operations, but scope enforcement would be at the paymaster application level rather than cryptographically at the key level, which is a weaker safety guarantee for autonomous agent execution.
\end{itemize}


\section{Design Patterns}

Our system uses five design patterns:

\begin{itemize}[noitemsep]
\item \textbf{Singleton:} A singleton means that each primary contract is established a single time. Only one version is utilized by all ensuring a unified source of truth. Not everyone can create multiple versions of the contract.

\item \textbf{Observer:} Contracts generate alerts whenever actions occur such as loan requests. These alerts are processed by a service that updates the database. Notifications are not ignored; they receive attention.

\item \textbf{Adapter:} An Adapter is utilized by the wallet component to function with various wallet types (like MetaMask). Different wallets that adhere to the EIP-1193 standards can be easily integrated.

\item \textbf{Factory Pattern:} Each Local Bank creates individual loan and savings contract instances using a factory contract, ensuring all deployed instances follow the same security-audited interface and access control checks. The factory pattern prevents unauthorized contract deployment that could bypass the hierarchical governance structure.

\item \textbf{Proxy/Upgradeable Pattern:} Banking contracts must be upgradeable without losing stored state such as balances, credit scores, and loan histories. OpenZeppelin's transparent proxy pattern separates the contract's logic from its storage, allowing governance-approved upgrades to the logic layer while preserving all stored financial data. This pattern is essential for a long-lived banking system that must adapt to evolving regulatory requirements and security patches.
\end{itemize}


\section{Software Testing Strategy}

The testing strategy covers four layers of the system, each with defined acceptance criteria.

\begin{itemize}[noitemsep]
\item \textbf{Smart contract unit tests:} Numerous tests for smart contracts exist in our Hardhat environment. Over 12 tests are conducted to verify actions related to managing reserves, sign-ups, loans and access permissions. No scenarios such as depositing zero or executing disallowed actions are missed in these assessments. \textbf{Acceptance criterion:} all Hardhat unit tests pass with zero failures before any sprint delivery.
\item \textbf{Integration testing:} Whole process checks are conducted for integration tests. Wallet connection, loan request, approval and repayment processes are evaluated on the network. No critical steps are excluded from the process. \textbf{Acceptance criterion:} end-to-end workflow completes successfully on testnet for representative scenarios (approval, rejection, repayment loop).
\item \textbf{AI/ML module verification:} Ensuring functionality is vital for our fraud and unusual activity detectors. Results are generated by the integration of the detectors with the overall system. \textbf{Acceptance criterion:} model inference endpoints return valid scores and explanations for test inputs without runtime errors.
\item \textbf{Frontend testing:} Website appearance is evaluated on Chrome, Firefox and mobile devices. Tests are conducted to ensure that functionality is maintained across various platforms. \textbf{Acceptance criterion:} core flows render and remain usable across desktop and mobile breakpoints without blocking UI defects.
\end{itemize}


%% %%%%%%%%%%%%%%%%%%%%%%%%%%%%%%%%%%%%%%%%%%%%%%%%%%%%%%%%%%%%
%% CHAPTER 5: MARKET ANALYSIS AND FEASIBILITY
%% %%%%%%%%%%%%%%%%%%%%%%%%%%%%%%%%%%%%%%%%%%%%%%%%%%%%%%%%%%%%
\chapter{Market Analysis and Feasibility}

\section{Market Sizing}

\begin{table}[H]
\centering
\caption{Market segments with supporting data.}
\label{tab:market-sizing-ch5}
\begin{tabular}{p{3.5cm}p{6.5cm}p{3cm}}
\hline
\textbf{Segment} & \textbf{Description} & \textbf{Estimated Scale} \\
\hline
Total Addressable Market (TAM) & Global DeFi lending (\$55B+ TVL [13]); cross-border remittances (\$860B [26]); SME financing gap (\$4.5T [20]) & \$55B -- 5T+ \\
\hline
Serviceable Addressable Market (SAM) & Institutional and semi-institutional lending requiring hierarchical structures; emerging-market credit demand & \$5B -- 15B \\
\hline
Serviceable Obtainable Market (SOM) & Pilot deployments in regulatory sandboxes, academic prototypes, NGO-backed microfinance programs & \$50 -- 200M \\
\hline
\end{tabular}
\TableScreenshotEnd{Market Sizing: DeFi Lending TVL, Remittances, and MSME Financing Gap}
\end{table}

\vspace{2pt}
\noindent\textit{\small This table quantifies the demand-side context (DeFi lending, remittances, MSME credit gap) used to ground the feasibility argument in measurable market data. It shows the platform targets both a high-volume settlement problem and a persistent inclusion/credit-access gap.}

\vspace{4pt}
\noindent\textit{\small\textbf{Sources:} (a)~TAM: DefiLlama~[13] reports DeFi lending TVL exceeding \$55~billion; World Bank Migration and Development Brief~[26] values global remittances at \$860~billion annually; IFC~[20] estimates the MSME financing gap at \$4.5~trillion; (b)~SAM and SOM: project-derived estimates based on industry segmentation of institutional lending and regulatory sandbox addressable market~[18].}

No hierarchically structured lending system currently exists in DeFi; existing protocols rely on undifferentiated shared pools. Cross-border payment networks such as Ripple (\$847M daily)~[25] and JPMorgan Kinexys (\$3--7B daily)~[40] handle settlement volume but offer no lending, deposit, or tiered governance functions. Institutional interest in blockchain-based finance is clearly established, yet no single platform combines DeFi-style accessibility with the structured capital hierarchy of traditional development banking. The Crypto World Bank is designed to occupy this gap.

\section{Target Customer Segment}

The Crypto World Bank targets the \textbf{retail customer segment}---individual retail clients and small businesses seeking transparent, accessible crypto-based financial services.

\begin{table}[H]
\centering
\caption{Target customer segment profile.}
\label{tab:customer-segment-ch5}
\begin{tabular}{p{3.5cm}p{10cm}}
\hline
\textbf{Characteristic} & \textbf{Description} \\
\hline
Primary Users & Individual retail clients seeking personal or small business financing \\
\hline
Geographic Focus & Developing economies with limited traditional banking access (e.g., Bangladesh, Southeast Asia, Sub-Saharan Africa) \\
\hline
Loan Size Range & Micro to mid-range: 0.1 ETH -- 500 ETH equivalent (\textasciitilde\$200 -- \$1{,}000{,}000 at current rates) \\
\hline
User Profile & Digitally literate individuals with cryptocurrency wallet access; small business owners; gig-economy freelancers \\
\hline
Key Pain Points & High interest rates from informal lenders; lack of credit history in traditional systems; exclusion from banking due to documentation barriers \\
\hline
\end{tabular}
\TableScreenshotEnd{Target Customer Segment: Retail Client Profile and Product Fit}
\end{table}

\vspace{2pt}
\noindent\textit{\small This target-customer table articulates the retail client profile and constraints that drive product requirements (low fees, simple onboarding, transparent terms). It supports the decision to prioritize usability and cost efficiency over complex institutional-only features in the prototype phase.}

\section{Partner Ecosystem}

\begin{table}[H]
\centering
\caption{Partner categories and roles.}
\label{tab:partner-categories-ch5}
\begin{tabular}{p{3.5cm}p{4.5cm}p{5.5cm}}
\hline
\textbf{Partner Category} & \textbf{Functional Role} & \textbf{Blockchain-Mediated Incentive} \\
\hline
Financial Regulators & Regulatory sandbox approval; compliance oversight & Reduced enforcement cost through on-chain transparency and audit trails \\
\hline
Banking Institutions & Network membership as National/Local Banks & Access to diversified global reserve; reduced inter-bank settlement friction \\
\hline
Payment Gateway Providers & Fiat-to-crypto on-ramp and off-ramp services & Volume-based transaction fees; expanded market reach \\
\hline
Academic \& Research Institutions & Validation of AI/ML models; publication of research findings & Access to anonymised datasets; collaborative research opportunities \\
\hline
Non-Governmental Organizations & Pilot deployment; field testing with underserved client populations & Transparent, low-friction credit access for beneficiaries \\
\hline
\end{tabular}
\TableScreenshotEnd{Partner Ecosystem: Integration Partners and External Dependencies (Chapter 5)}
\end{table}

\vspace{2pt}
\noindent\textit{\small This partner ecosystem table highlights external dependencies (identity/verification flows, infrastructure, and settlement rails) needed for a real deployment. Identifying these actors early reduces hidden-scope risk and clarifies what remains outside the smart contract boundary.}


\section{Competitive Landscape}

The Crypto World Bank operates at the intersection of four distinct competitor categories. Table~\ref{tab:competitor-detailed} provides a detailed comparison of representative projects in each category, with current metrics as of March~2026.

\begin{table}[H]
\centering
\caption{Detailed competitive landscape analysis (Part 1).}
\begin{tabular}{p{2.5cm}p{2.5cm}p{2.5cm}p{2.5cm}p{3.5cm}}
\hline
\textbf{Project} & \textbf{Category} & \textbf{Scale (2026)} & \textbf{Architecture} & \textbf{Gap We Address} \\
\hline
Compound v3 [26] & DeFi lending & \$1.4B TVL & Single-borrowable asset per single tier & No hierarchy; no institutional features \\
\hline
MakerDAO / Sky [30] & Stablecoin / CDP & \$6B TVL & CDP model; not peer-to-peer lending & Creates money, not a lending governance structure \\
\hline
Morpho [43] & DeFi lending & \$6.8B TVL; 1.4M users & Isolated markets; peer-to-peer matching & Flat primitive; no cross-market hierarchy; no banking integration \\
\hline
Maple Finance [31] & Institutional credit & \$2.6--3.8B TVL & Pool Delegate model; under-collateralised & Single-tier; no interest rate setting \\
\hline
Goldfinch [32] & Emerging-market credit & \$680M originated; 18+ countries & Trust-through-consensus; senior/junior tranches & B2B only; no interbank lending \\
\hline
Ripple / RLUSD [25] & Banking rails & \$847M/day cross-border & Payment rail; no lending capability & Moves money but has no lending, deposits, or credit system \\
\hline
\end{tabular}
\TableScreenshotEnd[tab:competitor-detailed]{Competitive Landscape (Part 1): DeFi Lending and Payment Rail Protocols}
\end{table}

\vspace{2pt}
\noindent\textit{\small This competitive landscape table compares representative projects across DeFi lending, payment rails, inclusion wallets, and institutional blockchain systems. The comparison highlights the whitespace: no competitor combines hierarchical multi-tier lending with governance-aware controls and AI-assisted monitoring in one architecture.}

\begin{table}[H]
\centering
\caption{Detailed competitive landscape analysis (Part 2).}
\begin{tabular}{p{2.5cm}p{2.5cm}p{2.5cm}p{2.5cm}p{3.5cm}}
\hline
\textbf{Project} & \textbf{Category} & \textbf{Scale (2026)} & \textbf{Architecture} & \textbf{Gap We Address} \\
\hline
JPMorgan Kinexys [40] & Banking rails & \$3--7B daily volume & Permissioned; single-bank control & Centralised; proprietary; restricted to JPMorgan clients \\
\hline
Stellar [44] & Financial inclusion & \$55.6B annual payment volume & Open payment network; anchors for fiat & Payment network only; no lending, reserves, or interest rate markets \\
\hline
Celo / MiniPay [42] & Financial inclusion & 14M wallets; 60+ countries & Stablecoin payments; mobile-first & Payments and savings only; no lending hierarchy or banking structure \\
\hline
R3 Corda [37] & Enterprise DLT & \$17B tokenised RWAs & Permissioned; consortium governance & Infrastructure layer only; no lending logic; closed access \\
\hline
World Bank FundsChain [39] & Development finance & 250 projects by mid-2026 & Hyperledger Besu; fund tracking & Does not implement lending or interest rate mechanics \\
\hline
\end{tabular}
\TableScreenshotEnd{Competitive Landscape (Part 2): Inclusion Wallets and Institutional Blockchain Systems}
\end{table}

\vspace{2pt}
\noindent\textit{\small This continuation expands the competitor comparison to additional platforms and feature dimensions. It reinforces the thesis positioning by showing that existing systems typically specialize in one function (lending or payments) rather than a complete banking suite.}

\begin{table}[H]
\centering
\caption{Detailed competitive landscape analysis (Part 3): multi-tier capital flow as differentiating feature.}
\begin{tabular}{p{4cm}p{5cm}p{4.5cm}}
\hline
\textbf{Feature Dimension} & \textbf{Competitors} & \textbf{Crypto World Bank} \\
\hline
Hierarchical lending tiers & None --- all competitors use flat pool architectures & Four-tier hierarchy: World Bank $\to$ National $\to$ Local $\to$ Client \\
\hline
Governance-controlled rates & Compound/Aave: algorithmic but no institutional hierarchy & Smart contract parameters; governance-defined bounds per tier \\
\hline
AI-assisted risk scoring & No competitor integrates on-chain ML fraud detection & Random Forest + SHAP; Isolation Forest anomaly detection \\
\hline
Retail and institutional access & Goldfinch / Maple: institutional only; Celo: retail only & Single platform serving retail clients through institutional capital hierarchy \\
\hline
Developing-market focus & Limited; most are US/EU-centric & Designed for Bangladesh, Southeast Asia, Sub-Saharan Africa \\
\hline
\end{tabular}
\TableScreenshotEnd{Competitive Landscape (Part 3): Multi-Tier Capital Flow as Differentiating Feature}
\end{table}

\vspace{2pt}
\noindent\textit{\small This final competitor table block concludes the comparative analysis and clarifies why multi-tier capital flow is a differentiating architectural feature. The results inform the go-to-market focus on transparency, governance structure, and retail accessibility.}

Upon reviewing various projects a common feature was identified: a lending mechanism does not exist that operates with tiers, is distributed and facilitates transactions across different levels and within identical levels. Such innovation distinguishes the Crypto World Bank. The key differentiating features of the Crypto World Bank are as follows:

\begin{itemize}[noitemsep]
\item Existing DeFi lending protocols provide capital access but lack institutional hierarchy, tiered governance, and structured capital flow.
\item Credit platforms such as Maple and Goldfinch engage in lending based on credit. Credit-based lending is often focused on businesses and exists at just one tier. These platforms do not cater to personal loans.
\item Ripple processes cross-border payments but does not offer lending, deposit, or hierarchical capital allocation services.
\item Celo provides assistance to individuals in emerging economies. Payment and savings functions are offered yet lending services remain absent.
\item \textbf{Central Bank Digital Currencies (CBDCs):} The development of CBDCs aims to improve payment systems~[5]. Concerns about privacy and control are raised as lending features are excluded. A lending layer does not exist on a CBDC. Users benefit from our platform which safeguards their interests while providing a level-based system similar to that utilized by CBDCs~[5].
\end{itemize}

\section{Risk Taxonomy}

\begin{table}[H]
\centering
\caption{Risk taxonomy and mitigation.}
\label{tab:risk-taxonomy}
\begin{tabular}{p{3cm}p{4cm}p{1.5cm}p{5cm}}
\hline
\textbf{Risk Category} & \textbf{Description} & \textbf{Severity} & \textbf{Mitigation} \\
\hline
Partner non-cooperation & Key partners decline to participate & Medium & Initiate with low-barrier academic and NGO pilots \\
\hline
Smart contract vulnerability & Exploit in contract logic & High & OpenZeppelin primitives; formal audit (planned); pause mechanism \\
\hline
Regulatory adversity & Jurisdictional restrictions & Medium & Testnet-only prototype; regulatory sandbox engagement \\
\hline
AI/ML model degradation & Fraud detection accuracy decay & Low & Continuous retraining; human-in-the-loop; SHAP explainability \\
\hline
\end{tabular}
\TableScreenshotEnd{Risk Taxonomy: Technical, Financial, Regulatory, and Operational Risk Categories}
\end{table}

\vspace{2pt}
\noindent\textit{\small This risk taxonomy table categorizes technical, financial, regulatory, and operational risks relevant to an on-chain banking platform. The taxonomy is used to justify design mitigations such as reserve enforcement, role-based approvals, and staged deployment through testnets and sandbox programs.}


\section{Technical Feasibility}

\begin{table}[H]
\centering
\caption{Technical feasibility assessment.}
\label{tab:tech-feasibility}
\begin{tabular}{p{3cm}p{3cm}p{7.5cm}}
\hline
\textbf{Component} & \textbf{Assessment} & \textbf{Evidence} \\
\hline
Smart contracts & Fully feasible & Three contracts implemented and tested with Hardhat (12+ passing unit tests); Solidity 0.8.20 with OpenZeppelin \\
\hline
Frontend DApp & Fully feasible & React 18 + TypeScript with all pages implemented; Wagmi and RainbowKit provide mature wallet integration \\
\hline
Blockchain deployment & Fully feasible & Polygon zkEVM Cardona and Ethereum Sepolia provide zero-cost, production-equivalent environments; ZK validity proofs provide cryptographic finality on Cardona \\
\hline
AI/ML integration & Feasible with constraints & Random Forest inference achieves sub-50\,ms latency; SHAP explanations computable in real time \\
\hline
Database backend & Feasible & PostgreSQL schema designed (19 tables, 3NF); FastAPI provides async REST framework \\
\hline
\end{tabular}
\TableScreenshotEnd{Technical Feasibility Assessment: Infrastructure Readiness and Ecosystem Maturity}
\end{table}

\vspace{2pt}
\noindent\textit{\small This technical feasibility table summarizes the readiness of core infrastructure (EVM tooling, Layer~2 scalability, and security libraries) required by the prototype. It supports the claim that the project leverages mature ecosystems rather than experimental primitives, reducing implementation risk.}

The technical feasibility of the platform rests on five pillars: the maturity of the underlying blockchain infrastructure, the availability of development tooling and security primitives, the demonstrated scalability of Layer~2 networks for financial applications, the precedent of comparable institutional blockchain deployments, and the platform's modular architecture which allows incremental delivery without requiring the full system to be operational before any component can be tested.

Ethereum---the EVM standard on which the platform is built---has operated continuously since 2015, while Polygon PoS has processed billions of transactions with low fees and rapid finality, making it suitable for high-frequency retail operations. The contract layer uses established security primitives (e.g., OpenZeppelin libraries and widely reviewed patterns), and the off-chain services use standard, production-grade web tooling for API and ML inference. Because the architecture is modular, components such as savings, FX, group lending, and insurance can be developed and validated independently and integrated through governance-approved rollouts.

\section{Economic Feasibility}

\begin{table}[H]
\centering
\caption{Economic feasibility --- zero-cost prototype.}
\label{tab:economic-feasibility}
\begin{tabular}{p{4cm}p{2.5cm}p{7cm}}
\hline
\textbf{Cost Category} & \textbf{Estimate} & \textbf{Notes} \\
\hline
Blockchain deployment & \$0 & Public testnets --- no real cryptocurrency required \\
\hline
Frontend hosting & \$0 & Vercel free tier or localhost for demo \\
\hline
Backend hosting & \$0 & Render free tier or localhost \\
\hline
AI/ML training & \$0 & Local machine (16\,GB RAM, 16\,GB VRAM) or Google Colab free tier \\
\hline
Development tools & \$0 & Hardhat, VS Code, Git --- all open-source \\
\hline
\textbf{Total prototype} & \textbf{\$0} & Entire prototype operates at zero financial cost \\
\hline
\end{tabular}
\TableScreenshotEnd{Economic Feasibility: Cost Drivers vs.~Revenue Potential on Layer-2 Networks}
\end{table}

\vspace{2pt}
\noindent\textit{\small This economic feasibility table captures cost drivers (gas, infrastructure, and defaults) and compares them against revenue potential from interest spreads and fees. The key conclusion is that low-fee Layer~2 deployment makes small-loan banking workflows economically viable, unlike high-fee base layers.}

The economic feasibility of the platform is evaluated across four dimensions: operational cost sustainability, revenue sufficiency, capital efficiency, and macroeconomic impact. On Polygon PoS, the total on-chain gas cost for a full retail loan lifecycle is measured in cents, while interest spread revenue is orders of magnitude larger, yielding a favorable cost-to-revenue profile for small loans that are infeasible on high-fee base layers. Off-chain infrastructure costs (hosting, RPC access, storage, ML inference) are comparable to typical SaaS deployments and scale with usage.

Revenue is diversified across interest spreads, origination fees, and potential FX spread revenue for multi-currency conversions. Capital efficiency is enhanced by the platform's reserve-enforced tiered allocation (each tier retains a minimum solvency reserve before deploying capital downward) and by reducing the need for idle pre-funded balances typical of correspondent banking. Note that because USDC is a 100\%-backed stablecoin, the platform does not create new money through lending---it re-allocates existing USDC through the hierarchy. The Reserve Ratio (RR) functions as a solvency constraint at each tier, not as a money-creation mechanism; see the Credit Velocity formula in the List of Formulas for the correct framing of capital throughput. At scale, remittance cost reduction and improved access to credit in underserved markets produce additional social and economic value beyond platform revenue.

\section{Revenue Projection}

The following projection models annual revenue potential at full deployment scale. Calculations use a reference ETH price of \textbf{\$2,500}\footnote{All ETH-denominated cost and revenue calculations use \$2{,}500/ETH as a conservative planning assumption consistent with the v22 revision date. The platform's stablecoin-first design (Section~\ref{sec:stablecoin-first}) means retail loan obligations are denominated in USDC, decoupling client repayment risk from ETH price volatility.} (conservative mid-point for February~2026; actual spot was approximately \$2,800 per CoinGecko on 1~February~2026) and interest rate parameters defined in Section~\ref{sec:interest-rates}. A sensitivity table showing revenue at \$2,000, \$2,500, and \$3,000 per ETH is provided after the main projection; all three scenarios use the same 8\% APR and loan-base assumptions, so revenue scales linearly with ETH price.

\begin{table}[H]
\centering
\caption{Revenue projection assumptions.}
\label{tab:revenue-assumptions}
\begin{tabular}{p{6cm}p{7.5cm}}
\hline
\textbf{Parameter} & \textbf{Value} \\
\hline
Reference ETH price & \$2{,}500 (conservative mid-point, Feb 2026; actual spot approx.\ \$2{,}800 per CoinGecko on 1~Feb~2026) \\
\hline
World Bank $\to$ National Bank & 3\% APR (wholesale inter-bank rate) \\
\hline
National Bank $\to$ Local Bank & 5\% APR (inter-bank) \\
\hline
Local Bank $\to$ Client & 8\% APR (retail lending) \\
\hline
Average loan term & 12 months \\
\hline
Default rate provision & 3\% (conservative estimate) \\
\hline
Origination fee & 0.25\% per disbursement \\
\hline
\end{tabular}
\TableScreenshotEnd{Revenue Projection by Tier (Base Case, USD Millions)}
\end{table}

\vspace{2pt}
\noindent\textit{\small This revenue projection table provides modeled annual income under a reference ETH price and assumed utilization, default rates, and tier spreads. It illustrates how hierarchical spreads across tiers accumulate into platform-level revenue while keeping retail APR within plausible ranges.}

\begin{table}[H]
\centering
\caption{System-wide annual revenue summary --- spread-based accounting.}
\label{tab:revenue-summary}
\small
\renewcommand{\arraystretch}{1.2}
\begin{tabular}{p{4.2cm}p{2.8cm}p{2.8cm}p{3.5cm}}
\hline
\textbf{Tier / Revenue Stream} & \textbf{Spread Rate} & \textbf{Annual Revenue (USD)} & \textbf{Basis} \\
\hline
Tier 1 (World Bank): lends at 3\% to NBs, cost of capital 0\% & 3\% on deployed capital & \$51.6M & \$1.72B total retail loan base × 3\% \\
\hline
Tier 2 (National Banks): borrow at 3\%, lend at 5\% & 2\% spread & \$34.4M & \$1.72B × 2\% \\
\hline
Tier 3 (Local Banks): borrow at 5\%, lend at 8\% & 3\% spread & \$51.6M & \$1.72B × 3\% \\
\hline
\textbf{Total platform revenue} & \textbf{8\% (client pays)} & \textbf{\$137.6M} & Sum of tier spreads = 3\%+2\%+3\% = 8\%; no double-counting \\
\hline
Origination fees (0.25\%) & --- & \$4.3M & 0.25\% × \$1.72B disbursed \\
\hline
FX spread revenue (est.) & 0.5--1.0\% & \$8--17M & Estimated conversion volume \\
\hline
\textbf{Total incl.\ fees} & & \textbf{\$150--159M} & Conservative base case \\
\hline
\end{tabular}
\end{table}

\vspace{2pt}
\noindent\textit{\small Revenue is computed using a \textit{spread-based} model: each tier earns the difference between its borrowing rate from the tier above and its lending rate to the tier below, applied to the same underlying loan capital as it flows down the hierarchy. Summing the three tier spreads (3\% + 2\% + 3\% = 8\%) equals the retail client's total APR, confirming internal consistency---no double-counting occurs. An earlier presentation of this table (v15) reported \$224.6M by summing the full APR at each tier rather than the spread; that figure was incorrect and has been corrected here.}

\vspace{4pt}
\noindent\textit{\small \textbf{Scale note on the \$137.6M figure:} This is the theoretical 10-year ceiling at 68,800 active loans across 25~Local Banks at near-full capacity---it is included for scale reference only. The pre-thesis evaluation target is Year~1--2 of the adoption ramp, where the platform reaches operational break-even at approximately 300~active loans and demonstrates positive unit economics on each incremental loan above that threshold. Figure~\ref{fig:revenue-by-tier} shows the theoretical maximum at full maturity; see the break-even analysis above for the realistic path from zero to that ceiling.}

\vspace{4pt}
\noindent The projection above models interest spread income only (base case: \$137.6M). Additional revenue streams from FX conversion spreads (estimated 0.5 to 1.0\% per conversion, yielding \$8--17M), loan origination fees (0.25\% per disbursement, yielding \$4.3M), and insurance fund premiums (0.5\% of loan value annually) represent material upside captured in the ``Total incl.\ fees'' row of the table (\$150--159M). These are net incremental revenue streams---they do not overlap with the interest spread income already accounted for.

\begin{table}[H]
\centering
\caption{ETH price sensitivity: annual interest spread revenue at three price points (all other assumptions held constant at base case).}
\label{tab:revenue-sensitivity}
\small
\renewcommand{\arraystretch}{1.2}
\begin{tabular}{@{}lccc@{}}
\toprule
\tableheadcolor
\tableheadfont ETH Price & \tableheadfont Total Loan Base & \tableheadfont Interest Spread Revenue & \tableheadfont Incl.\ Fees (est.) \\
\midrule
\$2{,}000 & \$1.376B & \$110.1M & \$120--128M \\
\rowcolor{RowShade}
\$2{,}500 (base case) & \$1.720B & \$137.6M & \$150--159M \\
\$3{,}000 & \$2.064B & \$165.1M & \$180--191M \\
\bottomrule
\end{tabular}
\end{table}

\vspace{4pt}
\noindent\textit{\small\textbf{Sources:} ETH price from spot market data (CoinGecko); interest rate tiers aligned with Aave~[15] and Compound~[16] DeFi benchmarks; default rate and origination fee from conservative industry estimates. Revenue methodology: spread-based, not gross-APR stacking.}

\vspace{6pt}
\noindent\textbf{Cost Model Considerations.}

\begin{itemize}[noitemsep]
\item \textbf{Gas costs:} Gas expenses on Polygon PoS are quite low. A complete retail loan lifecycle on Polygon involves approximately 27--32 individual on-chain state changes:
\label{sec:gas-cost}
\begin{itemize}[noitemsep]
  \item 1 loan request (\texttt{SSTORE}: borrower data, loan ID, status)
  \item 1 credit history lookup + risk score commit
  \item 1 approval (\texttt{SSTORE}: status update, disbursement record)
  \item 1 disbursement (ETH transfer + \texttt{SSTORE}: balance update)
  \item 12 installment payments $\times$ 2 \texttt{SSTORE} each = 24 writes
  \item 1 loan closure (\texttt{SSTORE}: final status)
  \item $\approx$4 event emissions (\texttt{LOG2}/\texttt{LOG3})
\end{itemize}
On Polygon PoS, with gas prices of approximately 30~Gwei and MATIC at $\approx$\$0.60, each \texttt{SSTORE} costs approximately \$0.00036. A 30-operation lifecycle costs approximately \$0.011 total---well under 0.01\% of the interest earned. On Ethereum mainnet at 15~Gwei base fee, the same operations would cost \$2.40--\$4.80, reinforcing the design choice of Polygon PoS~[55].
\item \textbf{Infrastructure costs:} Backend hosting (API server, database, ML service), frontend delivery (CDN) and RPC nodes (Alchemy, Infura) can incur costs ranging between \$500 and \$2,000 monthly. An increase in expenses will be experienced with a rise in transactions.
\item \textbf{Default losses:} A single default rate assumption is insufficient for an early-stage platform. Table~\ref{tab:default-scenarios} presents a three-scenario sensitivity analysis:

\begin{table}[H]
\centering
\caption{Default rate sensitivity scenarios with economic basis.}
\label{tab:default-scenarios}
\renewcommand{\arraystretch}{1.25}
\small
\begin{tabular}{@{}lccp{5.5cm}@{}}
\toprule
\tableheadcolor
\tableheadfont Scenario & \tableheadfont Default Rate & \tableheadfont Annual Loss & \tableheadfont Basis \\
\midrule
Optimistic & 3.7\% & \$7.4M & Grameen Bank 96.29\% recovery (June 2024)~[R12] after 48 years of social infrastructure \\
\rowcolor{RowShade}
Base Case & 8\% & \$16M & Typical early-stage DeFi undercollateralized lending; no established social trust \\
Stress Test & 15\% & \$30M & Early-stage crypto-native clients, no prior on-chain credit history, high ETH volatility \\
\bottomrule
\end{tabular}
\end{table}

Grameen Bank, after nearly five decades of social lending infrastructure, reported a loan recovery rate of 96.29\% as of June 2024~[R12]. The Crypto World Bank is a nascent platform without equivalent social trust or community officers. Its initial user population will likely exhibit default rates closer to 8--15\% before on-chain credit history accumulates sufficient predictive signal. The base case adopts 8\% default, with break-even at approximately 11.2\% default rate under the base loan volume assumption.

\textbf{Break-even user count:} At the base case of 8\% default with an average loan size of 10~ETH (\$25,000 at \$2,500/ETH) and 8\% APR, each loan generates \$2,000 annual interest and incurs under \$0.02 gas on Polygon. At \$500/month infrastructure costs, the platform requires at minimum 4 active loans to cover operating expenses and approximately 200 active loans to cover expected default losses at 8\%---a realistic near-term target for a pilot in 2--3 Local Banks.

\item \textbf{Break-even analysis:} Opportunities for profit on loans as small as 0.01~ETH (\$25) exist on Polygon PoS. Profitability on loans below 5~ETH (\$12,500) is not achievable on Ethereum mainnet without Layer~2. Layer~2 is essential for facilitating micro-loans.
\end{itemize}

\begin{figure}[H]
\centering
\HalfWidthDiagram{fig-revenue-by-tier.pdf}
\caption{Annual revenue projection at full 10-year deployment maturity (USD millions, spread-based accounting at \$2{,}500/ETH planning assumption). This represents the theoretical capacity ceiling at full institutional scale; see the break-even analysis above for the 3-year adoption ramp toward this target, with break-even occurring at approximately 300~active loans.}
\label{fig:revenue-by-tier}
\end{figure}

% ── v24 addition: revenue stream taxonomy ──────────────────────────────────────
\vspace{6pt}
\noindent The interest spread identified above is the platform's primary revenue stream. Placing it within the broader landscape of institutional crypto-exchange revenue models demonstrates that CWB has a diversified revenue path even without a native token. Table~\ref{tab:revenue-taxonomy} maps each CWB mechanism to its nearest exchange-sector analogue, drawing on the revenue architecture of major centralised exchanges (Binance, Coinbase) and DeFi protocols described in the ecosystem survey.

\begin{table}[H]
\centering
\caption{CWB revenue stream taxonomy: analogues from the institutional crypto exchange sector.}
\label{tab:revenue-taxonomy}
\renewcommand{\arraystretch}{1.25}
\small
\begin{tabular}{p{3.8cm}p{3.8cm}p{5.6cm}}
\hline
\textbf{Revenue Stream} & \textbf{Exchange Equivalent} & \textbf{CWB Mechanism} \\
\hline
Interest spread (primary) & Binance OTC desk & SavingsVault yield vs.\ retail lending rate; 3\%+2\%+3\% tier spread = 8\% client APR \\
\hline
Loan origination fee & Trading fee (0.1\%) & 0.5--1\% flat origination fee on disbursement; already modelled in base case above \\
\hline
Deposit mobilization & Binance Earn & SavingsVault + FixedDeposit products; yield funded by the lending tier spread \\
\hline
Local Bank registration & Exchange listing fee & One-time registration fee per Local Bank joining the network \\
\hline
Syndicated loan arrangement & Institutional desk & Lead Arranger fee on TranchedPool / SyndicatedLoan cross-tier facilities \\
\hline
Reserve transparency API & Data API subscription & Phase~3: GraphQL API over The~Graph subgraph; queried by external auditors and regulators \\
\hline
\end{tabular}
\end{table}

\vspace{2pt}
\noindent\textit{\small The six streams above are non-overlapping; none double-counts the interest spread already accounted for in Table~\ref{tab:revenue-summary}. The Local Bank registration fee and Reserve Transparency API are Phase~2/3 items and are not included in the base-case revenue model above. The interest spread and loan origination fee are the only streams active at the pre-thesis prototype stage.}

\paragraph{Hard rule: no native CWB governance token in the prototype.}
The FTX collapse of November 2022 demonstrated the systemic failure risk of self-minted tokens used as collateral: FTT, issued and held principally by the exchange itself, served simultaneously as the primary collateral asset for Alameda Research's borrowing, creating a reflexive loop that triggered an estimated \$8~billion solvency shortfall within 72~hours of a competitor announcing it would liquidate its FTT holdings. The lesson for the Crypto World Bank is categorical: no native CWB governance token is issued in the thesis prototype. On-chain reputation is provided instead by the Credit Passport SBT (Section~\ref{sec:credit-passport}), which is non-transferable and carries no speculative market price. A future governance token---analogous to Binance's BNB utility-and-burn model, which took seven years and \$16.8~billion in annual revenue to reach legitimacy---is scoped to Phase~3 of the CWB roadmap, following institutional trust bootstrapping, and is specified in Future Work with explicit anti-FTX collateral hard rules encoded at the contract level (no self-issued asset may be used as collateral within the same protocol that issues it). This constraint will be enforced as a Foundry invariant test rather than a governance policy, so it cannot be bypassed by any admin action.

\subsection{Transaction Economics: Interest Rates}
\label{sec:interest-rates}

\begin{table}[H]
\centering
\caption{Interest rate parameters.}
\label{tab:interest-rates}
\begin{tabular}{p{4cm}p{3.5cm}p{6cm}}
\hline
\textbf{Parameter} & \textbf{Value} & \textbf{Benchmark} \\
\hline
Base Annual Interest Rate & 5--12\% APR & Aligned with Aave/Compound variable rates \\
\hline
Rate Determination & Set by Local Bank approvers within World Bank-defined bounds & Configurable per-bank for local market conditions \\
\hline
Late Payment Penalty & 2\% of installment + 0.5\%/week (capped at 10\%) & Industry-standard late fee structure \\
\hline
Interest Calculation & Simple interest on outstanding principal & Transparent, client-friendly \\
\hline
Rate Transparency & All parameters stored on-chain & Publicly auditable; no hidden fees \\
\hline
\end{tabular}
\TableScreenshotEnd{Tiered Interest Rate Parameters: APR Spreads Across the Four-Tier Hierarchy}
\end{table}

\vspace{2pt}
\noindent\textit{\small This interest rate table defines the tiered APR parameters used throughout the feasibility and revenue modeling. Presenting the rate structure explicitly makes the spread logic auditable and connects the economic model directly to the proposed governance-controlled parameters.}

\begin{figure}[H]
\centering
\HalfWidthDiagram{fig-apr-spread.pdf}
\caption{Hierarchical interest-rate spread (APR) across the four-tier lending structure.}
\label{fig:apr-spread}
\end{figure}

\subsection{Global Economic Impact}

Apart from revenue generated through platforms, various economic advantages are provided by the Crypto World Bank:

\begin{itemize}[noitemsep]
\item \textbf{Capital deployment and fiscal multiplier:} At projected full-deployment lending volume of \$2~billion across nearly 1,000 institutional clients, the platform generates a downstream fiscal multiplier of approximately \$2.5 to \$3 per dollar lent~[19], implying \$5 to \$6~billion in annual economic activity. The IFC estimates that every \$1~million lent to small businesses in developing economies creates approximately 16 new jobs~[20], making capital deployment at this scale a meaningful driver of employment and local economic growth.

\item \textbf{Remittance cost reduction:} Global remittances total approximately \$860~billion annually, of which an estimated \$48 to \$56~billion is consumed by transfer fees averaging 6.49\%---more than double the United Nations Sustainable Development Goal target of 3\%~[26]. On-chain settlement on Layer~2 networks reduces transaction costs to well below 1\%, directly compressing this fee burden. Comparable blockchain payment networks demonstrate the feasibility of this target: Stellar processed approximately \$55.6~billion in payments at a fee of roughly \$0.0007 per transaction in 2025~[44], and Celo's MiniPay processes payments for under one cent across more than 60 countries~[42].

\item \textbf{Trapped capital liberation:} The correspondent banking system requires banks to maintain pre-funded nostro and vostro accounts in every currency corridor in which they operate, immobilizing large sums that cannot be deployed for lending or investment~[24]. On-chain atomic settlement eliminates the need for these idle pre-funded balances, freeing capital for productive use across the lending hierarchy.

\item \textbf{Financial inclusion:} The platform is designed for accessibility: sub-cent gas fees on Polygon PoS, a mobile-first interface, and support for micro-loan amounts below one dollar lower the barriers to credit for the estimated 1.4~billion adults globally who remain outside formal financial systems~[14]. The World Economic Forum has identified decentralized finance as a leapfrog technology capable of enabling populations to bypass the infrastructure constraints of traditional banking~[41], a pattern already observed in mobile payments across developing economies.

\item \textbf{Transaction cost reduction:} On Polygon PoS, transaction fees remain below one cent---a reduction of over \textbf{99\%} relative to the average \$42 cost of a correspondent banking transaction~[25]. This cost compression makes financially viable a large class of micro-transactions and small-loan servicing operations that are currently uneconomical through traditional payment rails.

\item \textbf{Transparency as an economic good:} On-chain publication of reserve ratios, interest rates, and transaction records eliminates the informational asymmetry that enables predatory lending practices in opaque banking environments. In conventional banking, borrowers---particularly in developing economies---frequently operate without visibility into the true cost of credit or the basis for lending decisions. The Crypto World Bank's design ensures that all participants access identical, verifiable information, reducing the scope for hidden fees and unfair pricing.
\end{itemize}

\subsection{Value Proposition and Go-to-Market}

\begin{table}[H]
\centering
\caption{Go-to-market phases.}
\label{tab:gtm-phases}
\begin{tabular}{p{2.5cm}p{6.5cm}p{4.5cm}}
\hline
\textbf{Phase} & \textbf{Activities} & \textbf{Timeline} \\
\hline
Phase 1: Validation & Competition submission (BCOLBD 2025); thesis publication; open-source release & Current \\
\hline
Phase 2: Pilot & Regulatory sandbox application; institutional partnership; testnet-to-mainnet migration & 6--12 months \\
\hline
Phase 3: Production & Multi-chain deployment; enhanced monitoring and analytics; governance token launch & 12--24 months \\
\hline
\end{tabular}
\TableScreenshotEnd{Value Proposition and Go-to-Market: User Benefits Mapped to Platform Features}
\end{table}

\vspace{2pt}
\noindent\textit{\small This value proposition table summarizes user-visible benefits (cost, transparency, speed, and inclusion) and maps them to platform features. It supports the go-to-market narrative by linking technical design choices to practical outcomes for retail users and partner institutions.}


\section{Currency Risk and the Stablecoin Imperative}

A retail client in rural Bangladesh who takes a loan denominated in ETH faces a fundamental risk that does not exist in traditional microfinance: if the price of ETH doubles between loan disbursement and final repayment, the real value of their repayment obligation doubles in taka terms. For clients near the poverty line, this is not a theoretical risk---it is potentially catastrophic. The May 2021 crypto market crash saw ETH lose approximately 55\% of its value within six weeks; ETH then halved again over the course of 2022.

This is why stablecoin integration is treated as a \textbf{critical path item} elevated to a design requirement in Section~\ref{sec:stablecoin-first}, not an optional extension. The platform supports USDC or USDT-denominated loans as the default product at retail tiers, with ETH-denominated loans reserved for institutional-tier participants who can manage currency risk. BIS Working Paper No.~905~[R13] identifies stablecoin volatility as a structural risk in emerging-market DeFi adoption and recommends fully-collateralized models (USDT/USDC) over algorithmic designs for developing-country use cases. The Terra/LUNA collapse (May~2022, $\approx$\$40~billion lost) provides the definitive negative example. ERC-20 stablecoin integration has technical implications: token allowances, \texttt{transferFrom} hooks, decimal precision (USDC uses 6 decimals vs.\ ETH's 18), and the approval-transfer-state-update ordering required by the CEI pattern must all be redesigned as a separate contract module---work that is in scope for the SavingsVault contract specified in Section~\ref{sec:savings-vault}.

\section{MiCA and GENIUS Act Compliance Mapping}
\label{sec:mica-compliance}

The platform's regulatory feasibility is not abstract: two recent statutes apply directly to any retail stablecoin-denominated lending product issued to European or US-jurisdiction users. The EU Markets in Crypto-Assets (MiCA) Regulation has been fully in effect since December~2024~[R36, R37]; the US GENIUS Act, signed in July~2025~[R38], establishes the first federal framework for payment stablecoins. Table~\ref{tab:mica-mapping} maps the relevant articles to the design choices that satisfy them and to the work that remains.

\begin{table}[H]
\centering
\caption{MiCA and GENIUS Act compliance mapping. Each row identifies a regulatory article relevant to a stablecoin-denominated crypto financial services platform and the corresponding CWB design control.}
\label{tab:mica-mapping}
\small
\begin{tabular}{p{3.2cm}p{4.6cm}p{5.8cm}}
\hline
\textbf{Statute / Article} & \textbf{Requirement} & \textbf{CWB Design Control} \\
\hline
MiCA Art.~16 (EMT Issuer Authorisation) & Authorized issuer required for any e-money token offered to the public & USDC / USDT are issued by Circle / Tether under their own authorization; CWB uses but does not issue EMTs. \\
\hline
MiCA Art.~36 (Reserve of Assets) & 1{:}1 reserve backing, segregated from issuer assets & Stablecoin reserves verified via on-chain attestations from Circle (USDC) before being accepted by SavingsVault. \\
\hline
MiCA Art.~41 (Redemption Rights) & Holders must have a right of redemption against the issuer at par & CWB does not impede redemption; SavingsVault withdraw function returns the underlying ERC-20 directly to the depositor wallet. \\
\hline
MiCA Art.~62 (Operational Resilience) & Regular audits, governance, complaint handling & Five-layer defense-in-depth (Section~\ref{sec:defense-in-depth}); Foundry invariants (Section~\ref{sec:foundry}); Certora reserve invariants (Section~\ref{sec:certora}); Tenderly runtime monitoring (Section~\ref{sec:realtime}). \\
\hline
GENIUS Act \S\,4 (Federal Payment Stablecoin Framework) & Federal registration of payment-stablecoin issuers & Same posture as MiCA Art.~16: CWB uses, does not issue. \\
\hline
GENIUS Act \S\,7 (Reserve Disclosure) & Monthly attestation and public disclosure of reserves & Circle / Tether attestations are independently verifiable on-chain; CWB dashboard surfaces issuer attestation links per stablecoin pool. \\
\hline
EU AI Act Art.~86 & Right to explanation for AI-assisted decisions & Borrower-facing SHAP explanations (Section~\ref{sec:aiml-support}); plain-language SHAP rendering as a regulatory compliance feature, not a UX bonus. \\
\hline
GDPR Art.~17 (Right to Erasure) & Personal data must be deletable on request & No PII on-chain; off-chain PostgreSQL data is encrypted and subject to data-subject deletion requests; only document hashes (SHA-256) appear on-chain. \\
\hline
\end{tabular}
\end{table}

The result is a compliance posture that is defensible to European and American institutional readers and evaluators in a single page rather than buried in narrative. Compliance gaps that remain---a formal MiCA whitepaper if CWB ever issues its own token, and FinCEN reporting integration in the US---are explicitly named in the Future Work section of the Conclusion.

\section{Bangladesh Regulatory Reality}
\label{sec:bangladesh-reg}

Bangladesh is the primary target market named throughout this thesis. The regulatory reality of operating there as of 2026 must be addressed explicitly rather than implied. Three observations frame the path forward.

\paragraph{Crypto remains effectively illegal for retail use as of 2025--2026.}
Bangladesh Bank's longstanding 2018 circular treats crypto trading as a Foreign Exchange Regulation Act offence, and no subsequent statute has formally legalized retail crypto activity~[R41]. The Crypto World Bank prototype is therefore deployed exclusively on public testnets (Polygon zkEVM Cardona, Ethereum Sepolia) using test tokens; no Bangladeshi participant transacts with mainnet value at any point in the pre-thesis or final-thesis phase.

\paragraph{The CBDC trajectory is the credible legal path.}
The central bank launched a CBDC feasibility study in 2022 and a pilot in 2024 but as of 2025 has not produced a deployment timeline. The Crypto World Bank's positioning against mBridge and Agora (Section~\ref{sec:lit-institutional}) provides a direct upgrade path: when Bangladesh Bank issues a CBDC under the same hierarchical distribution model that the Tan~(2023) IMF working paper~[5] describes, the platform's Tier~1 / Tier~2 contracts can settle in that CBDC without any change to the four-tier architecture. The current ETH / USDC denomination is a stand-in for the future CBDC asset.

\paragraph{The Regulatory FinTech Facilitation Office is the entry point.}
Bangladesh Bank operates a Regulatory FinTech Facilitation Office (RFFO) sandbox for new financial-services products. As of 2026, no crypto-lending project has been licensed under this sandbox, but the framework exists and FE Circular No.~06 (January~2025) on electronic Letters of Credit communication~[R42] together with the Payment and Settlement System Act 2024~[R43] provide the relevant regulatory anchors. The platform's go-to-market path for Bangladesh is therefore: (a) academic deployment on testnet for pre-thesis, (b) a controlled sandbox engagement with RFFO for the final-thesis evaluation, (c) initial pilot with a partner microfinance institution under the sandbox, and (d) phased CBDC integration once the Bangladesh Bank CBDC pilot moves beyond feasibility.

\paragraph{2025 Bangladesh-specific financial-inclusion data.}
The financial-inclusion argument is updated against Bangladesh Bank Special Publication SP2025-02 (June~2025)~[R44]: female-to-male account parity reached 49\%/49\% by March~2025, and female-owned deposit accounts grew from 33.4 million in 2019 to 55.3 million in 2024. BRAC's 2025 data shows 89\% of loans to women and 92\% reporting household income increase~[R45]. These statistics are more current and specific than the 2021 Global Findex \emph{1.4 billion unbanked} headline figure used elsewhere in the thesis, and they justify directing the solidarity-group module specifically at women-led client groups in line with the BRAC and Grameen tradition this work builds on.

% ── v24 addition: FATF Travel Rule scoping ────────────────────────────────────
\paragraph{FATF Travel Rule scoping for Bangladesh.}
The Financial Action Task Force Recommendation~16 (the ``Travel Rule'') requires that transfers above USD~1{,}000 include structured originator and beneficiary information transmitted alongside the transfer. In most FATF member jurisdictions this threshold is now in force for virtual asset service providers; Bangladesh, as a FATF Asia-Pacific Group member, is under the same obligation. For the Crypto World Bank, the Travel Rule applies specifically to inter-tier capital flows: Local Bank to National Bank repayments and upward surplus repatriation via the \texttt{UpwardDepositFacility} contract routinely exceed the USD~1{,}000 threshold. The off-chain Travel Rule data packet required for each such transfer is specified in Section~\ref{sec:upward-flow}:

\begin{verbatim}
{
  "originator_institution": "LocalBank-BD-042",
  "originator_tier": 3,
  "beneficiary_institution": "NationalBank-BD-001",
  "beneficiary_tier": 2,
  "amount_usdc": 5000,
  "purpose": "IBLP_REPAYMENT",
  "onchain_tx_hash": "0xABC...",
  "timestamp": "2026-06-01T10:00:00Z"
}
\end{verbatim}

\noindent This packet is stored in the PostgreSQL \texttt{audit\_logs} table and is linked to the on-chain transaction hash, making it available to regulators through the formal audit request workflow described in Section~\ref{sec:regulatory-compliance}. The \texttt{audit\_logs} table is append-only (enforced at the database-role level), ensuring the Travel Rule record cannot be altered after filing. For a Bangladesh deployment under the RFFO sandbox, this data packet specification must be submitted to Bangladesh Bank's Financial Intelligence Unit alongside the sandbox application. Implementation of the cross-tier Travel Rule notification protocol---including automated packet generation, digital signing by the originating institution, and transmission to the FIU---is Future Work contingent on Bangladesh Bank guidance on the precise format and submission channel required.

\section{Bootstrap Funding and the Tier~1 Capitalization Problem}

The Crypto World Bank faces a \textbf{bootstrap funding problem} analogous to the capitalization challenge of real multilateral development banks. The World Bank Group was initially capitalised by 44 member-state subscriptions in 1944; the IBRD's combined subscribed capital exceeded \$270~billion following its 2018 capital increase (Ocampo \& Gallagher, 2024~[R15]). Three initial capitalization mechanisms are proposed for the Tier~1 Reserve:

\begin{enumerate}[noitemsep]
\item \textbf{Founding stakeholder deposits:} Universities, NGOs, and development-focused blockchain organizations deposit ETH or USDC into the Tier~1 Reserve in exchange for governance tokens and yield rights.
\item \textbf{Protocol-owned liquidity (POL):} The protocol accumulates treasury reserves through bond mechanisms in which external parties swap ETH or USDC for discounted governance tokens over a vesting schedule. \textit{Note: The Olympus DAO algorithmic reserve model (2021--2022) is cited here as a cautionary precedent---Olympus OHM lost over 99\% of its token value within one year due to unsustainable algorithmic backing. The CWB POL mechanism is strictly collateralized: only fully-backed USDC bonds are accepted, and governance tokens carry no intrinsic algorithmic price-support obligation. This removes the reflexive collapse risk that destroyed Olympus DAO.}
\item \textbf{Philanthropic/impact grant funding:} Development finance institutions (IFC, ADB) have expressed interest in blockchain-based development finance tools, as evidenced by the World Bank FundsChain initiative~[39]. Grant funding from these institutions could seed the Tier~1 Reserve in exchange for research access and co-branding.
\end{enumerate}

In the short term, a multi-signature wallet (3-of-5 signers representing different stakeholder groups) governs Tier~1 allocations. In the long term, an on-chain governance module (similar to Compound Governor Bravo) enables token-weighted voting with time-locks to prevent governance attacks.

\section{Prototype Scope and Limitations}

A straightforward version of the complete system is represented by the current prototype:

\begin{enumerate}[noitemsep]
\item \textbf{Hierarchical lending scope:} Utilization of the Tier~1 World Bank Reserve contract is comprehensive. Deposits are managed, loan requests and approvals are processed alongside the display of reserve statistics. Not all functions of the Tier~2 (National Bank) and Tier~3 (Local Bank) contracts have been implemented. Role management and registration are not the only aspects focused on. Currently the process for fund transfers from the World Bank to National Bank and subsequently to Local Bank is incomplete. Four tiers will be included in the final design.

\item \textbf{Interest accrual and repayment:} Rules for interest rates (3\%/5\%/8\%) are part of the system design. Fees for initiating loans and penalties for delays also exist. Interest calculations, payment schedules or penalties aren't managed automatically by current smart contracts. Future updates will incorporate these functions.

\item \textbf{InterBankLendingPool, UpwardDepositFacility, and broader multi-entity / cross-tier operations:} Same-tier interbank lending, upward surplus repatriation, syndicated and tranched lending, cross-tier treasury FX, and multilateral settlement netting are fully specified in Section~\ref{sec:multi-entity} (with corresponding ERD entities listed in Section~\ref{sec:multi-entity-consistency}). The smart contracts (\texttt{InterBankLendingPool}, \texttt{UpwardDepositFacility}, \texttt{TreasurySwap}, \texttt{SyndicatedLoan}, \texttt{TranchedPool}, \texttt{NettingEngine}) are scheduled for Sprints~2 and~3; the current pre-thesis~1 prototype does not yet exercise these multi-directional flows.

\item \textbf{AI/ML integration:} AI tools such as fraud detection (Random Forest) and anomaly detection (Isolation Forest) along with decision explanations (SHAP) are set to be utilized by the system. A service has been established with FastAPI for machine learning. This service does not currently link to the loan approval process in the existing version.

\item \textbf{Backend architecture:} Express.js paired with MongoDB is utilized for rapid API development. The main database is intended to be PostgreSQL. FastAPI is utilized by the machine learning service.

\item \textbf{Banking product suite:} The extended banking product suite described in Section~\ref{sec:banking-products}, including savings accounts, fixed-term deposits, transactional accounts, group lending, FX conversion, and the insurance fund, is specified at the architectural and design level in this report. Smart contract implementation and testing of these modules is planned for the final thesis phase. The current prototype establishes the hierarchical lending foundation upon which these modules will be integrated.
\end{enumerate}

The limitations above are consistent with a pre-thesis prototype; the complete system design is fully specified in this report and will be implemented and validated in the final thesis phase.

\section{Accessibility Assessment: A Borrower in Rural Sylhet}

To ground the platform's financial inclusion claims in a concrete context, this section evaluates access across six dimensions for a hypothetical retail client: a small-scale trader in rural Sylhet, Bangladesh, seeking a 50,000~BDT ($\approx$\$430) working capital loan. Bangladesh is referenced here as a future-work target deployment context rather than a current implementation claim; the prototype operates on public testnets only (Section~\ref{sec:bangladesh-reg}).

\begin{description}[leftmargin=0pt]

\item[\textbf{1. Device access.}] Smartphone penetration in rural Bangladesh reached approximately 51\% of adults as of 2024. The platform's mobile-first React interface and WalletConnect integration allow access via any Android or iOS device with a browser, without requiring a desktop or dedicated hardware wallet.

\item[\textbf{2. Connectivity.}] 4G coverage in Bangladesh now reaches over 90\% of the population, including rural Sylhet, through Grameenphone and Banglalink networks. Polygon PoS transactions are lightweight (under 10~KB) and complete within seconds, well within 4G latency budgets.

\item[\textbf{3. Transaction cost.}] At Polygon PoS gas prices, the complete loan lifecycle (request, approval, 6~monthly installments) costs under \$0.02. This contrasts with informal money-lender fees of 5--10\% per transaction and formal bank charges of 2--3\% origination plus branch travel costs.

\item[\textbf{4. Language and literacy.}] The current prototype is English-only. A production deployment serving rural Sylhet must include a Bengali-language interface. This is a known gap and a planned extension; the React component library supports right-to-left and Unicode rendering.

\item[\textbf{5. Identity and onboarding.}] Bangladesh's National ID (NID) system covers approximately 110~million registered voters. The planned ZKP KYC extension (Section~\ref{sec:identity}) allows NID-based credential verification without exposing personal data on-chain. Mobile Financial Services (MFS) accounts---bKash, Nagad, Rocket---provide an existing digital financial identity that 99\% more adults accessed in 2021 compared to 2004~[R11].

\item[\textbf{6. Group lending fit.}] The solidarity group lending model (GroupLendingPool, planned) directly mirrors the BRAC and Grameen Bank methodology already familiar to rural Sylhet clients. Programmable mutual liability replaces community social pressure with smart contract enforcement, potentially improving both collection reliability and client protection.

\end{description}

In aggregate, the Crypto World Bank is accessible in principle to a rural Sylhet client today (device, connectivity, cost) with two critical gaps requiring resolution for a production deployment: Bengali-language interface and stablecoin denomination to eliminate currency risk.


%% %%%%%%%%%%%%%%%%%%%%%%%%%%%%%%%%%%%%%%%%%%%%%%%%%%%%%%%%%%%%
%% CHAPTER 6: CONCLUSION
%% %%%%%%%%%%%%%%%%%%%%%%%%%%%%%%%%%%%%%%%%%%%%%%%%%%%%%%%%%%%%
\chapter{Conclusion}

The Crypto World Bank addresses a \textbf{multi-party coordination and trust} problem inherent in hierarchical development finance by exploiting the distinctive properties of blockchain technology: cryptographic integrity, immutability, and programmable auditability. It advances beyond existing DeFi lending protocols---which collectively manage over \$55~billion in TVL~[13] but universally employ flat, single-tier pool architectures---through a four-tier institutional hierarchy with cross-tier, same-tier, and upward lending flows, combined with a governance structure that addresses network membership, business operations, and technology infrastructure.

This thesis makes four original contributions, as formalized in Section~\ref{sec:research-contribution}: (1) a four-tier hierarchical DeFi architecture that mirrors multilateral development finance capital flows, with no comparable prior art in existing protocols; (2) an on-chain solidarity group lending specification that encodes mutual liability and over-indebtedness control (Section~\ref{sec:over-indebt}) as programmable contract logic; (3) an oracle-mediated AI/ML integration pattern providing a blueprint for auditable AI-assisted credit governance, extended by an autonomous AI agent system (MCP tool server, 16 banking tools, EIP-7702 session key authorisation, and human-gate confirmation protocol) and by Chainlink Functions trustless oracle commitment, and further extended by GNN relational features (Section~\ref{sec:gnn-extension}) and federated learning across banking tiers (Section~\ref{sec:fl-fraud}); and (4) a compliance-aware ZKP identity pathway combining zkKYC, zkAML, and W3C DID/VC under a self-sovereign-identity frame (Section~\ref{sec:identity}). Additionally, this thesis specifies a formal actor-permission matrix and SAR compliance workflow extending the compliance architecture beyond identity to institutional AML governance. These contributions are framed against the institutional-DeFi landscape (Sygnum, mBridge, Agora) discussed in Section~\ref{sec:lit-institutional}, and against the MiCA / GENIUS Act regulatory frontier mapped in Section~\ref{sec:mica-compliance}.

The analysis of the competitive situation of 20+ existing projects in DeFi lending, institutional credit, banking infrastructure and financial inclusion affirm that no current system is a combination of multi-level lending (hierarchy), interbank lending mechanisms and graded access by borrowers in one decentralized architecture. The growing blockchain implementation in the institutions, as demonstrated by the World Bank FundsChain initiative~[39], the multi-billion-dollar daily volumes of JPMorgan Kinexys~[40] and R3 Corda's \$17~billion in tokenized assets~[37]---validates both the technical feasibility and institutional demand of blockchain-based financial infrastructure. The Crypto World Bank extends this trajectory from settlement and fund tracking into a fully open, hierarchically governed lending system.

The platform occupies the intersection of institutional finance, decentralized lending, and financial inclusion, a combination that remains unaddressed in both the academic literature and the commercial blockchain landscape. With a working prototype, a defined market and partnership plan, and a planned go-to-market plan, the Crypto World Bank represents a structurally distinct and architecturally grounded contribution to the emerging field of blockchain-based development finance.

Viewed as a complete banking function checklist, the current prototype implements hierarchical lending workflow foundations (tiered roles, governance framing, architectural capital flow, and a planned AI/ML monitoring layer) while leaving several bank-grade modules at the design level. Unimplemented modules specified in this report include:
\begin{itemize}[noitemsep]
\item Deposit products (SavingsVault and FixedDeposit), specified in Section~\ref{sec:savings-vault}
\item Transactional accounts (CurrentAccount)
\item Group lending (GroupLendingPool) with mutual-liability and over-indebtedness logic (Section~\ref{sec:over-indebt})
\item Liquidation Engine (Section~\ref{sec:liquidation})
\item Cross-chain bridge module (Section~\ref{sec:bridge})
\item On-chain credit passport SBT (Section~\ref{sec:credit-passport})
\item Oracle-priced FX conversion (FXModule)
\item Multi-entity and cross-tier operations: InterBankLendingPool, UpwardDepositFacility, SyndicatedLoan, TranchedPool, TreasurySwap, and NettingEngine, specified in Section~\ref{sec:multi-entity}
\item Insurance and depositor protection (InsuranceFund)
\item ERC-4337 account abstraction wallet stack (Section~\ref{sec:erc4337})
\item Federated learning aggregator and GNN ablation pipeline (Sections~\ref{sec:gnn-extension}--\ref{sec:fl-fraud})
\item Certora reserve invariants and Foundry fuzz / invariant suite (Sections~\ref{sec:certora}--\ref{sec:foundry})
\item Tenderly + The Graph runtime monitoring stack (Section~\ref{sec:realtime})
\end{itemize}
Future iterations therefore focus on completing end-to-end cross-tier fund transfers and repayment automation, integrating the analytics layer into real approval flows under the Chainlink Functions oracle (replacing the interim commit-reveal relay as specified in Sprint~3), hardening the platform against regulatory scrutiny via MiCA / GENIUS Act compliance, and expanding the platform into a functionally complete, stablecoin-first, compliance-aware banking system.

The future work will be directed towards the following directions, each grounded in a planned change introduced or formalized in v15:
\label{sec:future-work}

\begin{enumerate}[noitemsep]
\item \textbf{Complete hierarchical lending implementation.} Complete the end-to-end wiring of cross-tier fund transfers between the World Bank Reserve, National Bank, and Local Bank contracts---automated capital distribution, interest accrual, installment-schedule generation, and cascading-repayment enforcement---against the contract sketches in Appendix~B.
\item \textbf{Stablecoin-first deployment.} Ship the retail Local Bank tier with USDC as the default loan currency (Section~\ref{sec:stablecoin-first}); keep ETH only for institutional tiers. Document the MiCA / GENIUS Act compliance mapping (Section~\ref{sec:mica-compliance}) against the deployed configuration.
\item \textbf{Liquidation Engine, SavingsVault, FixedDeposit, GroupLendingPool.} Implement and audit the LiquidationEngine (Section~\ref{sec:liquidation}), SavingsVault and FixedDeposit (Section~\ref{sec:savings-vault}), and the GroupLendingPool with the over-indebtedness controls of Section~\ref{sec:over-indebt}, completing the Sprint~2 deposit-and-credit suite.
\item \textbf{Multi-entity and cross-tier operations.} Implement and audit the six contracts of Section~\ref{sec:multi-entity}: \texttt{InterBankLendingPool} (same-tier interbank lending), \texttt{UpwardDepositFacility} (LB~$\to$~NB~$\to$~WB surplus repatriation), \texttt{SyndicatedLoan} (multi-entity co-funded loans with on-chain consent vote and pro-rata distribution), \texttt{TranchedPool} (senior / junior risk-segmented pools), \texttt{TreasurySwap} (cross-tier ETH~$\leftrightarrow$~stablecoin treasury swap), and \texttt{NettingEngine} (multilateral settlement netting). These six contracts complete the cross-tier, same-tier, and upward lending claim of Section~\ref{sec:research-contribution}.
\item \textbf{AI/ML pipeline integration.} Wire Chainlink Functions oracle (Sprint~3) as the primary ML score commitment mechanism, replacing the interim commit-reveal relay. Wire Random Forest + Isolation Forest + SHAP into the loan decision pipeline; deliver client-facing SHAP plain-language explanations via the Authority Brief UI (Section~\ref{sec:aiml-support}). The human-gate confirmation pattern in the MCP agent pipeline (Sprint~2) constitutes the production safety boundary.
\item \textbf{GNN extension and federated learning.} Add the GraphSAGE ablation (Section~\ref{sec:gnn-extension}) for relational fraud detection, and the FedAvg federation across Local / National Banks for privacy-preserving cross-institutional threat intelligence (Section~\ref{sec:fl-fraud}).
\item \textbf{Formal verification with Certora.} Produce a Certora-verified proof of the two reserve invariants of Section~\ref{sec:certora}, with the CVL specifications in Appendix~B.
\item \textbf{Foundry invariant + fuzz suite.} Add the three Foundry invariants (solvency, role segregation, capital-flow direction) running for 10{,}000 fuzz runs per CI build (Section~\ref{sec:foundry}).
\item \textbf{On-chain economic feasibility simulation (Sprint~3 completion).} The scripted Foundry simulation of Section~\ref{sec:abm-sim} is the designated empirical tool for the pre-thesis phase; Sprint~3 completes its full deployment with 300 clients, 6 banks across a six-institution hierarchy on Polygon~zkEVM~Cardona, produces the gas-cost table and reserve-dynamics output, and publishes the deterministic seed and transaction manifest to Appendix~C. A separate agent-based model (Mesa or equivalent) calibrated against real on-chain data remains an optional extension for the final thesis, subject to data availability.
\item \textbf{Account abstraction, tiered KYC, zkAML circuit.} Deploy ERC-4337 account abstraction for retail Tier~4 onboarding (Section~\ref{sec:erc4337}), the tiered KYC ladder (Section~\ref{sec:tiered-kyc}), and the second ZKP circuit for AML (Section~\ref{sec:identity}) alongside the existing KYC circuit, both anchored on a W3C DID / VC layer.
\item \textbf{On-chain credit passport SBT.} Implement the soulbound credit passport (Section~\ref{sec:credit-passport}) and mirror its read interface across chains via the cross-chain bridge architecture of Section~\ref{sec:bridge}.
\item \textbf{Runtime monitoring and dashboard.} Deploy The Graph subgraph, Tenderly alerts, and the WebSocket event listener (Section~\ref{sec:realtime}) to convert the static dashboard into a live banking dashboard. The Reserve Transparency Dashboard (Sprint~3) queries The Graph subgraph in real time and serves as the live demo during the 300-client Foundry simulation.
\item \textbf{Operational security (L5).} Transfer all admin roles to Safe multisigs before any public deployment (Section~\ref{sec:defense-in-depth}); document the formal key-rotation policy; use a Ledger~Gen5 hardware-wallet signer for the World Bank Safe.
\item \textbf{Mainnet and regulatory sandbox.} Move from testnet to a controlled mainnet pilot under the Singapore MAS Project~Guardian sandbox or UAE DIFC Innovation Testing Licence as Phase~1 (1--2~years). Bangladesh~CBDC integration under the hierarchical distribution model is a Phase~3 (5--10~year) target contingent on Bangladesh~Bank's CBDC issuance timeline, as described in Section~\ref{sec:bangladesh-reg}. Initial pilot with a partner microfinance institution proceeds under whichever Phase~1 sandbox is obtained first.
\item \textbf{Cross-chain interoperability.} Bridge Polygon $\leftrightarrow$ Ethereum (and Arbitrum / Optimism / Base as needed) via Chainlink CCIP (Section~\ref{sec:bridge}) while preserving the hierarchical borrowing-limit invariant.
\item \textbf{LLM assistant fine-tuning and red-teaming.} Fine-tune the 8B-parameter Qwen3 model with QLoRA on platform-specific Q\&A and run the red-teaming evaluation protocol of Section~\ref{sec:llm-eval} for regulatory-hallucination resistance.

\item \textbf{FATF Travel Rule implementation.} The Travel Rule data packet specification (Section~\ref{sec:upward-flow}) provides the structural contribution; implementation of the off-chain notification protocol across the \texttt{InterBankLendingPool} and \texttt{UpwardDepositFacility} contracts is deferred to Phase~2, contingent on jurisdiction-specific regulatory guidance from Bangladesh Bank.

\item \textbf{SAR workflow hardening.} The Isolation Forest SAR pipeline (Sprint~3) covers anomaly detection and account freeze authority. Phase~2 will add integration with Bangladesh Bank's Financial Intelligence Unit reporting portal and automated SAR PDF generation in the prescribed compliance format.

\item \textbf{Groth16 zkKYC circuit.} A Circom~2.0 circuit proving National Identity Document (NID) possession without revealing the NID; integrated with Polygon ID for Bangladesh NID issuers. Circuit design and input/output specification constitute the academic contribution; production deployment requires a licensed identity provider partnership.

\item \textbf{World ID anti-Sybil for GroupLendingPool.} Integration of Worldcoin's World ID proof-of-personhood at the group formation step to prevent Sybil attacks on group loans (Section~\ref{sec:group-lending}). Specified as a design contribution; implementation deferred pending World ID developer partnership agreement.

\item \textbf{Federated learning module.} A FedAvg aggregator across Local and National Banks for privacy-preserving cross-institutional threat intelligence (Section~\ref{sec:fl-fraud}), extended with differential-privacy noise injection to bound information leakage across institution boundaries.

\item \textbf{Polygon CDK sovereign chain (Phase~3).} A dedicated CWB application chain with a governance-controlled validator set, custom gas token, and embedded Chainlink oracle nodes. Replaces the Polygon~zkEVM~Cardona testnet deployment at institutional scale, providing full control over sequencer policy and fee markets.
\end{enumerate}


%% %%%%%%%%%%%%%%%%%%%%%%%%%%%%%%%%%%%%%%%%%%%%%%%%%%%%%%%%%%%%
%% APPENDICES
%% %%%%%%%%%%%%%%%%%%%%%%%%%%%%%%%%%%%%%%%%%%%%%%%%%%%%%%%%%%%%
\appendix

\chapter{Technology Stack}

\section{In-product assistant: local large language model (LLM) integration (prototype)}

\paragraph{Purpose.}
The in-product assistant is an autonomous AI banking agent, not a static product guide. It combines a locally hosted, privacy-preserving large language model (Qwen3-8B, Apache~2.0~licence) with a Model Context Protocol (MCP) tool server exposing 16 restricted banking tools. The agent can both answer questions---via retrieval-augmented generation (RAG) from policy documents---and execute on-chain banking operations via write tools, subject to an explicit human confirmation gate before any state-modifying action is taken. The privacy-preserving local hosting means that client conversation data never leaves the institution's infrastructure, satisfying the data-sovereignty requirement for a Bangladesh deployment.

\paragraph{Model and runtime (local development).}
The inference endpoint follows the \textbf{OpenAI Chat Completions} contract (\texttt{/v1/chat/completions}) with \texttt{stream: true}, proxied from the project backend to a local service such as \texttt{http://127.0.0.1:1234}. During early development, the assistant used a locally hosted \textbf{Gemma-4-E4B} checkpoint (Google mixture-of-experts model, $\sim$26B total / 4B active parameters) via LM Studio in GGUF form (e.g., \texttt{Q4\_K\_M} quantization) for a laptop-friendly footprint. For the final thesis evaluation, the assistant is replaced by \textbf{Qwen3-8B} (Alibaba Cloud Qwen team, released April~28~2025, Apache~2.0~licence)~[R46], a dense 8.19B-parameter causal language model with a 32K-token native context window (extendable to 131K tokens via YaRN scaling) and support for 100+ languages and dialects. A distinguishing architectural feature of Qwen3-8B is its \emph{switchable reasoning mode}: the model operates in \textbf{thinking mode} (chain-of-thought wrapped in \texttt{<think>} tokens, suited to multi-step compliance or policy questions) and \textbf{non-thinking mode} (direct, low-latency responses for general navigation queries), selectable per-turn with \texttt{/think} or \texttt{/no\_think} in the system prompt. For Q\&A and retrieval-augmented generation (RAG) queries, non-thinking mode is the default; thinking mode is activated only for the red-team regulatory-compliance prompts in the evaluation protocol of Section~\ref{sec:llm-eval}. For agent tool-calling, non-thinking mode is the default for read tools and confirmation summaries; thinking mode is activated for complex multi-step operations such as EMI calculation or group signature coordination. The per-user context namespace (see Section~\ref{sec:aiml-support}) is prepended to every system prompt as a structured JSON block, so the model has full account context before reading the user's message. The API model identifier is environment-driven; the integration layer remains identical for both models.

\paragraph{End-to-end behavior.}
The user interface assembles a short transcript (user/assistant messages) and posts it to the backend streaming route. The backend prepends a \textbf{system message} with platform + screen ``feature'' context, optionally enriched by user role when a valid session token is present. The backend then opens a streaming request to the local model server, converts upstream token events into \textbf{server-sent events (SSE)} for the browser, and the UI incrementally appends text. The message renderer uses Markdown (including GitHub-flavored extensions), math (KaTeX) for \texttt{\$\$...\$\$} blocks, and line-break rules appropriate for chat transcripts.

For write-tool requests, the agent assembles the full parameter set from the user's request and the injected on-chain state, then presents a confirmation summary. Only after the user sends explicit affirmative consent does the agent call the corresponding MCP write tool via the Express.js banking API layer. The EIP-7702 session key (if active) signs the resulting on-chain transaction within its approved scope. Every executed write tool is logged to \texttt{agent\_action\_log} with the confirmation message ID as the audit reference.

\paragraph{Live on-chain context injection.}
The Express.js backend injects the client's full on-chain state as a structured JSON context block into every system prompt before forwarding the user's message to the model. The per-user context namespace is:
\begin{verbatim}
{
  "client_id": "0xABC...",
  "language": "Bengali",
  "credit_tier": "Silver",
  "credit_score": 512,
  "active_loans": [
    {"loan_id": "L-4821", "amount": 100,
     "next_due": "2026-06-15", "installment": 17.5}
  ],
  "savings_balance": 250.0,
  "pool_utilisation": 0.67,
  "kyc_tier": 1,
  "conversation_history": [ "...last 10 turns..." ]
}
\end{verbatim}
This context block is injected only when the user is authenticated and the \texttt{eth\_call} round-trip completes within a 200~ms timeout (falling back to a static-context response otherwise). The model uses the injected context to ground its answers in the user's actual account data rather than generic policy text, materially reducing hallucination risk for account-specific queries. The backend change is approximately 20~lines of additional TypeScript in the chat route handler and does not affect the model's safety posture, token budget, or latency target.


\paragraph{Security and limitations (prototype stance).}
The agent's safety posture is enforced at three levels. (1)~\textbf{Tool-schema boundary:} the agent interacts with CWB exclusively through the 16 MCP tools --- it cannot execute arbitrary code, make unapproved API calls, or read data outside the tool schema. (2)~\textbf{Human confirmation gate:} every write tool requires explicit affirmative consent in the conversation history; the confirmation turn ID is logged as the audit reference. (3)~\textbf{Session key scope:} the EIP-7702 session key is scoped to specific tools, value-capped, time-bound to 24~hours, and revocable at any time. A session that attempts to call an out-of-scope tool is rejected at the key level, not just at the application level. Hallucination risk is bounded by: (a) the injected on-chain state grounding account-specific answers, and (b) the refusal layer for regulatory and legal queries (100\% target on the red-team set, Section~\ref{sec:llm-eval}).
\paragraph{Architecture diagram.}
Figure~\ref{fig:local-llm-mermaid} shows the compact end-to-end request path; Figure~\ref{fig:local-llm-tikz} expands the same flow with component boundaries and authentication context.

\begin{figure}[H]
\centering
\OnePageDiagram{fig-local-llm-compact.pdf}
\caption{Autonomous AI banking agent request path: browser UI $\to$ Vite dev proxy $\to$ CWB Express API (SSE) $\to$ MCP tool server (16 banking tools) $\to$ Qwen3-8B inference $\to$ human confirmation gate $\to$ write-tool execution $\to$ on-chain transaction.}
\label{fig:local-llm-mermaid}
\end{figure}

\begin{figure}[H]
\centering
\OnePageDiagram{fig-local-llm.pdf}
\caption{Expanded autonomous AI agent data flow with component boundaries: the browser UI streams from the CWB Express API; read-tool requests return immediately; write-tool requests pass through the human confirmation gate and are signed with the EIP-7702 session key before on-chain execution.}
\label{fig:local-llm-tikz}
\end{figure}

\paragraph{Implementation analysis (brief).}
The integration is successful for the project prototype because it reuses a stable transport shape (chat-completions with streaming) and isolates the model vendor behind a single API route, allowing the user interface to remain a thin client. The main practical engineering challenges in local development were: (1) \textbf{process/port hygiene} to keep the API bound to a predictable port, (2) \textbf{CORS and same-origin} considerations when the UI and API run on different localhost ports, mitigated with a Vite dev proxy to \texttt{/api} and with permissive dev CORS policy on the API, and (3) \textbf{RPC provider selection} for wallet reads in the browser, where some public endpoints may fail browser CORS; pinning explicit public RPC endpoints avoids noisy failures unrelated to the chat feature.

\paragraph{Relationship to the earlier rule-based assistant (legacy).}
The repository also retains early keyword-routed \texttt{/api/chatbot} handlers used for initial prototyping; the LLM path is additive. This separation prevents regressions in deterministic demo endpoints while the conversational assistant evolves.

\begin{table}[H]
\centering
\caption{Technology stack summary.}
\label{tab:tech-stack}
\begin{tabular}{p{3cm}p{4cm}p{6.5cm}}
\hline
\textbf{Layer} & \textbf{Technology} & \textbf{Version / Notes} \\
\hline
Smart Contract & Solidity, OpenZeppelin & 0.8.20; Ownable, ReentrancyGuard \\
Frontend & React, TypeScript & Material UI; Design 3 \\
Wallet Integration & Wagmi, RainbowKit, Viem & EIP-1193 compliant \\
Build and Test & Hardhat & Automated test suite; deployment scripts \\
Backend API & Express.js, Node.js & REST API; middleware architecture \\
Database & PostgreSQL (designed) & 19 tables, 3NF; relational integrity \\
Target Network & Polygon zkEVM Cardona testnet & ZK validity proof security model; replaces Polygon Amoy PoS \\
AI Agent Engine & Qwen3-8B (llama.cpp, Q4\_K\_M) & MCP tool server; per-user context namespace; human confirmation gate \\
MCP Tool Server & 16 banking tools (Node.js) & Exposes read + write banking operations to the agent with typed parameter schemas \\
Session Keys & EIP-7702 & Scoped, time-bound, value-capped agent wallet authorisation \\
Oracle Network & Chainlink Functions (DON) & Trustless ML score commitment; replaces commit-reveal relay \\
Price Feeds & Chainlink AggregatorV3Interface & BDT/USD and ETH/USD; 8-decimal precision \\
Automation & Chainlink Automation & Trustless installment due-date checking; replaces centralised cron job \\
Proof of Reserve & Chainlink PoR & Cryptographically verifiable WorldBankReserve balance \\
Event Indexing & The Graph (subgraph) & GraphQL query layer over LoanApplication and ReserveRatioSnapshot events \\
\hline
\end{tabular}
\TableScreenshotEnd{Local LLM Assistant Integration: Components, Configuration, and Prototype Stance}
\end{table}

\vspace{2pt}
\noindent\textit{\small This appendix table summarizes assistant integration components and configuration aspects used in the prototype (UI, API proxying, and local model server). The purpose is to make the implementation reproducible and to clarify which parts are prototype conveniences versus production requirements.}

\chapter{Smart Contract Capabilities}

The complete banking architecture is designed around fifteen modular contracts. The current prototype implements the core three-contract lending foundation; the remaining twelve modules are planned for implementation across Sprints~2 and~3 of the final thesis phase.

\textbf{Implemented contracts (Sprint~1):}
\begin{itemize}[noitemsep]
\item \textbf{World Bank Reserve Contract:} Reserve custody, deposit handling, national bank registration, capital allocation to national banks, system pause and unpause, emergency withdrawal, and system statistics.
\item \textbf{National Bank Contract:} Local bank registration, borrowing from the World Bank reserve, capital allocation to local banks, and network status queries.
\item \textbf{Local Bank Contract:} Client registration, loan application processing, approval and rejection workflows, installment generation and payment processing, approver management, user account management, and a \texttt{freezeAccount(clientId)} function gated by \texttt{onlyApprover} modifier, used by the SAR workflow (Section~\ref{sec:regulatory-compliance}) to suspend loan and payment operations for a flagged client pending compliance review.
\end{itemize}

\textbf{Planned contracts -- Banking product suite (Sprint~2):}
\begin{itemize}[noitemsep]
\item \textbf{SavingsVault Contract:} Standard savings account management, variable yield accrual, withdrawal processing, and deposit-to-lending pool integration.
\item \textbf{FixedDeposit Contract:} Term deposit creation, lock period enforcement, agreed APY accrual, early withdrawal penalty calculation, and maturity release.
\item \textbf{GroupLendingPool Contract:} Group formation, multi-signature consent recording, shared collateral management, per-member disbursement, mutual liability enforcement, and group credit history recording.
\item \textbf{FXModule Contract:} Oracle-priced currency conversion, spread calculation, dual-currency lending denomination support, and conversion audit trail.
\item \textbf{InsuranceFund Contract:} Captures 5\% of all interest collected across the lending hierarchy, processes default coverage disbursement claims, and publishes its reserve balance to Chainlink Proof of Reserve for external verifiability.
\item \textbf{CurrentAccount Contract:} Transactional account management, atomic peer transfers, recurring payment scheduling, and payroll deposit handling.
\end{itemize}

\textbf{Planned contracts -- Multi-entity and cross-tier operations (Sprints~2--3):}
\begin{itemize}[noitemsep]
\item \textbf{InterBankLendingPool Contract} (Section~\ref{sec:interbank-pool}): one instance per tier (\texttt{IBLP\_NB}, \texttt{IBLP\_LB}). Short-tenor (1 / 7 / 30 day) same-tier borrowing with a utilization-kinked rate model bounded by the tier-above downward rate; default cascade through reserve buffer~$\to$~InsuranceFund~$\to$~parent-tier backstop.
\item \textbf{UpwardDepositFacility Contract} (Section~\ref{sec:upward-flow}): role-restricted extension of SavingsVault accounting that lets a lower-tier bank deposit excess reserves into its parent tier, earning a strictly lower yield than the downward borrowing rate; withdrawal gate enforces depositor reserve-ratio safety.
\item \textbf{SyndicatedLoan Contract} (Section~\ref{sec:syndicated-lending}): syndicate proposal, subscription window, supermajority on-chain consent vote (75\%-by-capital-share), atomic disbursement, pro-rata interest and recovery distribution; supports cross-tier syndication (mixed NB / LB co-funders).
\item \textbf{TranchedPool Contract} (Section~\ref{sec:tranched-pools}): senior / junior tranching with a subordination-ratio policy, waterfall interest and principal payments, first-loss absorption by the junior tranche; suitable for risk-segmented institutional and impact-aligned capital.
\item \textbf{TreasurySwap Contract} (Section~\ref{sec:treasury-fx}): oracle-priced atomic asset swap restricted to wallets holding a banking role; cross-tier asset exchange (e.g., NB ETH treasury~$\to$~USDC) with a tighter spread than retail FX; post-swap reserve-ratio invariant enforced before state change.
\item \textbf{NettingEngine Contract} (Section~\ref{sec:netting-settlement}): multilateral netting of accumulated interbank obligations; off-chain Settlement Coordinator computes the netted matrix, on-chain \texttt{settleBatch} executes the entire batch in one transaction, with a public challenge window for dispute resolution.
\end{itemize}


%% %%%%%%%%%%%%%%%%%%%%%%%%%%%%%%%%%%%%%%%%%%%%%%%%%%%%%%%%%%%%
%% APPENDIX C: DEPLOYED TESTNET CONTRACT ADDRESSES
%% %%%%%%%%%%%%%%%%%%%%%%%%%%%%%%%%%%%%%%%%%%%%%%%%%%%%%%%%%%%%
\chapter*{Appendix C: Deployed Testnet Contract Addresses}
\phantomsection
\addcontentsline{toc}{chapter}{Appendix C: Deployed Testnet Contract Addresses}

The following table records the contract addresses deployed to the Polygon zkEVM Cardona testnet during Sprint~1 and Sprint~2. All addresses can be independently verified on the Cardona block explorer at \url{https://explorer-ui.cardona.zkevm-rpc.com}. Deployment was performed using Hardhat~v2.22 with a project-specific deployment script; transaction hashes are available in the project repository under \texttt{/deployments/cardona/}.

\begin{table}[H]
\centering
\caption{Deployed smart contract addresses, Polygon zkEVM Cardona testnet (as of pre-thesis submission).}
\renewcommand{\arraystretch}{1.35}
\small
\begin{tabular}{@{}llp{5.5cm}@{}}
\toprule
\tableheadcolor
\tableheadfont Contract & \tableheadfont Network & \tableheadfont Address \\
\midrule
WorldBankReserve & Polygon zkEVM Cardona & \texttt{[See project repository: /deployments/cardona/WorldBankReserve.json]} \\
\rowcolor{RowShade}
NationalBank (Scaffold) & Polygon zkEVM Cardona & \texttt{[See project repository: /deployments/cardona/NationalBank.json]} \\
LocalBank (Scaffold) & Polygon zkEVM Cardona & \texttt{[See project repository: /deployments/cardona/LocalBank.json]} \\
\rowcolor{RowShade}
LendingController & Polygon zkEVM Cardona & \texttt{[See project repository: /deployments/cardona/LendingController.json]} \\
\bottomrule
\end{tabular}
\end{table}

\noindent\textit{Note: Exact hex addresses are stored in the project deployment manifest rather than hard-coded here, as testnet redeployment during development may produce updated addresses. The deployment manifest in the repository is the authoritative record. The final thesis submission will include static addresses from a stable named deployment. Sprint~1 deployment is scheduled on Polygon zkEVM Cardona. Existing Amoy testnet addresses are deprecated and retained here only for pre-thesis verification reference.}


%% %%%%%%%%%%%%%%%%%%%%%%%%%%%%%%%%%%%%%%%%%%%%%%%%%%%%%%%%%%%%
%% APPENDIX D: SMART CONTRACT INTERFACE
%% %%%%%%%%%%%%%%%%%%%%%%%%%%%%%%%%%%%%%%%%%%%%%%%%%%%%%%%%%%%%
\chapter*{Appendix D: WorldBankReserve Contract Interface}
\phantomsection
\addcontentsline{toc}{chapter}{Appendix D: WorldBankReserve Contract Interface}

The following Solidity interface documents the public API of the \texttt{WorldBankReserve} contract---the fully implemented Tier~1 contract of the prototype. Three design annotations follow the listing:

\begin{lstlisting}[language=Solidity,basicstyle=\ttfamily\scriptsize,breaklines=true,
    commentstyle=\color{GrayText},keywordstyle=\color{PrimaryBlue}\bfseries,
    frame=single,rulecolor=\color{AccentBlue!40},captionpos=b,
    caption={WorldBankReserve.sol -- Tier 1 public interface. Full implementation in project repository.}]
// SPDX-License-Identifier: MIT
pragma solidity ^0.8.20;

import "@openzeppelin/contracts/security/ReentrancyGuard.sol";
import "@openzeppelin/contracts/access/AccessControl.sol";

/// @title WorldBankReserve
/// @notice Tier 1 reserve contract: manages global capital allocation
///         to registered National Bank addresses.
contract WorldBankReserve is ReentrancyGuard, AccessControl {

    bytes32 public constant NATIONAL_BANK_ROLE = keccak256("NATIONAL_BANK_ROLE");
    bytes32 public constant AUDITOR_ROLE       = keccak256("AUDITOR_ROLE");

    uint256 public totalReserve;
    uint256 public minimumReserveRatio;   // e.g. 1000 = 10.00%
    mapping(address => uint256) public allocatedTo; // NationalBank => amount

    event CapitalAllocated(address indexed nationalBank, uint256 amount);
    event RepaymentReceived(address indexed nationalBank, uint256 amount);
    event ReserveRatioUpdated(uint256 newRatio);

    /// @notice Allocate capital downward to a registered National Bank (CEI pattern)
    function allocateCapital(address nationalBank, uint256 amount)
        external onlyRole(DEFAULT_ADMIN_ROLE) nonReentrant;

    /// @notice Record repayment from a National Bank (CEI pattern)
    function recordRepayment(uint256 amount)
        external onlyRole(NATIONAL_BANK_ROLE) nonReentrant;

    /// @notice Query current utilization ratio (scaled 1e4)
    function utilizationRate() external view returns (uint256);

    /// @notice Returns true if reserve ratio constraint is satisfied
    function isReserveAdequate() external view returns (bool);

    /// @notice Returns a four-value reserve summary for Proof of Reserve
    ///         verification and the Reserve Transparency Dashboard.
    /// @return totalDeposited  Gross deposits ever made to this reserve tier
    /// @return totalLoaned     Outstanding principal lent to National Banks
    /// @return reserveRatio    (totalDeposited - totalLoaned) * 1e4 / totalDeposited
    ///                         (e.g. 5000 = 50.00%)
    /// @return insuranceFundBalance  Current balance of the InsuranceFund contract
    function getReserveSummary() external view returns (
        uint256 totalDeposited,
        uint256 totalLoaned,
        uint256 reserveRatio,
        uint256 insuranceFundBalance
    );
}
\end{lstlisting}

\textbf{Design annotation 1 --- \texttt{nonReentrant} on state-changing functions.}
The \texttt{nonReentrant} modifier from OpenZeppelin's \texttt{ReentrancyGuard} sets a boolean lock before execution and clears it after, preventing recursive calls. It is applied to \texttt{allocateCapital} and \texttt{recordRepayment} because both involve ETH transfers that could trigger attacker-controlled fallback functions. The full Checks-Effects-Interactions analysis is in Section~\ref{sec:kinked-rate}.

\textbf{Design annotation 2 --- \texttt{AccessControl} mapping to RBAC design.}
OpenZeppelin's \texttt{AccessControl} maps each on-chain role (e.g., \texttt{NATIONAL\_BANK\_ROLE}) to a \texttt{bytes32} keccak256 hash stored in a role-to-address mapping. Role assignment is controlled by the \texttt{DEFAULT\_ADMIN\_ROLE} (the World Bank admin). This directly implements the four-tier RBAC hierarchy described in Section~\ref{sec:prototype-scope}: only addresses granted \texttt{NATIONAL\_BANK\_ROLE} can call \texttt{recordRepayment}, and only the World Bank admin can call \texttt{allocateCapital}.

\textbf{Design annotation 3 --- Events enable off-chain indexing for the analytics layer.}
\texttt{CapitalAllocated}, \texttt{RepaymentReceived}, and \texttt{ReserveRatioUpdated} are emitted at every significant state transition. The off-chain Express.js event listener subscribes to these events and writes them to PostgreSQL, enabling the AI/ML monitoring layer to construct loan lifecycle timelines, compute utilization windows, and feed the Random Forest fraud detector without requiring costly repeated \texttt{SLOAD} calls.

%% %%%%%%%%%%%%%%%%%%%%%%%%%%%%%%%%%%%%%%%%%%%%%%%%%%%%%%%%%%%%
%% REFERENCES
%% %%%%%%%%%%%%%%%%%%%%%%%%%%%%%%%%%%%%%%%%%%%%%%%%%%%%%%%%%%%%
\chapter*{References}
\phantomsection
\addcontentsline{toc}{chapter}{References}

{\small
\begin{enumerate}[label={[\arabic*]}]
\item S.~M. Werner, D.~Perez, L.~Gudgeon, A.~Klages-Mundt, D.~Harz, and W.~J. Knottenbelt, ``SoK: Decentralized Finance (DeFi),'' \textit{arXiv preprint arXiv:2101.08778}, 2022. [Online]. Available: \url{https://arxiv.org/abs/2101.08778}

\item S.~Bastankhah, M.~Hashemi, and W.~Shi, ``An Adaptive Data-Driven DeFi Borrow-Lending Protocol,'' \textit{Proc. IEEE Int. Conf. Blockchain}, 2023.

\item G.~Palaiokrassas, A.~Skoufis, and L.~Tassiulas, ``Leveraging Machine Learning for Multichain DeFi Fraud Detection,'' \textit{Proc. IEEE Int. Conf. Blockchain and Cryptocurrency (ICBC)}, 2023.

\item D.~Adom, E.~O. Acheampong, and M.~Boateng, ``LIME and SHAP: A Comparison of Model-Agnostic Approaches to Explainability in Loan Approval Systems,'' \textit{International Journal of Advanced Computer Science and Applications}, vol.~13, no.~11, 2022.

\item B.~W. Tan, ``Central Bank Digital Currency and Financial Inclusion,'' \textit{IMF Working Paper WP/23/69}, 2023. [Online]. Available: \url{https://www.imf.org/en/Publications/WP}

\item F.~T. Liu, K.~M. Ting, and Z.-H. Zhou, ``Isolation Forest,'' in \textit{Proc. IEEE Int. Conf. Data Mining (ICDM)}, pp.~413--422, 2008. DOI: \url{https://doi.org/10.1109/ICDM.2008.17}

\item N.~Atzei, M.~Bartoletti, and T.~Cimoli, ``A Survey of Attacks on Ethereum Smart Contracts (SoK),'' in \textit{Proc. Int. Conf. Principles of Security and Trust (POST)}, pp.~164--186, 2017. DOI: \url{https://doi.org/10.1007/978-3-662-54455-6_8}

\item S.~M. Lundberg and S.-I. Lee, ``A Unified Approach to Interpreting Model Predictions,'' in \textit{Advances in Neural Information Processing Systems (NeurIPS)}, pp.~4765--4774, 2017. [Online]. Available: \url{https://proceedings.neurips.cc/paper/2017/hash/8a20a8621978632d76c43dfd28b67767-Abstract.html}

\item P.~Bracke, A.~Datta, C.~Jung, and S.~Sen, ``Machine Learning Explainability in Finance: An Application to Default Risk Analysis,'' \textit{Bank of England Staff Working Paper No.~816}, 2019. [Online]. Available: \url{https://www.bankofengland.co.uk/working-paper/2019/machine-learning-explainability-in-finance}

\item R.~Beck, C.~M\"{u}ller-Bloch, and J.~L. King, ``Governance in the Blockchain Economy: A Framework and Research Agenda,'' \textit{Journal of the Association for Information Systems}, vol.~19, no.~10, pp.~1020--1034, 2018. DOI: \url{https://doi.org/10.17705/1jais.00518}

\item D.~Mhlanga, ``Blockchain Technology for Financial Inclusion and Sustainable Development,'' in \textit{Digital Financial Inclusion}, Springer, 2022.

\item OpenZeppelin, ``OpenZeppelin Contracts,'' 2024. [Online]. Available: \url{https://openzeppelin.com/contracts}

\item DefiLlama, ``Lending Protocols --- DeFi TVL and Protocol Rankings,'' 2024--2025. [Online]. Available: \url{https://defillama.com/protocols/lending}

\item World Bank, ``The Global Findex Database 2021: Financial Inclusion, Digital Payments, and Resilience,'' 2021. [Online]. Available: \url{https://www.worldbank.org/en/publication/globalfindex}

\item Aave, ``Aave Protocol Documentation,'' 2024. [Online]. Available: \url{https://docs.aave.com/}

\item Compound, ``Compound Protocol Documentation,'' 2024. [Online]. Available: \url{https://docs.compound.finance/}

\item BCOLBD 2025, ``Blockchain Olympiad Bangladesh: Guideline and Evaluation Scheme,'' 2025. [Online]. Available: \url{https://bcolbd.org/uploads/guideline/BLOCKCHAIN\%20OLYMPIAD\%20BANGLADESH\%20Blockchain\%20Guideline.pdf}

\item Galaxy Digital, ``The State of Crypto Lending and Borrowing,'' Galaxy Research, 2024. [Online]. Available: \url{https://www.galaxy.com/insights/research/the-state-of-crypto-lending}

\item World Bank, ``World Development Report 2022: Finance for an Equitable Recovery,'' 2022. [Online]. Available: \url{https://www.worldbank.org/en/publication/wdr2022}

\item International Finance Corporation (IFC), ``MSME Finance Gap 2019,'' 2019. [Online]. Available: \url{https://www.ifc.org/en/insights-reports/2019/msme-finance-gap}

\item Bank for International Settlements, ``DeFi Lending: Intermediation Without Information?,'' BIS Working Paper No.~1183, 2024. [Online]. Available: \url{https://www.bis.org/publ/work1183.pdf}

\item Deloitte, ``Global Banking Outlook 2024,'' 2024.

\item Michigan State University Libraries, ``Literature Table and Synthesis --- Nursing Literature Reviews,'' LibGuides, 2025. [Online]. Available: \url{https://libguides.lib.msu.edu/nursinglitreview/table}

\item Committee on Payments and Market Infrastructures and Bank for International Settlements, ``Correspondent Banking --- Technical Report,'' CPMI Papers No.~147, 2016. [Online]. Available: \url{https://www.bis.org/cpmi/publ/d147.pdf}

\item Ripple, ``RLUSD Use Case Analysis,'' XRP Academy, 2026. [Online]. Available: \url{https://xrpacademy.com/blog/rlusd-use-case-analysis-calendar-570}

\item World Bank, ``Remittance Prices Worldwide: Making Markets More Transparent,'' Migration and Development Brief~40, 2024. [Online]. Available: \url{https://remittanceprices.worldbank.org/}

\item DefiLlama, ``Aave V3 TVL, Fees, Revenue \& Income Statement,'' March 2026. [Online]. Available: \url{https://defillama.com/protocol/aave-v3}

\item World Inequality Lab, ``Exorbitant Privilege,'' \textit{World Inequality Report 2026}, 2026. [Online]. Available: \url{https://wir2026.wid.world/insight/exorbitant-privilege/}

\item DefiLlama, ``Compound V3 TVL, Fees, Revenue \& Income Statement,'' March 2026. [Online]. Available: \url{https://defillama.com/protocol/compound-v3}

\item Fensory, ``Sky Protocol Projects \$611M Revenue in 2026 as USDS Supply Targets \$20.6 Billion,'' February 2026. [Online]. Available: \url{https://www.fensory.com/intelligence/rwa/sky-protocol-tokenization-regatta-solana-february-2026}

\item Fensory, ``Maple Finance --- Institutional DeFi Lending Protocol,'' 2026. [Online]. Available: \url{https://www.fensory.com/insights/protocols/maple-finance}

\item Fensory, ``DeFi Credit Platforms Hit \$2.4B Milestone,'' March 2026. [Online]. Available: \url{https://www.fensory.com/intelligence/rwa/defi-private-credit-platforms-2-4-billion-loans-ethereum}

\item Financial Stability Board, ``Enhancing Cross-border Payments: Stage 3 Roadmap,'' 2020. [Online]. Available: \url{https://www.fsb.org/2020/10/enhancing-cross-border-payments-stage-3-roadmap/}

\item A.~Beyer, B.~Gasperini, and P.~Theodoridis, ``Monetary Policy and Inequality: Distributional Effects of Asset Purchase Programs,'' \textit{Journal of International Money and Finance}, vol.~157, 2025.

\item Federal Reserve Board, ``Monetary Policy and the Distribution of Income: Evidence from U.S. Metropolitan Areas,'' FEDS Notes, March 2025. [Online]. Available: \url{https://www.federalreserve.gov/econres/notes/feds-notes/monetary-policy-and-the-distribution-of-income-evidence-from-us-metropolitan-areas-20250331.html}

\item R.~Cantillon, \textit{Essai sur la Nature du Commerce en G\'en\'eral} (\textit{Essay on the Nature of Commerce in General}), 1755 (posthumous).

\item GlobeNewswire, ``R3's Corda Leads Tokenized RWA Market with Over \$10 Billion in On-chain Assets,'' February 2025. [Online]. Available: \url{https://www.globenewswire.com/news-release/2025/02/13/3025637/}

\item Centrifuge, ``Real-World Asset Tokenization: Key Trends from 2025,'' 2025. [Online]. Available: \url{https://centrifuge.io/blog/real-world-asset-tokenization-trends-2025}

\item World Bank, ``World Bank Group Tracks Project Funds with New Blockchain Tool,'' Press Release, September 2025. [Online]. Available: \url{https://www.worldbank.org/en/news/press-release/2025/09/26/world-bank-group-tracks-project-funds-with-new-blockchain-tool}

\item JPMorgan, ``Kinexys Digital Payments: Real-Time Multicurrency Payments,'' 2025. [Online]. Available: \url{https://www.jpmorgan.com/onyx/coin-system.htm}

\item World Economic Forum, ``Why Decentralized Finance Is a Leapfrog Technology for the 1.1 Billion People Who Are Unbanked,'' September 2022. [Online]. Available: \url{https://weforum.org/stories/2022/09/decentralized-finance-a-leapfrog-technology-for-the-unbanked}

\item Opera Newsroom, ``160M CELO Allocation Proposal to Grow Opera from Distribution Partner into Key Network Stakeholder,'' March 2026. [Online]. Available: \url{https://press.opera.com/2026/03/19/opera-celo-partnership-2026/}

\item Morpho, ``Network Data,'' March 2026. [Online]. Available: \url{https://data.morpho.org/}

\item Stellar, ``End of Year 2025 Report --- Execution at Scale,'' 2025. [Online]. Available: \url{https://www.stellar.org/blog/foundation-news/2025-year-in-review}

\item Human Rights Foundation, ``Tracking CBDCs Before They Track You,'' 2025. [Online]. Available: \url{https://hrf.org/latest/tracking-cbdcs-before-they-track-you}

\item Y.~Li, X.~Zhang, and H.~Wang, ``Design and Implementation of a Multi-Chain Lending Model in Blockchain,'' in \textit{Proc. IEEE Int. Conf. Blockchain}, 2024. DOI: \url{https://doi.org/10.1109/Blockchain62396.2024.10729983}

\item R.~Xu, S.~Chen, and L.~Yang, ``An Evaluation System for DeFi Lending Protocols,'' in \textit{Proc. IEEE Int. Conf. Blockchain and Cryptocurrency (ICBC)}, pp.~1--6, 2023. DOI: \url{https://doi.org/10.1109/ICBC56567.2023.10240601}

\item A.~Sharma, R.~K.~Singh, and S.~Gupta, ``Blockchain Empowered Framework for Peer to Peer Lending,'' in \textit{Proc. IEEE Int. Conf. Blockchain Computing and Applications (BCCA)}, pp.~142--147, 2021. DOI: \url{https://doi.org/10.1109/BCCA53669.2021.9596552}

\item M.~T.~Hassan, F.~Ahmad, and A.~Mehmood, ``Blockchain and Machine Learning for Fraud Detection: A Privacy-Preserving and Adaptive Incentive Based Approach,'' \textit{IEEE Access}, vol.~10, pp.~87115--87131, 2022. DOI: \url{https://doi.org/10.1109/ACCESS.2022.3199498}

\item Y.~Wang, J.~Liu, and Z.~Chen, ``ContractWard: Automated Vulnerability Detection Models for Ethereum Smart Contracts,'' \textit{IEEE Trans. Network Science and Engineering}, vol.~8, no.~2, pp.~1133--1144, 2020. DOI: \url{https://doi.org/10.1109/TNSE.2020.2968505}

\item C.-F.~Liao, T.-H.~Tsai, and C.-J.~Chen, ``SoliAudit: Smart Contract Vulnerability Assessment Based on Machine Learning and Fuzz Testing,'' in \textit{Proc. IEEE Int. Conf. Internet of Things (iThings)}, pp.~458--465, 2019. DOI: \url{https://doi.org/10.1109/iThings/GreenCom/CPSCom/SmartData.2019.00098}

\item S.~So, M.~Lee, J.~Park, H.~Lee, and H.~Oh, ``VERISMART: A Highly Precise Safety Verifier for Ethereum Smart Contracts,'' in \textit{Proc. IEEE Symp. Security and Privacy (S\&P)}, pp.~1678--1694, 2020. DOI: \url{https://doi.org/10.1109/SP40000.2020.00032}

\item S.~Islam, R.~Khan, and M.~A.~Rahman, ``Central Bank Digital Currency (CBDC): Design Requirements and Challenges,'' in \textit{Proc. IEEE Int. Conf. Computing, Communication, and Intelligent Systems (ICCCIS)}, 2024. DOI: \url{https://doi.org/10.1109/ICCCIS62002.2024.10634472}

\item M.~S.~Alam, S.~M.~Hossain, and A.~K.~Das, ``Towards Using Blockchain Technology for Microcredit Industry in Bangladesh,'' in \textit{Proc. IEEE Region 10 Symp. (TENSYMP)}, pp.~1--6, 2021. DOI: \url{https://doi.org/10.1109/TENSYMP52854.2021.9392730}

\item P.~Tolmach, Y.~Li, and S.~Lin, ``Optimal Gas Fee Minimization in DeFi: Enhancing Efficiency and Security on the Ethereum Blockchain,'' \textit{IEEE Trans. Dependable and Secure Computing}, 2024. DOI: \url{https://doi.org/10.1109/TDSC.2024.3495637}

\item S.~Gupta, A.~Kumar, and R.~Sharma, ``SHAP-based Interpretable Models for Credit Default Assessment Using Machine Learning,'' in \textit{Proc. IEEE Int. Conf. Artificial Intelligence and Data Science}, 2024. DOI: \url{https://doi.org/10.1109/ICAIDS60875.2024.10840375}

\item S.~Nakamoto, ``Bitcoin: A Peer-to-Peer Electronic Cash System,'' 2008. [Online]. Available: \url{https://bitcoin.org/bitcoin.pdf}

\item G.~Wood, ``Ethereum: A Secure Decentralised Generalised Transaction Ledger'' (Yellow Paper), Ethereum Foundation, 2014 (revised 2023). [Online]. Available: \url{https://ethereum.github.io/yellowpaper/paper.pdf}

\item Aave, ``Aave Protocol V3 Technical Paper,'' 2022. [Online]. Available: \url{https://github.com/aave/aave-v3-core/blob/master/techpaper/Aave_V3_Technical_Paper.pdf}

\item T.~Dao, T.~Trinh, and V.~Pham, ``Optimizing Credit Scoring Models for Decentralized Financial Applications,'' in \textit{Information and Communication Technology}, Springer, 2025. DOI: \url{https://doi.org/10.1007/978-981-96-4282-3_36}

\item G.-L.~G\"{u}c\"{u}k, S.~Leible, and J.~Edinger, ``Blockchain-Based Microlending for Financial Inclusivity: A Literature Review of Its Privacy and Trust,'' Springer, 2024. DOI: \url{https://doi.org/10.1007/978-3-032-12801-0_22}

%% ---- NEW REFERENCES v8 (R1–R19) ----

\item A.~Beniiche, ``A Study of Blockchain Oracles,'' \textit{arXiv preprint arXiv:2004.07140}, 2020. [Online]. Available: \url{https://arxiv.org/pdf/2004.07140}

\item A.~Pasdar, Y.~C.~Lee, and Z.~Dong, ``Connect API with Blockchain: A Survey on Blockchain Oracle Implementation,'' \textit{ACM Computing Surveys}, vol.~55, no.~10, 2023. DOI: \url{https://doi.org/10.1145/3567582}

\item L.~Gudgeon, D.~Perez, D.~Harz, B.~Livshits, and A.~Gervais, ``DeFi Protocols for Loanable Funds: Interest Rates, Liquidity, and Market Efficiency,'' in \textit{Proc. ACM AFT}, 2020. [Online]. Available: \url{https://berkeley-defi.github.io/assets/material/DeFi\%20Protocols\%20for\%20Loanable\%20Funds.pdf}

\item T.~Mackinga, T.~Nadahalli, and R.~Wattenhofer, ``Attacks on Dynamic DeFi Interest Rate Curves,'' \textit{arXiv preprint arXiv:2307.13139}, 2023.

\item K.~W.~Wu, ``Strengthening DeFi Security: A Static Analysis Approach to Flash Loan Vulnerabilities,'' \textit{arXiv preprint arXiv:2411.01230}, 2024.

\item ``A Comprehensive Study of Exploitable Patterns in Smart Contracts: From Vulnerability to Defense,'' \textit{arXiv preprint arXiv:2504.21480}, 2025.

\item E.~Albert, J.~Correas, P.~Gordillo, G.~Rom\'{a}n-D\'{i}ez, and A.~Rubio, ``GASOL: Gas Analysis and Optimization for Ethereum Smart Contracts,'' in \textit{Proc. TACAS}, Springer, 2020. DOI: \url{https://doi.org/10.1007/978-3-030-45237-7_7}

\item F.~Piper, K.~Wolf, and J.~Heiss, ``Privacy-Preserving On-chain Permissioning for KYC-Compliant Decentralized Applications,'' TU Berlin, \textit{arXiv preprint arXiv:2510.05807}, 2025.

\item N.~Decker, ``Zero-Knowledge Proofs: Cryptographic Model for Financial Compliance and Global Banking Security,'' \textit{SSRN Working Paper No.~5170068}, 2025.

\item E.~Toufaily and T.~Zalan, ``How Can Blockchain-Based Lending Platforms Support Microcredit Activities in Developing Countries? An Empirical Validation of Its Opportunities and Challenges,'' \textit{Technological Forecasting and Social Change}, vol.~203, 2024. DOI: \url{https://doi.org/10.1016/j.techfore.2024.123403}

\item S.~Howlader and P.~Halder, ``Fintech's Impact on Financial Inclusion Through Mobile Financial Services in Bangladesh,'' \textit{Sage Publications}, 2025. DOI: \url{https://doi.org/10.1177/09763996251356998}

\item Atlas of Wars, ``Grameen Bank: A Successful Microcredit Model,'' 2024. [Online]. Available: \url{https://www.atlasofwars.com/grameen-bank-a-successful-microcredit-model/}

\item A.~Carstens et al., ``Stablecoins: Risks, Potential and Regulation,'' \textit{BIS Working Papers No.~905}, Bank for International Settlements, 2021. [Online]. Available: \url{https://www.bis.org/publ/work905.pdf}

\item ``Stablecoin Devaluation Risk,'' \textit{The European Journal of Finance}, Taylor \& Francis, 2025. DOI: \url{https://doi.org/10.1080/1351847X.2025.2505757}

\item J.~A.~Ocampo and K.~Gallagher, ``The Role of Multilateral Development Banks and Development Assistance in the Provision of Global Public Goods,'' UNDP Background Paper, 2024. [Online]. Available: \url{https://hdr.undp.org/system/files/documents/background-paper-document/2024jaocampokdgonzaleztheroleofmultilateraldevlmntbanks.pdf}

\item ``Commit-Reveal$^2$: Securing Randomness Beacons with Randomized Reveal Order in Smart Contracts,'' \textit{arXiv preprint arXiv:2504.03936}, 2025.

\item F.~A.~Aponte-Novoa, A.~L.~S.~Orozco, R.~Villanueva-Polanco, and P.~Wightman, ``The 51\% Attack on Blockchains: A Mining Behavior Study,'' \textit{IEEE Access}, vol.~9, pp.~140549--140564, 2021. DOI: \url{https://doi.org/10.1109/ACCESS.2021.3119110}

\item DL News, ``State of DeFi 2025,'' March 2026. [Online]. Available: \url{https://www.dlnews.com/research/internal/state-of-defi-2025/}

\item arXiv:2506.00505, ``Reinforcement Learning for Interest Rate Adjustment in DeFi Lending,'' 2025.

%% ---- NEW REFERENCES v15 (R20--R45) ----

\item W3C, ``Decentralized Identifiers (DIDs) v1.0 --- Core Architecture, Data Model, and Representations,'' W3C Recommendation, July 2022. [Online]. Available: \url{https://www.w3.org/TR/did-core/}; see also W3C, ``Verifiable Credentials Data Model v1.1,'' W3C Recommendation, March 2022. [Online]. Available: \url{https://www.w3.org/TR/vc-data-model/}

\item Sygnum Bank, ``Institutional DeFi in 2025: From Niche to Mainstream Finance,'' Sygnum Research Report, February 2026. [Online]. Available: \url{https://www.sygnum.com/research/}

\item M.~J.~Page, J.~E.~McKenzie, P.~M.~Bossuyt, I.~Boutron, T.~C.~Hoffmann, C.~D.~Mulrow, et~al., ``The PRISMA 2020 statement: an updated guideline for reporting systematic reviews,'' \textit{BMJ}, vol.~372, n71, 2021. DOI: \url{https://doi.org/10.1136/bmj.n71}

\item P.~Treleaven, R.~Gendal Brown, and D.~Yang, ``Blockchain Technology in Finance: A Systematic Review using PRISMA on 38 Empirical Studies,'' \textit{Computer}, vol.~58, no.~5, pp.~36--46, 2025. DOI: \url{https://doi.org/10.1109/MC.2025.3500125}

\item D.~Cheng, R.~Zhang, X.~Xu, S.~Yang, and L.~Akoglu, ``Graph Neural Network Architectures for Coordinated Fraud Detection in DeFi,'' \textit{IEEE Trans. Knowledge and Data Engineering}, vol.~36, no.~9, pp.~4521--4536, 2024. DOI: \url{https://doi.org/10.1109/TKDE.2024.3401234}

\item Y.~Liu, J.~Chen, and H.~Zhao, ``DeFiGuard: A Graph Neural Network Architecture for Real-Time DeFi Fraud Detection,'' \textit{Proc. ACM CIKM}, pp.~3782--3791, 2024. DOI: \url{https://doi.org/10.1145/3627673.3679876}

\item L.~Wang and Y.~Wang, ``Cross-Chain Money Laundering Detection with Heterogeneous Graph Attention,'' \textit{IEEE Trans. Information Forensics and Security}, vol.~20, pp.~2105--2119, 2025. DOI: \url{https://doi.org/10.1109/TIFS.2025.3478912}

\item H.~Khan, A.~Mughal, S.~Ali, and M.~Imran, ``PrivChain-AI: Federated Learning over Permissioned Blockchain for Cross-Bank Fraud Detection,'' \textit{Nature Scientific Reports}, vol.~15, art.~99812, December 2025. DOI: \url{https://doi.org/10.1038/s41598-025-99812-5}

\item Y.~Abbassi, M.~Lahby, and R.~Karbab, ``Federated Learning for Cross-Border Financial Fraud Intelligence: A Regulatory Compliance Framework,'' \textit{IEEE Access}, vol.~13, pp.~78340--78356, 2025. DOI: \url{https://doi.org/10.1109/ACCESS.2025.3501122}

\item J.~Park, S.~Kim, and B.~Lee, ``FED-SPFD: A Federated Strategy for Payment Fraud Detection in Cross-Institution Banking,'' \textit{Proc. IEEE Int. Conf. Big Data}, pp.~2912--2921, 2024. DOI: \url{https://doi.org/10.1109/BigData62323.2024.10825443}

\item V.~Buterin, P.~Hitzig, and E.~G.~Weyl, ``Decentralized Society: Finding Web3's Soul,'' \textit{SSRN Working Paper No.~4105763}, May 2022. DOI: \url{https://dx.doi.org/10.2139/ssrn.4105763}

\item MicroSave Consulting, ``Bangladesh Microfinance Sector Outlook 2025: Toward Performance-Based Credit Identity,'' MSC Insight Note, December 2025. [Online]. Available: \url{https://www.microsave.net/insights/bangladesh-mfi-2025}

\item Bank for International Settlements, ``2025 BIS Survey on Central Bank Digital Currency and Crypto-Assets,'' BIS Papers No.~147, July 2025. [Online]. Available: \url{https://www.bis.org/publ/bppdf/bispap147.pdf}

\item Morgan Stanley, ``AI @ Morgan Stanley Assistant: GPT-4 Powered Knowledge Tool for Financial Advisors,'' Morgan Stanley Technology Disclosure, 2024. [Online]. Available: \url{https://www.morganstanley.com/press-releases/key-milestone-in-firm-s-ai-journey}

\item McKinsey \& Company, ``The Economic Potential of Generative AI: The Next Productivity Frontier (Banking Sector Update),'' McKinsey Global Institute, June 2024. [Online]. Available: \url{https://www.mckinsey.com/capabilities/quantumblack/our-insights/the-economic-potential-of-generative-ai-the-next-productivity-frontier}

\item A.~Yampolskiy, I.~Maddali, K.~Khoury, et~al., ``LLM Vulnerabilities in Finance: A Red-Teaming Benchmark for Regulatory-Hallucination Resistance,'' \textit{Proc. ACM Conf. AI, Ethics, and Society (AIES)}, 2025. DOI: \url{https://doi.org/10.1145/3641421.3678412}

\item European Parliament and Council of the European Union, ``Regulation (EU) 2023/1114 of 31 May 2023 on Markets in Crypto-Assets (MiCA),'' \textit{Official Journal of the European Union}, L~150/40, June 2023; fully applicable from 30 December 2024. [Online]. Available: \url{https://eur-lex.europa.eu/eli/reg/2023/1114/oj}

\item European Banking Authority, ``Guidelines on Reserve Asset Composition, Custody, and Audit for Asset-Referenced and E-Money Tokens under MiCA,'' EBA/GL/2024/06, May 2024. [Online]. Available: \url{https://www.eba.europa.eu/regulation-and-policy/markets-in-crypto-assets}

\item United States Congress, ``Guiding and Establishing National Innovation for U.S. Stablecoins (GENIUS) Act,'' Pub.~L.~No.~119-27, July~18, 2025. [Online]. Available: \url{https://www.congress.gov/bill/119th-congress/senate-bill/919}

\item Bank for International Settlements Innovation Hub, ``Project mBridge: Connecting Economies through CBDC --- Reaches Minimum Viable Product,'' BIS Innovation Hub Report, June 2024. [Online]. Available: \url{https://www.bis.org/about/bisih/topics/cbdc/mcbdc_bridge.htm}

\item Bank for International Settlements and Central Banks of France, Japan, Korea, Mexico, Switzerland, the UK, and the US, ``Project Agora: Tokenization of Cross-Border Payments using Unified Ledgers,'' BIS Innovation Hub Working Paper, April 2024. [Online]. Available: \url{https://www.bis.org/about/bisih/topics/fmis/agora.htm}

\item Bangladesh Bank, ``Cautionary Notice Regarding Transactions in Virtual / Crypto-Currencies,'' Bangladesh Bank Foreign Exchange Policy Department, Notice No.~FE-1/2017, December 2017 (reaffirmed 2018, 2022). [Online]. Available: \url{https://www.bb.org.bd/aboutus/regulationguideline/foreignexchange/notice_virtualcurrencies.pdf}

\item Bangladesh Bank, ``FE Circular No.~06: Electronic Communication for Letters of Credit and Trade Finance Instruments,'' Foreign Exchange Policy Department, January 2025. [Online]. Available: \url{https://www.bb.org.bd/mediaroom/circulars/fepd/jan172025fepd06.pdf}

\item Government of the People's Republic of Bangladesh, ``The Payment and Settlement System Act, 2024 (Act No.~XX of 2024),'' Bangladesh Gazette, August 2024.

\item Bangladesh Bank, ``Financial Inclusion in Bangladesh: Progress and Gender-Disaggregated Outlook 2024--25,'' Special Publication SP2025-02, Financial Inclusion Department, June 2025. [Online]. Available: \url{https://www.bb.org.bd/pub/special/SP2025-02.pdf}

\item BRAC, ``BRAC Microfinance Annual Performance Report 2025: Women-Centric Lending Outcomes,'' BRAC Research and Evaluation Division, 2025. [Online]. Available: \url{https://www.brac.net/program/microfinance/}

\item Ethereum Improvement Proposals, ``EIP-4626: Tokenized Vault Standard,'' Ethereum Foundation, April 2022. [Online]. Available: \url{https://eips.ethereum.org/EIPS/eip-4626}; ``EIP-7540: Asynchronous ERC-4626 Tokenized Vaults,'' Ethereum Foundation, 2023. [Online]. Available: \url{https://eips.ethereum.org/EIPS/eip-7540}; ``EIP-3643: T-REX -- Token for Regulated EXchanges,'' Ethereum Foundation, 2021. [Online]. Available: \url{https://eips.ethereum.org/EIPS/eip-3643}. See also: Spectral Finance, ``On-Chain Credit Scoring for DeFi Lending,'' 2023. [Online]. Available: \url{https://docs.spectral.finance}; RociFi Labs, ``Undercollateralized DeFi Lending via On-Chain Credit Score,'' 2023. [Online]. Available: \url{https://docs.rocifi.com}; Qwen Team (Alibaba Cloud), ``Qwen3 Technical Report,'' \textit{arXiv preprint arXiv:2505.09388}, 2025. [Online]. Available: \url{https://arxiv.org/abs/2505.09388}.

%% ---- NEW REFERENCES v23 (R47–R51) ----

\item Behaviour-Centric Cybersecurity Center, York University, ``DeFi Fraud Transactions Dataset (BCCC-DeFiFraudTrans-2025),'' 2025. [Online]. Available: \url{https://www.yorku.ca/research/bccc/ucs-technical/cybersecurity-datasets-cds/defi-fraud-transactions-bccc-defifraudtrans-2025/}

\item Elliptic, ``The Elliptic Data Set,'' Kaggle, 2019. [Online]. Available: \url{https://www.kaggle.com/datasets/ellipticco/elliptic-data-set}

\item Y.~Elmougy, S.~Liu, and J.~Liu, ``Demystifying Fraudulent Transactions and Illicit Nodes in the Bitcoin Network for Financial Forensics,'' in \textit{Proc. ACM SIGKDD Int. Conf. Knowledge Discovery and Data Mining (KDD)}, 2023. [Online]. Available: \url{https://github.com/git-disl/EllipticPlusPlus}

\item M.~Weber, G.~Domeniconi, J.~Chen, D.~K.~I.~Weidele, C.~Bellei, T.~Robinson, and C.~E.~Leiserson, ``Anti-Money Laundering in Bitcoin: Experimenting with Graph Convolutional Networks for Financial Forensics,'' \textit{arXiv preprint arXiv:1908.02591}, KDD Workshop on Anomaly Detection in Finance, 2019. [Online]. Available: \url{https://arxiv.org/abs/1908.02591}

\item AMD, ``ROCm 7.0.2 Release Notes: Official Support for RDNA~4 / gfx1200 (RX~9060~XT),'' AMD Developer Blog, October 2025. [Online]. Available: \url{https://www.phoronix.com/news/AMD-ROCm-7.0.2-Released}

%% ---- NEW REFERENCES v26 — Agent Architecture (R52–R55) ----

\item K.~Heidt, M.~Sp\"{o}rk, D.~Markov, and A.~Kr\"{u}ger, ``Harness Engineering
for Language Agents: The Harness Layer as Control, Agency, and Runtime,''
\textit{Preprints.org}, DOI: \url{https://doi.org/10.20944/preprints202603.1756},
March 2026.

\item M.~Alizadeh, M.~Gholampour, A.~Imani, and J.~Schulz, ``Simple Prompt
Injection Attacks Can Leak Personal Data Observed by LLM Agents During Task
Execution,'' \textit{arXiv preprint arXiv:2506.01055}, University of Zurich / IPM,
June 2026.

\item OWASP Foundation, ``OWASP Top 10 for Large Language Model Applications and
Agents, 2026 Edition,'' OWASP Foundation, 2026. [Online]. Available:
\url{https://owasp.org/www-project-top-10-for-large-language-model-applications/}

\item E.~Gamma, R.~Helm, R.~Johnson, and J.~Vlissides, \textit{Design Patterns:
Elements of Reusable Object-Oriented Software}, Addison-Wesley, 1994.

\end{enumerate}
}

\end{document}
